% !TEX program = lualatex
% SirixDB: A Bitemporal Database System — in-depth technical report
% Compile with:  latexmk -lualatex -interaction=nonstopmode SIRIX_ARCHITECTURE_REPORT.tex
% (biblatex needs biber; latexmk runs lualatex -> biber -> lualatex -> lualatex automatically)
% PDF/A-2b export: gs -dPDFA=2 -dBATCH -dNOPAUSE -dNOSAFER -sColorConversionStrategy=RGB
%   -sDEVICE=pdfwrite -dPDFACompatibilityPolicy=1 -sOutputFile=<out>.pdf /tmp/PDFA_def.ps <in>.pdf
\documentclass[11pt,a4paper,oneside]{report}

%============================ Fonts (lualatex) ============================
\usepackage{fontspec}
\setmainfont{TeX Gyre Pagella}
\setsansfont{TeX Gyre Heros}[Scale=0.92]
\setmonofont{DejaVu Sans Mono}[Scale=0.80]
\usepackage{microtype}

%============================ Page geometry ==============================
\usepackage[a4paper,top=2.6cm,bottom=2.6cm,inner=2.35cm,outer=2.35cm]{geometry}
\setlength{\parskip}{0.45em}
\setlength{\parindent}{0pt}
\linespread{1.05}
% Let TeX break before long unbreakable \texttt identifiers rather than overflowing.
\setlength{\emergencystretch}{3em}
\tolerance=2500
\hbadness=2500

%============================ Math & theorems ============================
\usepackage{amsmath}
\usepackage{amssymb}
\usepackage{amsthm}

%============================ Tables & lists ============================
\usepackage{array}
\usepackage{booktabs}
\usepackage{tabularx}
\usepackage{longtable}
\usepackage{xltabular}
\usepackage{alphalph}
\usepackage{ragged2e}
\usepackage{enumitem}
\setlist{topsep=2pt,itemsep=1pt,parsep=0pt,leftmargin=1.4em}

%============================ Graphics ============================
\usepackage{graphicx}
\graphicspath{{figures/}}
\usepackage{xcolor}
\definecolor{ember}{HTML}{B5532A}
\definecolor{slate}{HTML}{233140}
\definecolor{inkblue}{HTML}{1F4E79}
\definecolor{boxrule}{HTML}{9AA5AD}
\definecolor{boxbg}{HTML}{F4F2EE}
\definecolor{layerbg}{HTML}{EAE6DF}
\definecolor{nodeold}{HTML}{6B7B8C}
\definecolor{nodenew}{HTML}{B5532A}
\usepackage{tikz}
\usetikzlibrary{positioning,fit,backgrounds,calc,shapes.geometric,arrows.meta,decorations.pathreplacing}

%============================ Verbatim diagrams / algorithms ============================
\usepackage{fancyvrb}
\usepackage{fvextra}
% Pure-literal verbatim (no commandchars): ASCII box-drawing passes through unchanged.
\fvset{frame=single,framerule=0.3pt,rulecolor=\color{boxrule},framesep=6pt,%
       xleftmargin=3pt,xrightmargin=3pt}
\DefineVerbatimEnvironment{diagram}{Verbatim}{fontsize=\footnotesize}
\DefineVerbatimEnvironment{diagramS}{Verbatim}{fontsize=\scriptsize}
\newcommand{\dcap}[1]{\par\vspace{1pt}\noindent{\small\sffamily\bfseries\color{slate}#1}\par\vspace{-2pt}}

%============================ Callout boxes ============================
\usepackage{tcolorbox}
\tcbuselibrary{skins,breakable}
\newtcolorbox{thesisbox}{enhanced,breakable,colback=ember!7,colframe=ember!65!black,
  boxrule=0.6pt,arc=2pt,left=10pt,right=10pt,top=7pt,bottom=7pt}
\newtcolorbox{honesty}{enhanced,breakable,colback=boxbg,colframe=black!28,
  boxrule=0.4pt,arc=1pt,left=8pt,right=8pt,top=5pt,bottom=5pt}

%============================ Sectioning style ============================
\usepackage{titlesec}
\titleformat{\chapter}[display]
  {\normalfont\sffamily\bfseries}
  {\color{ember}\large\MakeUppercase{\chaptertitlename}\ \thechapter}
  {6pt}{\Huge\color{slate}}
\titlespacing*{\chapter}{0pt}{4pt}{20pt}
\titleformat{\section}{\normalfont\Large\sffamily\bfseries\color{slate}}{\thesection}{0.6em}{}
\titleformat{\subsection}{\normalfont\large\sffamily\bfseries\color{slate!92}}{\thesubsection}{0.6em}{}
\titleformat{\subsubsection}{\normalfont\sffamily\bfseries\color{slate!85}}{\thesubsubsection}{0.6em}{}

%============================ Captions, headers ============================
\usepackage{caption}
\captionsetup{font=small,labelfont={bf,sf},labelsep=period}
\usepackage{fancyhdr}
\pagestyle{fancy}
\fancyhf{}
\renewcommand{\headrulewidth}{0.2pt}
\fancyhead[L]{\small\sffamily\nouppercase{\leftmark}}
\fancyhead[R]{\small\sffamily\thepage}
\renewcommand{\chaptermark}[1]{\markboth{\thechapter.\ #1}{}}
\usepackage{needspace}

%============================ Theorem styles ============================
\newtheoremstyle{sirixthm}{9pt}{9pt}{\normalfont}{0pt}
  {\bfseries\sffamily\color{slate}}{.}{.5em}{}
\theoremstyle{sirixthm}
\newtheorem{property}{Property}[chapter]
\newtheorem{theorem}[property]{Theorem}

%============================ Hyperref (load late) ============================
\usepackage[bookmarks=true,bookmarksnumbered=true]{hyperref}
\hypersetup{
  colorlinks=true, linkcolor=slate, citecolor=ember, urlcolor=inkblue,
  pdftitle={SirixDB: A Bitemporal Database System},
  pdfauthor={The SirixDB Project},
  pdfsubject={Architecture, Algorithms, and Evaluation}}

% Bibliography (biblatex + biber)
\usepackage[backend=biber,style=numeric-comp,sorting=nyt,giveninits=true,maxbibnames=8]{biblatex}
\addbibresource{references.bib}
\defbibnote{refsnote}{\noindent\textit{Indicative; the works this design and document draw upon.}\par\medskip}

% Inline-code helper (camelCase identifiers need no escaping; escape _ # & % { } by hand).
\newcommand{\code}[1]{\texttt{#1}}
% Path helper: breaks cleanly at / and . , renders _ literally (no escaping needed).
\usepackage{url}
\DeclareUrlCommand\pp{\urlstyle{tt}}
\Urlmuskip=0mu plus 2mu\relax  % stretchable glue so long paths break instead of overflowing

\begin{document}

%================================ TITLE PAGE ================================
\begin{titlepage}
\centering
\vspace*{1.2cm}
{\color{ember}\sffamily\large\bfseries AN IN-DEPTH TECHNICAL REPORT}\\[2pt]
{\color{boxrule}\rule{0.5\linewidth}{0.4pt}}\\[1.6cm]
{\sffamily\bfseries\color{slate}\fontsize{30}{34}\selectfont SirixDB\par}
\vspace{0.45cm}
{\sffamily\bfseries\color{slate}\LARGE A Bitemporal\\[4pt] Database System\par}
\vspace{0.9cm}
{\large\itshape Architecture, Algorithms, and Evaluation\par}
\vspace{1.7cm}
{\color{boxrule}\rule{0.7\linewidth}{0.4pt}}\\[1.0cm]
{\large An in-depth technical report synthesised from the source tree and design\\
documents of \code{sirix-core}, \code{sirix-enterprise}, \code{brackit},\\
and the SirixDB Web GUI.\par}
\vfill
{\sffamily\large\bfseries\color{slate} Johannes Lichtenberger\par}
\vspace{0.3cm}
{\large The SirixDB Project\par}
\vspace{0.7cm}
{\color{boxrule} Working draft \textperiodcentered\ June 2026}\\[2pt]
\vspace*{0.6cm}
\end{titlepage}

%================================ ABSTRACT ================================
\chapter*{Abstract}
\addcontentsline{toc}{chapter}{Abstract}
\markboth{Abstract}{}

SirixDB is a bitemporal, versioned document store for JSON and XML, built on a single idea:
\textbf{immutability with structural sharing}. Every write transaction publishes a new,
immutable \emph{revision} of a resource by copying only the pages on the path from a change
to the tree's root (copy-on-write); every unchanged page is shared by reference with all
prior revisions. Revisions order data in \emph{transaction time}, and resources may also
carry user-defined \emph{valid-time} bounds, so the store is bitemporal, not merely
versioned. Most of the remaining properties follow from the immutability decision alone:
time-travel reads are direct traversals instead of reconstructions, durability is
append-only with no separate write-ahead log, snapshot isolation on the read path is
lock-free, and the design adds sub-document field history, change tracking, and hash-guided
semantic diffs.

This document develops the architecture bottom-up across fourteen chapters: the
copy-on-write page trie and node model (Ch.~2); the sliding-snapshot page-versioning
algorithm and its reconstruction correctness (Ch.~3); the self-describing,
checksum-protected on-disk format, its dual-beacon durable-commit protocol, the
evolution of that protocol from seven sync calls per commit to one explicit fsync,
and the crash-recovery contract validated under SIGKILL fault injection (Ch.~4); the
Height-Optimised Trie (HOT) used for index and intra-page navigation, with its
invariant catalogue and the engineering campaign that reduced a reproducer's
violations from 127 to 1 (Ch.~5); the path summary with incremental statistics and
the full secondary-index family of value, name, DeweyID, columnar, and HNSW vector indexes
(Ch.~6); the brackit JSONiq/XQuery engine and its
cost-based, index-aware optimiser (Ch.~7); the columnar/SIMD vectorised execution
path, with its differential-testing discipline and deliberately narrow, fail-closed scope (Ch.~8);
the enterprise io\_uring, S3, and Kafka back-ends, including the falsification of the
io\_uring linked-commit chain by the W1 benchmark (Ch.~9); the embedded and REST APIs, covering the single-writer/$N$-reader transaction model and the Vert.x server (Ch.~10);
and the SolidJS web front-end (Ch.~11). Chapter~12 presents a rigorous, same-machine evaluation
against
PostgreSQL that PostgreSQL \emph{wins} in its home regime, and explains precisely why
and where each system's advantages lie. Chapter~13 situates the work; Chapter~14
concludes and lists the open frontier.

\vspace{0.4cm}
\begin{honesty}
\small This is an uncommitted working document. It privileges honesty over polish:
where a subsystem is partially implemented, aspirational, or carries known technical
debt, this is stated explicitly. Benchmark numbers are reproduced from the project's
own \code{docs/COMPARISON\_POSTGRES.md} and \code{docs/BENCHMARKS.md} with their
caveats intact.
\end{honesty}

%================================ TOC ================================
\microtypesetup{protrusion=false}
\tableofcontents
\microtypesetup{protrusion=true}
\cleardoublepage

%================================ NOTATION ================================
\chapter*{Notation and Abbreviations}
\addcontentsline{toc}{chapter}{Notation and Abbreviations}
\markboth{Notation and Abbreviations}{}

\begin{minipage}[t]{0.47\linewidth}
\small
{\sffamily\bfseries\color{slate} Symbols}\par\vspace{3pt}
\begin{tabularx}{\linewidth}{@{}l >{\RaggedRight}X@{}}
\toprule
$N$ & a revision number\\
$R$ & sliding-snapshot window (\code{revisionsToRestore}, default 3)\\
$n$ & document size (pages or nodes)\\
$k$ & pages touched by a write\\
$V,\,P$ & values written; distinct paths\\
$\beta$ & a HOT discriminative bit\\
$O(\cdot)$ & asymptotic cost\\
$\lceil\cdot\rceil$ & ceiling\\
$\subseteq$ & sub-mask / subset (HOT routing)\\
$\circ$ & function composition\\
\bottomrule
\end{tabularx}
\end{minipage}\hfill
\begin{minipage}[t]{0.49\linewidth}
\small
{\sffamily\bfseries\color{slate} Abbreviations}\par\vspace{3pt}
\begin{tabularx}{\linewidth}{@{}l >{\RaggedRight}X@{}}
\toprule
CoW & copy-on-write\\
TIL & Transaction Intent Log\\
HOT & Height-Optimised Trie\\
PCR & path-class record\\
CAS & content-and-structure (index)\\
PAX & Partition Attributes Across\\
PEXT/PDEP & parallel bit extract / deposit\\
FSST & fast static symbol table\\
HNSW & Hierarchical Navigable Small World\\
SPI & service-provider interface\\
CDC & change data capture\\
FFM & Foreign Function \& Memory API\\
SQE/CQE & io\_uring submission / completion entry\\
FUA & force unit access\\
W1--W6 & the evaluation workloads (Ch.~\ref{ch:eval})\\
\bottomrule
\end{tabularx}
\end{minipage}

\vspace{0.6cm}
\noindent\small Conceptual terms (beacon, fragment, path class, sliding snapshot, \ldots) are
defined in the Glossary, \S\ref{app:glossary}.
\cleardoublepage

%================================================================
\chapter{Introduction: Claim, Principles, and Contributions}
\label{ch:intro}

\section{The problem}

Most databases overwrite data in place. Where they keep history at all, it is usually
bolted on: a write-ahead log that checkpointing consumes and recycles, audit tables
driven by triggers, or temporal-table machinery sitting on top of a mutable heap. In all
of these, reading the past means rebuilding it, whether by replaying a log, filtering on
a validity interval, or joining against a history table.

This matters more than it once did. Audit and compliance regimes increasingly demand a provable record
of what a dataset said and when; debugging a production incident often means asking what a document
looked like at the moment things went wrong; analytics over how data \emph{changed}, not merely its
current state, underpins everything from fraud detection to product metrics; and reproducible machine
learning needs the exact training snapshot rather than today's mutated table. In each case the question
is temporal, and a store that overwrites in place can answer it only by reconstructing the past from a
side channel that was never designed to be queried as data. The premise of this report is that history
is better treated as a first-class property of the store than as something bolted on after the fact.

SirixDB inverts that arrangement. A resource (a JSON or XML document of arbitrary size)
is stored as an append-only sequence of immutable page fragments, and every commit
publishes a new revision that shares all unchanged structure with earlier ones. So a
request for the document as of revision~17, or as of last Tuesday at 14:00, is not a
reconstruction. The snapshot is already on disk, and a plain traversal reaches it without
a temporal predicate or a join.

\section{The central claim}

\begin{thesisbox}
\itshape A disciplined commitment to immutability and copy-on-write, executed through a
high-fan-out page trie, a sliding-snapshot page-versioning algorithm, and a
self-describing dual-beacon on-disk format, yields a document store whose temporal
capabilities (time travel, sub-document history, lock-free snapshot reads, hash-guided
semantic diffs) are intrinsic rather than emulated, at a per-version storage and write
cost of $O(\text{change})$ plus fixed metadata instead of $O(\text{document})$.
\end{thesisbox}

This is not a claim that SirixDB is universally faster. Chapter~\ref{ch:eval} shows it is
not, in the small-document regime where a relational engine is at home. What matters here
is the shape of the cost and what the design makes possible. SirixDB pays a fixed
per-revision metadata floor and runs a deliberately strong durability protocol, and in
exchange it gets a class of temporal operations that a \code{jsonb} column with a trigger
cannot express at any document size.

\section{Lineage}

SirixDB sits at the confluence of two lines of native-database research. Most of its ideas were first
worked out for XML, and SirixDB's contribution is to carry them over to JSON.

The storage model descends from \textbf{Treetank}, a versioned XML store built at the University of
Konstanz~\cite{treetank}. Its founding idea, a single large persistent tree that retains every version
by copying only what changed, is Marc Kramis's~\cite{kramis} (an earlier prototype was named
\emph{Idefix}). Sebastian Graf's doctoral work~\cite{graf} then developed the node-granular versioning
and the \emph{sliding snapshot} page-versioning strategy that SirixDB still uses, and SirixDB is the
continuation of that system, re-engineered and extended to JSON. The persistent-tree heritage places it
in the wider tradition of immutability with structural sharing (Okasaki~\cite{okasaki}) and of
copy-on-write filesystems (WAFL~\cite{wafl}, ZFS, Btrfs).

The second line is the native-XML-database work of Theo H\"arder's group at TU Kaiserslautern, built
around the XTC system, and two of its results are load-bearing here. The \textbf{node labeling} is an
ORDPATH-style scheme (O'Neil et al.~\cite{ordpath}) of the kind H\"arder, Haustein, Mathis and Wagner
reconsidered and refined into stable path labeling identifiers~\cite{haerderlabels}: hierarchical,
insert-friendly identifiers whose compressed binary form yields document order under a plain byte
comparison (\S\ref{sec:deweyids}). And the \textbf{query engine} is \emph{brackit}, the
XQuery~3.1/JSONiq processor that came out of Sebastian B\"achle's doctoral work in that same
group~\cite{baechle}; SirixDB embeds it and layers a cost-based, index-aware optimiser on top
(Ch.~\ref{ch:query}).

The through-line is that versioned node storage, DeweyID labeling, and a set-oriented XQuery engine were
all native-XML ideas, and SirixDB makes them first-class for JSON and JSONiq on a modern off-heap,
flyweight-node engine. Its trie and index machinery additionally draw on recent main-memory index
research (ART~\cite{leisart} and the Height-Optimised Trie, or HOT, of Binna et al.~\cite{binna}).
Chapter~\ref{ch:related} develops these connections.

\section{Design principles}
\label{sec:principles}

The chapters that follow describe many mechanisms, but the system is best read as the disciplined
application of five commitments. Most of its properties are corollaries of these principles rather than
independent features, and most of its limitations are the points at which a principle is deliberately
pushed to its boundary. Naming them here is the report's \emph{red thread}: each later chapter develops one
or more of them and states what it costs.

\begin{enumerate}[label=\textbf{P\arabic*.},leftmargin=2.6em]
\item \textbf{Immutability with structural sharing.} Nothing is overwritten. A commit copies only the
  pages on the path from a change to the root and shares everything else by reference (Ch.~\ref{ch:pds}).
  This is the root commitment from which the rest follow: it makes every revision a consistent, lock-free
  snapshot (Property~\ref{prop:snapshot}) and reduces durability to pure append.

\item \textbf{Bounded cost, not free history.} Retaining all history is made affordable by bounding
  \emph{both} sides of the trade --- a read reconstructs a page from at most $R$ fragments, and a write
  preserves each record at most once per window pass (the sliding snapshot, Properties~\ref{prop:bounded}
  and~\ref{prop:amort}). Per-version cost is therefore $O(\text{change})$, not $O(\text{document})$
  (Ch.~\ref{ch:versioning}; measured in Ch.~\ref{ch:eval}).

\item \textbf{Separation of the logical from the physical.} One logical model admits many physical
  realisations: a query plan runs serially, in parallel, morsel-driven, or vectorised without being
  rewritten (physical data independence, Property~\ref{prop:pdi}), and a resource moves between storage
  backends byte-for-byte (Property~\ref{prop:backendid}). New capability joins at a service-provider seam,
  not by forking the engine --- the \emph{separation of concerns} that the Kaiserslautern lineage above
  made its method~\cite{baechle}.

\item \textbf{A store that verifies itself.} The format carries enough to check its own integrity: a
  self-describing on-disk layout, per-page and per-beacon checksums, two redundant beacons that survive a
  torn write (Property~\ref{prop:torn}), and rolling node hashes that serve both integrity and the
  hash-guided semantic diff (Ch.~\ref{ch:ondisk}). Crash recovery is a contract exercised under SIGKILL
  fault injection, not an assumption.

\item \textbf{Measured, not asserted.} Every performance claim in this report is a number from a named
  harness under stated conditions, and where the design loses --- the small-document write regime, the
  single-writer ceiling --- it is shown losing (Ch.~\ref{ch:eval}; \S\ref{sec:limitations}). Scope is
  drawn plainly rather than blurred; the honesty boxes throughout are part of the method, not a disclaimer
  appended to it.
\end{enumerate}

\noindent The rest of the report can be held against these five: each subsequent chapter is, in effect,
one principle worked out in full --- its mechanism, its proof or its test, and the price it pays. The
\textbf{contributions} below are what that working-out produced.

\section{Contributions documented here}

\begin{enumerate}[label=\arabic*.,leftmargin=2em]
\item A complete account of the \textbf{copy-on-write page trie} and the four
  \textbf{page-versioning strategies}, with the sliding-snapshot reconstruction
  algorithm and its correctness argument (Ch.~\ref{ch:pds}--\ref{ch:versioning}).
\item A self-describing, checksum-protected \textbf{on-disk format} with a
  \textbf{dual-beacon commit protocol}, and a documented evolution of that protocol
  (from 7 sync calls per commit down to 1 explicit fsync) that includes the discovery and
  elimination of a quadratic \code{access(2)} syscall pathology, all re-validated under
  power-loss and SIGKILL fault injection (Ch.~\ref{ch:ondisk}).
\item A production \textbf{Height-Optimised Trie} adapted to multi-key leaves, with a
  16-invariant catalogue and an engineering campaign reducing a reproducer's violations
  from 127 to a single, reads-correct, marginal case (Ch.~\ref{ch:hot}).
\item An \textbf{incremental path summary} with deferred per-path statistics
  (HyperLogLog cardinality, Roaring page bitmaps) that turns \emph{N} value writes into
  $O(\text{distinct-paths})$ copy-on-write operations (Ch.~\ref{ch:pathsummary}).
\item An \textbf{index-aware, cost-based query optimiser} over an embedded
  JSONiq/XQuery engine, including a path-rewrite that must respect node-kind to avoid a
  destructive-replace correctness bug (Ch.~\ref{ch:query}).
\item A \textbf{columnar/SIMD execution} path with frame-of-reference and dictionary
  encodings, validated by interpreter-differential testing. It is functional for a
  deliberately narrow, fail-closed set of query shapes (Ch.~\ref{ch:vector}).
\item A pluggable \textbf{enterprise I/O} family (io\_uring via the Java FFM API with no
  JNI; S3; Kafka), and the finding that an io\_uring linked-commit chain is actually
  slower than the serial path for single-resource commits (Ch.~\ref{ch:enterprise}).
\item An \textbf{evaluation} against PostgreSQL in which the relational engine wins in
  the small-document regime, with a precise account of why (Ch.~\ref{ch:eval}).
\end{enumerate}

\section{Scope and reading guide}

This report is an architecture and algorithms study, not a tutorial or an API manual. It explains how
SirixDB is built and why, grounds every mechanism in the source, and measures the result honestly; it
does not aim to teach JSONiq, to document every method, or to cover the internals of the embedded brackit
engine (cited where relevant). Distributed consensus, the GraalVM native-image build, and the full
operational runbook are out of scope beyond the deployment appendix.

The chapters build bottom-up, each depending on the one before it, so a linear reading is intended,
though a reader after a specific subsystem can enter at its chapter and follow the cross-references.
Three conventions recur. \emph{Worked examples}, set in boxed monospace, trace a concrete mechanism end
to end. \emph{Caveat} notes mark the honest boundary of a claim, where something ships partially or a
result is machine-specific. And \emph{properties} state the invariants the design relies on. The
appendices collect a class map, a glossary, the byte-level on-disk format, configuration and function
references, a JSONiq cookbook, and the deployment and build details.

\section{The shape of the system}

Figure~\ref{fig:stack} situates the layers treated by the remaining chapters. A query
arriving from the SolidJS GUI or a brackit client descends through the JSONiq/XQuery
engine and optimiser, into a transactional node layer (cursors, path summary, indexes,
hashing), onto the copy-on-write page trie and its sliding-snapshot versioning, and
finally onto a pluggable storage SPI whose dual-beacon commit protocol makes a revision
durable.

The book follows that stack from the bottom up. Chapter~\ref{ch:pds} builds the persistent
structure everything else rests on: a copy-on-write page trie whose pages are never
overwritten. Chapter~\ref{ch:versioning} makes that structure cheap to version, and
Chapter~\ref{ch:ondisk} makes it durable and crash-safe. Only once storage is settled do
the higher layers appear, each depending on the one below it: the navigation trie of
Chapter~\ref{ch:hot}, the path summary and secondary indexes of
Chapter~\ref{ch:pathsummary}, and the query planner and executor of
Chapters~\ref{ch:query} and~\ref{ch:vector}. Chapter~\ref{ch:enterprise} swaps in
alternative I/O backends, and Chapter~\ref{ch:eval} measures the whole against PostgreSQL.
Every layer earns its place by serving the claim above: a revision that costs
$O(\text{change})$ to write and a plain traversal to read back.

\begin{figure}[t]
\centering
\begin{tikzpicture}[font=\sffamily\small,
  layer/.style={draw=slate!55,fill=layerbg,rounded corners=2pt,minimum width=10.6cm,
                minimum height=1.0cm,align=center,inner sep=3pt},
  tag/.style={font=\sffamily\footnotesize\bfseries,color=ember},
  client/.style={draw=slate!55,fill=white,rounded corners=2pt,minimum width=3.1cm,
                 minimum height=0.85cm,align=center}]
  \node[client] (gui)  {Web GUI\\(SolidJS)};
  \node[client,right=0.5cm of gui] (cli) {brackit\\clients};
  \node[layer,below=0.7cm of gui.south east,anchor=north,xshift=-0.0cm] (rest)
        {REST API (Vert.x)\ \textbullet\ JSONiq/XQuery {+} cost-based optimiser};
  \node[layer,below=0.28cm of rest] (vec) {Vectorised columnar executor (SIMD)};
  \node[layer,below=0.28cm of vec]  (node)
        {Transactional node layer\\cursors \textbullet\ path summary \textbullet\ indexes \textbullet\ hashing};
  \node[layer,below=0.28cm of node] (page)
        {Page layer: CoW trie \textbullet\ sliding snapshot \textbullet\ HOT navigation};
  \node[layer,below=0.28cm of page] (sto)
        {Storage SPI: dual-beacon commit \textbullet\ recovery\\FileChannel \textbullet\ mmap \textbullet\ io\_uring \textbullet\ S3 \textbullet\ Kafka};
  % chapter tags
  \node[tag,right=0.2cm of rest.east]  {Ch.~7,\,10};
  \node[tag,right=0.2cm of vec.east]   {Ch.~8};
  \node[tag,right=0.2cm of node.east]  {Ch.~6};
  \node[tag,right=0.2cm of page.east]  {Ch.~2--3, 5};
  \node[tag,right=0.2cm of sto.east]   {Ch.~4, 9};
  \draw[-{Latex[length=2mm]},slate] (gui.south) -- ([xshift=-2.6cm]rest.north);
  \draw[-{Latex[length=2mm]},slate] (cli.south) -- ([xshift=0.6cm]rest.north);
\end{tikzpicture}
\caption{The layered architecture. Each layer is treated by the indicated chapters;
control flows top-to-bottom on a query, bottom-to-top on a commit.}
\label{fig:stack}
\end{figure}

\section{Modules and cross-cutting structure}

SirixDB is a Gradle multi-module build. The engine itself is \code{sirix-core}: the node and
transaction API, the page layer (the copy-on-write trie and its versioning), the storage SPI and its
backends, the caches, the secondary indexes, and the diff machinery. The query language lives in
\code{sirix-query}, which binds SirixDB into \textbf{brackit}, a separate XQuery~3.1/JSONiq engine
(groupId \code{io.sirix}) that supplies the parser, the data model, the function library, and the update
facility; the binding adds SirixDB's cost-based optimiser (Ch.~\ref{ch:query}) and the
\code{sdb:}/\allowbreak\code{jn:} functions. \code{sirix-rest-api} is the Vert.x HTTP/2 server
(Ch.~\ref{ch:rest}). Smaller modules wrap the core for other audiences: a Kotlin/coroutine API, a
command-line tool, examples and JMH benchmarks, and a Model Context Protocol server that exposes the
store to AI agents. The advanced I/O backends of Ch.~\ref{ch:enterprise} (io\_uring locally, S3 for
object storage, and a Kafka-based replication path) attach to this same storage SPI, with the
distributed pieces supplied by a separate enterprise layer on top of the open-source core.

\dcap{The open-source modules and the brackit engine}
\begin{diagramS}
sirix-core         storage engine: node / page / storage layers, caches, indexes, diff
sirix-query        brackit binding: cost-based optimiser, jn:/xml:/sdb: function libraries, CLI
sirix-rest-api     Vert.x HTTP/2 server: OAuth2/JWT, streaming shredder, metrics
sirix-kotlin-api   idiomatic Kotlin / coroutine wrapper
sirix-kotlin-cli   standalone command-line tool
sirix-mcp          Model Context Protocol server (AI-agent access)
sirix-examples / sirix-benchmarks
brackit            the XQuery 3.1 / JSONiq engine (a separate repo, used by sirix-query)
sirix-enterprise   distributed backends (S3, Kafka replication) layered on the storage SPI
\end{diagramS}

The runtime is three layers with narrow interfaces between them, and most of this report is a walk down
through them. At the top the \emph{node layer} (\code{NodeReadOnlyTrx}, \code{NodeTrx},
\code{NodeCursor}, \code{Axis}) presents a cursor over one immutable revision. Beneath it the \emph{page
layer} (\code{StorageEngineReader} and \code{StorageEngineWriter} over the trie of Ch.~\ref{ch:pds})
turns node keys into pages. At the bottom the \emph{storage SPI} (\code{IOStorage}, \code{Reader},
\code{Writer}) turns pages into bytes. Each seam is a real extension point. The storage backend is
chosen per resource (\code{FILE\_CHANNEL}, \code{MEMORY\_MAPPED}, \code{IO\_URING}, in-memory, or an
external provider such as S3 discovered through \code{ServiceLoader}); the bytes pass through a
configurable \code{ByteHandler} pipeline (Deflate or LZ4 compression, optional Tink encryption); and the
versioning strategy, hash mode, and secondary-index backend are likewise per-resource choices fixed in
the \code{ResourceConfiguration}.

Two concerns cut across all three layers. \textbf{Concurrency} is single-writer, many-reader per
resource: a process-wide registry hands out one write permit per resource path, while readers take no
lock at all because they observe immutable snapshots (Ch.~\ref{ch:rest}). Commit parallelises page
serialisation across a fork-join pool, and the REST layer runs on Vert.x event loops with blocking work
pushed to a worker pool. \textbf{Caching} is centralised in one JVM-wide buffer manager, sized like a
relational engine's shared buffer pool: a large record-page cache for leaf data, a fragment cache for
versioning reconstruction, and smaller caches for metadata pages, names, path summaries, and revision
roots, all evicted by background clock-sweeper threads that respect the oldest revision any live reader
still needs. Above them sit the query planner's plan and statistics caches (Ch.~\ref{ch:query}).

\section{Concepts and terminology}

A reader new to versioned stores may find a few recurring ideas worth stating up front; the glossary
(Appendix~\ref{app:glossary}) collects the rest.

\textbf{Immutability and structural sharing.} A \emph{persistent} data structure, in the
functional-programming sense, never destroys an old version when it produces a new one; the two share
everything that did not change. SirixDB applies this at page granularity: a write copies only the pages
it touches and shares the rest by reference, so every past revision remains a complete, readable tree.

\textbf{Copy-on-write and the root swap.} Because pages are immutable, a change cannot be written in
place; it produces new pages up to a new root, and the resource advances by atomically swapping a single
pointer to that root. This is the discipline copy-on-write filesystems use, lifted to a document store.

\textbf{Node identity and labeling.} Every node has a stable key that persists across revisions and,
optionally, a hierarchical \emph{DeweyID} that encodes its position. A good labeling scheme lets a node
be inserted without renumbering its neighbours, and lets document order and ancestor tests be decided by
comparing labels rather than walking the tree.

\textbf{Page versioning.} Storing a fresh full copy of a page on every change is simple but wasteful;
storing only deltas saves space but slows reads. A \emph{page-versioning strategy} chooses where on that
trade-off to sit, and SirixDB's sliding snapshot is the one that bounds the read cost and the write
spikes at the same time.

\textbf{Bitemporality.} A bitemporal store distinguishes \emph{transaction time} (when the system
recorded a fact) from \emph{valid time} (when the fact holds in the modeled world). SirixDB's revisions
provide the first inherently, and its valid-time bounds add the second.

\textbf{Row versus columnar layout.} Records stored field-by-field (row-oriented) are cheap to read and
write whole; the same data stored column-by-column is cheap to scan and compress one field at a time.
SirixDB is row-oriented at its core, with an optional columnar projection for analytical scans.

%================================================================
\chapter{The Persistent Data Structure}
\label{ch:pds}

\section{The copy-on-write page trie}

The cost claim of the previous chapter rests on one structure: a tree of pages that is
never overwritten in place, so that a write copies only the pages it actually touches.
This chapter describes that tree from the apex down.

The physical hierarchy is a multi-level trie that maps a 64-bit logical \emph{node key} to
a physical page. Its apex is the \textbf{\code{UberPage}} (\code{io/sirix/page/UberPage.java}),
a singleton resource descriptor that holds the reference to the current
\textbf{\code{RevisionRootPage}}. Each \code{RevisionRootPage}
(\code{io/sirix/page/RevisionRootPage.java}) anchors one immutable revision and carries
the revision number, commit timestamp, the maximum allocated node key, and references to
the document's primary node index and the auxiliary structures (path summary, name
page, and the CAS/path/name/DeweyID/projection-index roots).

\dcap{The trie spine: uber page \textrightarrow\ revision root \textrightarrow\ indirect levels \textrightarrow\ leaf}
\begin{diagram}
                  UberPage  (one per resource; pointer swapped atomically at commit)
                     │  getRevisionRootReference()
                     ▼
              RevisionRootPage  (revision N: revNo, timestamp, maxNodeKey,
                     │           refs → {document index, pathSummary, name,
                     │                   CAS, path, name, deweyID, projection})
                     ▼
              IndirectPage  ── 1024 references ──┐    (height grows with maxNodeKey)
                     ▼                           ▼
              IndirectPage                   IndirectPage
                     ▼
              KeyValueLeafPage  (≤ 1024 records, off-heap MemorySegment, slot bitmap)
\end{diagram}

From the document-index reference, a chain of \textbf{\code{IndirectPage}} nodes
(\pp{io/sirix/page/IndirectPage.java}) decomposes the node key into 10-bit slices, each level
fanning out by up to 1024. An \code{IndirectPage} is not a fixed 1024-slot array. It is an
\emph{adaptive} reference container that promotes through three layouts as it fills. It
begins as a \code{ReferencesPage4} (a linear list of up to four \code{(offset, ref)} pairs),
grows into a \code{BitmapReferencesPage} (a 1024-bit occupancy bitmap over a dense list, mapped
sparse-to-dense by a popcount prefix-sum), and becomes a \code{FullReferencesPage} only once about
1008 of the 1024 slots are used. A sparse trie level therefore costs a few references rather than
8\,KB. Navigation extracts successive 10-bit offsets via
the exponent ladder $[70,60,50,40,30,20,10,0]$:

\dcap{Algorithm 2.1 --- Trie navigation: nodeKey \textrightarrow\ leaf}
\begin{diagram}
  input: rootRef, pageKey         // pageKey = nodeKey >> log2(recordsPerPage)
  page ← deref(rootRef)
  level ← pageKey
  for exp in [70,60,50,40,30,20,10,0][8 - height ..]:   // only the live levels
      offset ← level >> exp
      level  ← level − (offset << exp)
      page   ← deref(page.getReference(offset))
  return page                      // a KeyValueLeafPage
\end{diagram}

\dcap{The same navigation worked: nodeKey 5\,000\,000 over a height-2 trie}
\begin{diagramS}
nodeKey 5 000 000
  pageKey = 5 000 000 >> 10 = 4882            slot = nodeKey & 1023 = 832
  4882 = 4·1024 + 786   →   indirect offsets [4, 786]

  RevisionRootPage ──▶ IndirectPage ─[4]─▶ IndirectPage ─[786]─▶ KeyValueLeafPage  (slot 832)
\end{diagramS}

The trie \textbf{height grows dynamically}. A document whose keys fit in 10 bits needs a
single indirect level, giving direct leaf access; larger key spaces add levels lazily. With
fan-out 1024 the depth is $\lceil \log_{1024}(\text{maxNodeKey}) \rceil$, at most seven
indirections for any realistic document, so node lookup is effectively constant-time
and small documents pay no structural overhead. Figure~\ref{fig:pagetrie} draws the addressing for the
worked key.

\begin{figure}[t]
\centering
\begin{tikzpicture}[font=\sffamily\small,>=Latex,
  pg/.style={draw=slate!60,fill=layerbg,rounded corners=2pt,align=center,inner sep=3pt,minimum height=0.85cm,minimum width=1.9cm},
  leaf/.style={draw=nodenew,fill=nodenew!12,rounded corners=2pt,align=center,inner sep=3pt,minimum height=0.85cm,minimum width=1.9cm},
  slice/.style={draw=ember!60,fill=ember!10,minimum width=1.9cm,minimum height=0.6cm,align=center,font=\scriptsize,inner sep=1pt},
  e/.style={->,slate}]
  % key slices, one above each level it indexes
  \node[font=\scriptsize\itshape,align=center] at (0.7,2.15) {nodeKey\\5\,000\,000};
  \node[slice] (s0) at (3.6,2.15) {level-0 = 4};
  \node[slice] (s1) at (6.5,2.15) {level-1 = 786};
  \node[slice,fill=nodenew!12,draw=nodenew!70] (sl) at (9.6,2.15) {slot = 832};
  \node[font=\scriptsize,anchor=west] at (10.7,2.15) {(10-bit fields)};
  % trie path
  \node[pg]   (rr) at (0.7,0) {Revision\\Root};
  \node[pg]   (i0) at (3.6,0) {Indirect\\level 0};
  \node[pg]   (i1) at (6.5,0) {Indirect\\level 1};
  \node[leaf] (lf) at (9.6,0) {Leaf\\(slot 832)};
  \draw[e] (rr) -- (i0);
  \draw[e] (i0) -- node[above,font=\scriptsize,color=ember!80!black]{[4]} (i1);
  \draw[e] (i1) -- node[above,font=\scriptsize,color=ember!80!black]{[786]} (lf);
  \draw[->,ember!60,dashed] (s0) -- (i0);
  \draw[->,ember!60,dashed] (s1) -- (i1);
  \draw[->,nodenew!70!black,dashed] (sl) -- (lf);
\end{tikzpicture}
\caption{Addressing a node by its 64-bit key. The key splits into 10-bit fields, one per trie level; each
field indexes a 1024-ary \code{IndirectPage} (here level-0 picks child 4, level-1 picks child 786), and the
low bits select the slot within the leaf (832). The height is
$\lceil\log_{1024}(\text{maxNodeKey})\rceil$ --- two indirect levels here --- so a lookup is a handful of
dereferences, effectively constant time, and a small document with a one-level trie reaches its leaves
directly.}
\label{fig:pagetrie}
\end{figure}

\section{The slotted leaf page}
\label{sec:slotted}

The terminal \textbf{\code{KeyValueLeafPage}} (\pp{io/sirix/page/KeyValueLeafPage.java}) is a single
contiguous off-heap \code{MemorySegment} in the PostgreSQL/LeanStore \emph{slotted-page} tradition,
laid out so that the in-memory and on-disk forms coincide (\code{PageLayout}):

\dcap{\code{KeyValueLeafPage} --- the slotted-page memory map}
\begin{diagramS}
KeyValueLeafPage — one contiguous off-heap MemorySegment   (in-memory == on-disk)
  0      32 B   header: recordPageKey, revision, populatedCount, heap watermarks, indexType
  32    128 B   slot bitmap (16 × long = 1024 bits; bit i ⇒ slot i populated)        on disk
  160   128 B   preservation bitmap (sliding-snapshot carry-forward)                 runtime only
  288  8192 B   slot directory (1024 × 8 B): [heapOffset:4][dataLength:3][kindId:1]   rebuilt on read
  8480   …      record heap (bump-allocated forward), one record =
                  [kindId:1][fieldOffsets:n×1][delta-varint fields][hash:8?][payload][deweyId][len:2]
\end{diagramS}

The 32-byte header records the page key, revision, populated count, and heap watermarks. The
128-byte \textbf{slot bitmap} (1024 bits) makes population $O(\text{populated})$ to iterate; a second
128-byte \textbf{preservation bitmap} is \emph{runtime-only}, never written, and drives the
sliding-snapshot carry-forward of Ch.~\ref{ch:versioning}. The 8\,KB \textbf{slot directory}
(1024 $\times$ 8~B) is \emph{rebuilt on read}, each entry packing
\code{[heapOffset:4][dataLength:3][kindId:1]}; the node kind is duplicated here so that dispatch is
branch-free without touching the heap. A \textbf{record} on the bump-allocated heap has the form
\code{[kindId:1][field-offset table][delta-varint fields][8\,B hash?][payload][deweyId][len:2]}. The
per-record one-byte offset table gives $O(1)$ random field access. Structural neighbour keys are
\textbf{delta-varint} encoded against the record's own key, and because siblings and children are
numerically close a pointer usually compresses to one or two bytes. The 8-byte hash is present only
when hashing is enabled. In-place mutation keeps the record width constant as long as the new delta
fits the same byte-width; otherwise it re-serialises and bumps the heap. A record above the spill
threshold moves out of line to an \code{OverflowPage} and is carried across fragment recombination
(Ch.~\ref{ch:versioning}) as a \code{PageReference}, so it is never lost. Two \code{MAX\_RECORD\_SIZE}
constants coexist, a 500-byte practical budget and a 150\,000-byte hard ceiling, which is a documented
refactor hazard. Because the directory and preservation bitmap are reconstructed at read time,
``self-describing on disk'' means precisely the header, slot bitmap, and heap; the rest is derived.

\section{The node model: kinds, fusion, and flyweights}

A \code{DataRecord} is one of some three dozen \code{NodeKind}s: the XML family (\code{ELEMENT},
\code{ATTRIBUTE}, \code{TEXT}, \ldots), the JSON family (\code{OBJECT}, \code{ARRAY}, the
\code{*\_VALUE} scalars, \code{JSON\_DOCUMENT}), the index and dictionary records, and the
\textbf{fused} \code{OBJECT\_NAMED\_\{OBJECT,ARRAY,STRING,NUMBER,BOOLEAN,NULL\}} kinds. Fusion is the
decisive space optimisation, and it is now mandatory. A JSON field \code{"k": v} is stored as a
\emph{single} record that plays both the object-key and the value role, instead of the historical
\code{OBJECT\_KEY} node plus a separate value child; that legacy pairing was removed outright. A
fused structural field still anchors a path-summary layer so that array and nested-object paths
resolve (Ch.~\ref{ch:pathsummary}). One consequence shows up on the wire: a fused key's children sit
as a bare array directly under its value, and that is the subtlety the REST serialiser and GUI
normaliser of Ch.~\ref{ch:rest}--\ref{ch:gui} absorb.

The node model is best read in two layers. The original \textbf{delegate} decomposition still names
a node's fields most clearly: a \code{NodeDelegate} (identity: node/parent keys, type,
previous/last-modified revision, a lazily decoded DeweyID), a \code{StructNodeDelegate} (first/last
child, left/right sibling, child and descendant counts), and a \code{ValueNodeDelegate} (a value byte
array, Deflate-compressed past a length threshold). That decomposition is now \emph{conceptual}: the
node families the page reader actually instantiates---\emph{both} the JSON kinds (\code{OBJECT},
\code{ARRAY}, the fused \code{OBJECT\_NAMED\_*}, the \code{*\_VALUE} scalars) and the XML kinds
(\code{ELEMENT}, \code{ATTRIBUTE}, \code{TEXT}, \code{NAMESPACE}, \code{COMMENT}, \code{PI})---all
implement \code{FlyweightNode} and carry \emph{no} delegate fields at all. They read and write their
fields directly against the page's \code{MemorySegment} through the per-record offset table, which is
what makes the slotted page's \emph{in-memory == on-disk} property pay off. This report treats
delegates as the conceptual data model and flyweights as its page-direct realisation.

Figure~\ref{fig:jsonencoding} makes the encoding concrete: a JSON object's fields become fused
\code{OBJECT\_NAMED\_*} records (one node holding the key \emph{and} its value), each with a stable node
key, and array elements are value nodes. The path summary (Ch.~\ref{ch:pathsummary}) then folds that
shape to one node per distinct path.

\begin{figure}[t]
\centering
\begin{tikzpicture}[font=\sffamily\footnotesize,>=Latex,
  box/.style={draw=slate!60,fill=layerbg,rounded corners=2pt,align=center,inner sep=4pt,minimum height=7mm},
  vbox/.style={draw=nodenew!70,fill=nodenew!12,rounded corners=2pt,align=center,inner sep=4pt,minimum height=7mm},
  key/.style={circle,draw=ember!80!black,fill=ember!12,text=ember!75!black,inner sep=0pt,minimum size=4.6mm,font=\scriptsize\bfseries},
  e/.style={draw=slate!55}]
  \node[font=\ttfamily\footnotesize,align=left,anchor=east] (json) at (-0.6,0)
    {\{\\\hspace{1.2ex}"name": "sirix",\\\hspace{1.2ex}"tags": ["json", "xml"]\\\}};
  \node[box] (obj) at (3.4,0) {OBJECT};
  \node[box] (name) at (1.5,-1.7) {OBJ\_NAMED\_STRING\\[-1pt]{\itshape "name" = "sirix"}};
  \node[box] (tags) at (5.3,-1.7) {OBJ\_NAMED\_ARRAY\\[-1pt]{\itshape "tags"}};
  \node[vbox] (v1) at (4.1,-3.5) {STRING\\[-1pt]{\itshape "json"}};
  \node[vbox] (v2) at (6.5,-3.5) {STRING\\[-1pt]{\itshape "xml"}};
  \foreach \n/\num in {obj/2,name/3,tags/4,v1/5,v2/6}
     \node[key] at (\n.north west) {\num};
  \draw[->,slate,thick] (json.east) -- (obj.west);
  \draw[e] (obj.south) -- (name.north);
  \draw[e] (obj.south) -- (tags.north);
  \draw[e] (tags.south) -- (v1.north);
  \draw[e] (tags.south) -- (v2.north);
\end{tikzpicture}
\caption{A JSON document and its SirixDB node encoding. Each object field is one \emph{fused}
\code{OBJECT\_NAMED\_*} record carrying the key \emph{and} its value (no separate key node); array
elements are value nodes. The circled numbers are the stable node keys, preserved across revisions.
Every distinct path collapses to a single \emph{path class} in the path summary
(Ch.~\ref{ch:pathsummary}); here the two array elements share the class \code{/tags/[]}.}
\label{fig:jsonencoding}
\end{figure}

Fusion is called ``the decisive space optimisation'' above; the slotted-page layout shows why, in bytes.
A single field \code{"name":"sirix"} costs, under the two encodings:

\dcap{Fusion in bytes: one field \texttt{"name":"sirix"}, legacy two-record vs fused one-record}
\begin{diagramS}
legacy (removed):  TWO records, TWO slot-directory entries, TWO node keys
  slot[i]    -> OBJECT_KEY "name"  : [kind][nodeKey][parentKey][firstChild -> value][lSib][rSib][nameKey]   ~32B
  slot[i+1]  -> STRING "sirix"     : [kind][nodeKey][parentKey][parent -> key][lSib][rSib][len][s i r i x]  ~32B
  to read the value: look up the key record, then CHASE firstChild to its value record  (2 records)

fused (mandatory): ONE record, ONE slot-directory entry, ONE node key
  slot[i]    -> OBJECT_NAMED_STRING: [kind][nodeKey][parentKey][lSib][rSib][nameKey "name"][len][s i r i x] ~40B
  to read the value: one slot lookup returns key AND value together                       (1 record)

saved per field: ~one whole record (~30B) + one slot entry + one node-key allocation + one pointer chase
at document scale (millions of fields) fusion roughly HALVES the structural record count
\end{diagramS}

\noindent The win is not only the $\sim$30 bytes of the eliminated intermediate node but a node key
(keys are a finite 64-bit budget tracked by \code{maxNodeKey}) and, on every read, a pointer chase: the
fused record returns key and value in one slot lookup. Because object fields are the most common shape in
real JSON, removing one record per field is close to halving the engine's structural record count, which
is why the pairing ``was removed outright'' rather than left as an option.

\section{Off-heap memory and the FFM frame pool}
\label{sec:offheap}

The slotted leaf page earns its ``in-memory equals on-disk'' property only because SirixDB keeps pages
\emph{off} the Java heap. A \code{KeyValueLeafPage} is one contiguous native \code{MemorySegment} reached
through the Foreign Function and Memory API (Project Panama), and its bytes are exactly the bytes written
to disk. Reading a node is therefore offset arithmetic into the segment with no object graph and no
deserialisation, and committing a page is handing its segment to the storage layer largely as-is. The
flyweight nodes of the previous section are thin cursors bound to a live segment; they hold no field
state of their own.

Native frames come from a small allocator hierarchy. \code{MemorySegmentAllocator} (with a
Linux-specialised \code{LinuxMemorySegmentAllocator}, chosen by \code{SegmentAllocators}) carves
page-sized frames out of a \emph{shared} FFM \code{Arena}, and a \code{FrameSlotAllocator} pools those
frames so that an evicted page's memory is reused rather than handed back to the operating system. The
whole buffer pool of \S\ref{sec:buffering} lives in this off-heap space, and that is the point: the
engine's largest memory consumer, the page cache, exerts \emph{no} pressure on the garbage collector,
and its size is a single budget dial rather than a function of heap tuning.

\dcap{Off-heap page frames from a shared FFM arena}
\begin{diagramS}
Java heap  │  off-heap native memory (FFM / Panama, a shared Arena)
           │
  thin node│   FrameSlotAllocator: a pool of reusable page-sized MemorySegment frames
  cursor ──┼──▶ KeyValueLeafPage = one MemorySegment   (its in-memory bytes are its on-disk image)
           │       holding the slot directory, slot bitmap, and record heap
           │   on eviction a frame returns to the pool; the clock-sweeper drives the reuse
\end{diagramS}

The price of going off-heap is manual lifecycle: a frame must be returned to the pool when its page is
evicted, which the clock-sweeper and the Transaction Intent Log handle between them; the shared
\code{Arena} is routed through \code{SharedArenas} (\code{Arena.ofShared} on HotSpot, \code{Arena.ofAuto}
in a native image), so a native image needs no shared-arena support. In exchange the design
gets zero-copy reads, GC-independent cache sizing, and the byte-for-byte identity between a page in
memory and the same page on disk that the rest of this report keeps relying on.

\section{Structural sharing}

When revision $N{+}1$ mutates a node, only the pages on its root-to-leaf path are copied;
sibling subtrees keep their old references. The \code{UberPage} is itself immutable:
a commit publishes a \emph{new} \code{RevisionRootPage} and atomically swaps the pointer
(Figure~\ref{fig:cow}).

\begin{figure}[t]
\centering
\begin{tikzpicture}[font=\sffamily\small,>=Latex,
  old/.style={draw=nodeold,fill=nodeold!12,circle,minimum size=7.5mm,inner sep=0pt},
  new/.style={draw=nodenew,fill=nodenew!14,circle,minimum size=7.5mm,inner sep=0pt,thick},
  rt/.style ={draw=nodeold,fill=nodeold!12,rounded corners=2pt,minimum size=7mm},
  rtn/.style={draw=nodenew,fill=nodenew!14,rounded corners=2pt,minimum size=7mm,thick}]
  % Revision N
  \node[rt] (rn) {$R_N$};
  \node[old] (a) [below left=0.9cm and 0.5cm of rn] {A};
  \node[old] (b) [below right=0.9cm and 0.5cm of rn] {B};
  \node[old] (c) [below=0.9cm of a] {C};
  \node[old] (d) [below=0.9cm of b] {D};
  \draw[->,nodeold] (rn)--(a); \draw[->,nodeold] (rn)--(b);
  \draw[->,nodeold] (a)--(c);  \draw[->,nodeold] (b)--(d);
  \node[align=center,below=0.2cm of c,font=\footnotesize\itshape] {Revision $N$};
  % Revision N+1
  \node[rtn] (rn1) [right=4.6cm of rn] {$R_{N{+}1}$};
  \node[new] (b1) [below right=0.9cm and 0.5cm of rn1] {B$'$};
  \node[new] (d1) [below=0.9cm of b1] {D$'$};
  \draw[->,nodenew,thick] (rn1)--(b1); \draw[->,nodenew,thick] (b1)--(d1);
  % shared edges back into A and C of revision N
  \draw[->,nodeold,dashed] (rn1) to[bend right=12] (a);
  \node[align=center,below=0.2cm of d1,font=\footnotesize\itshape,xshift=-0.6cm]
       {Revision $N{+}1$ (modify a leaf under B)};
  \node[font=\scriptsize\itshape,color=nodenew,above right=0.1cm and -0.3cm of d1]
       {B$'$,\,D$'$ copied};
  \node[font=\scriptsize\itshape,color=nodeold,below=0.05cm of a]
       {A,\,C shared};
\end{tikzpicture}
\caption{Copy-on-write structural sharing. Mutating a leaf under \textsf{B} copies only
the root-to-leaf path (\textsf{B$'$}, \textsf{D$'$} and a new root); the unchanged
subtree \textsf{A\,\textrightarrow\,C} is shared by reference with revision $N$.}
\label{fig:cow}
\end{figure}

\begin{figure}[t]
\centering
\begin{tikzpicture}[font=\sffamily\scriptsize,>=Latex,
  r/.style={draw=slate!60,fill=layerbg,circle,minimum size=7.5mm,inner sep=0pt},
  rn/.style={draw=nodenew,fill=nodenew!14,circle,minimum size=7.5mm,inner sep=0pt,thick}]
  \node[r] (r1) at (0,0)   {$r_1$};
  \node[r] (r2) at (1.7,0) {$r_2$};
  \node[r] (r3) at (3.4,0) {$r_3$};
  \node[r] (r4) at (5.1,0) {$r_4$};
  \node[rn] (r5) at (3.4,-1.5) {$r_5$};
  \draw[->,slate] (r1)--(r2); \draw[->,slate] (r2)--(r3); \draw[->,slate] (r3)--(r4);
  \draw[->,nodenew,thick] (r2) to[bend right=22] (r5);
  \node[font=\scriptsize\itshape,color=nodenew!55!black,anchor=west] at (3.8,-1.5) {revertTo($r_2$), edit, commit};
  \node[font=\scriptsize,anchor=west] at (5.4,0) {latest};
\end{tikzpicture}
\caption{Editing the past branches rather than overwrites. \code{revertTo($r_2$)} followed by a commit
produces a new revision $r_5$ whose content starts from $r_2$, while $r_3$ and $r_4$ are retained. A
resource's history is therefore a directed acyclic graph of immutable revisions, not a single line;
\code{truncateTo} is the only operation that removes revisions.}
\label{fig:revdag}
\end{figure}

For a write touching $k$ pages of an $n$-page document, pages written $= O(k + \text{height})$
and storage growth $= O(\text{actual delta})$, not $O(n)$. This is the cost-shape claim
of Chapter~\ref{ch:intro} made concrete. Every record additionally carries
\code{previousRevision} and \code{lastModifiedRevision}, enabling temporal predicates
(``in which revisions did node 12345 exist or change?'') to be answered structurally.

\begin{property}[Snapshot immutability]\label{prop:snapshot}
For any committed revision $N$, the set of pages reachable from $\textit{Root}_N$ is fixed
for all time; no later commit mutates a page in that set. Hence a read transaction bound
to revision $N$ observes a consistent snapshot with no locks.
\end{property}
This follows directly from copy-on-write: a mutation produces new pages and never
overwrites a page referenced by an already-published root. Figure~\ref{fig:mvcc} puts that property in
time: it is what lets many readers run concurrently, each on its own revision, with no lock between any of
them and the writer.

\begin{figure}[t]
\centering
\begin{tikzpicture}[font=\sffamily\small,>=Latex,
  rev/.style={draw=ember!80!black,very thick},
  rdr/.style={draw=slate!60,fill=layerbg,rounded corners=3pt,minimum height=0.5cm,align=center,inner sep=2pt}]
  \draw[->,slate!70,thick] (0,0) -- (12.6,0) node[below right,font=\scriptsize\itshape]{time};
  \node[font=\scriptsize\itshape,slate,anchor=west] at (0,3.0) {writer: one permit, commits in series};
  \foreach \x/\r in {2/1,5.5/2,9/3}{
     \draw[rev] (\x,0.12) -- (\x,2.5);
     \node[font=\scriptsize,color=ember!80!black,above] at (\x,2.5) {commit r\r};
  }
  \node[rdr,minimum width=4.4cm] (ra) at (4.5,-0.7)  {$R_a$ reads r1};
  \node[rdr,minimum width=4.0cm] (rb) at (8.1,-1.45) {$R_b$ reads r2};
  \node[rdr,minimum width=2.6cm] (rc) at (11.0,-2.2) {$R_c$ reads r3};
  \draw[->,slate,dashed] (2,0)   |- (ra.west);
  \draw[->,slate,dashed] (5.5,0) |- (rb.west);
  \draw[->,slate,dashed] (9,0)   |- (rc.west);
\end{tikzpicture}
\caption{Snapshot isolation in time. The single writer commits revisions in series (\code{r1}, \code{r2},
\code{r3}); each reader binds to one committed revision and observes its immutable page set for its whole
lifetime, so \code{R}$_a$ keeps seeing \code{r1} even after \code{r2} and \code{r3} commit. No reader takes
a lock and none blocks the writer; the only serialisation in the picture is the writer's single permit
between commits (Property~\ref{prop:snapshot}).}
\label{fig:mvcc}
\end{figure}

A concrete trace makes the ``$O(k+\text{height})$ written, the rest shared'' claim tangible. Take the
height-2 trie of \S\ref{sec:slotted}'s navigation example and update one node.

\dcap{One node update, revision N $\to$ N+1: what copies and what is shared}
\begin{diagramS}
update node key 5 000 000   (from §2.1: pageKey 4882 → indirect offsets [4, 786], slot 832)

COPIED — only the root-to-leaf path (newest-wins):
   KeyValueLeafPage @[4][786]   slot 832 takes the new value                       1 leaf
   IndirectPage (level 1) @[4]  its reference to that leaf re-pointed              1 indirect
   IndirectPage (level 0)       its reference to level 1 re-pointed                1 indirect
   RevisionRootPage N+1         fresh, points at the copied level-0 indirect       1 revroot
   UberPage                     pointer swapped atomically N → N+1                 1 uber
SHARED with revision N (identical PageReferences, not rewritten):
   every other leaf, every other indirect subtree, the path summary, and every index not on this path

written ≈ 3 path pages + revroot + uber = O(height);  read @ N still sees the OLD value, because
RevisionRootPage N still points at the original leaf — never mutated (Property~\ref{prop:snapshot}).
\end{diagramS}

Put numbers on it. A document of $5\times10^{6}$ nodes fills about $5\,000\,000/1024 \approx 4\,900$ leaf
pages under five level-1 indirects and one level-0 indirect, on the order of tens of megabytes on disk.
The single-node commit above rewrote \emph{five} pages, on the order of tens of kilobytes --- roughly
\textbf{0.1\%} of the document. The remaining $\sim$99.9\% of pages are not copied, not even read: revision
$N{+}1$'s root simply re-uses revision $N$'s \code{PageReference}s for every subtree off the changed path.
A second commit on an \emph{unrelated} branch shows the sharing compounding rather than decaying:

\dcap{A second update on a different branch: the revisions form a sharing DAG}
\begin{diagramS}
commit B @ N+2:  update node key 1 000 000   (pageKey 976 -> indirect offsets [0, 976], slot 832)

COPIED for B:  its own leaf @[0][976], the level-1 indirect @[0], the level-0 indirect, revroot, uber
SHARED by N+2 with N+1 (untouched by B, so the SAME page objects):
   the leaf and level-1 indirect @[4] that commit A had copied  ── A's work is inherited, not redone
   every other leaf and subtree                                  ── ~99.9% of the tree

result: revisions N, N+1, N+2 are three roots over ONE page DAG.
   total growth for the two commits  ≈ 10 pages (~tens of KiB), NOT 2 x (tens of MiB)
   storage cost = O(sum of deltas),  never O(commits x document)
\end{diagramS}

\noindent The level-0 indirect and the revision-root are copied by \emph{every} commit (they sit on every
root-to-leaf path), which is the fixed $O(\text{height})$ floor; everything below the divergence point is
shared. This is the structural reason the evaluation finds that retaining a hundred revisions costs about
one percent over the live data (\S\ref{sec:crossover}): each revision adds its delta to a shared DAG, and
nothing recopies what did not change.

\section{DeweyIDs}
\label{sec:deweyids}

When a resource is built with \code{ResourceConfiguration.Builder\#useDeweyIDs(true)}, every node also
carries a \textbf{\code{SirixDeweyID}} (\pp{io/sirix/node/SirixDeweyID.java}): an outline-style label
such as \code{1.17.33}, one \emph{division} per level from the root. This is the insert-friendly node
labeling SirixDB inherits from the native-XML world (the ORDPATH~\cite{ordpath} and
SPLID~\cite{haerderlabels} schemes), chosen so that document order, ancestor and descendant tests, and
sibling tests all reduce to arithmetic on the label rather than a walk of the tree. In memory a DeweyID
is an \code{int[]} of division values plus a level; the byte form is computed lazily and cached.

\begin{figure}[t]
\centering
\begin{tikzpicture}[font=\sffamily\scriptsize,>=Latex,
  n/.style={draw=slate!60,fill=layerbg,rounded corners=2pt,align=center,inner sep=3pt,minimum height=6mm},
  e/.style={draw=slate!55}]
  \node[n] (root) at (4,0)      {document\\[-1pt]{\ttfamily\color{ember!80!black}1}};
  \node[n] (b0)   at (1.8,-1.7) {book{[}0{]}\\[-1pt]{\ttfamily\color{ember!80!black}1.17}};
  \node[n] (b1)   at (6.2,-1.7) {book{[}1{]}\\[-1pt]{\ttfamily\color{ember!80!black}1.33}};
  \node[n] (t0)   at (0.4,-3.5) {title "A"\\[-1pt]{\ttfamily\color{ember!80!black}1.17.17}};
  \node[n] (p0)   at (2.7,-3.5) {price 12\\[-1pt]{\ttfamily\color{ember!80!black}1.17.33}};
  \node[n] (t1)   at (5.0,-3.5) {title "B"\\[-1pt]{\ttfamily\color{ember!80!black}1.33.17}};
  \node[n] (p1)   at (7.3,-3.5) {price 8\\[-1pt]{\ttfamily\color{ember!80!black}1.33.33}};
  \draw[e](root)--(b0); \draw[e](root)--(b1);
  \draw[e](b0)--(t0); \draw[e](b0)--(p0); \draw[e](b1)--(t1); \draw[e](b1)--(p1);
\end{tikzpicture}
\caption{A document tree labelled with DeweyIDs. Each child extends its parent's label by one division,
and siblings are spaced by the default distance of 16, so \code{1.17} and \code{1.33} leave room for a
later \code{1.25} between them with no relabelling. A depth-first, left-to-right reading of the tree is
exactly ascending DeweyID order, which is why an unsigned byte comparison of two encoded labels recovers
document order.}
\label{fig:deweytree}
\end{figure}

\textbf{The encoding.} Each division is written with a prefix-free, tier-based variable-length code, so
that a byte-by-byte unsigned comparison of two encoded labels reproduces their numeric, level-by-level
order. Five tiers cover the 31-bit division range; the prefix (\code{0}, \code{10}, \ldots, \code{1111})
selects the tier and a fixed-width suffix carries the value within it.

\dcap{DeweyID division encoding: prefix-free tiers (with division \texttt{17} worked through)}
\begin{diagramS}
tier  prefix  suffix bits   value range (half-open)          encoded size
 0     0          7         [0, 127)                          1 byte
 1     10        14         [127, 16 511)                     2 bytes
 2     110       21         [16 511, 2 113 663)               3 bytes
 3     1110      28         [2 113 663, 270 549 119)          4 bytes
 4     1111      31         [270 549 119, 2 147 483 647]      5 bytes (35 bits, unaligned)

worked example:  division 17  →  tier 0  →  prefix 0  ++  7-bit (17+1)  =  0 0010010  =  0x12
(tier-0 suffix stores value+1 so no all-zero byte occurs; the root division 1 is implicit)
encode(d₁ … dₙ) = the byte strings of d₁, d₂, … concatenated
\end{diagramS}

The tier bounds are \emph{exclusive} upper limits, and a 2026-01 fix made them so: a value uses tier
\emph{i} when it is strictly below that tier's bound, so 127 belongs to tier 1 rather than overflowing
tier 0's seven bits. The verification test \code{DeweyIDEncodingVerificationTest} round-trips every tier
boundary, and a separate \code{DeweyIDEncoder} performs the same encoding with table lookups and
multi-byte writes, asserted byte-identical to the reference algorithm.

\begin{theorem}[Order preservation]\label{thm:dewey}
For all DeweyIDs $a,b$: $a.\textsf{compareTo}(b) < 0 \iff
\textsf{compareUnsigned}(\textsf{encode}(a),\textsf{encode}(b)) < 0$.
\end{theorem}
The proof rests on two facts: the tier prefixes are prefix-free and increase with the tier
(\code{0} < \code{10} < \code{110} < \code{1110} < \code{1111}), and within a tier the suffix is a
fixed-width big-endian integer that is monotone in the value. Together they make the encoded bytes a
faithful, order-preserving image of the division array. (The repository also carries a longer
\pp{docs/DEWEYID\_HOT\_INDEX\_FORMAL\_PROOF.md} as a design-and-proof artifact; the guarantee the
running system actually relies on is the one these tests check.)

\textbf{Insert-friendly labeling.} The scheme is ORDPATH's parity trick. Odd division values are real
node steps and even values (and zero) are \emph{carets}, intermediate steps that do not count as a
level. New siblings are spaced a fixed \emph{distance} apart (the default is 16, an even number), which
leaves room to insert later without touching any existing label.

\dcap{Careting-in: insertion without relabelling (default distance 16)}
\begin{diagram}
three children of node 1, spaced by 16:     1.17     1.33     1.49      (odd = real node)

append one more on the right:   last 1.49 + 16             →  1.65
wedge between 1.17 and 1.33:    odd midpoint               →  1.25
wedge between adjacent 1.17 and 1.19 (no odd in between):
        descend with an even caret that adds no level      →  1.18.17
the labels 1.17, 1.33, 1.49 are never rewritten
\end{diagram}

Appending a sibling adds the distance to the trailing division. Wedging between two siblings takes the
odd midpoint when one exists, and otherwise descends one level, appending an even caret followed by a
fresh child division. Because a caret is even it does not change the node's level, so \code{1.18.17}
sits between \code{1.17} and \code{1.19} in document order while staying a child of node \code{1}. No
existing label is ever renumbered, which is what makes sub-document inserts cheap and keeps node
identity stable across edits.

\begin{property}[Insertion without relabelling]\label{prop:norelabel}
Inserting a node between two existing siblings assigns it a fresh DeweyID and changes \emph{no} existing
label; the document order, ancestor, and sibling relationships of every other node are preserved.
\end{property}

\textbf{Structural axes.} Because a child's label extends its parent's, the structural relationships are
prefix arithmetic on the division array, with no tree traversal.

\dcap{Ancestor, sibling, and order tests as prefix arithmetic}
\begin{diagramS}
a = 1.17.5     b = 1.17.5.33     c = 1.17.49     d = 1.33

a ancestor of b : a's divisions are a prefix of b's                              → yes  (O(depth))
a sibling of c   : same level, share 1.17, differ only in the last odd division   → yes
d follows a      : first differing division 33 > 17 decides document order        → yes
parent(b) = a    : drop the trailing division (and any even carets)
\end{diagramS}

\code{isAncestorOf}, \code{isDescendantOf}, \code{isSiblingOf}, \code{getParent} and the
preceding/\allowbreak following tests are all elementwise prefix checks over the \code{int[]}, so
ancestor and descendant queries are $O(\text{depth})$ and document-order comparison is a single scan.
These tests power the bitwise structural axes without ever dereferencing a sibling or parent pointer.

\textbf{On disk and as an index.} A DeweyID lives at the tail of its record, and it is stored
\emph{delta-encoded against the previous sibling on the page}, not as a fixed-width field
(\code{NodeSerializerImpl}). The first record on a page writes a one-byte length followed by the full
label bytes; each later record writes only how many leading bytes it shares with its predecessor plus
the divergent suffix, so a run of siblings (whose labels differ only in the last division or two) costs
a byte or two each. Decoding is lazy: a node keeps the raw bytes and materialises the \code{int[]} form
only when something asks for it (\code{NodeDelegate.getDeweyID}), so a bulk scan that never needs
structure never pays to decode. The labels also back a dedicated \code{DEWEYID\_TO\_RECORDID} index
(\code{IndexType.DEWEYID\_TO\_RECORDID}, a per-revision page tree mapping node keys to their DeweyIDs);
the order-preserving encoding above is what lets DeweyID-keyed range scans return nodes already in
document order.

\section{Node hashing}

Controlled by \code{HashType} (\code{io/sirix/access/trx/node/HashType.java}), SirixDB
maintains a \textbf{rolling Merkle hash} over the node tree. Each node stores a 64-bit
hash computed (in \code{node/json/*.java}) over its identity, kind, structural pointers,
name key, and aggregated child/descendant counts:

\dcap{Per-node hash inputs}
\begin{diagram}
hash(node) = XXH3( nodeKey || parentKey || kind || childCount || descendantCount
                   || leftSibling || rightSibling || firstChild [|| lastChild] || nameKey )
\end{diagram}

On mutation the node's hash is recomputed and the change propagates up the ancestor chain
to the root, forming a Merkle tree whose root digests the whole document. The payoff is
\textbf{hash-guided diffing}: two revisions are compared by descending only into subtrees
whose hashes differ, never deserialising unchanged nodes. This is the basis of the
sub-millisecond semantic diff measured in Ch.~\ref{ch:eval}. The three modes differ in how
much work the recompute does. \code{ROLLING} is a Horner-rule polynomial updated up the
ancestor chain alone: each ancestor folds in the delta as
$\text{hash} - \text{old}\cdot\textsc{prime} + \text{new}\cdot\textsc{prime}$
(\textsc{prime} $= 77081$) and bumps its descendant count in the same pass, so a node's hash digests
its structural pointers and its child/descendant \emph{counts} rather than child content directly.
\code{POSTORDER} is heavier, re-reading every child on the path to the root (and, for XML, folding
in attribute and namespace hashes). \code{NONE} disables hashing and additionally omits the hash and
descendant-count fields from the record, trading diff and integrity capability for about 13\,\% more
commit throughput (measured, \S\ref{sec:w1}).

A trace shows the \code{ROLLING} update touching only one root-to-leaf path.

\dcap{Rolling-hash propagation for one value change (\textsc{prime} $=77081$)}
\begin{diagramS}
change  book[0].price : 12 → 19.99

  recompute the leaf:  h_price' = XXH3(record with the new value)
  walk UP the ancestor spine, folding the child delta with the Horner update:
    book[0]:  h ← h − h_price·P + h_price'·P            (P = 77081; subtract old, add new contribution)
    books[]:  h ← h − h_book·P  + h_book'·P
    root:     h ← h − h_books·P + h_books'·P
  each ancestor bumps its descendant count in the SAME pass → one O(depth) walk, no sibling re-read

  untouched siblings (title, book[1], …) keep their hashes unchanged
  ⇒ a later jn:diff prunes every one of them with SAMEHASH, descending only this path (\S\ref{sec:diff})
\end{diagramS}

The arithmetic is one machine word wide. Working the \code{book[0]} step with illustrative leaf hashes
makes the cost concrete (real XXH3 values are full 64-bit; small ones here keep the multiply checkable):

\dcap{One Horner step in hex: $h \leftarrow h - h_{old}\cdot P + h_{new}\cdot P$ (mod $2^{64}$, $P=77081$)}
\begin{diagramS}
  parent book[0] hash, before:   H      = 0xA1B2C3D4_E5F60718
  price leaf hash:    old = 1000,  new = 1500      (illustrative; real values are 64-bit)
  delta = (new - old) * P = 500 * 77081 = 38 540 500 = 0x024C14D4
  H' = H + delta (mod 2^64)             = 0xA1B2C3D4_E8421BEC
       high half 0xA1B2C3D4 unchanged here only because this delta is small;
       a real 64-bit delta carries through and mixes the entire word
\end{diagramS}

\noindent The recompute is therefore one 64-bit multiply, one add and one subtract per ancestor --- no
big-integer, no child re-read. So the hash a node stores digests its structural pointers and its
child/descendant \emph{counts}, and keeping the whole-document Merkle root current after an edit is one
walk up the spine, the same $O(\text{depth})$ the copy-on-write page copy already pays. That shared cost
is why hashing is current by the time a diff needs it, without a separate pass.

\section{Names, large values, and string compression}
\label{sec:names}

Three mechanisms keep records small without losing fidelity.

\textbf{Name interning.} Element and attribute names, and JSON object keys, repeat enormously, so SirixDB
stores each distinct name once. A name is interned to an integer key in a per-resource \code{NamePage}
dictionary, and a node holds only that key. The dictionary is itself copy-on-write and shared across
revisions, and it is cached (the names cache of \S\ref{sec:buffering}); the path summary and the name and
path indexes key on these integer name keys rather than on strings, which is part of why those structures
stay compact.

\textbf{Overflow records.} Most records sit inline on the bump-allocated heap of their
\code{KeyValueLeafPage}, but a value above a spill threshold would bloat the page and hurt the slot
density that makes scans fast. Such a record moves out of line into an \code{OverflowPage} and is
referenced from its slot by a \code{PageReference}. The sliding-snapshot combination carries that
reference across newest-wins without decoding the large value (Ch.~\ref{ch:versioning}), so a big blob is
never re-serialised by a version merge.

\textbf{String compression.} Within a page, repeated string content is compressed with a per-page
\textbf{FSST} symbol table, a fast static dictionary that replaces common substrings with short codes.
The table lives with the page; during a multi-fragment sliding-snapshot combine each compressed slot is
decoded through \emph{its own} fragment's table and rewritten, so the merged page carries no stale
dictionary and the next commit rebuilds one coherent table. This is a correctness requirement, not just a
space optimisation, because a slot's codes are meaningless against the wrong table.

A concrete table makes both the saving and the hazard visible. Suppose a page's strings are dominated by a
few substrings; the FSST builder assigns each a one-byte code:

\dcap{A per-page FSST table, and why a fragment merge must re-encode}
\begin{diagramS}
page strings:  "item-000123"  "item-000456"  "active"  "inactive"
symbol table (THIS page):   0x01 -> "item-00"    0x02 -> "active"    0x03 -> "in"

encode "active":    0x02                       6 bytes  -> 1 byte
encode "inactive":  0x03 0x02                  8 bytes  -> 2 bytes
encode "item-000123": 0x01 '0' '0' '0' '1' '2' '3'   (the "item-00" prefix folds to one code)

now a sliding-snapshot combine pulls in a fragment B whose table assigns  0x01 -> "category":
   a RAW byte-copy of A's 0x01 slot would decode as "category..." under B's table  -> CORRUPTION
   so each slot is decoded through ITS OWN fragment's table, then re-encoded under the merged page's table
\end{diagramS}

\noindent The codes are page-local by construction --- \code{0x01} is whatever \emph{this} page's table
says --- so the only safe merge decodes every incoming slot through the table it was written against
before the combined page builds a fresh dictionary. That is the precise sense in which carrying the table
``with the page'' is a correctness invariant and not merely a place to keep it.

\section{XML and JSON: one engine, two node families}
\label{sec:xmljson}

Everything in this report applies to XML resources as much as to JSON ones, because the two share a
single engine and differ only in their node kinds. A transaction is one \code{AbstractNodeTrxImpl}
skeleton with the node family bound as a type parameter, so the page layer, the commit pipeline, the
sliding snapshot, the transaction registries, the write semaphore, and the caches are identical for
both; only the node-kind operations differ. \textbf{XML} adds the markup machinery: \code{ELEMENT},
\code{ATTRIBUTE}, \code{TEXT}, comment, and processing-instruction nodes, with attributes and namespaces
attached to their element and reached by \code{moveToAttribute}/\allowbreak\code{moveToNamespace}.
\textbf{JSON} instead has objects, arrays, the scalar value kinds, and the fused \code{OBJECT\_NAMED\_*}
records that pack a key and its value into one node. The query bridge mirrors the split: JSON items
surface through \code{JsonDBObject}/\allowbreak\code{JsonDBArray} and XML through the element/attribute item
model, but both drive the same write transaction underneath. The report leads with JSON because that is
SirixDB's present focus, yet none of the storage, versioning, durability, or indexing machinery is in
any way JSON-specific.

%================================================================
\chapter{Versioning: The Sliding Snapshot}
\label{ch:versioning}

\section{The strategy space}
\label{sec:strategyspace}

Copy-on-write as described so far would rewrite a whole leaf page for every single-record
change. Avoiding that is the job of versioning. The page layer reconstructs a logical page
from a chain of \textbf{page fragments} linked by \code{PageFragmentKey}s, under one of four
strategies (\code{io/sirix/settings/VersioningType.java}).

\begin{table}[h]
\small\setlength{\tabcolsep}{5pt}
\renewcommand{\arraystretch}{1.25}
\begin{tabularx}{\linewidth}{@{}l X c X@{}}
\toprule
\textbf{Strategy} & \textbf{On-disk per commit} & \textbf{Read} & \textbf{Write spikes / storage}\\
\midrule
\textbf{FULL} & whole page & $1$ & none; high, $O(\text{page})$ per change\\
\textbf{INCREMENTAL} & delta; full every $R$ revisions & $\leq R$ & periodic full rewrite; low between spikes\\
\textbf{DIFFERENTIAL} & delta-from-last-full; full every $R$ & $\leq 2$ & periodic full rewrite; delta grows toward each dump\\
\textbf{SLIDING\_SNAPSHOT} \textit{(default)} & sparse delta {+} preserved tail & $\leq R$ & \textbf{none}; bounded, even\\
\bottomrule
\end{tabularx}
\end{table}

FULL is simplest but write-heavy. INCREMENTAL keeps each delta small (the change since the
previous revision) but re-emits a whole page every $R$ revisions, a $6\times$ write spike for a
one-record change in a 1024-record page, and a read may replay up to $R$ fragments. DIFFERENTIAL
trades the other way: a read touches at most two fragments (the latest differential plus the last
full dump), but each differential re-stores every record changed since that dump, so the per-commit
delta grows toward the next full dump, which it still rewrites every $R$ revisions. SLIDING\_SNAPSHOT
combines INCREMENTAL's small deltas with DIFFERENTIAL's bounded reads while eliminating the periodic
full-dump spikes both of them pay. It is the strategy that defines the system.

The four strategies are easiest to compare on one concrete write sequence to a single page.

\dcap{The same six commits under all four strategies (4-slot page, window $R=3$)}
\begin{diagramS}
write sequence:  r1 write a,b,c,d   r2 change a   r3 change b   r4 change c   r5 change a   r6 change d

strategy        what each commit writes  (r1   r2   r3    r4*    r5   r6)        read @ r6 touches
FULL            {abcd}{abcd}{abcd}{abcd}{abcd}{abcd}                            1   (latest full page)
INCREMENTAL     {abcd}  {a}  {b}  {abcd}  {a}  {d}     (* full dump every R)    ≤ 3 (replay since r4 dump)
DIFFERENTIAL    {abcd}  {a}  {ab} {abcd}  {a}  {ad}    (delta since last dump)  2   (last dump + latest)
SLIDING (def.)  {abcd}  {a}  {b}  {cd}    {a}  {bd}    (preserve the tail)      ≤ 3, and NO dump spike

SLIDING preserves a slot only when it leaves the window uncovered:
   r4 writes {c} and preserves d (r1 leaves window {r2,r3,r4}; d not re-covered) → {c,d}
   r6 writes {d} and preserves b (r3 leaves window {r4,r5,r6}; b not re-covered) → {b,d}
DIFFERENTIAL's r3 re-stores both a and b (everything changed since the r1 dump); SLIDING's r3 writes only {b}.
\end{diagramS}

Reading down the ``writes'' column shows the trade in miniature: FULL never grows its read cost but pays a
whole page every commit; INCREMENTAL and DIFFERENTIAL keep commits small but both rewrite the entire page
at \code{r4} (the dump), the write spike; and SLIDING\_SNAPSHOT writes only the changed slot plus the
occasional one-slot preservation, so it has small commits \emph{and} bounded reads \emph{and} no spike,
at the cost of the Phase-2 preservation bookkeeping the next section details.

\section{The sliding-snapshot algorithm}
\label{sec:slidealg}

Let $R = \textit{revisionsToRestore}$ (the window; default~3). On
modification, the page is reconstructed and the \emph{next} fragment is prepared in two
phases:

\dcap{Algorithm 3.1 --- \texttt{SLIDING\_SNAPSHOT.combineRecordPagesForModification}}
\begin{diagram}
  input: fragment chain F = [f₀ (newest) … f_{m}], window R
  complete ← empty page;  inWindow ← 1024-bit zero bitmap
  // Phase 1 — in-window merge (newest wins)
  for f in f₀ … f_{min(R-1, m)}:               // newest R-1 fragments
      for slot s populated in f:
          mark s in inWindow
          if s not in complete:  complete[s] ← f[s]     // newer shadows older
  // Phase 2 — out-of-window preservation
  if |chain| == R:                              // an oldest fragment is about to fall out
      f_old ← the R-th fragment
      for slot s populated in f_old:
          if s not in complete:  complete[s] ← f_old[s]          // fill a gap
          if s not in inWindow:  modifying.markSlotForPreservation(s)  // carry forward
  return (complete /*read view*/, modifying /*sparse page written at commit*/)
\end{diagram}

Phase~1 merges the newest $R{-}1$ fragments, newer slots shadowing older. Phase~2
inspects the single fragment about to leave the window: any slot it holds that \emph{no
in-window fragment covers} is \textbf{marked for preservation}, copied forward lazily
so the next on-disk fragment remains self-sufficient within the window.

\begin{property}[Bounded reconstruction]\label{prop:bounded}
A logical page is fully reconstructable from at most $R$ fragments at any revision.
\end{property}
Each record leaving the window is preserved into the new fragment exactly when no newer
fragment already covers it, so the union of any $R$ consecutive fragments covers every
live slot.

\begin{property}[Even amortisation]\label{prop:amort}
Each record is preserved at most once per window pass.
\end{property}
Preservation is triggered only for slots uncovered in-window and is not repeated while the
record stays covered, spreading the carry-forward cost across revisions ($\approx 1/R$
extra writes per revision) rather than concentrating it in periodic full rewrites.

A short trace makes the mechanism concrete.

\dcap{A worked sliding-snapshot trace (window $R = 3$) for one logical page}
\begin{diagramS}
one logical page with slots a, b, c, d. Each commit writes a sparse fragment of only the changed slots,
plus any slot whose last writer is leaving the R-window and is not re-covered inside it.

  r1  write a,b,c,d   → fragment {a,b,c,d}
  r2  change a        → {a}
  r3  change b        → {b}
  r4  change c        → {c} ; r1 leaves window {r2,r3,r4}, its d uncovered → preserve d → {c,d}
  r5  change a        → {a} ; r2 leaves window {r3,r4,r5}, its a re-covered by r5 → nothing extra
  r6  change d        → {d} ; r3 leaves window {r4,r5,r6}, its b uncovered → preserve b → {b,d}

  read @ r6 = combine last 3 fragments {r6,r5,r4}: a(r5) b(r6) c(r4) d(r6)   (all slots, ≤ R reads)
\end{diagramS}

Every read touches at most $R$ fragments and every slot is always present in the window, while the
extra preservation writes (d at r4, b at r6) are spread out instead of bunched into a periodic full
rewrite. That is the whole point of the strategy in one page.

The implementation handles two special cases concretely. \textbf{Overflow records} are not decoded
inline during combination; their \code{PageReference} is carried across, newest-wins, so a large
value is never re-serialised by a merge. \textbf{FSST symbol tables} (the per-page string-compression
dictionary) are inherited by reference on a single-fragment combine, but on a multi-fragment combine
each compressed slot is decoded through \emph{its own fragment's} table and rewritten uncompressed, so
the merged page carries no table and the next commit rebuilds one coherent dictionary. That is a
correctness invariant, not merely an optimisation.

Algorithm~3.1 is the \emph{write} path (\code{combineRecordPagesForModification}); a symmetric
\code{combineRecordPages} serves reads, and a third method \code{getRevisionRoots} decides which
revision roots to dereference. For \code{SLIDING\_SNAPSHOT} the read path is structurally identical to
\code{INCREMENTAL}, a newest-wins fill over the last $R$ fragments, and the entire distinction
lives in the write-side preservation of Phase~2. \code{DIFFERENTIAL} instead dereferences only
$\{\text{previous}, \text{last full dump}\}$, while \code{INCREMENTAL}/\allowbreak\code{SLIDING\_SNAPSHOT}
dereference the last $R$ consecutive revisions (Figure~\ref{fig:slide}; the default $R=3$).

The read paths differ in shape, and tracing the \emph{same} write history under two strategies makes the
contrast concrete. Take the worked trace above (r1 writes a,b,c,d; r2--r6 each change one slot) and read
at \code{r6}:

\dcap{Read @ r6 under DIFFERENTIAL vs SLIDING\_SNAPSHOT (same write history)}
\begin{diagramS}
DIFFERENTIAL  (dereferences {last full dump, previous delta}):
  fragments read:  r1 dump {a,b,c,d}  +  r6 delta {a,b,c,d}       -> 2 reads, ALWAYS
  but the r6 delta is every slot changed since the r1 dump, so it has grown to all four slots and keeps
  bloating until the next full dump -> few reads, but write/storage amplification climbs over time

SLIDING_SNAPSHOT  (dereferences the last R=3 revisions):
  fragments read:  r6 {b,d}  +  r5 {a}  +  r4 {c,d}               -> 3 reads (= R), FOREVER
  newest-wins merge:  a<-r5   b<-r6   c<-r4   d<-r6               each fragment stays small and sparse

trade: DIFFERENTIAL pins reads at 2 but lets the delta grow; SLIDING bounds BOTH (<= R reads, small
fragments) -- exactly why §\ref{sec:versioningbench} measures SLIDING smallest while DIFFERENTIAL's
storage sits near FULL (its delta approaches a whole re-dump)
\end{diagramS}

\section{Fragment linkage, reconstruction, and caching}
\label{sec:caching}

\begin{figure}[t]
\centering
\begin{tikzpicture}[font=\sffamily\small,>=Latex,
  frag/.style={draw=slate!55,fill=layerbg,minimum width=1.45cm,minimum height=0.95cm,align=center},
  newfrag/.style={draw=nodenew,fill=nodenew!12,thick,minimum width=1.45cm,minimum height=0.95cm,align=center}]
  \node[frag] (f1) {$r_{N\text{-}4}$};
  \node[frag,right=2mm of f1] (f2) {$r_{N\text{-}3}$};
  \node[frag,right=2mm of f2] (f3) {$r_{N\text{-}2}$};
  \node[frag,right=2mm of f3] (f4) {$r_{N\text{-}1}$};
  \node[newfrag,right=2mm of f4] (f5) {$r_{N}$};
  \draw[decorate,decoration={brace,amplitude=5pt,mirror},slate]
    ([yshift=-3pt]f3.south west) -- ([yshift=-3pt]f5.south east)
    node[midway,below=7pt,align=center]{reconstruction window $R=3$\\(a read combines $\le R$ fragments)};
  \draw[->,nodenew,thick] (f2.north) to[bend left=26]
    node[midway,above,font=\footnotesize]{uncovered slots $\Rightarrow$ preserved} (f5.north);
  \node[font=\footnotesize\itshape,above=4mm of f1.north,anchor=south west,xshift=-2mm]{older fragments (shareable)};
\end{tikzpicture}
\caption{The sliding snapshot. A read reconstructs a page from at most $R$ consecutive fragments
(here $R=3$). When a fragment leaving the window holds a slot that no newer in-window fragment
covers, that slot is preserved into the newest fragment $r_N$ at write time, which bounds
reconstruction to $R$ fragments while spreading carry-forward evenly across revisions.}
\label{fig:slide}
\end{figure}

A logical leaf is a \code{PageReference} whose \code{key} is the on-disk offset of the \emph{newest}
fragment, plus an ordered list of \code{PageFragmentKey}s for the older ones; each fragment key is a
four-field record \code{(revision, key, databaseId, resourceId)}. Reconstruction
(\code{getPageFragments}) loads fragment~0, returns early if it is already a complete page or there
are no older fragments, and otherwise fires one asynchronous \code{readPage} per fragment key, joins
them, and \textbf{sorts by revision descending} so the merge sees newest-first. Two caches sit in
front: a \code{RecordPageFragmentCache} of \emph{raw} single fragments and a \code{RecordPageCache} of
\emph{combined} pages, both keyed by the offset-bearing \code{PageReference} (database id, resource
id, log key, and offset), deliberately \emph{not} the revision, so a fragment is shared across all
revisions that reference it. Because reconstruction is roughly half of read CPU, a
\textbf{single-fragment fast path} skips \code{combineRecordPages} entirely whenever the newest
fragment is already complete or no historical fragments exist, loading straight into the
combined-page cache.

A fragment key resolves to bytes in three hops. On a cache miss for
\code{PageFragmentKey (revision=5, key=65536, databaseId=42, resourceId=17)}:

\dcap{Resolving one PageFragmentKey to a page on disk}
\begin{diagramS}
key = (revision 5, key 65536, databaseId 42, resourceId 17)

  1. (databaseId 42, resourceId 17)  -> the resource's file handles { sirix.data, sirix.revisions }
  2. key 65536                       -> the byte OFFSET of this fragment inside sirix.data
  3. at offset 65536:  [u32 len][page payload ...][u64 XXH3]
                       read len, read payload, verify XXH3 — mismatch => torn/bit-rot => fail closed

  result -> a raw single fragment, inserted into RecordPageFragmentCache keyed by the PageReference
            (db id, resource id, log key, offset) — NOT the revision, so every revision citing it shares it
\end{diagramS}

\noindent The revision field selects \emph{which} historical fragment, but the cache key deliberately
omits it: two revisions that both reference the unchanged fragment hit the same cache entry, which is the
in-memory face of the on-disk structural sharing of \S\ref{sec:caching}.

\section{The Transaction Intent Log}

Uncommitted state lives in the \textbf{Transaction Intent Log (TIL)} as a list of
\code{PageContainer}s, each holding a \code{complete} reconstruction (the read view during
modification) and a sparse \code{modified} page (what is serialised at commit); for
\code{FULL} the two are the same instance. \code{TransactionIntentLog.put} takes \emph{exclusive}
ownership of a reference: it evicts that reference from the record and fragment caches, assigns the
container a log-slot key, and stamps a generation tag so an in-flight asynchronous commit and the
live writer never alias. Reads during a write resolve through the TIL first (returning the
container's \code{complete} page), which is why a read-during-write sees a consistent in-progress
view without exposing partial state. Commit serialises the modified pages in parallel, the point
at which each page's deferred preservation copy (\S\ref{sec:slidealg}) is finally materialised, then writes them to
the append-only data file and links them into the new \code{RevisionRootPage}. \textbf{Rollback} is
simply discarding the TIL (closing each \code{PageContainer}), with no disk writes and the
on-reference fragment chains left untouched.

A three-operation transaction shows the TIL mediating every read and write:

\dcap{The TIL across a three-operation write transaction}
\begin{diagramS}
beginNodeTrx()        acquire the write semaphore;  TIL empty
write node A (slot 5) TIL.put(leafRef_A): container{ complete = reconstructed page, modified = {slot 5} }
                      ref evicted from record/fragment caches; log-slot key + generation tag stamped
write node B (slot 1) TIL.put(leafRef_B): a second container, same complete/modified shape
read  node A          resolves through the TIL FIRST -> returns container_A.complete (sees slot 5's NEW
                      value), never the bare on-disk page; so a read-during-write is consistent, not partial
commit()              serialize each container.modified in PARALLEL (preservation copies materialise here),
                      append to sirix.data, link into RevisionRootPage r+1, clear the TIL, release semaphore
   -- or --
rollback()            close each PageContainer; NO disk writes; on-reference fragment chains untouched
\end{diagramS}

\noindent The \code{complete}/\allowbreak\code{modified} split is what makes both properties hold at once
(Figure~\ref{fig:til}):
reads during the write see \code{complete} (a full, consistent in-progress page), while only the sparse
\code{modified} page is serialised at commit --- and because nothing reached disk until commit, rollback
is just dropping the in-memory list.

\begin{figure}[t]
\centering
\begin{tikzpicture}[font=\sffamily\small,>=Latex,
  io/.style={draw=slate!70,fill=slate!12,rounded corners=2pt,align=center,inner sep=3pt,
             minimum height=1.05cm,minimum width=2.0cm,font=\footnotesize},
  comp/.style={draw=slate!55,fill=layerbg,align=center,inner sep=2pt,
               minimum height=0.72cm,minimum width=1.55cm,font=\scriptsize},
  modf/.style={draw=ember!75,fill=ember!12,align=center,inner sep=2pt,
               minimum height=0.72cm,minimum width=1.2cm,font=\scriptsize,text=ember!85!black},
  cstep/.style={draw=slate!55,fill=white,rounded corners=2pt,align=center,inner sep=3pt,
                minimum height=0.82cm,text width=3.0cm,font=\footnotesize},
  rb/.style={draw=nodeold,fill=nodeold!12,rounded corners=2pt,align=center,inner sep=3pt,
             font=\scriptsize,text=nodeold!55!black,text width=3.0cm},
  e/.style={->,slate},
  ce/.style={->,ember}]

  % writer
  \node[io] (w) at (0,2.05) {write txn\\\code{put} A (slot 5)\\\code{put} B (slot 1)};

  % TIL: two PageContainers, each split complete | modified
  \node[font=\footnotesize\bfseries,text=slate!80!black] at (3.85,3.35) {Transaction Intent Log};
  \node[font=\scriptsize\itshape,text=slate!70!black] at (3.2,2.98) {read view};
  \node[font=\scriptsize\itshape,text=ember!75!black] at (4.55,2.98) {commit view};
  \node[comp] (cA) at (3.2,2.45) {complete\\(page A)};
  \node[modf] (mA) at (4.55,2.45) {mod\\\{5\}};
  \node[comp] (cB) at (3.2,1.62) {complete\\(page B)};
  \node[modf] (mB) at (4.55,1.62) {mod\\\{1\}};
  \draw[rounded corners=3pt,draw=slate!45] (2.38,1.18) rectangle (5.28,2.86);

  \draw[e] (w.east) -- node[above,font=\scriptsize,text=slate!70!black]{\code{put}} (2.38,2.05);

  % commit path: serialise modified, append, link  (spaced 1.15cm apart so the boxes never overlap)
  \node[cstep] (ser)  at (7.7,2.55) {serialise \code{modified} ($\parallel$)};
  \node[cstep] (app)  at (7.7,1.40) {append to \code{sirix.data}};
  \node[cstep] (link) at (7.7,0.25) {link into \code{RevisionRoot} $r{+}1$, then beacon};
  \draw[ce] (mA.east) to[out=0,in=180] node[above,font=\scriptsize,text=ember!80!black,pos=0.55]{commit} (ser.west);
  \draw[ce] (mB.east) to[out=0,in=195] (ser.west);
  \draw[ce] (ser)--(app); \draw[ce] (app)--(link);

  % rollback path (kept clearly left of the commit chain and below the TIL box)
  \node[rb] (rb) at (3.55,0.0) {rollback: discard TIL --- no disk I/O, fragment chains untouched};
  \draw[->,nodeold] (3.4,1.18) -- node[right,font=\scriptsize,text=nodeold!55!black,inner sep=1.5pt]{or} (rb.north);
\end{tikzpicture}
\caption{The write path through the Transaction Intent Log. Each mutation puts a \code{PageContainer}
holding a full \code{complete} page (the consistent read view during the write) and a sparse
\code{modified} page (only the touched slots). Reads during the write resolve through the TIL first and
see \code{complete}; \emph{commit} serialises just the \code{modified} pages in parallel --- the moment
each page's deferred preservation copy (\S\ref{sec:slidealg}) materialises --- appends them to
\code{sirix.data}, and links them into \code{RevisionRootPage} $r{+}1$ before the dual-beacon barrier
(Figure~\ref{fig:committimeline}). \emph{Rollback} simply drops the in-memory list: nothing reached disk,
so the on-reference fragment chains are untouched.}
\label{fig:til}
\end{figure}

%================================================================
\chapter{The On-Disk Format and the Durable-Commit Protocol}
\label{ch:ondisk}

\noindent\textit{(\code{docs/DISK\_FORMAT.md}; \code{io/sirix/io/*};
\code{io/sirix/page/PagePersister.java}.)}

\section{File layout}

Chapters~\ref{ch:pds} and~\ref{ch:versioning} built a versioned page tree in memory; this chapter
is what makes that tree durable on disk and crash-safe. A resource's
\code{FILE\_CHANNEL}/\allowbreak\code{MEMORY\_MAPPED} storage is two files. Both open
with an identical 64-byte \textbf{superblock} (\code{io/sirix/io/Superblock.java})
validated fail-fast at open: the 8 ASCII magic bytes \code{SIRIXDB!}, layout version~0, a
role byte (\code{DATA}/\allowbreak\code{REVISIONS}), an endianness-check word (\code{0x01020304}),
geometry (beacon-slot size, primary-beacon offset, content-start), and an XXH3-64 CRC at
byte~56 covering the \emph{leading 56 bytes} (it cannot checksum itself). \code{SuperblockValidator}
memoises validated paths per-JVM and fails in a fixed order (too-short, all-zero, bad magic,
wrong version, endianness mismatch, role mismatch, CRC mismatch), with the all-zero case carrying
a distinct ``interrupted first commit / lost header'' message routed to recovery (\S\ref{sec:recovery}).

\dcap{The two-file on-disk layout}
\begin{diagram}
sirix.data
  0       Superblock (64 B): magic, version, role=DATA, endianness, geometry, XXH3 CRC
  64..4096  sparse hole (written only at first commit)
  4096    PRIMARY beacon slot   [u32 len][UberPage payload][u64 XXH3][zero pad]   (4 KiB block)
  8192    SECONDARY beacon slot [identical copy — a SEPARATE 4 KiB block]         (4 KiB block)
  12288   DATA_REGION_START — append-only page records: [u32 len][payload], 8-aligned
          …───────────────────────────────────────────────────────────────────▶ (grows)

sirix.revisions
  0       Superblock (role=REVISIONS)
  4096    REVISIONS_RECORDS_START — fixed 32-byte slots, slot N at 4096 + N*32:
          [u64 rootPageOffset][u64 epochMillis][u64 checksum][u64 rootPageHash]
\end{diagram}

Endianness regimes: \emph{every} multi-byte record field is pinned \textbf{little-endian} ---
superblocks, beacon checksums, revision records, and the page-record length prefixes and payload
primitives alike --- so a resource is byte-identical across host architectures. Two fields stay
\textbf{big-endian} by design, where byte order \emph{is} sort order: the HOT trie keys
(Ch.~\ref{ch:hot}) and the canonical page-hash word; the superblock's endianness word is retained as a
defensive gate against a foreign-endian host.
The \textbf{deterministic} revision-slot layout (slot $N$ is always at $4096 + N{\times}32$,
never appended at file-size) means a failed commit cannot shift later slots, and recovery
can derive the \emph{logical} write frontier from the durable revision graph instead of
from \code{fileChannel.size()}, the property the preallocated-commit mode
(\S\ref{sec:prealloc}) depends on. The fourth revision-record field is the XXH3-64 of the
revision-root page's compressed body, the only integrity anchor on the otherwise
parent-less \code{readRevisionRootPage} path. The checksum covers 24 bytes when a hash is
present and 16 for legacy records, so older resources open without false corruption errors.

A byte-level look at one beacon slot shows what ``torn'' means and how the framing catches it. The primary
slot at offset~4096 holds the latest uber page as \code{[u32 len][payload][u64 XXH3]} inside its 4\,KiB
block:

\dcap{A beacon slot before, during, and after committing revision N+1}
\begin{diagramS}
offset 4096 (primary), before the commit — holds revision N's uber page:
  4096: [ len = 312 ]                         (u32, little-endian)
  4100: < 312 bytes of UberPage payload >
  4412: [ XXH3 = 0x9F3A..C1 ]                  (u64 over those 312 payload bytes)
  4420 .. 8192: zero pad        (the SEPARATE secondary copy lives at offset 8192)

commit N+1 overwrites the slot IN PLACE with [len=320][new payload][new XXH3], AFTER the data tail has
been fdatasync'd (tail-before-beacon ordering, Property~\ref{prop:torn}):

SIGKILL after only ~256 of the ~332 new bytes reach the platter — a TORN beacon:
  reopen reads [len = 320], reads 320 payload bytes, recomputes XXH3
  but the stored-XXH3 field was in the un-written tail -> still old/zero -> recomputed != stored -> INVALID
  => primary rejected; fall back to SECONDARY (offset 8192), which still names revision N, intact
  later, repairBeaconSlotsAfterTruncate heals the torn primary by copying the good slot over it
\end{diagramS}

\noindent The two copies occupy \emph{separate} 4\,KiB blocks (4096 and 8192) precisely so a single torn
block can damage only one: a one-block atomic write is the device's unit of failure, so splitting the
beacon across two blocks guarantees a surviving copy. That is the byte-level content of the dual-beacon
recovery table of \S\ref{sec:recovery}.

\section{The storage SPI and its backends}
\label{sec:storagespi}

Everything above the bytes talks to storage through one small service interface, so the on-disk format
of this chapter is really the format of \emph{one} backend among several. The page layer reads and
writes pages only through \code{IOStorage}, which hands out a \code{Reader} and a \code{Writer}. The
\code{Reader} fetches a page given its \code{PageReference} (and the uber-page and revision-root pages);
the \code{Writer} adds \code{write}, the \code{writeUberPageReference} that drives the durable-commit
protocol below, and \code{truncate}. Between the page bytes and the backend sits a configurable
\code{ByteHandlerPipeline} for compression and optional encryption, the per-page codec described later
in this chapter.

\dcap{The storage SPI and its backends (\texttt{IOStorage} / \texttt{Reader} / \texttt{Writer})}
\begin{diagramS}
page layer ──(PageReference: read / write a page)──▶ IOStorage ──▶ ByteHandlerPipeline ──▶ backend
                                                     createReader() / createWriter()   (compress / encrypt)

backends (StorageType, chosen per resource in the ResourceConfiguration):
  IN_MEMORY      RAM only, for tests and ephemeral resources
  FILE           RandomAccessFile
  FILE_CHANNEL   the default; the two-file, dual-beacon layout the rest of this chapter describes
  MEMORY_MAPPED  memory-mapped files
  IO_URING       Linux io_uring asynchronous I/O (Ch. 9)
  S3             object storage, supplied by an external StorageProvider (ServiceLoader)
\end{diagramS}

The backend is a per-resource choice, fixed in the \code{ResourceConfiguration} as a \code{StorageType}.
\code{FILE\_CHANNEL} is the default and the one whose two-file, dual-beacon layout the rest of this
chapter describes; \code{MEMORY\_MAPPED} maps the same files; \code{IN\_MEMORY} keeps everything in RAM
for tests; \code{IO\_URING} swaps in Linux asynchronous I/O (Ch.~\ref{ch:enterprise}). The list is open:
\code{StorageType.getStorage} consults external \code{StorageProvider} implementations through a
\code{ServiceLoader} before falling back to the built-ins, which is the seam an enterprise layer uses to
add an \code{S3} object-store backend without touching the core. Because the page layer only ever sees
\code{IOStorage}, nothing above it changes when the backend does.

\textbf{How a backend is found.} The seam is a \code{ServiceLoader} SPI. A \code{StorageProvider}
declares a \code{getName()}, an \code{isAvailable()} (so a provider whose native library is missing
removes itself), an \code{isEnterprise()} marker, and a \code{getPriority()}; \code{StorageProviders}
loads every provider on the classpath, indexes them by name, and resolves a name collision in favour of
the highest priority. \code{StorageType.getStorage} consults that registry first and only then falls back
to the built-in \code{getInstance}, so \code{IO\_URING} and \code{S3} resolve to real implementations
when the \code{sirix-enterprise-core}/\allowbreak\code{-s3} modules are present and otherwise throw a message that
names the missing module rather than failing obscurely. The built-in three need no provider.

\textbf{What each built-in backend actually does.} The default \code{FILE\_CHANNEL}
(\pp{io/sirix/io/filechannel/}) is not one channel but \emph{three}, each with the durability its data
needs: a \emph{buffered} data channel for the bulk page tail, a \emph{write-through} (\code{O\_SYNC})
revisions channel so a 32-byte revision record is durable the instant it returns, and a write-through
(\code{O\_DSYNC}) beacon channel for the two uber-page slots. Reads go through a pooled direct
\code{ByteBuffer} (128\,KiB by default, \code{-Dsirix.filechannel.bufferBytes}) to avoid an allocation per
page. \code{MEMORY\_MAPPED} (\pp{io/sirix/io/memorymapped/}, \code{MMStorage}\,+\,\code{MMFileReader})
reads by \emph{slicing a mapped \code{MemorySegment}} with no byte-array copy and hints the kernel with
\code{madvise(MADV\_SEQUENTIAL)} for scans, but \textbf{reuses \code{FileChannelWriter} verbatim for
writes}, so it inherits exactly the dual-beacon commit; its one subtlety is generational arena
management, when a commit grows the file a new mapping generation is created and the old one is unmapped
only once every reader still holding it has released its reference count (on HotSpot via
\code{Arena.ofShared}, on GraalVM via \code{Arena.ofAuto}). \code{IN\_MEMORY} (\code{RAMStorage}) keeps
pages and revision roots in two \code{ConcurrentHashMap}s and is for tests and ephemeral resources only,
it persists nothing.

\dcap{The built-in backends, by durability and copy behaviour}
\begin{diagramS}
backend        read path                       write path                durability
FILE_CHANNEL   pread → pooled 128KiB buffer     3 channels (data buf,     fsync + dual beacon
               (default)                         revisions O_SYNC,         (O_SYNC / O_DSYNC)
                                                 beacon O_DSYNC)
MEMORY_MAPPED  slice mapped MemorySegment       FileChannelWriter         same as FILE_CHANNEL
               (zero-copy) + madvise SEQUENTIAL  (reused verbatim)         (writes go via FileChannel)
IN_MEMORY      ConcurrentHashMap get            ConcurrentHashMap put     NONE (tests / ephemeral)
\end{diagramS}

\textbf{The shared read contract.} Whatever the backend, \code{AbstractReader} enforces one integrity
discipline: a page's body carries an XXH3-64 hash, computed over the \emph{compressed} bytes before the
byte-handler pipeline runs, and \code{verifyChecksumIfNeeded} checks it on read (gated by
\code{verifyChecksumsOnRead} so a latency-critical deployment can opt out). The revision record carries
the same hash for the parentless \code{RevisionRootPage}, which is how a time-travel open validates the
root it just seeked to. The page layer above never learns which backend served the bytes or whether they
were compressed, encrypted, mapped, or read through a 128\,KiB buffer.

\section{The commit protocol and its crash contract}
\label{sec:commit}

The dual-beacon protocol writes the new uber page to two separate 4~KiB blocks, after the
page tail is durable:

\dcap{Algorithm 4.1 --- Commit (default profile, \texttt{FileChannelWriter.writeUberPageReference})}
\begin{diagram}
  flush buffered page tail (ends with the new RevisionRootPage)
  write superblocks if first commit
  force(false)  data                     # WRITE-AHEAD BARRIER (tail durable)
  # revision record already written via O_SYNC channel, durable at write-return
  write SECONDARY beacon  (buffered to data channel)
  write PRIMARY beacon    (buffered to data channel)
  force(false)  data                     # COMMIT-END BARRIER (both beacons durable)
\end{diagram}

Figure~\ref{fig:committimeline} draws the protocol as a timeline, which turns the crash contract into a
matter of \emph{where} on it a failure lands rather than a table to memorise.

\begin{figure}[t]
\centering
\begin{tikzpicture}[font=\sffamily\small,>=Latex,
  ph/.style={draw=slate!55,fill=layerbg,rounded corners=2pt,minimum height=0.8cm,minimum width=1.9cm,align=center,inner sep=2pt},
  ack/.style={draw=nodenew,fill=nodenew!12,rounded corners=2pt,minimum height=0.8cm,minimum width=2.2cm,align=center,inner sep=2pt},
  fs/.style={draw=ember!80!black,very thick,->}]
  \draw[->,slate!70,thick] (-0.6,0) -- (12.4,0) node[below right,font=\scriptsize\itshape]{time};
  \node[ph]  (tail) at (1.1,0) {append\\page tail};
  \node[ph]  (sec)  at (5.3,0) {write\\SECONDARY};
  \node[ph]  (pri)  at (7.4,0) {write\\PRIMARY};
  \node[ack] (ack)  at (11.0,0) {revision\\acknowledged};
  \draw[fs] (3.1,0.85) -- (3.1,-0.85);
  \node[font=\scriptsize,align=center,color=ember!80!black] at (3.1,1.3) {force(data)\\write-ahead\\barrier};
  \draw[fs] (9.2,0.85) -- (9.2,-0.85);
  \node[font=\scriptsize,align=center,color=ember!80!black] at (9.2,1.3) {force(data)\\commit-end\\barrier};
  \draw[decorate,decoration={brace,amplitude=4pt,mirror},slate] (-0.4,-1.15) -- (3.0,-1.15)
     node[midway,below=4pt,font=\scriptsize,align=center,text width=3.0cm]{crash here\\$\Rightarrow$ previous revision (never acknowledged)};
  \draw[decorate,decoration={brace,amplitude=4pt,mirror},slate] (3.2,-1.15) -- (9.1,-1.15)
     node[midway,below=4pt,font=\scriptsize,align=center,text width=5.0cm]{tail durable, beacons being written $\Rightarrow$ previous revision, or fall back to the intact beacon if one block tore};
  \draw[decorate,decoration={brace,amplitude=4pt,mirror},nodenew!70!black] (9.3,-1.15) -- (12.0,-1.15)
     node[midway,below=4pt,font=\scriptsize,align=center,text width=2.6cm]{crash here\\$\Rightarrow$ new revision (acknowledged)};
\end{tikzpicture}
\caption{The default durable-commit timeline. The page tail (ending in the new revision-root) is appended
and made durable by the \emph{write-ahead barrier} before either beacon is written; the two beacon copies
then go to separate 4\,KiB blocks, and the \emph{commit-end barrier} acknowledges the revision.
Acknowledgement is a single durable event, so a crash falls cleanly on one side of it: the brace under each
window gives the outcome on reopen, and only the middle window needs the dual-copy, checksum-validated
fall-back of Property~\ref{prop:torn}.}
\label{fig:committimeline}
\end{figure}

\textbf{The crash contract.} Because the tail is fsynced \emph{before} any beacon is
written, and the two beacon copies are in separate blocks:

\begin{table}[h]
\small\setlength{\tabcolsep}{5pt}
\renewcommand{\arraystretch}{1.2}
\begin{tabularx}{\linewidth}{@{}l c c X@{}}
\toprule
\textbf{Crash point} & \textbf{PRIMARY} & \textbf{SECONDARY} & \textbf{Outcome on reopen}\\
\midrule
before write-ahead barrier & old & old & previous revision (commit never acknowledged) --- correct\\
after barrier, before beacons & old & old & previous revision --- correct\\
mid-beacon writes (one torn) & new/torn & old/new & reader validates \code{[len][payload][XXH3]}, falls back to the intact copy, naming a \emph{durable} tail\\
after commit-end barrier & new & new & new revision --- acknowledged\\
\bottomrule
\end{tabularx}
\end{table}

\begin{property}[Torn-block safety]\label{prop:torn}
Whichever beacon copy survives a torn-block crash names a durable tail, because the tail is
fsynced before any beacon is written and the two copies occupy distinct filesystem blocks.
\end{property}
One subtle residual case is a crash between the secondary and primary writes that leaves the
secondary advertising a \emph{truncated-away} revision. It is healed by
\code{repairBeaconSlotsAfterTruncate}, which copies the good slot over the stale one so both
agree.

\subsection*{Where a commit spends its time}
\label{sec:commitphases}

That is the commit's correctness anatomy; its \emph{cost} anatomy is worth measuring too, because it
explains the per-commit latencies reported in the evaluation. The writer instruments its own commit into
named phases --- \code{serialize} (encode the copied pages, with optional compression, into the off-heap
buffers and append them to the data file), \code{recursive} (relink the copied page-trie spine up to the
new revision root), \code{force} (the \code{fdatasync} barrier over the data tail), and \code{fsync} (the
beacon writes). Capturing those phases for a bulk ingest commit at three document sizes (lean
configuration, warmed JVM) isolates where the time goes.

\begin{table}[h]
\centering\footnotesize\setlength{\tabcolsep}{7pt}\renewcommand{\arraystretch}{1.25}
\begin{tabular}{@{}l r r r r@{}}
\toprule
\textbf{Ingest commit} & \textbf{serialize} & \textbf{relink} & \textbf{durability barrier} & \textbf{total}\\
\midrule
1\,MB  & 33\,ms (83\%)  & 3\,ms  & 3\,ms  & 40\,ms\\
4\,MB  & 122\,ms (83\%) & 9\,ms  & 16\,ms & 147\,ms\\
16\,MB & 185\,ms (84\%) & 15\,ms & 19\,ms & 221\,ms\\
\bottomrule
\end{tabular}
\caption{Commit-phase breakdown of a bulk ingest at three sizes (\emph{durability barrier} sums the data
\code{fdatasync} and the two beacon writes). Page serialisation dominates at a steady $\sim$83\%; the
durability barrier is a small, slowly-growing tail.}
\label{tab:commitphases}
\end{table}

Two things stand out. Page \textbf{serialisation dominates} at a near-constant $\sim$83\% of commit time
across a $16\times$ size range: a copy-on-write commit's cost is mostly the CPU to encode the copied
pages, not the disk barrier. And the \textbf{durability barrier is a small, slowly-growing tail}
(3--19\,ms) --- which is the reverse of where the cost lands for a \emph{tiny} per-update commit, where
serialisation is negligible and that same $\sim$10--15\,ms barrier becomes the floor. That is exactly the
shape the large-document crossover sees from the outside (\S\ref{sec:crossover}): SirixDB's
$\sim$8--10\,ms per single-field commit (full configuration) is barrier-dominated --- one changed page is
cheap to serialise --- and flat in document size, whereas a whole-document rewrite pays serialisation
proportional to the document. The phase breakdown is the inside view of that flat line.

\subsection*{Commit latency under JMH: barrier floor, then serialisation}

The phase split above is the engine's own instrumentation; an independent, statistically controlled view
comes from \textbf{JMH} (the Java Microbenchmark Harness), which forks fresh JVMs and runs warmup
iterations before timing, so the cold-start and ordering biases a hand-rolled loop is prone to cannot
creep in. A \code{CommitLatencyBenchmark} (in the project's \pp{sirix-benchmarks} module) times one durable
commit on a 1\,MB lean resource (hashing off, no path summary, no DeweyIDs) as the number of changed slots
per commit sweeps $1\to512$ (3 forks, 3 warmup {+} 8 measurement iterations, 24 samples per row):

\begin{table}[h]
\centering\footnotesize\setlength{\tabcolsep}{8pt}\renewcommand{\arraystretch}{1.25}
\begin{tabular}{@{}r r r@{}}
\toprule
\textbf{changed slots / commit} & \textbf{latency (ms/op)} & \textbf{$\pm$99.9\% CI}\\
\midrule
1   & \textbf{4.30} & 0.28\\
8   & 9.22 & 0.63\\
64  & 37.0 & 4.4\\
512 & 47.6 & 11.8\\
\bottomrule
\end{tabular}
\caption{JMH commit-latency sweep (\code{Mode.AverageTime}, mean of 24 samples). A single-field commit sits
near the durability-barrier floor; latency then climbs with the number of leaf \emph{pages} the change
touches, and saturates once scattered changes span the document's $\sim$78 leaf pages.}
\label{tab:jmhcommit}
\end{table}

The curve settles what bounds a commit. A \textbf{single-field commit is barrier-bound}: at 4.3\,ms it is
essentially the fsync floor, because serialising one changed page is negligible. As the change grows,
latency rises with the number of \emph{distinct leaf pages} touched --- serialisation, not the barrier ---
but \emph{sub-linearly} in the slot count, because scattered changes soon hit every one of the document's
$\sim$78 leaf pages: 64 and 512 changes both rewrite nearly all of them, so their latencies ($37$ and
$48$\,ms) are close. \textbf{Commit latency is therefore barrier-bound when you change little and
serialisation-bound when you change a lot}, with the crossover where the change starts touching most pages;
the bulk-ingest $\sim$83\,\% serialisation figure is simply the far-right end of this same curve. The run
also retroactively explains the write-scaling correction of \S\ref{sec:concurrency}: the cold single-thread
baseline that inflated that result is exactly the artifact JMH's warmup exists to remove, which is why the
project carries its benchmarks as JMH harnesses and why the figures here are quoted with confidence
intervals.

\section{Recovery and integrity}
\label{sec:recovery}

Recovery (\pp{io/sirix/io/AbstractReader.java}, \pp{FileChannelWriter.java},
\pp{InterruptedFirstCommitRecovery.java}) rests on three pillars: \textbf{dual-beacon
redundancy} (validate bounds {+} XXH3, fall back PRIMARY\,\textrightarrow\,SECONDARY);
\textbf{checksummed records and page bodies} (a torn record or page is caught before it is
trusted); and \textbf{beacon repair} after truncation. The \textbf{interrupted-first-commit}
gap (a crash before a resource's header/beacons exist, leaving a non-empty file with a
zero header) is healed \emph{conservatively} by \code{InterruptedFirstCommitRecovery}: it
re-initialises the resource empty only if the on-disk bytes \emph{prove} that at most the
empty bootstrap revision was ever acknowledged (all-zero or checksum-valid superblock;
beacons either zero or advertising only revision~0; and no checksum-valid revision record);
otherwise it rethrows the original failure untouched.

\dcap{Dual-beacon recovery: which uber page wins after a crash}
\begin{diagramS}
on open, validate both 4 KiB beacon slots (length bounds + XXH3):

  PRIMARY(4096)     SECONDARY(8192)    resolution
  rev N  ok         rev N  ok          clean: use rev N
  rev N  BAD        rev N-1  ok        torn primary: fall back to SECONDARY (rev N-1)
  rev N  ok         torn               use PRIMARY (rev N), then repair SECONDARY
  all-zero          all-zero           interrupted first commit → re-init empty, only if proven safe

a .commit marker present on open ⇒ truncateTo the last good revision before a writer is accepted
\end{diagramS}

The whole protocol is exercised by \textbf{\code{CrashRecoveryInjectionTest}}
(\code{io/sirix/crash/}), an opt-in (\code{-Dsirix.crash.run=true}) harness that forks a
writer, \code{SIGKILL}s it at random instants in its commit loop, and verifies on reopen
that: the database opens; every acknowledged commit is intact; all pages deserialise; the
truncate-recovery path runs when the \code{.commit} marker is present; and the resource
accepts a new writer. A complementary power-loss simulator models a stricter metadata-split
durability than POSIX fdatasync, which is \emph{why} (see \S\ref{sec:evolution}) the two
data barriers that cover the tail append remain full fsyncs deliberately even when the rest
of the protocol relaxes.

Walking the kill point through the protocol shows the contract holding at every instant.

\dcap{What reopen sees for a \texttt{SIGKILL} at each phase of committing revision N+1}
\begin{diagramS}
SIGKILL while ...                              reopen resolves to ...
  building N+1 in the TIL (nothing on disk)      revision N   (commit never reached disk)
  after the tail fsync, before any beacon        revision N   (both beacons still name N; N+1 unacknowledged)
  mid primary-beacon write (torn)                validate [len][payload][XXH3]: primary bad → fall back to secondary (N)
  after secondary, before primary                primary=N, secondary=N+1 → repairBeaconSlots heals both to agree
  after the commit-end barrier                   revision N+1 (both beacons name N+1; acknowledged, durable)

at every kill point the harness asserts: the database opens; every acknowledged commit is intact;
all pages deserialise; a .commit marker triggers truncateTo the last good revision; a fresh writer is accepted.
\end{diagramS}

The key property is that the acknowledgement boundary is a single durable event, the commit-end
barrier, so a crash is always cleanly on one side of it: either N+1 was acknowledged and is fully
present, or it was not and the resource is exactly revision N. There is no in-between state in which a
half-written N+1 is visible, which is what the dual beacons and the tail-before-beacon ordering buy.

\section{The commit protocol's evolution (an engineering case study)}
\label{sec:evolution}

The protocol reproduced in \S\ref{sec:commit} is the endpoint of a measured journey
(\code{docs/BENCHMARKS.md}). The instructive part is how \emph{much} of the per-commit cost
was incidental rather than fundamental.

\begin{enumerate}[label=\arabic*.,leftmargin=2em]
\item \textbf{Seven sync calls per commit.} The original protocol issued 5~\code{fsync} {+}
  2~\code{fdatasync}. A redundant \code{forceAll} (\code{t3}) ran while the tail was still
  buffered and covered strictly less than \code{writeUberPageReference}'s internal barrier.
  Removing it and reducing the acknowledge barrier to a data-only \code{fdatasync} cut the
  count to \textbf{four} (fsync data write-ahead; fsync revisions; fdatasync data
  beacon-order; fsync data acknowledge).
\item \textbf{A quadratic \code{access(2)} pathology.} Wall-clock profiling of a
  10\,000-commit build showed the late phase dominated by \code{access(2)}. A syscall census
  revealed \textbf{50\,196\,928} \code{access()} calls over the build ($\sim$50~M of them
  \code{ENOENT}), versus 520\,k for 1\,000 commits: a perfect $\Sigma i \approx N^2/2$.
  Root cause: \code{initializeIndexController} probed \code{revision.xml},
  \code{(revision-1).xml}, \ldots, \code{0.xml} with one \code{Files.exists} per step to find
  the latest index definitions, and a new controller is created \emph{per commit}. With no
  secondary indexes (the default), no file ever exists, so every commit walked the entire
  history. Fix: a single directory listing picking the max-numbered file $\leq$ revision (an
  empty directory short-circuits). This, plus an \textbf{amortised}
  \code{RevisionIndex.withNewRevision} (capacity-doubling shared arrays with a deferred
  Eytzinger rebuild, replacing a full $O(\text{size})$ array copy per commit), turned the
  build from \textbf{48.4~s} (4.84~ms/commit, declining to $\sim$150 commits/s at depth
  10\,k) into \textbf{20.5~s} (2.05~ms/commit, flat $\sim$570 commits/s), with the decline
  \emph{eliminated}, not merely reduced.
\item \textbf{Toward one explicit sync.} The ``one explicit sync'' design was then
  implemented: the revisions record goes through an \code{O\_SYNC} channel (durable incl.\
  size at write-return) and both beacons through an \code{O\_DSYNC} channel (in-place;
  write-return gives secondary-before-primary ordering and makes the primary's return the
  commit acknowledge), leaving a \emph{single} explicit \code{fdatasync} for the data tail.
  On consumer NVMe with FUA this measured at \textbf{parity} with the 4-sync protocol (three
  serialised write-through round-trips cost about what the saved flushes did), with the
  structural win of an explicit durable-on-return contract; gains are expected on server
  stacks where FUA is materially cheaper than a cache flush. All power-loss and SIGKILL gates
  re-validated green.
\end{enumerate}

The lesson, that most of a commit's apparent cost was redundant syscalls and a hidden
quadratic and not the durability barriers themselves, directly motivated the preallocated
profile below.

\section{The preallocated, journal-free commit profile}
\label{sec:prealloc}

Two \textbf{opt-in} flags (default \emph{off} on \code{main}) configure a low-latency
commit path: \code{sirix.commit.preallocated} and \code{sirix.commit.bufferedBeacons}. With both
off (the shipped default) the protocol is byte-identical to the legacy grow-and-\code{force(true)}
path; the release-candidate work that flips them on by default is validated (below) but, at the
time of writing, uncommitted.

Under preallocation, the data and revisions files are grown ahead of the write frontier in
\textbf{geometrically increasing chunks} (\code{sirix.commit.preallocInitialChunkBytes}
64~KiB, doubling to a \code{preallocChunkBytes} 64~MiB cap). A low-commit resource therefore
ends only a chunk or two past its real data, with a fast first commit; a high-commit-rate
resource ramps to the cap and amortises the one-time per-chunk fsync over many commits.
Because \code{i\_size} is then stable across commits, the write-ahead barrier becomes a
\textbf{journal-free \code{fdatasync}} instead of a growing-file \code{fsync} (which on
ext4/xfs forces a metadata-journal commit). The logical write frontier is \textbf{derived
from the durable revision graph}: read the primary/secondary beacon for the last
committed revision, then its offset from the revisions file. It is never taken from the
preallocation-inflated \code{fileChannel.size()}.

\code{bufferedBeacons} writes the two beacon copies buffered and makes them durable with a
single commit-end \code{fdatasync} instead of two O\_DSYNC (FUA) writes: one fewer
device round-trip per commit, preserving both the write-ahead ordering and the two-copy
redundancy.

A \emph{true} single-\code{fdatasync} design (folding tail and beacons into one barrier) was
evaluated and \textbf{rejected}. Its worst case is a single unordered flush in which both
beacon blocks reach the device before the tail. That leaves \emph{no surviving prior beacon}
and would demand prior-uber reconstruction or a ping-pong scheme; and ping-pong sacrifices
the dual-beacon redundancy, so a single bit-rotted beacon block could lose an
\emph{already-acknowledged} revision. The chosen flags therefore change only the \emph{I/O
shape}, never the durability model. The default-on configuration has been validated under the SIGKILL
injection harness (25/25 random-timing kills, no acknowledged-data loss) and across the full
unit suite (9\,255 tests; the only failures were file-size assertions in tests that assumed a
growing file, since preallocated commits write in place). \code{MEMORY\_MAPPED} uses its own
mmap/msync path, ignores the flags, and remains format-compatible (it reads preallocated
files via the beacon; the trailing zeros are unreachable).

\section{Page serialisation, compression, and the revision index}

Serialisation is dispatched by \code{PagePersister}: every page on disk begins with a one-byte
\textbf{page-kind tag} (\code{UBERPAGE}=3, \code{REVISIONROOTPAGE}=5, \code{INDIRECTPAGE}=4,
\code{KEYVALUELEAFPAGE}=1, the CAS/path/name pages, \code{HOT\_LEAF\_PAGE}=12, \ldots) followed by
a one-byte \code{BinaryEncodingVersion} (\code{V0}); deserialisation reads the tag and routes to
the matching reader. A page's child references serialise through whichever of the three adaptive
reference-container layouts it currently holds, each with its own sub-tag (\code{ReferencesPage4}=0,
\code{BitmapReferencesPage}=1, \code{FullReferencesPage}=2); one reference is
\code{[u8 fragmentCount][per-fragment: u32 rev, u64 key]\ldots[u64 key][u8 hashPresent][u64 XXH3?]},
where the fragment list is exactly the sliding-snapshot chain of Ch.~\ref{ch:versioning}. The
database and resource ids are \emph{not} stored; they are re-injected after read from the
opening context, a deliberate saving across an indirect page's up-to-1024 references.

Compression is \textbf{per page, inside the leaf format}, not an outer stream. The configurable
\code{ByteHandlerPipeline} defaults to identity (the \code{-Dsirix.compression} default is
\code{none}; \code{lz4} inserts a native FFI handler); the real work is a per-page \textbf{codec
bake-off} that encodes the leaf's compacted body with \code{ZeroRunByteCodec}, \code{ByteRunCodec},
and \code{SirixLZ77Codec} and writes the smallest, tagged by a one-byte codec id, or uses native
LZ4 directly when present. Each page body is checksummed with \textbf{XXH3-64} over the
\emph{compressed} bytes and stored in the parent \code{PageReference}; native segments are hashed
zero-allocation by address.

That inner codec is wrapped, when configured, by the \emph{outer} \code{ByteHandlerPipeline}, an ordered
chain of \code{ByteHandler}s applied to the whole page payload. Three handlers ship.
\code{DeflateCompressor} is the classic stream-based zlib codec. \code{FFILz4Compressor} is the modern
one: a Foreign-Function binding to native LZ4 that compresses and decompresses directly over
\code{MemorySegment}s with no intermediate byte array, in a \code{FAST} mode by default or a
\code{HIGH\_COMPRESSION} mode for storage-bound, read-heavy resources, drawing its decompression buffers
from a fixed pool sized to twice the core count (so it stays well-behaved under millions of virtual
threads) and refusing to keep a compressed result that does not save at least a few percent.
\code{Encryptor} wraps Google Tink's \code{StreamingAead} for per-page authenticated encryption. The
handlers compose in order, compress then encrypt, and the pipeline advertises zero-copy
(\code{supportsMemorySegments}) only when \emph{every} handler in it does, so a pure-LZ4 pipeline stays
allocation-free end to end while adding stream-only Deflate or encryption deliberately forces a
byte-array round-trip. (The encryption key currently lives in cleartext in \code{encryptionKey.json}, a
limitation the code flags in place.)

\dcap{Two compression layers: the inner leaf codec inside the outer page pipeline}
\begin{diagramS}
leaf body ─▶ INNER codec bake-off (pick smallest)        ─▶ tagged compacted body
              { ZeroRunByteCodec | ByteRunCodec | SirixLZ77Codec | native LZ4 }
page payload ─▶ OUTER ByteHandlerPipeline (default: identity)
              { DeflateCompressor | FFILz4Compressor(FAST|HIGH) | Encryptor(Tink AEAD) }   compose: compress ▶ encrypt
              zero-copy over MemorySegments iff EVERY handler supports it (pure-LZ4 does; Deflate/Encryptor force a copy)
\end{diagramS}

A concrete leaf shows the bake-off choosing (sizes illustrative): a compacted body is encoded by every
inner codec, the smallest kept, then the outer pipeline runs if configured.

\dcap{One leaf body through the two layers, with sizes}
\begin{diagramS}
compacted leaf body: 10 240 bytes  (slot directory + FSST-coded key/value heap)

INNER bake-off  (encode with each, keep the smallest, tag with its 1-byte codec id):
   ZeroRunByteCodec -> 8 200 B    ByteRunCodec -> 7 800 B    SirixLZ77Codec -> 6 500 B  <- smallest kept
   tagged result: [codec = SirixLZ77][6 500 B]      XXH3-64 taken over THESE compressed bytes

OUTER ByteHandlerPipeline  (here FFILz4 in FAST mode over the whole payload):
   6 500 B -> 4 200 B   saves > a few percent, so kept; otherwise the raw 6 500 B is stored unchanged

on-disk frame: [u8 page-kind][u8 V0][u8 pipeline tag][4 200 B payload]   the parent ref carries the XXH3
\end{diagramS}

\noindent The two layers are independent: the inner bake-off always runs (its codecs are tiny and
content-aware), while the outer pipeline is opt-in and refuses any result that fails to save a few percent,
so enabling LZ4 never enlarges an incompressible page. The checksum is taken over the \emph{inner}
compressed bytes, before the outer layer, so integrity is verified on exactly the bytes the leaf reader
will decode.

A measured ingest shows the two layers compounding. Shredding a 4\,MB document (the inner bake-off always
on), the on-disk size with the outer pipeline set to identity, FFI-LZ4 \code{FAST}, and
\code{HIGH\_COMPRESSION} --- size is deterministic, so the figures are exact, not sampled:

\dcap{Outer-pipeline compression on a 4\,MB document (inner bake-off always on)}
\begin{diagramS}
outer pipeline       on-disk      vs identity
  identity (none)     2.88 MB      1.00x          (4 MB JSON already down to 2.88 MB from the inner bake-off)
  FFI-LZ4 FAST        1.49 MB      1.94x
  FFI-LZ4 HIGH        1.49 MB      1.94x          ~2.7x overall vs the raw 4 MB JSON
\end{diagramS}

\noindent The outer LZ4 layer roughly halves what the inner bake-off leaves ($1.94\times$), and \code{FAST}
already reaches that ratio for this repetitive document, so \code{HIGH\_COMPRESSION} buys no further shrink
here (it earns its keep on less compressible, read-heavy data). LZ4 also ingests \emph{no slower} than
identity in the same run --- writing fewer bytes to disk offsets the compression CPU.

The naive conclusion would be ``so enable LZ4 by default,'' but the architecture argues the opposite, and
it is worth being precise about why. A general block compressor makes the on-disk bytes \emph{not} the
in-memory bytes, so the engine's defining \textbf{in-memory~$=$~on-disk} property (\S\ref{sec:offheap}) is
forfeited: an LZ4 page can no longer be read as a slice of its \code{MemorySegment} or SIMD-scanned in
place, but must be \emph{decompressed on every load} before a single node is touched. This is exactly the
trade the TU~Munich \emph{Data Blocks} line~\cite{datablocks} argues a storage engine should refuse: lean
on \textbf{lightweight, byte-addressable, scan-over-compressed} codecs --- frame-of-reference, dictionary,
bit-packing, the very toolkit of the \emph{inner} bake-off and the PAX columns of Chapter~\ref{ch:vector}
--- and do \emph{not} reach for heavyweight compression that must be undone before the data can be
accessed. So the honest reading inverts the naive one: the architecturally-aligned layer is the inner
bake-off, which preserves direct access while already recovering $\sim$1.4$\times$ here (4\,MB raw to
2.88\,MB); the outer LZ4 is a deliberate, storage-bound-\emph{only} escape hatch that buys ratio by
surrendering the engine's core property, not a default the design's own philosophy would choose. The
better direction --- strengthening the inner codecs toward a full \emph{scan-over-compressed} capability,
so the engine never \emph{needs} block compression --- has largely shipped for numeric/columnar data:
\code{NumberRegionCompact} (Frame-of-Reference + bit-packing) and the SIMD region kernels evaluate
predicates and aggregates directly over the bit-packed bytes in place; the narrower remaining gaps are
noted in \S\ref{sec:futuredirections}.

Two in-memory structures accelerate the revisions file. \code{RevisionFileData} (an
\code{offset}/\allowbreak\code{timestamp}/\allowbreak\code{pageHash} record) is cached per revision in an asynchronous
cache and bulk-loaded with a single \code{pread}, giving the position-independent random-revision
opens measured in \S\ref{sec:histpath}. \code{RevisionIndex} is an \textbf{Eytzinger-laid-out}
(BFS-order, cache-optimal) timestamp\,\textrightarrow\,revision search structure for
time-travel-by-instant; its \code{withNewRevision} is amortised $O(1)$, a capacity-doubling
append with a \emph{deferred} Eytzinger rebuild. This is the fix that flattened the
large-history commit-rate decay of \S\ref{sec:evolution}, replacing a full $O(n)$ rebuild per
commit.

A lookup makes the cache-optimal layout concrete. Seven commits at timestamps
$1000,1015,1020,1050,1100,1200,1300$ (revisions 1--7) lay out in BFS order, and a time-travel query
predecessor-searches it:

\dcap{Eytzinger time-travel: find the revision live at instant T=1120}
\begin{diagramS}
sorted:   ts  1000 1015 1020 1050 1100 1200 1300     (rev 1..7)
Eytzinger eyt[1..7]  (BFS / implicit-BST order; children of k are 2k and 2k+1):
          eyt = [ 1050, 1015, 1200, 1000, 1020, 1100, 1300 ]
                    1     2     3     4     5     6     7

predecessor search for T=1120  (the largest ts <= T):
  k=1   eyt[1]=1050 <= 1120  -> candidate rev4 ; go right   k=3
  k=3   eyt[3]=1200 >  1120  -> go left                     k=6
  k=6   eyt[6]=1100 <= 1120  -> candidate rev5 ; go right   k=13 > 7 -> stop
  answer:  revision 5 (ts 1100), the revision live at instant 1120

cache win: each level (eyt[1] | eyt[2..3] | eyt[4..7]) is CONTIGUOUS in memory, so the next probe usually
sits in the already-prefetched cache line; a classic sorted-array binary search jumps to scattered n/2, n/4
\end{diagramS}

\noindent The Eytzinger order is what buys the speed: a textbook binary search over a sorted array is
$O(\log n)$ comparisons but with a cache miss at almost every step, whereas the BFS layout keeps each
descent within prefetched lines, turning a time-travel-by-instant lookup into a handful of in-cache probes.

\medskip\noindent\textbf{The two axes, together.} The revisions file and the page trie are the two
halves of every historical read, and they are \emph{independent} coordinates. Opening revision $r$ is a
direct index into the revisions file --- its 32-byte record sits at byte $4096 + r\cdot 32$, so \emph{any}
revision is one array probe away, never a reconstruction --- and the \code{offset} that record yields seeks
straight to that revision's \code{RevisionRootPage} in \code{sirix.data}. From there the copy-on-write
node-key trie of Chapter~\ref{ch:pds} (Figure~\ref{fig:pagetrie}) takes over \emph{within} the immutable
snapshot. The fourth field, a hash of the root page, lets the read verify it landed on an intact root
before trusting a single byte. Figure~\ref{fig:twoaxis} draws the whole resolution.

\begin{figure}[t]
\centering
\begin{tikzpicture}[font=\sffamily\small,>=Latex,
  tcell/.style={draw=slate!55,fill=slate!8,minimum width=1.0cm,minimum height=0.58cm,
                inner sep=1pt,font=\footnotesize,align=center},
  thi/.style={draw=ember!80,fill=ember!15,minimum width=1.0cm,minimum height=0.58cm,
              inner sep=1pt,font=\footnotesize,align=center,text=ember!85!black},
  fld/.style={draw=boxrule,fill=layerbg,minimum width=1.18cm,minimum height=0.6cm,
              inner sep=1pt,font=\scriptsize,align=center},
  fldhi/.style={draw=ember!80,fill=ember!12,minimum width=1.18cm,minimum height=0.6cm,
                inner sep=1pt,font=\scriptsize,align=center,text=ember!85!black},
  pg/.style={draw=slate!60,fill=layerbg,rounded corners=2pt,align=center,inner sep=2pt,
             minimum height=0.82cm,minimum width=1.85cm,font=\footnotesize},
  leaf/.style={draw=nodenew,fill=nodenew!12,rounded corners=2pt,align=center,inner sep=2pt,
               minimum height=0.82cm,minimum width=1.85cm,font=\footnotesize},
  e/.style={->,slate},
  be/.style={->,ember,thick},
  de/.style={->,boxrule,densely dashed}]

  % --- TIME AXIS: the sirix.revisions side-file, direct-indexed by revision ---
  \node[font=\footnotesize\itshape,text=slate!75!black,align=center]
       at (2.35,5.15) {time axis --- \code{sirix.revisions}: fixed 32-byte records, revision $r$ at byte $4096+r\cdot 32$};
  \node[tcell] (c0) at (0.00,4.35) {rec 0};
  \node[tcell] (c1) at (1.05,4.35) {rec 1};
  \node[tcell] (c2) at (2.10,4.35) {$\cdots$};
  \node[thi]   (cr) at (3.15,4.35) {rec $r$};
  \node[tcell] (c4) at (4.20,4.35) {$\cdots$};

  % exploded 32-byte record of revision r
  \node[fldhi] (f1) at (1.61,2.95) {data\\offset};
  \node[fld]   (f2) at (2.79,2.95) {time\\stamp};
  \node[fld]   (f3) at (3.97,2.95) {check\\sum};
  \node[fld]   (f4) at (5.15,2.95) {page\\hash};
  \draw[de] (cr.south) -- (3.15,3.55);
  \node[font=\scriptsize,text=boxrule!50!black] at (3.38,3.62) {32 B};
  \draw[de] (3.15,3.55) -- (f1.north);
  \draw[de] (3.15,3.55) -- (f4.north);

  % --- bridge: dataOffset seeks into sirix.data, lands on the revision root ---
  \draw[be] (f1.south) -- node[right,font=\scriptsize,text=ember!85!black,inner sep=2pt]
            {seek \code{offset} in \code{sirix.data}} (1.61,1.72);

  % --- SPACE AXIS: the node-key trie within revision r (cf. fig:pagetrie) ---
  \node[font=\footnotesize\itshape,text=slate!75!black,align=center]
       at (6.9,2.15) {space axis --- node-key trie\\within revision $r$ (Fig.~\ref{fig:pagetrie})};
  \node[pg]   (rr) at (1.61,1.30) {Revision\\Root $[r]$};
  \node[pg]   (i0) at (4.30,1.30) {Indirect};
  \node[pg]   (i1) at (6.85,1.30) {Indirect};
  \node[leaf] (lf) at (9.45,1.30) {Leaf\\slot $k\bmod 1024$};
  \draw[e] (rr) -- node[above,font=\scriptsize,text=ember!80!black]{[off]} (i0);
  \draw[e] (i0) -- node[above,font=\scriptsize,text=ember!80!black]{[off]} (i1);
  \draw[e] (i1) -- node[above,font=\scriptsize,text=ember!80!black]{[off]} (lf);
\end{tikzpicture}
\caption{Resolving a historical address $(\text{revision }r,\ \text{nodeKey }k)$ along two independent
axes. The \emph{time axis} is a direct array index into the \code{sirix.revisions} side-file: record $r$
lives at byte $4096+r\cdot 32$ and carries the revision's \code{RevisionRootPage} offset, commit timestamp,
a record checksum, and the root-page hash --- so opening any revision is one probe, not a replay. The
\code{offset} seeks into \code{sirix.data} to that \code{RevisionRootPage}; from there the \emph{space axis}
is the copy-on-write node-key trie of Figure~\ref{fig:pagetrie}, walked \emph{inside} the immutable
snapshot. \code{RevisionIndex} (the Eytzinger structure above) supplies the by-instant entry into the same
time axis.}
\label{fig:twoaxis}
\end{figure}

\section{Buffer management and caching}
\label{sec:buffering}

Above the storage backend sits one \textbf{JVM-wide buffer pool} (\code{BufferManager}, a single
instance held by \code{Databases}), modeled on a relational engine's shared buffer pool rather than a
per-resource cache. All resources draw from the same pool, and per-resource isolation comes from
\emph{keying}: every cached page is addressed by an offset-bearing \code{PageReference} that embeds the
database and resource ids, so pages can be evicted database by database or resource by resource without
separate pools. Pages live off-heap as \code{MemorySegment} frames from an FFM-arena allocator, so the
cache footprint does not count against the Java heap and the in-memory page stays byte-identical to its
on-disk image (Ch.~\ref{ch:pds}).

The pool is really several caches, each holding one kind of thing and sized as a fraction of the
off-heap budget (overridable with \code{-Dsirix.cache.*}):

\dcap{The buffer pool: what is cached, and roughly how much}
\begin{diagramS}
cache                        holds                                             rough share / cap
  record-page cache          KeyValueLeafPages (the primary data)              ~50 % of budget
  record-page-fragment cache single fragments, for sliding-snapshot rebuild    ~18.75 %
  page cache (Caffeine)      metadata pages: uber, revision-root, indirect, …   ~6.25 %
  revision-root cache        RevisionRootPage by (resource, revision)           5 000 entries
  index-node cache           secondary-index tree nodes                         50 000 entries
  names cache                NamePage dictionaries                              500 entries
  path-summary cache         PathSummary structures                            20 entries
\end{diagramS}

\begin{figure}[t]
\centering
\begin{tikzpicture}[font=\sffamily\scriptsize,>=Latex,
  c/.style={draw=slate!55,fill=layerbg,rounded corners=2pt,minimum width=2.8cm,minimum height=0.75cm,align=center},
  sw/.style={draw=ember!70!black,fill=ember!10,rounded corners=2pt,align=center}]
  \node[font=\bfseries\color{ember}] at (3.4,2.55) {JVM-wide off-heap buffer pool};
  \node[c] (rp) at (0,1.55)   {record-page\\$\sim$50\%};
  \node[c] (fr) at (3.4,1.55) {fragment\\$\sim$19\%};
  \node[c] (pc) at (6.8,1.55) {metadata page\\$\sim$6\%};
  \node[c] (rr) at (0,0.45)   {revision-root};
  \node[c] (ix) at (3.4,0.45) {index-node};
  \node[c] (nm) at (6.8,0.45) {names + path-summary};
  \node[sw] (cs) at (3.4,-1.15) {clock-sweeper threads\\(second chance; never evict below the\\oldest live reader's revision)};
  \draw[->,ember!70!black,dashed] (cs.north) -- (ix.south);
  \draw[->,ember!70!black,dashed] (cs.north) to[bend left=11] (rp.south);
  \draw[->,ember!70!black,dashed] (cs.north) to[bend right=11] (pc.south);
\end{tikzpicture}
\caption{The buffer pool is one off-heap region shared by all resources, partitioned into caches sized
as fractions of the budget and evicted by background clock-sweeper threads. A sweeper skips any page a
live reader has \emph{guarded} (\S\ref{sec:reclaim}), so the lock-free read guarantee holds under memory
pressure.}
\label{fig:bufferpool}
\end{figure}

Eviction runs in the background, PostgreSQL-bgwriter style. A set of \textbf{clock-sweeper} threads
(one per sharded page cache) wake every 100\,ms, scan about a tenth of their entries per cycle from a
persistent clock hand, and give each frame a second chance, clearing a \emph{hot} bit on the first pass
and reclaiming the frame on the second; a synchronous pressure listener evicts on the spot if the
off-heap allocator is exhausted. What keeps this safe under a concurrent reader is the subject of the
next section. Because pages are keyed by offset and not by revision, a fragment shared by many revisions
is cached once and served to all of them.

A few sweep cycles show the second-chance rule, where \emph{hot} means ``touched since the hand last
passed'':

\dcap{Clock-sweep eviction: the hot bit grants one reprieve}
\begin{diagramS}
6 pages, clock hand starts at A.   hot bits: A1 B0 C1 D0 E0 F0   (A and C were recently read)

cycle 1  hand A -> D, visits A,B,C:
   A hot=1 -> clear to 0, KEEP (second chance)   B hot=0 -> RECLAIM   C hot=1 -> clear to 0, KEEP
   (between cycles a reader touches E -> E hot=1)
cycle 2  hand D -> A, visits D,E,F:
   D hot=0 -> RECLAIM   E hot=1 -> clear to 0, KEEP (its recent use saved it)   F hot=0 -> RECLAIM
cycle 3  hand at A, visits the survivors A,C,E:
   A hot=0 -> RECLAIM   C hot=0 -> RECLAIM   E hot=0 -> RECLAIM   (none re-touched since cycle 1/2)

a guarded page (a live reader is inside it, §\ref{sec:reclaim}) is SKIPPED on every visit, hot bit aside
\end{diagramS}

\noindent A page survives a pass only if it was used since the hand last swept it, so frequently-read
pages keep resetting their hot bit and stay resident while cold pages fall on the second visit --- an
$O(1)$-per-page approximation of LRU that needs no per-access bookkeeping, only a bit flip on read.

\section{Safe reclamation: page guards and the epoch watermark}
\label{sec:reclaim}

A background sweeper that frees off-heap frames and a lock-free reader that holds raw pointers into them
are on a collision course: nothing in the type system stops the sweeper from recycling a
\code{MemorySegment} while a reader is mid-traversal. SirixDB resolves this with a \textbf{page guard},
the LeanStore/Umbra discipline~\cite{leanstore} made explicit. Before a reader touches a leaf it calls
\code{acquireGuard}, which increments an \code{AtomicInteger guardCount} and records the page's current
\emph{frame version}; the work happens inside a \code{PageGuard} \code{AutoCloseable} that releases the
count on close. The sweeper's rule is then simple: \textbf{a page with \code{guardCount > 0} is skipped},
so no frame is ever reclaimed while a reader is inside it.

\dcap{The page-guard lifecycle: pin, read, validate-on-release}
\begin{diagramS}
reader:   try (PageGuard g = page.acquireGuard()) {        // guardCount++ ; record versionAtFix
              … traverse the leaf's MemorySegment …
          }                                                 // close: guardCount-- ; check version

sweeper:  if (page.hot)        clear hot bit ; keep            // second chance
          else if (guardCount>0) skip                         // a reader is inside — never reclaim
          else                 close frame ; version++ ; uncharge weight

race caught: if versionAtFix != page.version on release → throw FrameReusedException → reader reloads
orphaned page (evicted while guarded): marked ORPHANED, torn down deterministically at the last release
\end{diagramS}

Two further mechanisms harden the same boundary. The guard captures the frame version \emph{at fix} and
re-checks it \emph{at release}; if a frame was somehow recycled in between, the mismatch throws a
\code{FrameReusedException} that turns into a cache miss and a clean reload rather than a read of stale
bytes, a belt-and-braces check on top of the skip rule. And a page removed from the cache while still
guarded is not freed but \textbf{marked orphaned}, then torn down deterministically when its last guard
releases, so eviction never races teardown.

A concrete interleaving shows the guard rule holding under sweeper pressure:

\dcap{A reader and the sweeper racing on one page P (guardCount + frame version)}
\begin{diagramS}
t0  reader R1:  acquireGuard(P)   guardCount 0->1 ; versionAtFix = 7     R1 now traversing P's segment
t1  sweeper:    visits P:  guardCount = 1 > 0  -> SKIP   (P stays live, its frame is not freed)
t2  reader R1:  still reading P   (the t1 skip kept its bytes valid)
t3  reader R1:  close guard       guardCount 1->0 ; version still 7 == 7  -> clean release
t4  sweeper:    visits P:  not hot, guardCount = 0  -> close frame ; version 7->8 ; uncharge weight

belt-and-braces, had a frame somehow been recycled mid-read:
   on release  versionAtFix 7 != P.version 8  ->  throw FrameReusedException  ->  R1 reloads (cache miss),
   never reading stale bytes
\end{diagramS}

\noindent The skip at \code{t1} is the whole safety argument: a frame is freed only at \code{t4}, after
the reader has left, so a lock-free reader's raw pointer is never dangling. The version re-check is the
second line of defence, converting any hypothetical violation into a clean reload rather than a corrupt
read.

Independently, a \code{RevisionEpochTracker} maintains the \textbf{oldest-active-revision watermark}:
each read transaction registers its revision in one of 65{,}536 packed \code{long} slots (an active bit
plus the revision number, written through a \code{VarHandle}), and \code{minActiveRevision} scans them
lock-free in tens of microseconds to find the lowest revision any reader still needs. A per-resource
sweeper consults this watermark to avoid reclaiming a fragment below it.

\begin{honesty}
\small\textbf{Caveat.} The watermark and the guard count are two protections of different maturity. The
\emph{guard count} is the one the system actually relies on: it is checked by every sweeper and is what
makes lock-free reads safe today. The \emph{epoch watermark} is consulted only by \emph{per-resource}
sweepers, and the sweepers are currently constructed in their global form, so in practice guards, not the
revision watermark, carry the correctness weight; the code's own header flags closing that gap (a
per-resource watermark registry keyed by \code{(databaseId, resourceId)}) as the intended fix. The
report says ``guarded,'' not ``below the watermark,'' wherever it describes what keeps a read safe.
\end{honesty}

\section{The off-heap frame-slot allocator}
\label{sec:allocator}

The frames the buffer pool hands out come from a custom off-heap allocator
(\code{FrameSlotAllocator}, the Linux default; \code{-Dsirix.allocator=pool} selects the legacy
pool, and Windows uses its own), again Umbra-shaped. It reserves \textbf{seven power-of-two size
classes} from 4\,KiB to 256\,KiB, and for each class reserves a large \emph{virtual} region (32\,GiB)
up front while committing physical pages only on demand (\code{MADV\_POPULATE\_WRITE} at allocation),
so a slot's address is stable for the JVM's lifetime but its RAM is paid for only when used. A slot is
recycled rather than unmapped on release (no \code{MADV\_DONTNEED}), trading a touch of resident memory
for a constant-cost, syscall-free recycle.

The interesting part is how a reader can trust a recycled slot. Each slot carries a \textbf{version
counter} with the even/odd discipline of a seqlock: an allocation or release bumps it through an odd
(``in progress'') value to the next even (``quiescent'') one, and a reader brackets its access with
\code{acquireVersion} before and \code{validateVersion} after, retrying if the version moved or is odd.
This is the optimistic, lock-free complement to the page guard: the guard is the coarse ``do not evict
me'' pin, and the slot version is the fine ``the bytes did not change under me'' check.

\dcap{Frame-slot allocation: size classes, lazy commit, seqlock versioning}
\begin{diagramS}
size classes:   4K  8K  16K  32K  64K  128K  256K        each: 32 GiB virtual reserved, RAM on demand
allocate(n):    reserve budget ; pop a free slot of the class ; version even→odd→even ; MADV_POPULATE_WRITE
release(slot):  version even→odd→even ; push to free stack    (no MADV_DONTNEED — pages stay resident)
reader:         v0 = acquireVersion ; if odd → retry ; read bytes ; if !validateVersion(v0) → retry
budget hit:     pressure listener → cache.evictUnderPressure() on every cache, then clear metadata pages
\end{diagramS}

\begin{figure}[t]
\centering
\begin{tikzpicture}[font=\sffamily\small,>=Latex,
  used/.style={draw=slate!55,fill=slate!20,minimum width=0.5cm,minimum height=0.4cm,inner sep=0},
  free/.style={draw=nodenew!60,fill=nodenew!16,minimum width=0.5cm,minimum height=0.4cm,inner sep=0},
  virt/.style={draw=slate!35,dashed,minimum width=0.5cm,minimum height=0.4cm,inner sep=0}]
  % all seven size classes, 4 KiB (top) down to 256 KiB (bottom)
  \foreach \cls/\y in {4K/6,8K/5,16K/4,32K/3,64K/2,128K/1,256K/0}{
     \node[font=\scriptsize,anchor=east] at (-0.15,0.5*\y) {class \cls};
     \node[used] at (0.45,0.5*\y){}; \node[used] at (0.97,0.5*\y){};
     \node[free] at (1.49,0.5*\y){};
     \foreach \i in {0,...,4}{ \node[virt] at (2.25+0.52*\i,0.5*\y){}; }
  }
  \draw[decorate,decoration={brace,amplitude=4pt},slate] (2.0,3.28) -- (4.6,3.28)
     node[midway,above=3pt,font=\scriptsize,align=center]{32\,GiB \emph{virtual} per class;\\RAM on first write};
  % legend
  \node[used,label={[font=\scriptsize]right:committed (in use)}]  at (6.4,3.0) {};
  \node[free,label={[font=\scriptsize]right:free-stack (recycled)}] at (6.4,2.5) {};
  \node[virt,label={[font=\scriptsize]right:reserved, uncommitted}] at (6.4,2.0) {};
  % the overall max memory: ONE shared physical budget, drawn down on demand
  \node[font=\scriptsize\bfseries,text=slate!80!black,anchor=west] at (6.4,1.25)
     {overall max memory $=$ one \emph{shared} physical budget};
  \draw[draw=slate!60,fill=white] (6.45,0.55) rectangle (10.95,0.95);
  \fill[ember!30] (6.45,0.55) rectangle (8.7,0.95);
  \node[font=\scriptsize,text=ember!82!black] at (7.575,0.75) {$\Sigma$ committed};
  \draw[<->,slate] (6.45,0.33) -- (10.95,0.33) node[midway,below,font=\scriptsize,inner sep=1.5pt]{$=$ max memory};
  \node[font=\scriptsize,anchor=north west,align=left,text width=5.7cm,text=slate!75!black] at (6.4,-0.05)
     {one global counter shared by \emph{all} classes, so the runtime mix is workload-driven: a single class can fill it (e.g.\ all 64\,KiB leaf pages during a bulk load), and the ceiling then evicts rather than failing.};
\end{tikzpicture}
\caption{The off-heap frame allocator. \textbf{All seven} power-of-two size classes (4\,KiB--256\,KiB) each
reserve a large 32\,GiB \emph{virtual} region up front but commit physical RAM only on first write
(\code{MADV\_POPULATE\_WRITE}), so a slot's address is stable for the JVM's life while memory is paid for
only when used; a per-class free-stack recycles released slots syscall-free (no \code{MADV\_DONTNEED}). The
configured \textbf{maximum memory is a single shared physical budget} --- one global counter --- drawn
down on demand: a request resolves to the smallest class that fits it, and because nothing is split
per-class the runtime distribution across classes is workload-driven, so a single class can consume the
entire budget (e.g.\ all 64\,KiB leaf pages during a bulk load). Reaching the ceiling triggers cache
eviction, not an allocation failure.}
\label{fig:allocator}
\end{figure}

A reuse cycle shows the seqlock catching a recycle under a live reader:

\dcap{A reader detecting slot reuse via the version seqlock}
\begin{diagramS}
slot #142 (64K class), version = 2 (even = quiescent), currently holds page P

reader R:   v0 = acquireVersion(142) = 2   even -> proceed ; reading P's bytes ...
            (concurrently) release(142):           version 2 -> 3 (odd, in progress) -> 4 (even)
            (concurrently) allocate(142) for page Q: version 4 -> 5 -> 6 ;  Q's bytes written in
reader R:   validateVersion(142): current 6  !=  v0 = 2   ->  RETRY (R re-fetches through the cache)
\end{diagramS}

\noindent Bumping on \emph{both} allocate and release is what makes the check sound: a reader that starts
mid-transition sees an odd value and retries immediately, and a reader that brackets a full
release\,+\,reallocate sees the version move by four and retries on validate --- so the only way to pass
both checks is for the slot to have been genuinely untouched for the whole read. The slot version is thus
the fine-grained ``the bytes did not change under me'' complement to the page guard's coarse ``do not
evict me'' pin.

When a size class nears its budget, the allocator fires a \textbf{pressure listener} that the buffer
manager wires straight to cache eviction, so back-pressure from the allocator becomes immediate eviction
rather than an allocation failure. A \code{KeyValueLeafPage} owns its slot's \code{MemorySegment} for its
cached lifetime and, on a copy-on-write modification, allocates a \emph{fresh} slot and copies into it
(Ch.~\ref{ch:pds}), so the old revision's frame is never mutated, the allocator-level expression of the
immutability the whole system rests on.

Two smaller cache tiers sit on either side of the shared pool. Each \emph{writer} keeps its uncommitted
state in the Transaction Intent Log (Ch.~\ref{ch:versioning}) plus a small bounded cache of the most
recently touched page containers, all dropped at commit or rollback. And the \emph{query} layer keeps a
plan cache of recently optimised query plans and a statistics catalogue of per-resource histograms and
cardinalities (with a time-to-live), both keyed so that a commit or an index change invalidates exactly
the affected entries (Ch.~\ref{ch:query}).

\section{An annotated superblock}
\label{app:hexdump}

The format and the commit protocol have been described in prose; one annotated example makes the file
header concrete (the full binary format is tabulated in Appendix~\ref{app:format}). Here is the 64-byte
superblock of a \code{sirix.data} file, byte for byte (little-endian multi-byte fields) --- the very bytes
the engine reads first on open, where a mismatch in the magic, the version, or the endianness word is a
hard refusal to proceed.

\dcap{The \texttt{sirix.data} superblock, 64 bytes, annotated}
\begin{diagramS}
offset  bytes (hex)                meaning
 0x00   53 49 52 49 58 44 42 21    magic  "SIRIXDB!"  (8 ASCII bytes)
 0x08   00 00 00 00                layout version = 0  (u32 LE)
 0x0C   00                         role = 0 (DATA;  REVISIONS file would be 1)
 0x0D   00 00 00                   reserved / flags
 0x10   04 03 02 01                endianness check word = 0x01020304  → reader confirms LE
 0x14   00 10 00 00                beacon-slot size = 4096   (u32 LE)
 0x18   00 10 00 00 00 00 00 00    primary-beacon offset = 4096   (u64 LE)
 0x20   00 30 00 00 00 00 00 00    content-start = 12288   (u64 LE)   (DATA_REGION_START)
 0x28   00 … 00  (16 bytes)        reserved (future resource UUID)
 0x38   xx xx xx xx xx xx xx xx    XXH3-64 over bytes [0x00, 0x38)   (self-checksum)

then:  0x1000 PRIMARY beacon slot (4 KiB) · 0x2000 SECONDARY beacon slot (4 KiB) · 0x3000 first page record
\end{diagramS}

Three of these fields are the whole open-time contract. The \textbf{magic} identifies the file; the
\textbf{role} byte stops a \code{sirix.data} file from ever being read as a \code{sirix.revisions} file or
the reverse; and the \textbf{endianness check word} \code{0x01020304} is stored as a raw integer, so a
reader that sees \code{04 03 02 01} knows it is little-endian and a reader on a big-endian host would see
\code{01 02 03 04} and could correct, the format is self-describing rather than host-assumed. The trailing
\textbf{XXH3} over the first 56 bytes catches a torn or bit-rotted superblock before any of those fields
is trusted. Everything after offset \code{0x3000} is the append-only data region whose growth the file-growth
trace of Appendix~\ref{app:format} walks.

\section{Failure modes and their handling}
\label{app:failures}

Those integrity fields are the engine's first line against corruption, and the chapter has developed the
rest in place; gathering them into one table shows the fault-tolerance story whole. Each row names a way
things can go wrong, how the system \emph{detects} it, and what it \emph{does}, with a cross-reference to
where it is developed.

{\small
\begin{xltabular}{\linewidth}{@{}>{\RaggedRight\arraybackslash}p{0.235\textwidth} >{\RaggedRight\arraybackslash}p{0.25\textwidth} >{\RaggedRight\arraybackslash}X@{}}
\toprule
\textbf{Failure mode} & \textbf{Detection} & \textbf{Handling}\\
\midrule
\endhead
Crash mid-beacon write (one torn) & per-slot \code{[len][payload][XXH3]} validation on open & fall back to the intact beacon copy; it names a durable tail (Prop.~\ref{prop:torn})\\
Crash before a beacon barrier & both beacons still name the old revision & the commit was never acknowledged; reopen on the previous revision\\
Interrupted commit (\code{.commit} marker present) & marker found on open & \code{truncateTo} the last good revision before accepting a writer\\
Interrupted \emph{first} commit (non-empty file, zero header) & a conservative on-disk byte-pattern proof & re-initialise empty \emph{only} if the bytes prove nothing past the bootstrap was acknowledged; else rethrow (\S\ref{sec:recovery})\\
Torn or bit-rotted page body & XXH3-64 over the compressed bytes mismatches & caught before the page is trusted; reload / surface as corruption (\S\ref{sec:recovery})\\
Frame recycled under a live reader & guard's \code{versionAtFix} $\neq$ version on release & throw \code{FrameReusedException} $\Rightarrow$ cache miss and clean reload (\S\ref{sec:reclaim})\\
Eviction races a live snapshot & sweeper checks \code{guardCount > 0} & a guarded page is skipped, never reclaimed; orphaned-while-guarded pages tear down at the last release\\
Off-heap budget exhausted & allocator reaches its budget & pressure listener fires synchronous cache eviction rather than failing the allocation (\S\ref{sec:allocator})\\
Streaming-shred producer outruns consumer & bounded channel full & pause the parser \emph{and} the HTTP request; resume at the low-water mark (Ch.~\ref{ch:rest})\\
Empty-channel lost wakeup (the alpha15 bug) & consumer about to suspend on an empty channel & resume the producer \emph{before} suspending, so a short tail cannot wedge it asleep\\
Second writer on a resource & \code{Semaphore} \code{tryAcquire(5\,s)} times out & fail fast with \code{SirixUsageException} rather than block indefinitely (Ch.~\ref{ch:rest})\\
Stale plan after an index create/drop & global index-schema version bumped & the plan cache key \code{queryText@version} misses, forcing a recompile (\S\ref{sec:optpipeline})\\
Optimiser path-kind mismatch (the alpha14 bug) & kind-agnostic \code{pathSummary.match} & bail to normal evaluation instead of a destructive empty-sequence rewrite (\S\ref{sec:diff})\\
Reboot before Keycloak is up & SirixDB launch gated on a readiness poll & wait for Keycloak before starting the backend, avoiding a refused-connection crash loop (\S\ref{app:deploy})\\
\bottomrule
\end{xltabular}}

\noindent A pattern runs through the column of handlings: \textbf{detect, then degrade to a known-good
state}, never proceed on unverified bytes. A checksum or a version counter catches the corruption, and
the response is always to fall back, the intact beacon, the previous revision, a clean reload, a
fail-fast exception, rather than to guess. That discipline is what lets the report claim durability
under \code{SIGKILL} and power-loss fault injection rather than merely assert it.

\chapter{HOT: The Height-Optimised Trie}
\label{ch:hot}

\noindent\textit{(\pp{docs/HOT_*.md}; \pp{io/sirix/page/HOTLeafPage.java},
\pp{HOTIndirectPage.java}; \pp{io/sirix/access/trx/page/HOTTrieReader.java},
\pp{HOTTrieWriter.java}, \pp{HOTRangeCursor.java}.)}

\section{Motivation}

With the storage layer of Chapter~\ref{ch:ondisk} settled, HOT is the trie that navigates
within and across those pages and backs the secondary indexes. It is a cache-conscious,
height-optimised trie used for index and intra-page navigation, replacing heavyweight
B${}^{+}$-tree-style indirection over flat arrays. Instead of one
tree level per bit, a HOT indirect node captures \emph{k} \textbf{discriminative bits}
(0--32) at chosen absolute bit positions and routes among 2--32 children by their values at
exactly those positions, packing multiple bits per level so a typical node fits in one or
two cache lines. It is the structure behind the DeweyID index
(Ch.~\ref{ch:pathsummary}) and the projection index (Ch.~\ref{ch:vector}), and is designed
generically for any prefix-free, bit-discriminable key serialisation.

\section{The structure: discriminative bits and the sparse path}

For each child slot, the stored partial key is the parallel-bit-extract of that child's
first key under the node's mask:

\begin{diagram}
stored[c] = PEXT(firstKey[c], mask)        // the "sparse path": only on-path bits kept
\end{diagram}

Worked example. Let \code{firstKey = 0b1101\_1010} and \code{mask = 0b1100\_0001}
(bits~0,~6,~7). \code{PEXT(0b1101\_1010, 0b1100\_0001)} gathers bits \{7,\,6,\,0\} $=$
\code{1,1,0}, giving stored partial \code{0b110}. Figure~\ref{fig:hot} shows how that sparse path routes a
lookup.

\begin{figure}[t]
\centering
\begin{tikzpicture}[font=\sffamily\small,>=Latex,
  nd/.style={draw=slate!60,fill=layerbg,rounded corners=2pt,align=center,inner sep=4pt},
  lf/.style={draw=nodenew,fill=nodenew!12,rounded corners=2pt,align=center,inner sep=3pt,minimum width=1.4cm},
  sl/.style={draw=ember!60,fill=ember!10,font=\scriptsize,inner sep=2pt,minimum width=0.85cm},
  e/.style={->,slate}]
  \node[nd] (root) at (3.5,2.2) {\textbf{HOTIndirectPage}\quad mask = discriminative bits \{7, 6\}};
  \node[sl] (c0) at (1.5,0.9) {00};
  \node[sl] (c1) at (3.5,0.9) {01};
  \node[sl] (c2) at (5.5,0.9) {11};
  \draw[e] (root) -- (c0); \draw[e] (root) -- (c1); \draw[e] (root) -- (c2);
  \node[lf] (l0) at (1.5,-0.6) {leaf $00\ldots$};
  \node[lf] (l1) at (3.5,-0.6) {leaf $01\ldots$};
  \node[lf] (l2) at (5.5,-0.6) {leaf $11\ldots$};
  \draw[e] (c0) -- (l0); \draw[e] (c1) -- (l1); \draw[e] (c2) -- (l2);
  \node[nd,fill=white,draw=ember!70,align=left] (q) at (9.9,1.0)
     {\textbf{lookup} K $=$ \code{0b01\ldots}\\[2pt]
      PEXT(K, mask) $= 01$\\[2pt]
      $\to$ slot \code{01}};
  \draw[->,ember!70,dashed] (q.west) .. controls (7.4,1.0) and (5.2,1.5) .. (c1.east);
\end{tikzpicture}
\caption{HOT routing by discriminative bits. An indirect node stores only the bit positions where its
children's keys diverge (its \emph{mask}); each child slot is labelled by the parallel-bit-extract of its
first key under that mask (its \emph{sparse path}, here 2 dense bits). A lookup extracts the same bits from
the query key and descends to the slot whose stored partial matches --- \code{PEXT(K, mask) = 01} routes to
the middle child. Only the discriminative bits are ever compared, which is what keeps a HOT node compact
and its routing branch-light.}
\label{fig:hot}
\end{figure}

Splitting a node to add a new discriminative bit $\beta$ is a $\textsf{compress} \circ
\textsf{expand}$ round-trip:

\dcap{Algorithm 5.1 --- Add discriminative bit $\beta$ to a node}
\begin{diagram}
  for each slot c:
      full  ← PDEP/Long.expand(stored[c], oldMask)   // restore old-width bits
  newMask ← oldMask | (1 << β)
  for each slot c:
      stored[c] ← PEXT/Long.compress(full_c possibly with β set, newMask)
\end{diagram}

\textbf{Invariant I-PDEP-PEXT} (\pp{docs/HOT_INVARIANTS_CATALOG.md}).
$\textsf{compress} \circ \textsf{expand}$ is the identity on the old partial-key space:
the round-trip is bijective, so routing is preserved across the split.

A concrete split makes that bijection visible. Take the same \code{oldMask = 0b1100\_0001} with three
children, and insert a key whose first disagreement with the \code{C} subtree falls on bit~3, forcing a
new discriminative bit $\beta=3$:

\dcap{A worked split: widening \texttt{oldMask}=0b1100\_0001 to admit bit $\beta$=3}
\begin{diagramS}
before:  oldMask = 1100_0001  (set bits 7,6,0)         3 children, dense 3-bit partials
  slot A  stored 000   --expand-->  full 0000_0000      (PDEP places dense bits into mask slots)
  slot B  stored 010   --expand-->  full 0100_0000      (dense bit 1  -> mask bit 6)
  slot C  stored 110   --expand-->  full 1100_0000      (dense bits 1,2 -> mask bits 6,7)

  newMask = oldMask | (1<<3) = 1100_1001   (set bits now 0,3,6,7)

after:   re-extract every full key under newMask (PEXT) -> a wider dense partial
  slot A  full 0000_0000  --PEXT-->  0000      old children keep routing,
  slot B  full 0100_0000  --PEXT-->  0100      each gains a 0 spliced in at beta's rank (dense pos 1)
  slot C  full 1100_0000  --PEXT-->  1100
  new     full 1100_1000  --PEXT-->  1110      the inserted key takes a fresh slot

invariant:  compress(expand(stored,oldMask),newMask) = stored with one 0-bit inserted at beta's rank
            so 010 -> 0100, 110 -> 1100 : a uniform relabelling, never a reordering
\end{diagramS}

\noindent Because $\beta$ is not in \code{oldMask}, the expanded full key carries a $0$ at position~3, and
re-extracting under \code{newMask} simply splices that $0$ into every old partial at $\beta$'s rank. The
relabelling is uniform across all prior children, so the submask test of Algorithm~5.2 still routes each
existing key to its original slot; only the newcomer, which alone carries a $1$ at bit~3, lands in the
fresh slot. That is what ``routing is preserved'' means operationally: a split rewrites the labels but not
the routing function.

Three node layouts encode the mask and partials in the page's \code{MemorySegment}:

\dcap{HOT indirect-node layouts}
\begin{diagramS}
BiNode (2 children):   [1B initialBytePos][8B bitMask(single bit)]                + child refs
SingleMask (3–32):     [1B initialBytePos][8B bitMask][N×4B partial][N×8B childRef]   (all bits in one 8B window)
MultiMask (bits spread):[extractionPositions[]][extractionMasks[]][partials in 2×32B slots] + child refs
\end{diagramS}

The read path descends to the highest-popcount slot whose stored partial is a subset of the
query's dense PEXT, then binary-searches the terminal leaf:

\dcap{Algorithm 5.2 --- \texttt{HOTTrieReader.getRecord(node, key)}}
\begin{diagram}
  if node is leaf: binary-search entries by key; return value-slice or ⊥
  densePK ← PEXT(key, node.mask)
  c ← argmax_{slots s : stored[s] is a submask of densePK} popcount(stored[s])
  return getRecord(child[c], key)
\end{diagram}

In practice these layouts are richer than the schematic above. A \code{HOTLeafPage} is
\textbf{prefix-compressed}: all keys in a leaf share a \code{commonPrefix}, and each of up to
$\code{MAX\_ENTRIES}=512$ entries stores only \code{[u16 suffixLen][suffix][u16 valueLen][value]};
unsigned-short lengths bound a key or value at 64\,KB, and an 8-byte big-endian word view lets
\code{Long.compareUnsigned} do lexicographic comparison. A second, \emph{in-leaf} PEXT index
accelerates leaves of $\le 32$ entries: differing bits of adjacent suffixes are folded into one dense
mask and the search becomes an AVX2 broadcast-compare of all entries in a single vector op, falling
back to branchless binary search. A \code{HOTIndirectPage} carries, for the single-mask case, an
\code{initialBytePos} and an 8-byte \code{bitMask} (the discriminative bits live in that one window)
plus a cached most-significant-bit index; the multi-mask case spreads bits across
\code{extractionPositions[]}/\allowbreak\code{extractionMasks[]}. Children are addressed by an \code{int[]} of
dense partial keys, a SIMD \code{SparsePartialKeys} search structure (byte/short/int width chosen by
bit count), and a \code{PageReference[]}. A ``BiNode'' is not a separate class; it is an ephemeral
two-child single-mask node produced by a split and integrated as a compound node, exactly as in the
Binna reference.

\section{Insert, split, and copy-on-write}

\begin{figure}[t]\centering
\begin{tikzpicture}[font=\sffamily\footnotesize,>=Latex,
  leaf/.style={draw=slate!55,fill=layerbg,rounded corners=1pt,minimum height=0.85cm,align=center,inner sep=3pt},
  full/.style={draw=nodenew,fill=nodenew!12,thick,rounded corners=1pt,minimum height=0.85cm,align=center,inner sep=3pt}]
  \node[full] (b) {full leaf\\$k_1\,k_2\,k_3\,k_4\,k_5$};
  \node[below=1mm of b,font=\scriptsize\itshape] {overflow};
  \node[draw=slate!70,fill=white,circle,inner sep=1.5pt,right=3.0cm of b] (par) {$\beta$};
  \draw[->,slate,thick] (b.east) -- (par.west) node[midway,above,font=\scriptsize]{split on $\beta$};
  \node[leaf,below left=9mm and 0mm of par] (l) {$k_1\,k_2$\\$\beta{=}0$};
  \node[leaf,below right=9mm and 0mm of par] (r) {$k_3\,k_4\,k_5$\\$\beta{=}1$};
  \draw[->,slate] (par) -- (l); \draw[->,slate] (par) -- (r);
\end{tikzpicture}
\caption{A leaf split. When a leaf overflows, the most-significant differing bit $\beta$ is
\emph{discovered from the keys}; entries partition by their value at $\beta$ into two leaves under a
new (ephemeral ``BiNode'') indirect node whose mask is exactly $\beta$, the partial keys recomputed
by the \code{compress}$\,\circ\,$\code{expand} round-trip of Algorithm~5.1.}
\label{fig:hotsplit}
\end{figure}

Insertion descends to a leaf; on overflow the leaf performs an \emph{atomic most-significant-bit
split} (Figure~\ref{fig:hotsplit}). The discriminative bit $\beta$ is \textbf{discovered from the
data, not chosen a priori}: \code{findMsdbWithNewKey} computes the most significant bit at which the
new key and its sorted neighbours differ, the split point is the first existing key with $\beta=1$,
and the new key joins whichever half its $\beta$ value selects. Encoding $\beta$ into the parent
(\code{addEntryWithPDep}) is the algorithmic core: if $\beta$'s byte falls outside the current 8-byte
window the node is \textbf{upgraded single-mask\,\textrightarrow\,multi-mask}; every existing partial
key is repositioned by \code{Long.expand} (the PDEP half of the \code{compress}$\,\circ\,$%
\code{expand} round-trip); and, crucially, the insert verifies that $\beta$ is \emph{constant
across every non-split sibling's subtree} before committing, falling back to a parent split if not.
This subtree-constancy check is invariant I5 (\S\ref{sec:hotinv}) made operational, and its
difficulty under multi-key leaves is the entire subject of the campaign below. Copy-on-write
preserves a node's \code{pageKey} identity across the deep-copy of its child references (invariant
I-CoW), and deletes are \textbf{tombstones}, not physical removals: a physical delete would let a
fragment merge (Ch.~\ref{ch:versioning}) resurrect the key from an older fragment.

\dcap{A worked HOT split: the discriminative bit is discovered from the data}
\begin{diagramS}
a leaf holds these 8-bit keys (binary), sorted:
   0100_0001   0100_0011   0100_0111         (all share the prefix bits 7,6 = 0,1)
insert 0110_0000:
   it agrees on bits 7,6 but diverges at bit 5  →  discriminative bit β = 5
   split the leaf on β:
     β = 0  →  { 0100_0001, 0100_0011, 0100_0111 }   (the old keys)
     β = 1  →  { 0110_0000 }                          (the new key)
   each child stores PEXT(firstKey, mask) as its partial; routing now tests bit 5
\end{diagramS}

\section{Lookup, range, and the reader walk-up}

Point lookup routes per level in \code{findChildSpanNode}: compute the dense partial key
\code{Long.compress(keyWord, bitMask)}, ask the SIMD \code{SparsePartialKeys} structure for the
candidate slots, and take the most specific (highest-popcount) subset match, but \emph{prefer} an
exact match when one exists, a deliberate deviation from pure Binna subset routing that avoids
delivering a key into a $\beta$-violating subtree. Range scans use \code{HOTRangeCursor}, which is
sibling-pointer-free (sibling pointers would be copy-on-write-hostile) and instead walks the reader's
parent stack; its \code{lowerBound} is a multi-phase \emph{walk-up}: descend to a candidate leaf and,
if the insertion point straddles a node boundary, climb while the discriminative bit is less
significant than the path's captured most-significant bits, then re-descend leftmost. A
fire-and-forget virtual-thread prefetch of the next sibling leaf overlaps I/O with iteration. This
walk-up is also what keeps reads correct under the one surviving invariant violation (\S\ref{sec:hotcampaign}).

\dcap{\texttt{HOTRangeCursor.lowerBound}: finding the successor without sibling pointers}
\begin{diagramS}
lowerBound(K):
  1 descend root → candidate leaf L by the point route (dense PEXT = Long.compress(keyWord, bitMask) per level)
  2 binary-search L for the first entry ≥ K
  3 if the insertion point is at L's right edge (it straddles a node boundary):
        climb the reader's PARENT STACK while the level's disc-bit < the path's captured MSBs
        then re-descend LEFTMOST into the next subtree → the true successor leaf
  4 iterate; a virtual-thread prefetch pulls the next sibling leaf in the background

no sibling pointers (CoW-hostile); the parent stack + disc-bit comparison reconstructs document order,
and is exactly the path that recovers a key even past the one surviving I8 straddle (§5.6).
\end{diagramS}

A keyed trace shows the walk-up finding a successor across a leaf boundary:

\dcap{A worked lowerBound straddle: the successor lies in the next subtree}
\begin{diagramS}
trie:  root mask {bit7, bit4}  ->  4 leaves (indexed by bits 7,4):
   L1 (0,0): {0000_0001, 0000_0010}      L2 (0,1): {0001_0000, 0001_0001}
   L3 (1,0): {1000_0000, 1000_0001}      L4 (1,1): {1001_0000}

lowerBound(K = 0000_1111):           (all keys sorted: 1,2,16,17,128,129,144 ; first >= 15 is 16)
  1 point route:  K has bit7=0, bit4=0 -> dense 00 -> descend to L1
  2 binary-search L1 {1,2} for first >= 15 -> none ; insertion point is at L1's RIGHT EDGE (straddle)
  3 climb the parent stack to root ; re-descend LEFTMOST into the next child after (0,0): the (0,1) child -> L2
  4 L2 leftmost = 0001_0000 = 16  ->  the successor ; iterate from here
result:  lowerBound = 16, found with NO sibling pointer — the parent stack alone reconstructed document order
\end{diagramS}

\noindent The straddle at step~2 is the whole reason a range cursor needs more than a leaf-local binary
search: the next key in document order lives in a \emph{different} leaf, and with sibling pointers banned
(they would have to be copy-on-write-rewritten on every change) the only way back to it is up the parent
stack and down the next subtree. That climb-then-redescend is $O(\text{height})$ per leaf boundary, paid
only at the $\sim$1-in-512 transitions between leaves, not per key.

\section{Invariants and the multi-key-leaf deviation}
\label{sec:hotinv}

The implementation tracks and validates sixteen invariants in full (soft ones are demoted to warnings
because the multi-key-leaf design makes them non-essential for correctness):

\begin{table}[h]
\footnotesize\setlength{\tabcolsep}{6pt}
\renewcommand{\arraystretch}{1.15}
\begin{tabularx}{\linewidth}{@{}l X@{}}
\toprule
\textbf{Inv.} & \textbf{Statement}\\
\midrule
I1 & leaf keys are unique within a leaf\\
I2 & leaf entries are lexicographically sorted (unsigned)\\
I3 & partial keys are unique \emph{(soft)}\\
I4 & the leftmost slot's partial key is 0\\
I5 ($\beta$-constancy) & every key in a slot's subtree has the slot's partial bits set\\
I6 & \code{descend(K)} lands at K's actual leaf \emph{(soft)}\\
I7 & partial keys are sorted \emph{(soft)}\\
I8 & children are sorted by first key \emph{(hard --- the one surviving violation)}\\
I9 & height $\leq 64$\\
I10 & fan-out $\in [2,32]$\\
I11 & trie condition: a parent's most-significant disc-bit precedes its child's disc-bit range\\
I12 & disc-bit monotone: no absolute bit is captured twice on a root-to-leaf path\\
I-Binna & sparse path: $\text{partial}[c] \subseteq \text{PEXT}(\text{firstKey}[c],\text{mask})$\\
I-PDEP & \code{compress}$\,\circ\,$\code{expand} is the identity on the old partial-key space\\
I-CoW & page-key identity is preserved across copy-on-write\\
I-MR & multi-revision read isolation (a revision-$N$ read sees only revision-$N$ state)\\
\bottomrule
\end{tabularx}
\end{table}

SirixDB deviates from the Binna reference by allowing \textbf{up to $\sim$512 keys per leaf},
which makes $\beta$-constancy (I5) non-trivial: a leaf can legitimately contain keys that
differ at an ancestor's discriminative bit (in a single-key leaf, every bit is trivially
constant). This routing subtlety is the source of essentially all the difficulty below.

\section{The hardening campaign}
\label{sec:hotcampaign}

A documented campaign (\pp{docs/HOT_CAMPAIGN_RESULTS.md}, \pp{HOT_OPERATIONS_INVARIANTS_MATRIX.md})
built a validator over a $35 \times 16$ \emph{operations}$\times$\emph{invariants} matrix and drove a
50\,000-insert reproducer's violations \textbf{from 127 to 1} ($99.2\,\%$), with \emph{zero} violations
on the 100\,K / 1\,M / 10\,M production runs. The instructive part is how stubbornly the last violation
resisted: closing one matrix cell reliably opened another. Tightening the I8 (children-sorted) gate on
an intermediate node took 1\,\textrightarrow\,5\,998; an ``Option-B'' reroute through a sibling subtree
took 1\,\textrightarrow\,28\,706; a force-split at the offending ancestor took
1\,\textrightarrow\,15\,237; an I11 gate took 5\,\textrightarrow\,2\,180. Tellingly, the residual
violations emerge from \emph{happy-path} writer operations: across 50\,000 inserts the recovery
counters fired \textbf{zero} times, so the bug is in the default path, not an exotic fallback. Routing is
structurally ambiguous (81\,\% of inserts route ambiguously; only 1.5\,\% actually split).

The single surviving I8 violation is \textbf{not a production-correctness defect}. Its nature is an
\emph{off-path straddle}: a multi-key leaf stores $\text{partial} = \text{PEXT}(\text{firstKey},\text{mask})$,
so a later-inserted \emph{smaller} key whose differing bit $\beta$ is 0 where the firstKey's is 1 makes
routing-order disagree with lex-order at one boundary: a bit constant across an older subtree but
variable across a newer one, which the stored-partial encoding cannot capture. Reads stay correct
regardless, because the reader's multi-phase walk-up and a lexicographic-descent fallback recover every
key (measured full-scan misses: zero). Eliminating it cleanly is a scoped routing-encoding rework
(\pp{docs/HOT_ROUTING_ENCODING_REWRITE.md}) with three candidates, each carrying a known cost:
\textbf{AND-encoding} the stored partial over the whole subtree (degenerates toward 0 as inserts
diversify), \textbf{proactive disc-bit extension} at insert time (the mask saturates at $\sim$32 bits,
forcing splits), or \textbf{recompute-on-CoW} from the subtree (correct but $O(\text{subtree})$ per
insert).

The DeweyID-HOT secondary index is backed by formal order-preservation and lookup-correctness
proofs (Theorem~\ref{thm:dewey} and \pp{docs/DEWEYID_HOT_INDEX_FORMAL_PROOF.md}), including
tier-boundary and deep-hierarchy edge cases.

\section{HOT as the index backend}

Beyond intra-page navigation, HOT is the \emph{default} backend for the secondary indexes of
Ch.~\ref{ch:pathsummary} (replacing the legacy red-black tree). A generic \code{HOTIndexWriter<K>}
serves CAS and NAME through a pluggable key serialiser, while \code{HOTLongIndexWriter} is a
primitive-\code{long} specialisation for PATH (no boxing). Index keys are made order-preserving by a
sign-flipped big-endian \code{long} encoding, so a HOT range scan yields results in key order; the
values (the \code{NodeReferences} posting lists) are stored as \textbf{chunked Roaring bitmaps} keyed
by a key prefix plus a chunk index, so a single-node update rewrites only one small chunk rather than
the whole posting list.

\begin{honesty}
\small\textbf{Caveat.} HOT is an area of active, documented engineering; the ``1
remaining marginal violation, production reads correct'' status is the project's own
characterisation, reproduced here as such.
\end{honesty}

%================================================================
\chapter{The Path Summary and Secondary Indexes}
\label{ch:pathsummary}

\section{The path summary}
\label{sec:pathsummary}

With navigation fast (Chapter~\ref{ch:hot}), the structures in this chapter are what let a query
skip a full scan: a compact catalogue of the resource's structure, and inverted indexes over its
values. The \textbf{path summary} (\pp{io/sirix/index/path/summary/}) is a compact tree of every
\emph{distinct} path class in a resource: all nodes sharing a path collapse to one \code{PathNode},
regardless of how many document nodes instantiate it. Figure~\ref{fig:pathsummary} carries one JSON
document through shredding to the summary it induces. \emph{Bare objects and leaf values never appear
in the summary}: anonymous \code{OBJECT} containers (the root and the two array elements) get no path
node, and the scalar values are not stored. What remains is one \code{PathNode} per distinct path,
each carrying a \emph{reference count} of how many document nodes traverse it. Here the two
\code{title} fields and the two \code{price} fields collapse to a single class apiece (count~2) beneath
the array's \code{[]} step (count~1), with the \code{price} class additionally accumulating per-path
\code{PathStats}.

\begin{figure}[t]
\centering
\resizebox{\textwidth}{!}{%
\begin{tikzpicture}[font=\sffamily\scriptsize,>=Latex,
  o/.style={draw=slate!55,fill=slate!8,rounded corners=1.5pt,align=center,inner sep=2.5pt,minimum height=4.4mm},
  n/.style={draw=slate!60,fill=layerbg,rounded corners=1.5pt,align=center,inner sep=2.5pt,minimum height=4.4mm},
  nv/.style={draw=nodenew!70,fill=nodenew!12,rounded corners=1.5pt,align=center,inner sep=2.5pt,minimum height=4.4mm},
  p/.style={draw=slate!65,fill=white,rounded corners=2pt,align=center,inner sep=2.5pt,minimum height=4.4mm},
  e/.style={draw=slate!50},
  m/.style={draw=ember!75!black,dashed,->,thin},
  big/.style={draw=ember!70!black,-{Latex[length=2.8mm]},line width=1.1pt}]
  % ===== column headers =====
  \node[font=\bfseries\color{ember}] at (0.0,1.55) {JSON source};
  \node[font=\bfseries\color{ember}] at (3.4,1.55) {SirixDB encoding};
  \node[font=\bfseries\color{ember}] at (9.7,1.55) {Path summary (one per path)};
  % ===== stage 1: JSON source =====
  \node[font=\ttfamily\scriptsize,align=left,anchor=north west] (json) at (-1.6,0.55)
    {\{"books":[\\\hspace{0.9ex}\{"title":"A","price":12\},\\\hspace{0.9ex}\{"title":"B","price":8\}]\}};
  % ===== stage 2: SirixDB node encoding (bare OBJECTs + fused named records) =====
  \node[o]  (eroot)  at (3.4,0)     {OBJECT};
  \node[n]  (ebooks) at (3.4,-1.5)  {books \textperiodcentered\ ARRAY};
  \node[o]  (eo0)    at (2.1,-3.0)  {OBJECT};
  \node[o]  (eo1)    at (4.7,-3.0)  {OBJECT};
  \node[n]  (et0)    at (1.4,-4.5)  {title "A"};
  \node[nv] (ep0)    at (2.7,-4.5)  {price 12};
  \node[n]  (et1)    at (4.1,-4.5)  {title "B"};
  \node[nv] (ep1)    at (5.4,-4.5)  {price 8};
  \draw[e](eroot)--(ebooks);\draw[e](ebooks)--(eo0);\draw[e](ebooks)--(eo1);
  \draw[e](eo0)--(et0);\draw[e](eo0)--(ep0);\draw[e](eo1)--(et1);\draw[e](eo1)--(ep1);
  % ===== stage 3: path summary (no objects, no values) =====
  \node[p] (proot)  at (9.7,0)     {/};
  \node[p] (pbooks) at (9.7,-1.5)  {books \textperiodcentered\ ref 1};
  \node[p] (parr)   at (9.7,-3.0)  {{[]} \textperiodcentered\ ref 1};
  \node[p] (ptitle) at (8.5,-4.5)  {title \textperiodcentered\ ref 2};
  \node[p] (pprice) at (10.9,-4.5) {price \textperiodcentered\ ref 2\\{\tiny +\,PathStats sketch}};
  \draw[e](proot)--(pbooks);\draw[e](pbooks)--(parr);\draw[e](parr)--(ptitle);\draw[e](parr)--(pprice);
  % ===== stage arrows: shred (JSON->encoding); the fused named array induces TWO path nodes =====
  \draw[big] (json.east) -- node[above=1pt,font=\itshape\scriptsize,color=ember!75!black]{shred} (eroot.west);
  \draw[big] (ebooks.east) -- node[above=1pt,font=\itshape\scriptsize,color=ember!75!black]{induces name} (pbooks.west);
  \draw[big] (ebooks.east) -- node[pos=0.58,below=1pt,font=\itshape\scriptsize,color=ember!75!black]{+\,{[]} step} (parr.west);
  % ===== detail annotations (routed clear of every label) =====
  \node[align=center,font=\tiny\itshape,color=ember!75!black] (bnote) at (-0.1,-2.95) {bare object:\\no path class};
  \draw[m] (bnote.east) -- (eo0.west);
  \draw[m,rounded corners=3pt] (ep1.south) -- ++(0,-0.8) -| (pprice.south);
  \node[font=\tiny\itshape,color=ember!75!black] at (7.8,-5.9) {leaf values fold into PathStats (not stored)};
\end{tikzpicture}%
}
\caption{From JSON to the path summary. \textbf{Shredding} maps each object field to one fused
\code{OBJECT\_NAMED\_*} record holding the key \emph{and} its value inline; array elements become bare
\code{OBJECT} nodes. The encoding then \textbf{induces} the path summary, one \code{PathNode} per
distinct path. A fused \code{OBJECT\_NAMED\_ARRAY} maps to \emph{two} path nodes: the field
\code{name} and the array's \code{[]} step (both reference count~1). Bare objects (the root and the two
array elements) get \emph{no} path node; their fields attach to that \code{[]} step. The two
\code{title} fields and the two \code{price} fields each
collapse to a single class (reference count~2) beneath the one \code{[]} step (count~1). \emph{Leaf
values are not stored in the summary}: the \code{price} class instead accumulates per-path
\code{PathStats} (count, sum, min/max, a HyperLogLog cardinality sketch, and a Roaring bitmap of the
leaf pages holding values).}
\label{fig:pathsummary}
\end{figure}

\code{PathNode} (\pp{PathNode.java}) is a single inlined heap object holding identity, full
struct-node state, naming keys, kind, a \code{references} count of document nodes using the
path, level, and optional \code{PathStats}. It replaces the legacy three-delegate pointer
chase with one load. The summary is maintained \textbf{incrementally} during writes:
\code{PathSummaryWriter.getPathNodeKey} resolves or creates the child path node in $O(1)$ via
a primitive \code{childLookupCache}, incrementing or initialising its \code{references};
deletion decrements, removing the path subtree only when the last user departs.

The construction is clearest stepped through. Shredding the document of Figure~\ref{fig:pathsummary}
drives a sequence of \code{getPathNodeKey} calls, each either creating a class or bumping an existing one's
reference count:

\dcap{Building the path summary as the document shreds (ref = reference count)}
\begin{diagramS}
event (depth-first shred)            childLookupCache      summary after
  books . ARRAY                       miss -> create        /books(1)
  [] step                             miss -> create        /books(1) /books/[](1)
  {..} elem 0  (bare OBJECT)          -- no path node --     (unchanged)
  title "A"                           miss -> create         .../[]/title(1)
  price 12                            miss -> create         .../[]/price(1)
  {..} elem 1  (bare OBJECT)          -- no path node --     (unchanged)
  title "B"                           HIT  -> ref++           .../[]/title(2)     <- no new node
  price 8                             HIT  -> ref++           .../[]/price(2)     <- no new node

end:  /books(1)  /books/[](1)  /books/[]/title(2)  /books/[]/price(2)      (= the figure)
\end{diagramS}

\noindent The second \code{title} and \code{price} create nothing: \code{childLookupCache} resolves the
existing \code{PathNode} in $O(1)$ and increments its count. The summary therefore grows only with the
number of \emph{distinct} paths, and each shred event is $O(1)$ --- which is what lets a 100\,M-row load
carry a handful-of-paths summary, and why those reference counts are the exact base cardinalities the
optimiser starts its propagation from (\S\ref{sec:joinorder}).

Per-path \textbf{statistics} (\code{PathStats}: count, sum, min/max, null count, a HyperLogLog
cardinality sketch, and a Roaring bitmap of the leaf pages holding values) are accumulated
cheaply. Value observations are \textbf{buffered} in a \code{pendingStats} map during the
transaction and applied in a single copy-on-write pass per path at commit
(\code{flushPendingStats}):

\begin{property}[Deferred-stats cost]\label{prop:defstats}
Maintaining per-path statistics over a transaction touching $V$ values across $P$ distinct
paths costs $O(P)$ copy-on-write operations, not $O(V)$.
\end{property}
This is decisive for analytical loads: 100~M rows over a few hundred paths incur a few hundred
CoW updates at commit, not 100~M.

\dcap{Deferred PathStats: 100M value writes, one CoW update}
\begin{diagramS}
a load writes 100 000 000 price values, all on the single path /store/book/[]/price   (P = 1 path)

during the transaction — each value folds into pendingStats[/…/price] in memory (no PathNode copy):
    count += 1 ;  sum += v ;  min/max ← v ;  HLL.add(hash(v)) ;  pageBitmap.set(leafPageOf(v))

at commit (flushPendingStats) — ONE copy-on-write update to the price PathNode applies the merged stats:
    count = 1e8,  sum,  min,  max,  HLL distinct-estimate,  Roaring bitmap of the value leaf pages

cost: O(P) = 1 CoW update, not O(V) = 1e8.  The optimiser later reads card and distinct straight off it.
\end{diagramS}

\section{The index catalogue and its lifecycle}
\label{sec:catalogue}

SirixDB carries a single, unified secondary-index subsystem (\pp{io/sirix/index/}). One
byte-tagged enum, \code{IndexType}, spans both the \emph{internal} trees that the engine
always maintains and the \emph{user-declarable} indexes:

\dcap{The \code{IndexType} catalogue --- internal trees and user indexes}
\begin{diagram}
internal (always present)                user-declarable
  REVISIONS(0)   revision → offset          PATH(5)                 path class → nodes
  DOCUMENT(1)    the primary node trie       CAS(6)                  (value,type,path) → nodes
  CHANGED_NODES(2)                           NAME(7)                 QName → nodes
  RECORD_TO_REVISIONS(3)  per-node history   DEWEYID_TO_RECORDID(8)  nodeKey ↔ DeweyID
  PATH_SUMMARY(4)                            VECTOR(9)               HNSW kNN over embeddings
                                             PROJECTION(10)          columnar covering index
\end{diagram}

A user index is declared by an \code{IndexDef} (\pp{IndexDef.java}, itself \code{Materializable}
to/from an XML \code{<index>} fragment). Its knobs include the type, the database flavour
(XML/JSON), an id, a set of indexed \code{Path<QNm>}s, \code{included}/\allowbreak\code{excluded} name-filter
sets, the CAS \code{unique} flag and typed \code{contentType}, and, for the newer index kinds,
vector parameters (dimension, distance, the HNSW $m$/\code{efConstruction}/\allowbreak\code{efSearch})
and the projection field list with aligned field types. Factory helpers live in \code{IndexDefs}
(\code{createCASIdxDef}, \code{createPathIdxDef}, \code{createFilteredNameIdxDef},
\code{createVectorIdxDef}, \code{createProjectionIdxDef}). The full set is a per-resource
\code{Indexes} catalogue (a copy-on-write set with a \code{dirty} flag).

\textbf{Persistence is one XML catalogue per revision}: at commit the definitions are serialised
to \pp{{resource}/indexes/{revision}.xml}, and on open the controller searches \emph{backward}
from the requested revision for the most recent \code{{n}.xml}: the same max-numbered-file
search whose naïve per-commit \code{Files.exists} walk was the $O(R^2)$ \code{access(2)} pathology
of \S\ref{sec:evolution}. Each index \emph{tree} is rooted in a fixed reference slot on the
\code{RevisionRootPage} (\code{PATH\_SUMMARY}, \code{NAME}, \code{CAS}, \code{PATH},
\code{DEWEYID}, \code{VECTOR}, \code{PROJECTION}) and materialised through a per-type page
(\code{PathPage}, \code{CASPage}, \code{NamePage}, \code{DeweyIDPage}, \code{VectorPage},
\code{ProjectionIndexPage}); multiple indexes of one type are distinguished by an index number
within that page. Every such page is copy-on-write versioned, so each historical revision keeps
its own independently readable indexes.

The lifecycle from declaration to steady state has three phases, and the cost falls entirely in the
middle one:

\dcap{An index's lifecycle: declare, build once, then maintain incrementally}
\begin{diagramS}
r5  resource has no user index.   declare:  IndexDefs.createCASIdxDef(/store/book/[]/price) -> id 7

build (the commit that creates it):
   IndexBuilder runs ONE DescendantAxis scan from the document root; the CASIndexListener visits each
   /store/book/[]/price and a CAS tree is constructed, rooted in the RevisionRootPage's CAS slot
   at commit r6:  definitions serialised to  store/indexes/6.xml   (now lists id 7)

steady state (every later write):
   insert a book with price 30  ->  the write transaction notifies CASIndexListener SYNCHRONOUSLY
   -> it writes (P_price, decimal, 30) -> nodeRef into the TIL, merged ATOMICALLY with the data at commit r7
   no rescan: the single build scan at r6 is the only full pass the index ever costs

open @ rev k:  the controller searches backward for the latest indexes/{<=k}.xml ; each revision's CAS page
               is CoW-versioned, so a historical read sees the index exactly as it stood then
\end{diagramS}

\noindent The contract is build-once, maintain-forever: a declaration pays one document scan at creation,
after which the per-type listener folds each insert and delete into the index inside the \emph{same} TIL
and the \emph{same} atomic commit as the data change. There is no background reindex and no separate index
log to replay --- the index and the document it covers are one durable revision, which is also why a
historical read of any revision finds a consistent index for exactly that revision.

\section{The value-and-structure indexes}

Four index kinds answer path and predicate queries over the document:

\begin{table}[h]
\small\setlength{\tabcolsep}{5pt}
\renewcommand{\arraystretch}{1.25}
\begin{tabularx}{\linewidth}{@{}l l X@{}}
\toprule
\textbf{Index} & \textbf{Key \textrightarrow\ value} & \textbf{Serves / notes}\\
\midrule
\textbf{PATH} & path-class record (PCR) \textrightarrow\ nodes & \code{//title}, structural existence; index-all when no paths declared\\
\textbf{CAS} & \code{(pathNodeKey, type, value)} \textrightarrow\ nodes & \code{[price > 50]} typed value predicates; ordered by path \emph{then} value, so each path's values form a contiguous range-scannable run\\
\textbf{NAME} & \code{QNm} \textrightarrow\ nodes & \code{*[name()='price']}; include/exclude name-set filters\\
\textbf{DeweyID} & nodeKey $\leftrightarrow$ DeweyID & document-order navigation, $O(1)$ ancestor/descendant tests\\
\bottomrule
\end{tabularx}
\end{table}

The \textbf{CAS} (content-and-structure) index is the workhorse for value predicates: its key
\code{CASValue} triples the typed atomic value, its \code{Type}, and the \code{pathNodeKey}, and
\code{compareTo} orders by path first and value second, making \code{CASFilterRange} a simple
inclusive/exclusive $[\min,\max]$ walk within a path's run. \emph{Caveat:} the \code{unique}
flag is declared, materialised, and round-tripped, but no code path enforces it; it is
metadata only, awaiting a uniqueness-checking consumer.

The \textbf{NAME} index (an optional inverted list over interned names) must not be confused with
the always-present \textbf{name dictionary} \code{Names} held in \code{NamePage}, which interns
every element/attribute/key name with a reference count and backs all node name storage. The
dictionary's reconstruction once grew under name churn (a high-water-mark scan); it is now
\textbf{O(live)} --- \code{Names} rebuilds from a persisted bitmap of live entry node-keys (the plan's
Approach~A, \pp{docs/NAME_DICTIONARY_RECONSTRUCTION_PLAN.md}), wired through \code{NamePage.getNames()}
and guarded by a differential test against the full-scan oracle; a documented invariant set
(I-N1\,\ldots\,I-N7) constrains it.

The \textbf{DeweyID} index stores \code{(nodeKey, SirixDeweyID)} records in a dedicated tree under
\code{DeweyIDPage} for the forward (key\,\textrightarrow\,label, document order) direction; the
reverse (a node's own label) is read straight off the node record, since every node already
carries its encoded DeweyID (Theorem~\ref{thm:dewey}).

\dcap{A worked NAME-index lookup}
\begin{diagram}
query:  jn:doc('db','catalog')//book          (all nodes named "book")

  the NAME index is keyed on the interned name key, not the string
  NameFilter resolves "book" → name key 42
        │  IndexFilterAxis scans the NAME tree for key 42
        ▼
  NodeReferences = { every node named "book" }, returned in document order
\end{diagram}

The CAS index turns a two-sided value predicate into a single ordered walk, which the NAME index's
single-key lookup cannot do.

\dcap{A worked CAS-index range scan: a two-sided value predicate as one walk}
\begin{diagramS}
query:  jn:doc('db','catalog').store.book[][?$$.price ge 10 and $$.price le 50]

  CAS index on /store/book/[]/price, keyed (pathNodeKey, type, value), ordered by PATH then VALUE
  resolve the price path → pathNodeKey P
  CASFilterRange walks the contiguous run for P, inclusive on both ends:
        (P, xs:decimal, 10.00)  ───────────▶  (P, xs:decimal, 50.00)
        |<──────────── every key in [10, 50], already in value order ───────────>|
  each key in range yields its NodeReferences (Roaring bitmap of node keys)
  ⇒ exactly the in-range price nodes; the AND of ge+le is ONE index walk, no document scan
\end{diagramS}

Because the key sorts on path \emph{before} value, every value of one field lives in a single
contiguous run, so the optimiser can fuse \code{>}\,/\,\code{>=} and \code{<}\,/\,\code{<=} on the same
field into one \code{BETWEEN} and serve it as a single bounded walk; a predicate on a \emph{different}
field is a different run entirely, never interleaved.

\section{Posting lists and the backing trees}

For PATH/CAS/NAME, the value at a key is a \code{NodeReferences} posting list backed by a
\textbf{\code{Roaring64Bitmap} of node keys}, a compressed multi-valued set rather than a pointer
chain. Each index can be built on either of two backends, selected by a \code{useHOT} switch
that every builder and listener threads through:

\begin{itemize}[leftmargin=1.4em]
\item \textbf{HOT} (the default): the Height-Optimised Trie of Ch.~\ref{ch:hot}, over the
  index's serialised key.
\item \textbf{Red-black tree} (\pp{index/redblacktree/}, the legacy/stability fallback): a
  BaseX-derived balanced BST that splits each entry into an \code{RBNodeKey} (tagged
  \code{PATHRB}/\allowbreak\code{CASRB}/\allowbreak\code{NAMERB}) pointing at a separate \code{RBNodeValue}.
\end{itemize}

\dcap{A secondary-index entry: key $\to$ posting list, over HOT or a red-black tree}
\begin{diagramS}
CAS index on /store/book/price (a value index):
   key (value, type, path-class)   →   NodeReferences = Roaring64Bitmap{ node keys }
   e.g. (12.5, double, PCR 7)      →   { 1042, 3310, 88914, … }

backing tree, chosen per index by useHOT:
   HOT (default)    Height-Optimised Trie over the serialised key (Ch. 5)
   RBTREE (legacy)  RBNodeKey (PATHRB / CASRB / NAMERB) ──▶ RBNodeValue

query: IndexFilterAxis scans the tree, applies the Filter set, emits NodeReferences
\end{diagramS}

At query time \code{IndexFilterAxis} scans the chosen tree, applies the declared \code{Filter}
set (a \code{PathFilter} resolving paths to PCRs, a \code{CASFilter}/\allowbreak\code{CASFilterRange}, or a
\code{NameFilter}), and yields each surviving entry's \code{NodeReferences}. A separate,
non-declared accelerator, \code{PageSkipRegistry}, opportunistically builds a
\code{nameKey}\,\textrightarrow\,Roaring-bitmap-of-pages map as a side effect of the first
analytical scan and invalidates it on a total-page-count mismatch; it is a scan heuristic, not a
persisted index. (The \pp{index/art/} \code{AdaptiveRadixTree} is a general-purpose in-memory
ordered map, \emph{not} wired as an index backend.)

A concrete query shows the path from predicate to result.

\dcap{A worked CAS-index point query}
\begin{diagramS}
query:  for $b in jn:doc('db','store').book[] where $b.price eq 12.5 return $b

  the optimiser rewrites the where-predicate to a CAS index scan on /store/book/[]/price
  CAS key = (12.5, double, PCR of that path)
        │  IndexFilterAxis scans the backing tree; CASFilter matches the exact key
        ▼
  NodeReferences = Roaring64Bitmap{ 1042, 3310, 88914 }     (the matching price nodes)
        │  move each to its owning book object (its parent)
        ▼
  result = the three matching book objects, in document order
\end{diagramS}

Two predicates on different fields compose by Roaring set algebra, with no document scan:

\dcap{Conjunctive predicates via Roaring posting-list AND}
\begin{diagramS}
query:  book[] where price in [10,50] AND category = "fiction"

  price CAS range scan [10,50] -> OR the per-value bitmaps along the run:
     10:{1001,1003,1102}  15:{1005,1050}  25:{1100,1101}  30:{1002,1101}  50:{2001}
     price_hits    = { 1001,1002,1003,1005,1050,1100,1101,1102,2001 }      (9 nodes)
  category CAS point "fiction":
     category_hits = { 1001,1050,1100,1101,2010 }                          (5 nodes)

  AND the two bitmaps (compressed-container intersection, no per-node work):
     price_hits AND category_hits = { 1001, 1050, 1100, 1101 }             (4 nodes satisfy BOTH)

cost: one CAS range-walk + one CAS point-lookup + one Roaring AND — never a pass over the books
\end{diagramS}

\noindent The range predicate \emph{unions} the bitmaps along the value run while the conjunction
\emph{intersects} across fields; both are compressed-container operations on \code{Roaring64Bitmap}s, so a
multi-dimensional filter costs a few set operations on posting lists rather than a scan, which is the
reason posting lists are bitmaps and not plain key lists.

\section{Columnar and vector indexes}

Two newer index kinds extend the family beyond posting lists.

The \textbf{projection index} (\code{PROJECTION}; \pp{index/projection/}) is a user-declared
\emph{columnar covering} index over a root path and an ordered field list (one row per matching
record, one bit-packed/dictionary-encoded column per field), feeding the SIMD scan of
Ch.~\ref{ch:vector}. It is wired through an interim \code{ProjectionIndexRegistry} that bypasses
the full listener lifecycle, and its builder requires a single-\code{pathNodeKey} root. A projection leaf is \textbf{sub-leaf chunked}: \code{ProjectionIndexHOTStorage} splits it into
fixed-size (4\,KiB) chunks, each stored as a \emph{separate} HOT slot keyed by
\code{(leafIndex, chunkIdx)}, so the HOT copy-on-write merge shares unchanged chunks across
revisions: a one-row change rewrites only the affected chunk, honouring the
\code{SLIDING\_SNAPSHOT} $O(\text{change})$-write contract rather than re-emitting the whole
$\sim$20~KB leaf. \emph{Caveat:} the remaining task~\#57 item is a \emph{further} optimisation,
lifting the chunk bytes out of the inline HOT values into dedicated chunk pages (a
\code{ChunkDirectory} of \code{chunkIdx}\,\textrightarrow\,\code{PageReference}), shrinking HOT slots
to $\sim$16-byte references so a leaf holds far more entries before splitting; it is required before a
public projection-index API commitment.

The \textbf{vector index} (\code{VECTOR}; \pp{index/vector/}) is the ``naming hazard'' of
\S\ref{ch:vector}: an \textbf{HNSW} (Hierarchical Navigable Small World) $k$-nearest-neighbour index
over float embeddings, \emph{not} the columnar projection. \code{HnswParams} fixes
$m{=}16$, \code{efConstruction}${=}200$, \code{efSearch}${=}50$, with \code{L2}/\allowbreak\code{COSINE}/%
\code{INNER\_PRODUCT} distances backed by SIMD kernels; the persistent layout is a vector tree whose
first key holds a metadata/entry-point node and whose remaining keys hold \code{VectorNode}s
(embedding {+} neighbour lists), with tombstone deletes. \emph{Caveat:} vector indexes are
populated \emph{explicitly}, not by document traversal (no traversal-driven builder), XML resources
do not support them, and the store keeps in-memory metadata; it is the most experimental member
of the family.

\dcap{HNSW search: descend the layers, then widen at the base}
\begin{diagram}
  layer 2   entry o───────────o            few nodes, long hops
                  │           │
  layer 1   o─────o─────o─────o─────o       more nodes, medium hops
                  │     │     │
  layer 0   o─o─o─o─o─o─o─o─o─o─o─o─o─o      every VectorNode (embedding + neighbour list)

  query: enter at the top, greedily hop to the closest neighbour per layer,
         then explore a beam of width efSearch at layer 0 → the k nearest
\end{diagram}

The build and query paths are the textbook HNSW algorithms over that persistent graph. An
\textbf{insert} draws a random top layer from a geometric distribution (\code{level}${=}\lfloor-\ln U
\cdot m_L\rfloor$ with $m_L{=}1/\ln m$, so most vectors live only at layer~0 and a thin tower rises
above), greedily descends from the entry point to that layer, and then, at each layer down to~0, runs a
width-\code{efConstruction} beam search and wires the new node to a \emph{diverse} neighbour set chosen
by Malkov's heuristic~\cite{hnsw}, neighbours that are closer to the new node than to one another, rather than simply
the nearest. Links are bidirectional, and a back-link that pushes a neighbour past its degree cap
(\code{mMax0}${=}2m$ at layer~0, $m$ above) re-runs the same heuristic to trim it. A \textbf{search}
greedily hops to the closest node per layer down to~0, then widens to an \code{efSearch} beam with
early termination once the candidate frontier is farther than the current $k$\textsuperscript{th} best.
Distances are the SIMD kernels (\code{FloatVector.SPECIES\_PREFERRED}, one pass accumulating dot product
and both norms for cosine). Because every \code{VectorNode} edit goes through the same
copy-on-write record path as any other index node, the graph is \emph{versioned for free}: a historical
read opens the graph as it stood at that revision and searches it with no special handling.

A concrete search makes the greedy-descend-then-widen strategy visible. Take 2D embeddings (real ones are
high-dimensional \code{float32}; 2D keeps the distances checkable) and query $q=(0.2, 0.8)$, seeking the
single nearest:

\dcap{An HNSW k=1 search: greedy hop per layer, then a layer-0 beam}
\begin{diagramS}
nodes:  A(0.1,0.9)  B(0.9,0.1)  C(0.3,0.7)  D(0.8,0.2)  E(0.15,0.85)  F(0.5,0.5)
towers: layer2 {B,F}   layer1 {B,C,F}   layer0 {all}        query q=(0.2,0.8)
L2(q, .):  A 0.14   B 0.99   C 0.14   D 0.85   E 0.071   F 0.42

layer 2   enter B (0.99) -> neighbour F (0.42) is closer -> hop F ; F is local-min -> descend at F
layer 1   at F (0.42)    -> neighbour C (0.14) is closer -> hop C ; C is local-min -> descend at C
layer 0   at C (0.14)    -> efSearch beam over C's neighbours -> E (0.071) closer -> hop E ; E local-min
result:   nearest = E (0.071)      — 3 hops, never a scan of all six nodes
\end{diagramS}

\noindent The thin upper towers do the long-range work --- the first two hops cross most of the space at
$O(\log n)$ cost --- and only the bottom layer, where every node lives, runs the wider \code{efSearch} beam
for accuracy. That descend-coarse-then-refine shape is the whole reason HNSW answers a $k$-NN query in
logarithmic hops instead of a linear scan.

\begin{honesty}
\small\textbf{Caveat.} The vector index is the family's most experimental member and its boundaries are
worth stating plainly. It is \textbf{not yet exposed in the query language}, there is no
\code{jn:vector-search} function, so a $k$-NN query today is a programmatic \code{VectorIndex.searchKnn}
call on a \code{StorageEngineReader}, not a JSONiq expression. Population is \textbf{explicit only} (the
builder is a deliberate no-op, so vectors are inserted by hand rather than harvested from a document
field), embeddings are stored as full-precision \code{float32} with \textbf{no quantisation}, deletes are
\textbf{tombstones} with no graph compaction, and the test coverage, while real (insert/search/recall
tests), exercises only small dimensions and dataset sizes. It demonstrates that an HNSW graph composes
cleanly with the versioned page trie; it is not a tuned production ANN engine.
\end{honesty}

\section{Maintenance}

A new index is populated once by \code{IndexBuilder}, a single \code{DescendantAxis} scan from the
document root that drives every registered builder's visitor. Thereafter the indexes are kept
current incrementally: the \code{IndexController} registers per-type \textbf{listeners}
(\code{PathIndexListener}, \code{CASIndexListener}, \code{NameIndexListener},
\code{VectorIndexListener}) that the write transaction notifies on every insert/delete. The hot
path uses a primitive \code{PathNodeKeyChangeListener} signature
(\code{(nodeKey, kind, pathNodeKey, name, value)}) to avoid materialising an immutable change
snapshot per node; each listener consults the path summary for path/type eligibility and routes
the change to \code{index(...)}/\allowbreak\code{remove(...)} on the chosen backend. Listeners are re-created
per write transaction, and the catalogue is re-serialised at commit.

\section{Change tracking: the history indexes}
\label{sec:changetracking}

Two of the engine's \emph{internal} indexes (\S6.2) turn a resource's change history into a direct
lookup rather than a tree comparison, both \textbf{on by default}. \code{RECORD\_TO\_REVISIONS} maps each
node to the ascending list of revisions in which it was inserted, modified, or deleted, stored as a
\code{RevisionReferencesNode} (an \code{int[]} of revision numbers), written on \emph{every} insert and
update by the write transaction and read back in one shot by \code{RecordRevisionsLookup}.
\code{CHANGED\_NODES} records, per revision, the set of changed subtree-roots, maintained as its own
copy-on-write trie hung off the \code{RevisionRootPage}. Independently, every node carries a
\code{previousRevision}/\allowbreak\code{lastModifiedRevision} pair, forming a \textbf{back-pointer chain}: from any
materialised node one walks straight back through exactly the revisions in which it changed, with no
comparison.

\dcap{Change tracking: two internal indexes plus a per-node back-pointer chain}
\begin{diagramS}
node X (key 12345), changed in revisions 3, 7, 12, 40:

  RECORD_TO_REVISIONS[X] = [3, 7, 12, 40]    a RevisionReferencesNode (int[] of revisions)
  back-pointers:  X@40 ──prev──▶ X@12 ──prev──▶ X@7 ──prev──▶ X@3    (walk back, no comparison)
  CHANGED_NODES@r = { changed subtree-roots }    a copy-on-write trie hung off RevisionRootPage_r
\end{diagramS}

\begin{property}[History-lookup cost]\label{prop:histcost}
For a node that changed in $k$ revisions, listing those revisions (\code{sdb:item-history}) costs $O(k)$
from the \code{RECORD\_TO\_REVISIONS} index, independent of the total number of revisions or the size of
the document.
\end{property}

The payoff is that the common temporal questions become \textbf{$O(\text{changes})$ index or pointer
lookups, not tree walks} (Figure~\ref{fig:histindex}). ``In which revisions did node $X$ change?'' (\code{sdb:item-history}) reads the
record's revision list and opens just those revisions in parallel virtual threads; the time-travel axes
(\code{jn:all-times}, \code{jn:past}, \code{jn:future}) use the same pointers to jump to a node's first
revision rather than opening a transaction per pre-creation revision (skipping, e.g., 989 useless
transaction starts for a node born at revision 990 of 1000). This is the sense in which \textbf{change
tracking can beat diffing}: it answers \emph{where} and \emph{when} something changed without
materialising or descending a single tree. It is complementary to the hash-guided \emph{diff}
(\S\ref{sec:diff}), which answers the different question of the \emph{exact transformation} between two
given revisions.

\dcap{Pre-creation skip: a node born late in a long history}
\begin{diagramS}
node Y first appears at revision 990 of a 1000-revision resource; later touched at 994 and 1000

  naive temporal scan:  open rtx@1, @2, … , @989  → each sees "Y absent"   (989 wasted opens)
                        then @990 found, continue @991 … @1000
  SirixDB:              read RECORD_TO_REVISIONS[Y] = [990, 994, 1000]      (one index lookup)
                        open ONLY those three revisions, in parallel virtual threads
     jn:first(Y) = Y@990     jn:past(Y) = {990, 994}     jn:all-times(Y) = {990, 994, 1000}

  cost = O(changes to Y), not O(revisions): the 989 pre-creation transaction starts never happen
\end{diagramS}

\begin{figure}[t]
\centering
\begin{tikzpicture}[font=\sffamily\small,>=Latex,
  io/.style={draw=slate!70,fill=slate!12,rounded corners=2pt,align=center,inner sep=3pt,
             minimum height=0.95cm,minimum width=1.85cm,font=\footnotesize},
  rev/.style={draw=ember!75,fill=ember!12,align=center,inner sep=2pt,
              minimum height=0.7cm,minimum width=1.05cm,font=\footnotesize,text=ember!85!black},
  res/.style={draw=slate!60,fill=slate!8,align=center,inner sep=2pt,
              minimum height=0.7cm,minimum width=1.05cm,font=\footnotesize},
  aside/.style={draw=nodeold,fill=nodeold!12,rounded corners=2pt,align=left,inner sep=4pt,
                font=\scriptsize,text=nodeold!55!black,minimum width=4.3cm},
  e/.style={->,slate},
  ce/.style={->,ember}]

  \node[font=\footnotesize\bfseries,text=slate!80!black,align=center] at (3.8,3.5)
       {\code{RECORD\_TO\_REVISIONS[Y]} --- revisions where Y changed (\code{int[]}, $O(k)$)};

  % the index array (a RevisionReferencesNode)
  \node[io]  (y)  at (0,2.7) {node Y\\(key 12345)};
  \node[rev] (r0) at (2.7,2.7) {990};
  \node[rev] (r1) at (3.8,2.7) {994};
  \node[rev] (r2) at (4.9,2.7) {1000};
  \draw[e] (y.east) -- node[above,font=\scriptsize,text=slate!70!black]{index} (r0.west);

  % opened node versions (only the changed revisions)
  \node[res] (v0) at (2.7,1.1) {Y@990};
  \node[res] (v1) at (3.8,1.1) {Y@994};
  \node[res] (v2) at (4.9,1.1) {Y@1000};
  \draw[ce] (r0)--(v0); \draw[ce] (r1)--(v1); \draw[ce] (r2)--(v2);
  \node[fill=white,inner sep=1.5pt,font=\scriptsize,text=ember!82!black] at (3.8,1.9)
       {open only these --- parallel virtual threads};

  % contrast: the naive scan the index avoids
  \node[aside] (n) at (8.95,2.55)
       {\textit{replaces} the naive scan:\\open rtx@1\,\ldots\,@989 $\to$ all ``Y absent''\\(989 wasted opens), then @990\,\ldots\,@1000};
\end{tikzpicture}
\caption{Sub-document history via the \code{RECORD\_TO\_REVISIONS} change-tracking index. For a node Y, the
index is a \code{RevisionReferencesNode} --- an \code{int[]} of exactly the revisions in which Y changed
--- so \code{sdb:item-history} is one $O(k)$ array read, never a per-revision tree walk (Property
\ref{prop:histcost}). The temporal axes are slices of that list: \code{jn:first}/\code{jn:last} are its ends,
\code{jn:past}/\code{jn:future} the prefix/suffix around a pivot revision, \code{jn:all-times} the whole
array. Only those revisions are opened --- in parallel virtual threads, each resolved through the time axis
of Figure~\ref{fig:twoaxis} --- so a node born at revision 990 of 1000 skips 989 pre-creation transaction
starts entirely.}
\label{fig:histindex}
\end{figure}

%================================================================
\chapter{Query Processing and Optimisation}
\label{ch:query}

\noindent\textit{(\code{brackit}; \pp{io/sirix/query/*}; \pp{io/sirix/diff/*};
\pp{docs/cost-based-optimizer-design.md}.)}

\section{Brackit: the JSONiq/XQuery engine}

With the indexes of Chapter~\ref{ch:pathsummary} in place, this chapter is how a query is
planned and rewritten to reach for them. SirixDB embeds \textbf{brackit}, a query engine from
B\"achle's doctoral work~\cite{baechle} whose three organising ideas, named in its title, are
\textbf{set orientation}, \textbf{physical data independence}, and \textbf{parallelism}. A query is parsed
into an \code{AST} whose nodes double as
the optimiser's annotation medium and the translator's blueprint; FLWOR expressions normalise
into tuple pipelines (for/let/where/group/order/return). The control-flow interface is a
\textbf{bind-evaluate (Volcano) iterator}: every \code{Expr.evaluate(QueryContext, Tuple)} returns a lazy
\code{Sequence} of \code{Item}s pulled on demand, so operators avoid materialisation. \emph{But the data
model underneath that interface is set-oriented, not one-tuple-at-a-time}: a pipeline operator is built
through \code{Operator.create(ctx, Tuple[] tuples, int len)}, a form that hands it a whole \emph{batch}
of tuples, and the pipeline physically transports tuples in bundles rather than singly. That separation,
a pull-iterator on top of set-oriented batch execution, is what \emph{physical data independence} buys:
the logical algebra never changes, yet the same pipeline can be run serially, in parallel, or through a
vectorised columnar kernel (\S\ref{sec:setoriented} below, and the SIMD path of
Chapter~\ref{ch:vector}); Figure~\ref{fig:pdi} and Property~\ref{prop:pdi} make that separation precise.
A \textbf{store SPI} lets back-ends plug in without touching the
compiler/runtime, and a parallel \textbf{vectorised-executor SPI} lets a back-end supply a columnar
kernel for the shapes it can accelerate. This is brackit as B\"achle frames it --- a \emph{retargetable
query compiler}, in which a host like SirixDB contributes physical optimisations (index-matching
rewrites, storage access, SIMD kernels) at fixed seams rather than by forking the engine.

\dcap{A FLWOR query compiled to a Volcano tuple pipeline}
\begin{diagramS}
for $b in jn:doc('db','books').store.book[]
let $p := $b.price   where $p gt 10   order by $p   return { title: $b.title, price: $p }

  Start ─▶ ForBind($b) ─▶ LetBind($p) ─▶ Select($p>10) ─▶ OrderBy($p) ─▶ Return{…}
  operators pull Tuples on demand (singly, or in batches when parallelised/vectorised); Sequences are
  lazy, so nothing materialises eagerly
\end{diagramS}

Each clause becomes one operator over a stream of \code{Tuple}s, and because every
\code{Sequence} is evaluated lazily, a query that selects a handful of items out of a 100~M-node
resource touches only the path it actually needs, never the whole document. The cost-based optimiser of
\S\ref{sec:optpipeline} rewrites this pipeline before it runs, turning the \code{where} predicate into an
index scan where a covering index exists.

\section{The compiler pipeline: from query text to plan}
\label{sec:compiler}

Before any of that runs, the query text passes through a four-phase compiler
(\code{CompileChain.compile}, \pp{io/brackit/query/compiler/}). Each phase rewrites a single mutable
\code{AST}, and SirixDB extends the chain rather than replacing it.

\dcap{The brackit compile chain (\texttt{CompileChain}), and where SirixDB hooks in}
\begin{diagramS}
query text
   │  parse      JsoniqParser (hand-written recursive descent over a Tokenizer; not ANTLR)
   ▼
parsed AST       nodes typed by XQ.* kind constants (FlowrExpr, PipeExpr, DerefExpr, InsertJsonExpr, …);
   │             a per-node properties map carries optimiser annotations
   │  analyze    Analyzer: resolve QNames/namespaces, build the symbol table of functions+variables,
   ▼             check cyclic declarations → emit one Target per compilation unit (main body, each UDF)
analysed AST
   │  optimize   a sequence of Stages, each a Walker run to a FIXPOINT (overrides SirixOptimizer ↓)
   ▼
optimised AST    snapshot kept (getOptimizedAST) → the GUI plan visualiser (Fig. plans)
   │  translate  Compiler dispatches on XQ kind; PipelineStrategy turns a PipeExpr into an Operator chain
   ▼
executable Expr / Operator tree
\end{diagramS}

\textbf{Parse.} \code{JsoniqParser} is a hand-written recursive-descent parser over a character-level
\code{Tokenizer}, covering XQuery~3.1 and the JSONiq dialect; there is no parser generator. It produces
an \code{AST} whose every node is tagged with an integer \emph{kind} from \code{XQ} (\code{FlowrExpr},
\code{ForClause}, \code{DerefExpr}, \code{ObjectConstructor}, \code{InsertJsonExpr}, and brackit's
set-oriented extensions \code{PipeExpr}, \code{Start}, \code{ForBind}, \code{Join}, \ldots), plus a
\code{properties} map the later phases use to annotate it. \textbf{Analyze.} The \code{Analyzer} does
semantic analysis, not rewriting: it resolves namespace prefixes and \code{QNm}s, builds a symbol table
from the prolog's function and variable declarations (with forward declarations for mutual recursion),
checks for cyclic initialisers, and slices the program into \code{Target}s, one per compilation unit.
\textbf{Optimise.} Each \code{Target} is handed to an \code{Optimizer}, a list of \code{Stage}s applied
in order; brackit's \code{TopDownOptimizer} runs \emph{Simplification} (constant folding, step merging),
\emph{Pipelining} (normalise FLWOR into the \code{Start\,$\to$\,bindings\,$\to$\,End} pipe), join
recognition, let-unnesting into left joins, and a finalisation pass. Each stage is a \code{Walker}, a
depth-first visitor re-run to a \emph{fixpoint}, so a rewrite that exposes another fires again in the
same pass. \textbf{Translate.} The \code{Compiler} walks the optimised AST dispatching on \code{XQ} kind,
and the pluggable \code{PipelineStrategy} compiles each \code{PipeExpr} into an \code{Operator} chain
(\S\ref{sec:physops}); the resulting \code{Expr}/\code{Operator} tree is what the Volcano iterator drives.

The seam SirixDB uses is exactly this stage list (Figure~\ref{fig:compilechain}). \code{SirixCompileChain} extends \code{CompileChain},
returning a \code{SirixOptimizer} (the ten cost-based stages of \S\ref{sec:optpipeline} appended to
brackit's base stages) and a translator whose \code{SequentialPipelineStrategy} consults the registered
\code{SirixVectorizedExecutor}. Because the chain keeps a copy of the AST after parsing and after
optimisation (\code{getParsedAST}/\allowbreak\code{getOptimizedAST}, serialisable to JSON), the GUI's
parsed-vs-optimised plan visualiser (Figure~\ref{fig:plans}) is a direct render of those two snapshots,
not a separate reconstruction.

A short trace shows one FLWOR through the phases.

\dcap{One query through the compile chain}
\begin{diagramS}
for $b in jn:doc('db','store').book[] where $b.price gt 50 return $b.title

parse →     FlowrExpr[ ForClause($b ← DerefExpr(book[])) , WhereClause(price>50) , ReturnClause($b.title) ]
analyze →   QNames resolved; $b scoped; one Target (the main body); no rewrite of shape
optimize →  Pipelining: FLWOR ⇒ PipeExpr[ Start → ForBind($b) → Select(price>50) → Return($b.title) → End ]
            CostBased (SirixOptimizer): annotate ForBind PREFER_INDEX ; IndexMatching → JsonCASStep(price≥50)
translate → PipeExpr ⇒ Operator chain: Start ▸ ForBind ▸ Select ▸ (Return expr) — driven by the Volcano iterator
\end{diagramS}

The parsed and optimised forms are the two the GUI shows side by side; the gap between them, a generic
\code{book[]} scan on the left and a CAS-index range scan on the right, is the whole work of the
optimiser made visible.

\begin{figure}[t]
\centering
\begin{tikzpicture}[font=\sffamily\small,>=Latex,
  art/.style={draw=slate!55,fill=slate!8,rounded corners=2pt,align=center,
              inner sep=4pt,minimum width=2.6cm,minimum height=0.62cm},
  io/.style={align=center,font=\sffamily\footnotesize\itshape,text=slate!70!black},
  ph/.style={fill=white,inner sep=1.6pt,font=\sffamily\footnotesize\itshape,text=slate!80!black},
  hook/.style={draw=nodenew,fill=nodenew!12,rounded corners=2pt,align=center,
               inner sep=4pt,font=\footnotesize,minimum width=3.7cm},
  gui/.style={draw=inkblue,fill=inkblue!8,rounded corners=2pt,align=center,
              inner sep=4pt,font=\footnotesize,minimum width=3.0cm,text=inkblue!88!black},
  e/.style={->,slate},
  he/.style={->,nodenew},
  ge/.style={->,inkblue,densely dashed}]

  % artifact spine (brackit's base compile chain)
  \node[io]  (qt) at (0,5.5) {query text};
  \node[art] (pa) at (0,4.4) {parsed AST};
  \node[art] (aa) at (0,3.1) {analysed AST};
  \node[art] (oa) at (0,1.8) {optimised AST};
  \node[io]  (op) at (0,0.45) {\code{Expr}/\code{Operator} tree\\(Volcano-driven)};

  \draw[e] (qt)--(pa) node[ph,midway] {parse};
  \draw[e] (pa)--(aa) node[ph,midway] {analyze};
  \draw[e] (aa)--(oa) node[ph,midway] {optimize};
  \draw[e] (oa)--(op) node[ph,midway] {translate};

  % SirixDB seams (right, rust)
  \node[hook] (h1) at (5.0,2.45) {\textbf{SirixOptimizer}\\ten cost-based stages appended\\to brackit's base stages};
  \node[hook] (h2) at (5.0,0.70) {\textbf{SequentialPipelineStrategy}\\routes each \code{PipeExpr} through\\the \code{SirixVectorizedExecutor}};
  \draw[he] (h1.west) to[out=180,in=0] (1.05,2.45);
  \draw[he] (h2.west) to[out=180,in=0] (1.05,1.12);

  % what the GUI consumes (left, inkblue, dashed)
  \node[gui] (g) at (-3.95,3.1) {GUI plan visualiser\\renders both snapshots\\\code{getParsedAST} /\\\code{getOptimizedAST}};
  \draw[ge] (pa.west) to[out=180,in=90]  (g.north);
  \draw[ge] (oa.west) to[out=180,in=-90] (g.south);
\end{tikzpicture}
\caption{The four-phase brackit compile chain (\code{CompileChain}) and the two seams where
\code{SirixCompileChain} splices in. The \code{SirixOptimizer}'s ten cost-based stages
(\S\ref{sec:optpipeline}) run \emph{after} brackit's base optimisation; at translation the
\code{SequentialPipelineStrategy} routes each \code{PipeExpr} through the
\code{SirixVectorizedExecutor}. The parsed and optimised AST snapshots are exactly what the GUI's
side-by-side plan visualiser (Figure~\ref{fig:plans}) renders --- the optimiser's work made visible.}
\label{fig:compilechain}
\end{figure}

\section{Set orientation, parallelism, and vectorised operators}
\label{sec:setoriented}

The pull iterator above is the \emph{interface}; the \emph{execution} is set-oriented, and that is what
lets the same logical pipeline be sped up three ways without rewriting it. All three live in brackit's
operator layer (\pp{io/brackit/query/operator/}) and are reached through the translator's pipeline
strategy.

\textbf{Batched tuple transport.} An \code{Operator} exposes two factory forms: \code{create(ctx, Tuple
single)} and \code{create(ctx, Tuple[] tuples, int len)}. The second is the set-oriented one: an operator
is fed a \emph{buffer} of tuples and processes the run, so the inter-operator hand-off is an array, not a
call per tuple. This is the concrete form of B\"achle's \emph{set orientation}: the algebra is defined
over tuple sets, and a single-tuple call is just the degenerate batch of one.

\textbf{Pipeline parallelism (\code{Parallelizer}).} To parallelise a pipeline, brackit splices in a
\code{Parallelizer} operator. Its cursor runs the upstream in a dedicated producer thread that fills
fixed-size \code{Tuple[]} buffers (2000 tuples each) drawn from a small free-list, and consumer threads
dequeue whole buffers and drive the downstream pipeline. Tuples thus cross the operator boundary in
bundles of up to 2000, decoupling producer from consumers and turning a serial pull chain into a
producer/consumer pipeline. A \code{BlockingParallelizer} variant gives the synchronising form for
order-sensitive consumers.

\dcap{Three set-oriented execution modes over one logical pipeline}
\begin{diagramS}
logical pipeline:   Start ─▶ ForBind ─▶ Select ─▶ (GroupBy | OrderBy) ─▶ Return     (unchanged)

serial            one Cursor chain; Tuple[] batches of 1 … N flow through on demand
Parallelizer      producer thread fills Tuple[2000] buffers ─▶ free-list hand-off ─▶ consumer threads
morsel-driven     MorselSource hands 2048-tuple Morsels to ForkJoinPool workers, each a pipeline CLONE;
                  pipeline breakers (GroupBy, OrderBy) sit OUTSIDE the morsel boundary and merge after
vectorised        non-breaker prefix feeds VectorizedGroupByExec: values land in reusable 2048-element
                  COLUMN arrays, aggregates update in place (count/sum/min/max), ~10 group entries
\end{diagramS}

\begin{figure}[t]
\centering
\begin{tikzpicture}[font=\sffamily\small,>=Latex,
  lp/.style={draw=slate!65,fill=slate!10,rounded corners=2pt,align=center,inner sep=4pt,
             minimum height=0.85cm,minimum width=8.6cm,font=\footnotesize},
  strat/.style={draw=ember!75,fill=ember!12,rounded corners=2pt,align=center,inner sep=3pt,
                font=\footnotesize,text=ember!85!black,minimum width=6.2cm},
  phys/.style={draw=slate!55,fill=layerbg,rounded corners=2pt,align=center,inner sep=3pt,
               minimum height=1.15cm,minimum width=3.0cm,font=\scriptsize},
  vec/.style={draw=nodenew,fill=nodenew!12,rounded corners=2pt,align=center,inner sep=3pt,
              minimum height=1.15cm,minimum width=3.0cm,font=\scriptsize},
  e/.style={->,slate},
  ee/.style={->,ember}]

  \node[lp] (lp) at (5.1,3.5)
       {\textbf{one logical \code{Operator} pipeline} (the optimised AST --- unchanged)\\
        Start $\to$ ForBind $\to$ Select $\to$ (GroupBy $\,|\,$ OrderBy) $\to$ Return};
  \node[strat] (s) at (5.1,2.25) {\code{SequentialPipelineStrategy} --- picks the physical form at \emph{translation}};
  \draw[ee] (lp) -- (s);

  \node[phys] (p1) at (0.55,0.7)  {\textbf{serial}\\\code{Cursor} chain\\\code{Tuple[]} $1\ldots N$};
  \node[phys] (p2) at (3.75,0.7)  {\textbf{\code{Parallelizer}}\\producer $\to$ 2000-tuple\\buffers $\to$ consumers};
  \node[phys] (p3) at (6.95,0.7)  {\textbf{morsel-driven}\\ForkJoinPool clones\\(Fig.~\ref{fig:morsel})};
  \node[vec]  (p4) at (10.15,0.7) {\textbf{vectorised}\\2048-col arrays\\$\to$ SIMD SPI (Ch.~\ref{ch:vector})};
  \draw[e] (s.south) -- (p1.north);
  \draw[e] (s.south) -- (p2.north);
  \draw[e] (s.south) -- (p3.north);
  \draw[e] (s.south) -- (p4.north);

  \node[font=\scriptsize\itshape,text=slate!72!black,align=center] at (5.35,-0.45)
       {the plan is compiled once and never rewritten; only the physical realisation is chosen --- and morsel $+$ vectorised compose};
\end{tikzpicture}
\caption{Physical data independence (Property~\ref{prop:pdi}), the central concern of B\"achle's
thesis~\cite{baechle}. A query is compiled \emph{once} into a single \code{Operator} pipeline; the
translator's pipeline strategy then chooses \emph{how} it runs --- a serial cursor chain, a
\code{Parallelizer} producer/consumer split, morsel-driven workers (Figure~\ref{fig:morsel}), or a
vectorised columnar kernel feeding the SIMD SPI of Chapter~\ref{ch:vector} --- without rewriting the plan
or recompiling the query. The same seam that admits these execution strategies admits new storage
backends (Property~\ref{prop:backendid}): one logical algebra, many physical realisations.}
\label{fig:pdi}
\end{figure}

\textbf{Morsel-driven parallelism (\code{MorselPipeline}).} The modern alternative to a producer thread is
morsel-driven scheduling (Leis et al.~\cite{morsel}). A \code{MorselPipeline} wraps an upstream pipeline whose output is
drawn by a serial \code{MorselSource}; $N$ ForkJoinPool workers each call \code{nextMorsel()} to grab a
\code{Morsel} (a fixed-capacity batch, \code{DEFAULT\_SIZE}~$=2048$ tuples) and run it through a
\emph{per-worker clone} of the pipeline. Because the per-morsel work is independent, the non-blocking
prefix of a pipeline parallelises trivially; \emph{pipeline breakers} (\code{GroupBy}, \code{OrderBy})
are kept outside the morsel boundary and reduce the per-worker partials afterward. The unit of work is a
morsel, so load balances across cores regardless of data skew.

\textbf{Vectorised columnar group-by (\code{VectorizedGroupByExec}).} For streaming group-by-aggregate,
brackit can abandon tuple-at-a-time entirely: it materialises the grouping key and the aggregated value
into two reusable \emph{columnar} arrays of 2048 elements and aggregates in place into a small hash of
group entries, so a batch of 2048 rows is summed, min/max'd, or counted as column operations rather than
2048 object-dispatched method calls. This is the engine-level analogue of, and the SPI seam for, the
SIMD numeric kernels of Chapter~\ref{ch:vector}: brackit owns the set-oriented, columnar \emph{framework}
and the \code{VectorizedExecutor} SPI, and SirixDB's \code{SirixVectorizedExecutor} plugs a genuine Java
Vector API (\code{LongVector}/\allowbreak\code{VectorMask}) implementation into it for the numeric, projection-indexed
shapes it can prove. The two layers compose: morsel parallelism splits the scan across workers, and each
worker's morsel is then evaluated by a vectorised kernel (Figure~\ref{fig:morsel}).

\begin{figure}[t]
\centering
\begin{tikzpicture}[font=\sffamily\small,>=Latex,
  box/.style={draw=slate!60,fill=layerbg,rounded corners=2pt,align=center,inner sep=3pt,minimum height=0.75cm},
  mor/.style={draw=ember!60,fill=ember!10,rounded corners=2pt,align=center,inner sep=3pt,minimum height=0.75cm},
  wk/.style={draw=slate!55,fill=slate!8,rounded corners=2pt,align=center,inner sep=3pt,minimum height=1.0cm,minimum width=3.7cm},
  brk/.style={draw=nodenew,fill=nodenew!12,rounded corners=2pt,align=center,inner sep=3pt,minimum height=1.3cm,minimum width=1.9cm},
  e/.style={->,slate}]
  \node[box] (scan) at (0,0) {columnar\\scan};
  \node[mor,right=0.5cm of scan] (m) {Morsels\\2048 tuples};
  \draw[e] (scan) -- (m);
  \node[wk] (w1) at (6.3,1.8)  {worker 1\\vectorised kernel\\(SIMD over 2048-col)};
  \node[wk] (w2) at (6.3,0.0)  {worker 2\\vectorised kernel};
  \node[wk] (w3) at (6.3,-1.8) {worker $n$\\vectorised kernel};
  \node[font=\scriptsize\itshape,slate] at (6.3,2.8) {ForkJoinPool: a pipeline clone per worker};
  \draw[e] (m.east) -- (w1.west); \draw[e] (m.east) -- (w2.west); \draw[e] (m.east) -- (w3.west);
  \node[brk] (brk) at (10.9,0) {breaker:\\group-by /\\order-by};
  \draw[e] (w1.east) -- (brk.north west); \draw[e] (w2.east) -- (brk.west); \draw[e] (w3.east) -- (brk.south west);
  \node[box,below=0.55cm of brk] (out) {result};
  \draw[e] (brk) -- (out);
\end{tikzpicture}
\caption{Morsel-driven execution with a vectorised kernel per worker (\S\ref{sec:setoriented},
Chapter~\ref{ch:vector}). The scan hands fixed 2048-tuple morsels to a pool of \code{ForkJoinPool} workers,
each an independent clone of the pipeline's pipelineable operators evaluated by a SIMD columnar kernel; the
\emph{breaker} (group-by or order-by) sits outside the morsel boundary and reduces the workers' partials.
The pipelineable prefix scales with cores; the breaker is the one synchronisation point.}
\label{fig:morsel}
\end{figure}

\begin{property}[Physical data independence]\label{prop:pdi}
One logical algebra admits many physical executions. A query is compiled once into a single
\code{Operator} pipeline; the translator's pipeline strategy then selects \emph{how} that pipeline runs ---
a serial cursor chain, a \code{Parallelizer} producer/consumer split, morsel-driven workers, or a
vectorised columnar kernel --- and which backend serves its scans, without rewriting the plan or
recompiling the query. New physical capability (a backend, an index-matching rewrite, a SIMD kernel) joins
at an SPI seam, not by forking the engine. This is the runtime companion of the backend-independent
\emph{format} of Property~\ref{prop:backendid}: there, one resource over many storage backends; here, one
plan over many execution strategies --- the two faces of B\"achle's \emph{retargetable} design.
\end{property}

\begin{honesty}
\small\textbf{Caveat.} These modes are opt-in and shape-restricted, not the default for every query. The
serial \code{Cursor} pipeline is what runs unless the translator's pipeline strategy selects a parallel
or vectorised operator, and the vectorised columnar path covers the same narrow, fail-closed set of
aggregate/group-by shapes described in Chapter~\ref{ch:vector}; anything outside it falls back to the
set-oriented interpreter, which is itself already batched. The point of this section is the
\emph{mechanism}, set orientation under a pull interface, that makes the speed-ups expressible at all,
not a claim that every query is parallel or vectorised.
\end{honesty}

\section{Physical operators: joins, grouping, and sorting}
\label{sec:physops}

The pipeline operators above are the easy, stateless half of the algebra. The other half, the operators
that must see more than one tuple at once to do their job, are where the set-oriented model earns its
keep and where the implementation has to decide what to do when the working set does not fit in memory.

\textbf{Joins.} brackit ships two join operators, and the optimiser of \S\ref{sec:joinorder} picks the
order in which they are applied. \code{NLJoin} is the nested-loop join: for each left tuple it re-opens
the right input and tests the join predicate pair by pair, an $O(\lvert L\rvert\cdot\lvert R\rvert)$ shape
that is the right choice only when the right input is tiny or the predicate is a non-equality. Its
\code{leftJoin} flag adds left-outer semantics by emitting an unmatched left tuple padded with nulls.
\code{TableJoin} is the workhorse \emph{hash join}: at cursor-creation it compares the two inputs' sizes
and \textbf{builds its in-memory table on the smaller side}, then streams the larger side and probes.
The table itself is a \code{MultiTypeJoinTable} that adapts to the comparison: an equality join
(\code{Cmp.eq}) uses a true \code{HashJoinTable}, while an order comparison falls back to a
\code{SortedJoinTable}, and heterogeneous key types are reconciled by the same numeric-promotion rules
the value comparisons use, so a join key that is \code{long} on one side and \code{double} on the other
still matches. \code{TableJoin} also supports the left-outer form; semi- and anti-joins are not a
distinct operator but fall out of the surrounding pipeline.

\dcap{The two join operators and when each is chosen}
\begin{diagramS}
NLJoin     nested loop, O(|L|·|R|)   re-open right per left tuple, test predicate   best when |R| tiny or non-equality
TableJoin  hash join                 build table on the SMALLER input, probe with the larger
           equality (Cmp.eq)  → HashJoinTable        order comparison → SortedJoinTable
           mixed key types reconciled by numeric promotion;  leftJoin flag → left-outer padding
\end{diagramS}

\textbf{Grouping.} \code{GroupBy} chooses one of three cursors by the \emph{shape} of its input.
\code{AllGroupBy} handles the no-group-key case, a single aggregate over the whole stream.
\code{SequentialGroupBy} assumes the input is already clustered by key (because a sort or an index
delivered it that way) and emits a group the moment the key changes, holding only one group's state at a
time. \code{HashGroupBy} is the general case: a \code{LinkedHashMap} from a hash-caching \code{Key} wrapper
to a \code{Grouping} accumulator, with a fast path that compares single grouping keys by a vectorised
\code{stringEquals} or a direct 64-bit integer equality rather than going through generic atomic
comparison. The accumulators are pluggable, so \code{count}, \code{sum}, \code{min}, \code{max}, and
\code{avg} are all just different \code{Grouping} implementations folded in as tuples arrive.

\textbf{Spilling.} Two operators bound their memory and spill to disk. \code{OrderBy} feeds its tuples to
a \code{TupleSort} with a memory budget of a quarter of the heap (\code{Runtime.maxMemory()/4}), and when
that budget is exceeded the sort writes sorted runs to disk and finishes as an external merge sort, so an
\code{order by} over data larger than memory still completes. \code{SpillableGroupBy} applies the same
quarter-heap budget to hash aggregation, as a radix hash-aggregation spill in the spirit of DuckDB and
PostgreSQL's \code{HashAgg}. It aggregates in memory until the table reaches the budget; thereafter a
tuple carrying a \emph{new} group is appended raw to one of 64 partition files chosen by a hash of its
grouping key, while a tuple whose group is \emph{already} resident keeps folding in at no extra memory.
A group therefore lives entirely in memory or entirely in one partition --- never split --- so once the
input drains, the resident groups are emitted and each partition is then re-aggregated, one at a time, by
the same routine; a partition that is itself too large repartitions again under a depth-rotated hash so
that the bits selecting the partition differ at each level. If nothing ever exceeds the budget no file is
created and the operator behaves exactly like \code{HashGroupBy}. The hot path is kept lean: on the common
existing-group hit it allocates no key wrapper (a single per-cursor lookup probe is reused and the key
hash is computed once, then reused both to probe the table and to choose the spill partition). That the
spilled and in-memory results are byte-identical, and that memory stays bounded under a deliberately tiny
budget, is checked end-to-end by \code{SpillableGroupBySpillTest}.

\begin{honesty}
\small\textbf{Residual.} The spill partitions by the 32-bit \code{hashCode} of the grouping key, so the
one shape it cannot bound is a single partition holding more \emph{distinct} keys than the budget admits
that \emph{also} collide in that hash. Depth-rotated repartitioning separates keys that merely share low
bits, but keys with identical 32-bit hashes cannot be split, and past a recursion-depth cap such a
partition is aggregated in memory. This is the standard limit of any hash-partition spiller and takes
adversarial input to reach; ordinary high-cardinality group-bys spill and complete.
\end{honesty}

\section{The SirixDB binding and temporal functions}

SirixDB implements that SPI as \code{JsonDBStore}/\allowbreak\code{JsonDBCollection}
(\pp{io/sirix/query/json/}), wrapping a Sirix \code{Database<JsonResourceSession>};
collections are time-aware (\code{getDocument(name, Instant)} opens a read-only transaction at
a point in time). Document nodes surface as a lazy \code{JsonDBItem} hierarchy wrapping a
cursor and materialising fields on access. The \code{jn:} family adds the temporal surface:
\code{jn:doc(db,res[,rev])}, the time-travel axes \code{jn:all-times}/\allowbreak\code{jn:past}/\allowbreak\code{jn:future},
\code{sdb:item-history}, and \code{jn:diff(rev1,rev2[,path])}, layering brackit's language directly
over the bitemporal store.

The temporal surface is far broader than those three functions. The \code{jn:} and \code{sdb:}
functions --- and the \code{xml:} functions for XML resources --- are \emph{added by SirixDB}, registered
by the \code{sirix-query} binding rather than being part of brackit: brackit supplies the XQuery/JSONiq
language and runtime, and SirixDB layers on the \code{jn:} (JSON), \code{xml:} (XML), and \code{sdb:}
(store and temporal metadata) function modules that bind it to the bitemporal store. The surface falls into
four groups, each \code{jn:} axis backed by a matching core temporal axis (\code{PastAxis},
\code{FutureAxis}, \ldots):

\dcap{The bitemporal function library}
\begin{diagramS}
opening / time-travel     jn:doc(coll,res[,rev]) · jn:open(coll,res,instant) · jn:open-revisions(…,from,to)
                          jn:valid-at(…,validTime) · jn:open-bitemporal(…) · jn:load · jn:store
temporal axes ($item)     jn:past · jn:future · jn:previous · jn:next · jn:first · jn:last
                          jn:first-existing · jn:last-existing · jn:all-times
history / diff            jn:diff(coll,res,r1,r2) · sdb:item-history($item)
node / revision metadata  sdb:nodekey · sdb:hash · sdb:path · sdb:revision · sdb:timestamp
                          sdb:author-id · sdb:author-name · sdb:child-count · sdb:descendant-count
                          sdb:valid-from · sdb:valid-to · sdb:is-deleted · sdb:select-item · sdb:select-parent
\end{diagramS}

\dcap{The temporal axes: one item navigated across the revisions it appears in}
\begin{diagram}
revisions where item I exists:   r3 ── r7 ── r12 ── r25 ── r40        ($item is at r12)

  jn:first(I)    = I@r3            jn:last(I) = I@r40
  jn:previous(I) = I@r7           jn:next(I) = I@r25
  jn:past(I)     = { r3, r7 }      jn:future(I) = { r25, r40 }
  jn:all-times(I) = { r3, r7, r12, r25, r40 }   (oldest to newest, via the change-tracking index)
\end{diagram}

So a single query pivots through time: \code{jn:all-times(jn:doc('db','res').field)} yields a field's
value in \emph{every} revision, \code{jn:open('db','res',\$instant)} opens a resource as-of a timestamp,
and \code{sdb:item-history(\$node)} returns the node in each revision that touched it (via the
change-tracking index of \S\ref{sec:changetracking}); \emph{valid-time} bounds are first-class too
(\code{sdb:valid-from}/\allowbreak\code{sdb:valid-to}, \code{jn:valid-at}, \code{jn:open-bitemporal}).

\begin{figure}[t]
\centering
\begin{tikzpicture}[font=\sffamily\footnotesize,>=Latex]
  \draw[->,slate] (0,0) -- (6.6,0) node[below right,font=\scriptsize]{valid time};
  \draw[->,slate] (0,0) -- (0,3.7) node[above left,font=\scriptsize]{transaction time};
  \foreach \x in {1,2,3,4,5} \draw[slate!22] (\x,0)--(\x,3.4);
  \foreach \y in {1,2,3} \draw[slate!22] (0,\y)--(6.3,\y);
  \fill[nodenew!18,draw=nodenew!70] (2,0.6) rectangle (5,1.7);
  \node[font=\scriptsize] at (3.5,1.15) {price 12 (as first recorded)};
  \fill[ember!14,draw=ember!70] (3,1.7) rectangle (5,3.1);
  \node[font=\scriptsize,color=ember!75!black] at (4,2.4) {price 19.99 (corrected)};
\end{tikzpicture}
\caption{Bitemporal data has two independent time axes. A fact occupies a rectangle: its
\emph{valid-time} extent (when it holds in the world) and its \emph{transaction-time} extent (the
revisions during which the system believed it). A later correction (the ember rectangle) supersedes part
of the original belief without erasing it, so a query can ask not only ``what was the price valid at
$t$'' but ``what did we \emph{believe} the price valid at $t$ to be, as of revision $r$''. SirixDB
provides the transaction axis through its revisions and the valid axis through valid-from/valid-to
bounds.}
\label{fig:bitemporal}
\end{figure}

\section{A worked bitemporal query: both axes at once}
\label{sec:bitemporal-worked}

The two-axis claim is easiest to believe from a trace. Configure a resource with valid-time paths
(\code{ResourceConfiguration.Builder.validTimePaths("validFrom","validTo")}, or the conventional-path
shorthand), so every top-level object carries an ISO-8601 \code{validFrom}/\allowbreak\code{validTo} pair that the
bitemporal functions read. Now record a price, then \emph{correct} it in a later revision, exactly the
two rectangles of Figure~\ref{fig:bitemporal}:

\dcap{Two revisions of a price fact: an original belief, then a correction}
\begin{diagramS}
resource "prices" (validTimePaths = validFrom/validTo); top level is an array of price facts

revision r5  committed at transaction time 2026-01-10
   [ { sku:"A", price:12,    validFrom:"2026-02-01", validTo:"2026-12-31" } ]

revision r9  committed at transaction time 2026-07-15   (a correction is recorded)
   [ { sku:"A", price:12,    validFrom:"2026-02-01", validTo:"2026-12-31" },
     { sku:"A", price:19.99, validFrom:"2026-06-01", validTo:"2026-12-31" } ]
\end{diagramS}

\noindent Two queries differ only in the \emph{transaction} axis; both apply the same valid-time
filter, an inclusive containment test \code{validFrom $\le$ t $\le$ validTo} evaluated by
\code{BitemporalFilterIter} over each top-level object:

\dcap{\texttt{valid-at} reads the latest belief; \texttt{open-bitemporal} reads a past belief}
\begin{diagramS}
Q1  jn:valid-at("mydb","prices", xs:dateTime("2026-08-01T00:00:00Z"))
      transaction axis → latest revision r9;   valid axis → 2026-08-01
      filter r9's array:  12    in [2026-02-01, 2026-12-31] ?  yes → keep
                          19.99 in [2026-06-01, 2026-12-31] ?  yes → keep
      result: BOTH facts (today we believe both intervals cover 2026-08-01)

Q2  jn:open-bitemporal("mydb","prices",
                       xs:dateTime("2026-03-01T00:00:00Z"),   ← transaction time
                       xs:dateTime("2026-08-01T00:00:00Z"))   ← valid time
      transaction axis → revision at-or-before 2026-03-01  = r5  (r9 is later, unseen)
      valid axis → 2026-08-01
      filter r5's array:  12 in [2026-02-01, 2026-12-31] ?  yes → keep
      result: ONLY {price:12}   (as of March we had not yet recorded the correction)
\end{diagramS}

The contrast is the whole point of bitemporality. \code{jn:valid-at} answers ``what do we \emph{now}
believe was valid on 2026-08-01'' and returns both the original and the corrected fact; the same valid
instant under \code{jn:open-bitemporal} with a March transaction time returns only the price~12 fact,
because the correction committed in r9 (July) did not exist in the snapshot a March-era reader would
have seen. The transaction axis is served by the immutable revision trie (the \code{Instant} resolves
to revision~r5 by the binary search of Property~\ref{prop:histcost}), and the valid axis by a streaming
filter over that revision's top-level items; neither rewrites the other's data. \code{sdb:valid-from}
and \code{sdb:valid-to} expose a selected fact's bounds as accessors for use inside an enclosing FLWOR,
so an application can compute interval overlaps, gaps, or as-of-then-vs-now differences without leaving
the query language. The filter keys on top-level objects with at least one bound present: an absent \code{validTo} is read
as ``valid from \code{validFrom} onward'' and an absent \code{validFrom} as ``valid up to \code{validTo}'';
a fact with neither bound, or one nested below the top level, is not selected.

\textbf{How valid time is stored, and an honest asymmetry.} The valid-time bounds are not a separate
on-disk structure: a resource declares two JSON paths
(the \code{ResourceConfiguration} builder's \code{validTimePaths("validFrom","validTo")}, or
\code{useConventionalValidTimePaths()} for the \code{\_validFrom}/\code{\_validTo} convention), the pair
is serialised into the resource config so it survives a reopen, and the bounds themselves are just
\emph{ordinary JSON fields} on the data nodes. That choice produces a real asymmetry between the two
axes, worth stating plainly: the transaction axis is indexed for free by the revision structure, while
the valid axis is accelerated only when its index is explicitly built for the resource (below):

\dcap{The two time axes have different cost shapes}
\begin{diagramS}
transaction axis (system belief)            valid axis (world truth)
  served by the immutable revision trie       served by validFrom/validTo JSON fields
  Instant→revision: O(log R) Eytzinger        stab at instant: O(h) RI-tree (scan fallback)
  indexed, position-independent               opt-in VALIDTIME interval index + CAS narrowing
\end{diagramS}

\begin{honesty}
\small\textbf{Status: both axes are now indexed.} The transaction axis uses the $O(\log R)$
commit-timestamp search of Property~\ref{prop:histcost}; the valid axis is served by a persistent
\code{VALIDTIME} \textbf{Relational-Interval-Tree}. \code{jn:valid-at} and \code{jn:open-bitemporal}
stab the RI-tree at the query instant in $O(h)$, streaming every record object whose interval contains
it, then \textbf{re-verify} each candidate against its exact \code{validFrom}/\code{validTo} bounds ---
provably the same set the linear scan returns (a monotone interval domain gives completeness; the
exact-instant recheck removes false positives). When no interval index exists but a \textbf{CAS index}
covers a valid-time path, the functions instead narrow candidates with a range scan and verify by
reading; absent both, they fall back to the original $O(n)$ scan of the top-level items. Open-ended
intervals are supported: an absent (or unparseable) \code{validTo} means ``valid from \code{validFrom}
onward'' and an absent \code{validFrom} means ``valid up to \code{validTo}'' (neither bound = no interval);
the interval-index writer maps a null bound to the domain min/max and registers the record identically, so
all paths agree. Two constraints still follow from the field-based design: only top-level records are
filtered (a fact nested deeper is not seen), and the exact bounds are parsed from ISO-8601 at verification
time. The bitemporal
\emph{model}, \emph{query surface}, \emph{and} valid-axis \emph{acceleration} are all complete --- the
earlier ``acceleration is future work'' caveat shipped in \code{1.0.0-beta3}.
\end{honesty}

\section{The binding internals: lazy items and the item factory}
\label{sec:bindinginternals}

The bridge between brackit's value model and SirixDB's node store is a thin layer of \emph{lazy} wrappers
(\pp{io/sirix/query/json/}). It is worth a look because it is what makes a query over a 100\,M-node
resource run on \emph{one} read transaction rather than opening one per item.

\dcap{From a Sirix node to a brackit json-item, by kind (\texttt{JsonItemFactory})}
\begin{diagramS}
rtx.getKind()                       wrapper (each carries only the nodeKey + the shared rtx + collection)
  OBJECT, OBJECT_NAMED_OBJECT    →  JsonDBObject     (fields materialised lazily on get(QNm), then cached)
  ARRAY,  OBJECT_NAMED_ARRAY     →  JsonDBArray      (elements materialised once on at(i)/values(), cached)
  *_STRING / STRING_VALUE        →  AtomicStrJsonDBItem
  *_BOOLEAN / BOOLEAN_VALUE      →  AtomicBooleanJsonDBItem
  *_NULL / NULL_VALUE            →  AtomicNullJsonDBItem
  *_NUMBER / NUMBER_VALUE        →  Int/Long → IntNumericJsonDBItem,  Double → DblNumericJsonDBItem,
                                    Float → FltNumericJsonDBItem,  BigDecimal/BigInteger → DecNumericJsonDBItem
\end{diagramS}

A \code{BasicJsonDBStore} implements brackit's store SPI, mapping each \code{Database} to a
\code{JsonDBCollection}; the collection is \emph{time-aware} (\code{getDocument(name, Instant)} opens a
read-only transaction at that point in time, \code{getDocument(name, rev)} at a revision), which is how
\code{jn:open} and \code{jn:doc(\ldots,rev)} resolve. From a cursor position the \code{JsonItemFactory}
dispatches on \code{NodeKind} to the wrapper above. The decisive properties are two. First, every wrapper
stores only the \textbf{node key} plus a reference to the \emph{shared} read transaction, and calls
\code{moveRtx()} to re-seek that one cursor before touching Sirix state, so a result set of a million
items rides on a single transaction. Second, structure is \textbf{materialised lazily}: \code{JsonDBObject.get(QNm)}
resolves a field only when asked, by a \code{FilterAxis} of a \code{ChildAxis} under a \code{JsonNameFilter}
(short-circuiting on the path summary when the key cannot exist), and caches it; \code{JsonDBArray} fills
and caches its element list on the first \code{at(i)}. The per-kind numeric wrappers are not cosmetic:
because \code{IntNumericJsonDBItem} implements \code{LonNumeric} and \code{DblNumericJsonDBItem} implements
\code{DblNumeric}, a document-sourced number dispatches to the right arithmetic and comparison path, the
fix behind the int-vs-double comparison correctness discussed for the node model (Ch.~\ref{ch:pds}).

A navigation trace shows the laziness and the one shared cursor together. Evaluate
\code{jn:doc('db','res').store.books[0].price} over a document rooted at node key~1, on a single read
transaction \code{R}:

\dcap{Lazy item navigation: one shared rtx, only the touched path read}
\begin{diagramS}
jn:doc(...)         -> JsonDBObject(nodeKey 1)        wraps the root; holds { nodeKey 1, rtx R }
  .store            get("store"): moveRtx(1); ChildAxis + JsonNameFilter -> nodeKey 2
                    -> JsonDBObject(2)         cached on the parent
  .books            get("books"): moveRtx(2); filtered child walk -> nodeKey 5  (NAMED_ARRAY)
                    -> JsonDBArray(5)          cached
  [0]             at(0): moveRtx(5); first child -> nodeKey 6  (the [0] object)
                    -> JsonDBObject(6)         element list cached
  .price            get("price"): moveRtx(6); filtered child walk -> nodeKey 9  (NAMED_NUMBER)
                    -> IntNumericJsonDBItem(9)

every step re-seeks the SAME cursor R (never a new transaction) and reads only this root-to-leaf path;
a million-item result set reuses R a million times — one transaction, not a million
\end{diagramS}

\textbf{The update bridge.} Writes run the same wrappers in reverse. A \code{JsonDBItem} lazily upgrades
to a writer through \code{getOrCreateWriteTrx} (taking the single write permit, and \code{revertTo}-ing
first if the target is an older revision, the branch of \S\ref{sec:worked}). A same-type leaf change is
an in-place \code{setStringValue}/\allowbreak\code{setNumberValue}; a type change is a
\code{replaceObjectRecordValue} followed by an \code{insertSubtree} of the new value; array and object
edits route through \code{JsonItemSequence.insert}, which dispatches atomic values to
\code{insert\ldots Value} calls and nested structures to \code{insertSubtreeAsLastChild} with auto-commit
suppressed. The commit decision lives one level up, in \code{SirixQueryContext.applyUpdates}: under the
\code{AUTO} strategy it maps each pending update's target back to its owning \code{JsonNodeTrx},
deduplicates by transaction id, and commits each once, then invalidates that resource's statistics, so
one updating query becomes one atomic revision per touched resource.

\section{The sequence and atomic type model}
\label{sec:typemodel}

Underneath the wrappers is brackit's JDM (JSON/XQuery Data Model, \pp{io/brackit/query/jdm/}), the type
lattice everything in the engine speaks. A \code{Sequence} is an ordered, possibly empty, possibly
lazy run of \code{Item}s; it is itself an \code{Expr}, so an expression and its value share a type. An
\code{Item} is a single value, either a \code{StructuredItem} (a node, or a JSON \code{Object}/\code{Array})
or an \code{Atomic}. Iteration is the \code{Iter} contract, \code{next()} returns items until null and
must be \code{close()}d, with an optional \code{split()} for the parallel operators of
\S\ref{sec:setoriented}, which is why a query over a huge sequence never materialises it.

\dcap{The brackit type lattice (JDM) and the numeric promotion hierarchy}
\begin{diagramS}
Sequence (ordered, lazy)  ─ iterate() → Iter.next()/close()/split()
  └ Item
      ├ StructuredItem ─ Node | json Object | json Array
      └ Atomic ─ Str, Bool, Null, DateTime/Date/Time/Dur…, QNm, B64/Hex, AnyURI, and Numeric:

numeric promotion interfaces (cmp / arithmetic dispatch on the INTERFACE, not the concrete class):
   Numeric
     ├ DblNumeric  → Dbl            ├ FltNumeric → Flt
     └ DecNumeric  → Dec            └ IntNumeric ← LonNumeric ← {Int32, Int64};  Int (arbitrary precision)
\end{diagramS}

The \code{Atomic} contract carries two comparisons: \code{cmp}, the value comparison with type promotion
(an \code{xs:integer} against an \code{xs:double} promotes), and \code{atomicCmp}, a stable \emph{total}
order used for \code{distinct} and indexing. The subtlety that bit the engine, and that the per-kind
\code{*NumericJsonDBItem} wrappers above exist to respect, is that numeric dispatch keys on the
\emph{interface} (\code{DblNumeric}, \code{IntNumeric}, \ldots), not the concrete class: a
document-sourced double wrapped so that it implements \code{DblNumeric} compares and computes correctly,
whereas matching only the concrete \code{Dbl} would have thrown \code{XPTY0004} on a perfectly valid
\code{3} vs \code{3.7} comparison. The lattice is shared verbatim between XML and JSON; the JSON additions
are only \code{Object}, \code{Array}, and the \code{Null} item.

\dcap{Why interface dispatch matters: \texttt{order by \$.rating} over mixed 3 and 3.7}
\begin{diagramS}
rating column holds 3 (an integer) on half the rows, 3.7 (a double) on the rest

concrete-class dispatch (the bug):  cmp(Int64 3, Dbl 3.7)
   Int64.cmp matches the operand's CONCRETE class → no branch for Dbl → throws XPTY0004
   ("cannot compare xs:integer with xs:double") — order by 500's on perfectly valid data

interface dispatch (the fix):       cmp(Int64 3, x where x instanceof DblNumeric)
   Int64.cmp tests the INTERFACE DblNumeric → promotes 3 → 3.0, compares 3.0 < 3.7 → ordered
   document-sourced doubles wrap as DblNumericJsonDBItem (implements DblNumeric, not the concrete Dbl)
\end{diagramS}

The example is not hypothetical: it is the demo's \code{rating} field, and the fix, dispatching
\code{Int32}/\allowbreak\code{Int64} arithmetic and comparison on the \code{DblNumeric}/\allowbreak\code{FltNumeric} interfaces
the way the base \code{Int} already did, lives in brackit and is mirrored by the per-kind
\code{*NumericJsonDBItem} wrappers (\S\ref{sec:bindinginternals}) so that a number read from a document
carries the right interface to the comparison.

The same interface dispatch governs \emph{arithmetic}, not only comparison, which a mixed-type aggregate
shows:

\dcap{Interface dispatch in arithmetic: \texttt{sum(\$.value)} over mixed 5 (int) and 3.5 (double)}
\begin{diagramS}
value holds 5 (Int32) on some rows, 3.5 (Dbl) on others

  add(Int32 5, x):  Int32.add tests  x instanceof DblNumeric  -> promote 5 -> 5.0 -> 5.0 + 3.5 = 8.5 (Dbl)
  the running sum WIDENS to Dbl the moment a double is seen; an all-integer column stays Int (no widening)
  concrete-class dispatch would seek Int32.add(Dbl), find no such branch, and throw XPTY0004 — the same
  trap as the ordering case, on the arithmetic path
\end{diagramS}

\noindent So the promotion lattice is not a comparison-only nicety: every binary numeric operator
(\code{+}, \code{*}, \code{cmp}, \ldots) widens along \code{Int $\to$ Dec $\to$ Dbl} by testing the
operand's interface, and the per-kind document wrappers exist precisely so a value read from storage
arrives carrying that interface.

\section{JSONiq updates and the commit model}

SirixDB is writable \emph{from the query language}: brackit implements the \textbf{JSONiq Update
Facility} and SirixDB applies the resulting operations to a \code{JsonNodeTrx}. Five statement forms are
supported (the XML XQuery Update Facility, with \code{insert}/\allowbreak\code{delete}/\allowbreak\code{rename}/\allowbreak\code{replace
node \ldots as first into | after | before \ldots}, running in parallel):

\dcap{JSONiq update statements}
\begin{diagramS}
insert json   {"c":3}             into $obj          insert json (4,5) into $arr at position 2
replace json value of  $obj.price with 19.99
delete json   $obj.field                             delete json $arr[1]
rename json   $obj.oldKey         as "newKey"
append json   $arr                into "x"
\end{diagramS}

Updates are \textbf{accumulated, not applied eagerly}: every mutating function marks itself
\code{isUpdating}, and the pending operations collect in the query context's update list. At end of
evaluation \code{applyUpdates} runs them against the in-memory items and, \textbf{when the commit
strategy is \code{AUTO}}, maps each target back to its owning \code{JsonNodeTrx} (deduplicated by
transaction) and commits, so one updating query becomes \emph{one atomic revision per touched
resource}, after which that resource's statistics and plan cache are invalidated. Under the
\code{EXPLICIT} strategy the caller commits via \code{sdb:commit(\$doc[, msg[, dateTime]])} (returning
the new revision number; \code{sdb:rollback(\$doc)} discards). This is exactly the batched-atomic-commit
model the GUI uses (Ch.~\ref{ch:gui}).

Underneath, this is brackit's standard \textbf{pending-update-list} (PUL) model from the XQuery Update
Facility (\pp{io/brackit/query/update/}), which SirixDB inherits unchanged up to the final apply.

\dcap{The two-phase update model (\texttt{UpdateList}); SirixDB binds only the last step}
\begin{diagramS}
Phase 1 — collect (during evaluation):
   each updating expression marks isUpdating, builds an UpdateOp primitive, and calls
   ctx.addPendingUpdate(op) — then returns the EMPTY sequence (an update yields no value)
   primitives: insertInto/Before/After, replaceNode, replaceValue, rename  (XML)
               insertIntoArray/Record, deleteArrayIndex/recordField, replaceArray/RecordValue, rename (JSON)

Phase 2 — applyUpdates() at end of query:
   1 sort the PUL by OpType priority   (replaceNode first … delete last)
   2 conflict check: two renames / two replaces of the SAME target → XQUF error
   3 collect delete targets; skip property updates to a to-be-deleted node (XQUF §3.2.1)
   4 apply each primitive in order → (SirixDB) maps each target to its JsonNodeTrx, commits once per trx
\end{diagramS}

The design point is that an updating expression has \emph{no} side effect when it runs: it appends a
primitive to the list and yields nothing, so order of evaluation never changes the outcome. The list is
then ordered by a fixed primitive priority, checked for conflicting edits to the same node (duplicate
renames or replaces raise a static-typed error), and applied; a property update to a node that another
primitive deletes is silently dropped, per the specification. Only the final step is SirixDB's: where
brackit's reference applies primitives to its in-memory DOM, SirixDB's \code{applyUpdates} routes each
primitive's target to the owning \code{JsonNodeTrx} (\S\ref{sec:bindinginternals}) and commits, turning a
spec-compliant PUL into one atomic SirixDB revision per resource.

A five-update query makes both phases, and the delete-last rule, concrete:

\dcap{A multi-update query: collected in eval order, applied in priority order}
\begin{diagramS}
one query, 5 updating expressions over the same object $o:
   insert json {"cost":50} into $o     replace json value of $o.price with 12    delete json $o.price
   rename json $o.id as "sku"          replace json value of $o.qty with 99

Phase 1 (collect, in evaluation order) — each yields () and appends a primitive:
   [ insert(cost), replaceValue(price,12), delete(price), rename(id->sku), replaceValue(qty,99) ]

Phase 2 (applyUpdates):
   1  sort by OpType priority:   replaceValue .. rename .. insert .. delete (LAST)
   2  collect delete targets = { price }
   3  SKIP replaceValue(price,12): its target is to-be-deleted (XQUF §3.2.1) — silently dropped
   4  apply:  replaceValue(qty,99) ; rename(id->sku) ; insert(cost,50) ; delete(price)
   (a SECOND rename of id, or a second replace of qty, would instead raise an XQUF conflict error)

   -> SirixDB maps every surviving target to $o's JsonNodeTrx and commits ONCE = one atomic revision
\end{diagramS}

\noindent Two things the trace shows that the prose only states: deletes are applied \emph{last} so that a
replace or insert can never touch an already-removed node, and a value update aimed at a field some other
primitive deletes is dropped rather than erroring --- both straight from XQUF~§3.2.1, and both invisible in
the source order because Phase~1 has no side effects.

\section{Diffing and the FMES diff-import}
\label{sec:diff}

Where change tracking (\S\ref{sec:changetracking}) answers \emph{where} and \emph{when}, the
\textbf{diff} answers \emph{what}: the exact transformation between two given revisions.
\code{jn:diff(rev1,rev2[,path])} returns a JSON edit script
(\code{\{database,resource,old-revision,new-revision,diffs:[\ldots]\}}) of insert/delete/update/replace
operations. The engine (\code{DiffFactory}) is \textbf{hash-guided}: it descends the two revision trees
in lock-step and, wherever a node's 64-bit rolling hash (Ch.~\ref{ch:pds}) matches, emits
\code{SAMEHASH} and \textbf{prunes the whole subtree}: an unchanged subtree is one hash comparison, never
descended, so deep structure below an unchanged node costs nothing (with \code{HashType.NONE} it degrades
to a full node-by-node walk). This is a \emph{vertical} prune; the lock-step walk still enumerates
\emph{siblings} child by child, so a wide unchanged fan-out is scanned even when little of it changed,
which the measured cost shape of \S\ref{sec:diffbench} makes precise and is one reason the stored-diff
short-circuit exists. For \emph{adjacent} revisions the engine additionally short-circuits to a per-commit
diff file precomputed at write time (default-on \code{storeDiffs}), reading the stored delta rather than
recomputing it. Change
tracking and diffing are thus complementary: the former is the cheap lookup for history and temporal
navigation, the latter the means to compute or export the precise edit script between two \emph{arbitrary}
revisions.

\dcap{Hash-guided diff: descend only where the rolling hashes differ}
\begin{diagramS}
rev 7                              rev 8 (only book[0].price changed)
   OBJECT  h=A                        OBJECT  h=A'      A≠A' → descend
     books h=B                          books h=B'      B≠B' → descend
       [0]   h=C                          [0]   h=C'    C≠C' → descend
         title h=D   ── SAMEHASH ──         title h=D   D=D  → prune subtree
         price h=E   ──  UPDATE   ──         price h=E'  E≠E' → emit update
       [1]   h=F     ── SAMEHASH ──         [1]   h=F    F=F  → prune the whole subtree

edit script = [ { update: price at /books/[0]/price } ]    cost ∝ the changed region, not document size
\end{diagramS}

\begin{figure}[t]
\centering
\begin{tikzpicture}[font=\sffamily\small,>=Latex,
  diff/.style={draw=ember!70,fill=ember!10,rounded corners=2pt,align=center,inner sep=3pt,minimum height=0.6cm},
  same/.style={draw=slate!45,fill=slate!8,rounded corners=2pt,align=center,inner sep=3pt,minimum height=0.6cm,text=slate!70!black},
  upd/.style={draw=nodenew,fill=nodenew!14,rounded corners=2pt,align=center,inner sep=3pt,minimum height=0.6cm},
  e/.style={->,slate}]
  \node[diff] (root)  at (4.0,3.0)  {root \quad $h{:}\,A\neq A'$ --- descend};
  \node[diff] (books) at (4.0,1.8)  {books \quad $B\neq B'$ --- descend};
  \node[diff] (b0)    at (2.1,0.6)  {book[0] \quad $C\neq C'$ --- descend};
  \node[same] (b1)    at (6.7,0.6)  {book[1] \quad $F{=}F$ --- SAMEHASH (prune)};
  \node[same] (title) at (0.9,-0.7) {title \quad $D{=}D$ --- prune};
  \node[upd]  (price) at (3.4,-0.7) {price \quad $E\neq E'$ --- UPDATE};
  \draw[e] (root)--(books); \draw[e] (books)--(b0); \draw[e] (books)--(b1);
  \draw[e] (b0)--(title); \draw[e] (b0)--(price);
\end{tikzpicture}
\caption{Hash-guided diff. The two revisions' trees are descended in lock-step, comparing each node's stored
64-bit rolling hash. Where the hashes differ (\emph{ember}) the walk descends; where they match
(\emph{grey}) it emits \code{SAMEHASH} and prunes the entire subtree unread, so \code{book[1]} and the
unchanged \code{title} cost one comparison each. Only the genuinely changed leaf, \code{price}, is emitted
as an \code{UPDATE}. The cost is thus the changed region's \emph{spine} plus the siblings enumerated along
it, never the whole document --- the shape \S\ref{sec:diffbench} measures.}
\label{fig:diff}
\end{figure}

Both the on-demand diff and the stored per-commit deltas share \textbf{one JSON change format}
(\code{JsonDiffSerializer}), the sense in which SirixDB does change tracking ``in our format'':

\dcap{The stored JSON change format (\texttt{JsonDiffSerializer})}
\begin{diagramS}
{ "database": "...", "resource": "...", "old-revision": N, "new-revision": M,
  "diffs": [
    { "insert":  { nodeKey, insertPositionNodeKey, insertPosition: "asFirstChild"|"asRightSibling",
                   path, deweyID, depth, type: "jsonFragment"|"string"|"number"|…, data } },
    { "delete":  { nodeKey, path, deweyID, depth } },
    { "update":  { nodeKey, path, deweyID, depth, name | (type, value) } },
    { "replace": { oldNodeKey, newNodeKey, path, deweyID, depth, type, data } } ] }
\end{diagramS}

A metadata envelope (\code{database}, \code{resource}, \code{old}/\allowbreak\code{new-revision}) wraps a
\code{diffs} array of single-key operation objects. Every operation carries the node key and, from the
path summary, the field \code{path} with array positions resolved to concrete indices
(\code{/store/book/[3]}); when DeweyIDs are stored it also carries the \code{deweyID} and tree
\code{depth}. An \code{insert} adds its \code{insertPosition} (\code{asFirstChild} or
\code{asRightSibling}) plus a \code{type}/\allowbreak\code{data} payload: a scalar value, or a
\code{jsonFragment} (the serialised subtree); fused \code{OBJECT\_NAMED\_*} records are emitted as a
\code{jsonFragment} so a diff exposes the same external shape as a legacy key{+}value pair. This is the
format \textbf{stored at commit} (default-on \code{storeDiffs}), returned by \code{jn:diff} and the REST
\code{/diff} endpoint, and merged client-side by the GUI (Ch.~\ref{ch:gui}), so ``what changed'' is
itself a versioned, queryable JSON document, not merely an in-memory computation.

\dcap{A concrete edit script: \texttt{jn:diff('db','store', 5, 6)}}
\begin{diagramS}
revisions 5 → 6: book[0].price changed 12 → 19.99, and a "discount" field was inserted

{ "database":"db", "resource":"store", "old-revision":5, "new-revision":6,
  "diffs": [
    { "update": { "nodeKey":7,  "path":"/store/book/[0]/price",
                  "type":"number", "value":19.99 } },
    { "insert": { "nodeKey":42, "insertPositionNodeKey":7, "insertPosition":"asRightSibling",
                  "path":"/store/book/[0]/discount", "type":"number", "data":0.1 } } ] }
\end{diagramS}

A third operation imports an \emph{external} document \emph{as} a diff. The \textbf{FMES} service
(\code{FMSEImport}, \pp{io/sirix/diff/algorithm/fmse/}) implements the Fast-Match\,/\,Edit-Script
tree-correction algorithm of Chawathe, Rajaraman, Garcia-Molina and Widom~\cite{chawathe}. Its job is to fold a
new version of a document into the store as a \emph{structural delta} that keeps the identity, the
stable node keys, of everything that did not change, instead of re-shredding the document and losing
that identity. The pipeline has three steps. First the new document is shredded into a throwaway
resource (\code{FMSEImport.shredder}), so both versions are now SirixDB trees with cursors. Second, a
\emph{matching} is computed between the stored (old) tree and the new one. Third, the algorithm walks
the new tree and applies a minimal edit script \textbf{in place against the stored tree} through an
ordinary write transaction, then commits, so the import lands as one new revision that \code{jn:diff}
can later reproduce.

\textbf{Phase~1, the fast match}, pairs nodes bottom-up and label-first. Nodes are first bucketed by
\code{NodeKind} (a \code{LabelFMSEVisitor} pass), and only nodes of the same kind are ever paired.
\emph{Leaves} are matched first: two leaves match when their values are similar enough, scored by a
Levenshtein ratio (\code{Levenshtein}, via \code{FMSENodeComparisonUtils.calculateRatio}, which
substitutes a character-frequency \emph{quick ratio} for strings longer than 50 characters) that must
exceed the threshold \code{FMESF}~$=0.5$. The JSON comparator scores typed values directly (numbers by
numeric value, booleans and nulls exactly) rather than stringifying them. Ties are broken by the path
summary, where equal root-to-leaf paths force a match and unequal-length paths forbid one, and for XML
by an optional \code{id}-attribute fast path that forces a match when ancestor ids agree. \emph{Inner}
nodes are then matched by structural overlap: two inner nodes match when the fraction of their
descendants already matched to one another, $\lvert\text{common}\rvert / \max(\lvert x\rvert,\lvert
y\rvert)$, reaches \code{FMESTHRESHOLD}~$=0.5$. Within each same-kind bucket the pairing itself is a
longest-common-subsequence pass (Myers'~\cite{myers} $O(ND)$ diff, \code{Util.longestCommonSubsequence}) followed by
a greedy sweep over whatever the LCS left unpaired. The result is a \code{Matching}: two directional
maps, old to new and back, plus a subtree-containment index that answers ``how many matched descendants
do these two share'' in time proportional to the subtree.

\textbf{Phase~2 builds and applies the edit script} in Chawathe's order. It iterates the new tree
breadth-first and, for each node $x$ under parent $y$, does exactly one of \textbf{insert} (when $x$ has
no partner, copy its subtree in with \code{copySubtreeAs\ldots}), \textbf{update} (when the partner's
value or key name differs, \code{setValue} or \code{setObjectKeyName}), or \textbf{move} (when $x$'s
partner sits under the wrong parent, \code{moveSubtreeTo\ldots}, which keeps the subtree's node keys).
After each node it \textbf{aligns children}: a second LCS, this time over the \emph{matched} children of
the pair, decides which children are already in relative order and which must be repositioned. A final
post-order pass over the \emph{old} tree (\code{DeleteFMSEVisitor}) removes every node that never found
a partner.

\dcap{FMES import: fast match, then a minimal, identity-preserving edit script}
\begin{diagram}
new document  ──shred──▶  throwaway resource ─┐
                                              ├──▶ fast match ──▶ Matching (old/new maps)
stored resource (old revision) ───────────────┘                       │
                                                                      ▼
  edit script, breadth-first over the new tree:
    no partner                     ──▶ insert   copySubtreeAs*
    value or key-name differs      ──▶ update   setValue / setObjectKeyName
    partner under the wrong parent ──▶ move     moveSubtreeTo*   (keeps node keys)
    children mis-ordered           ──▶ align    LCS over matched children, move the rest
  then, post-order over the old tree:
    no partner                     ──▶ delete
                                                                      │
                                                           wtx.commit() ──▶ one new revision
\end{diagram}

Identity is preserved exactly where it should be: a matched node is updated or moved in place and keeps
its node key, an inserted subtree is fresh material copied from the new tree, and a deletion is old
material with no partner. \code{FMES} is wired to the function \code{xml:import(\$coll, \$res, \$doc)}
(registered under both the \code{xml:} and \code{sdb:} libraries) and is reachable from a small CLI as
well. A full JSON port exists (\code{JsonFMSE}, sharing its matching core \code{FMSEAlgorithm} with the
XML driver), but it is not yet bound to any import function, so as a user-facing feature FMES is
XML-only today. Two implementation notes are worth recording. The paper's ``in-order'' traversal is
realised as \emph{post-order}, the code noting that the paper's wording is wrong for non-binary trees.
And the JSON delete visitor does not yet handle the two fused structural record kinds
(\code{OBJECT\_NAMED\_OBJECT} and \code{OBJECT\_NAMED\_ARRAY}), so an unmatched named container is
descended into rather than removed; that gap is one reason the JSON path remains test-only.

\dcap{A worked FMES import: an edited document becomes a structural delta}
\begin{diagramS}
stored (old)          new document        FMES decision
  <a>                   <a>                 match a↔a (same kind, overlapping descendants)
   <b id=1>x</b>         <b id=1>x</b>      match b↔b via the id fast-path → unchanged
   <c>y</c>              <c>Y</c>           match c↔c, value differs → UPDATE c := "Y"
                        <d>z</d>           no partner → INSERT <d>z</d> right of c
   <e>w</e>                                no partner in new → DELETE <e>

edit script: update(c), insert(d), delete(e).  b keeps its key; the import is one structural delta
\end{diagramS}

\section{Index-aware path rewriting (and a correctness bug worth studying)}

The optimiser's centrepiece is the path-index rewrite (\code{AbstractJsonPathWalker} and
subclasses \code{JsonPathStep}, \code{JsonCASStep}, \code{JsonObjectKeyNameStep}). It extracts
the JSON access path of a \code{jn:doc(...)[...]} expression, matches it against the path
summary, and rewrites the subtree into an \code{IndexExpr} backed by a PATH or CAS index when
applicable.

A subtle, historically real bug: for \code{.store.book[][?\$\$.price gt 10]}, the rewrite
fuses the scalar predicate leaf (\code{price}) into the path. The path-summary lookup may
search for \code{price} as an \code{OBJECT\_NAMED\_OBJECT} and find nothing, even though
\code{price} exists as an \code{OBJECT\_NAMED\_NUMBER}. Treating ``no match of the assumed
kind'' as ``the path is empty'' once \textbf{destructively replaced a live subtree with the
empty sequence}, silently dropping results. The fix checks segment existence \emph{regardless
of kind} before any such replacement. The lesson generalises: index rewrites must distinguish
``this path genuinely yields nothing'' from ``my kind assumption was wrong'', or they become
silent-wrong-result generators.

Rewrites are gated by a cost flag so the optimiser can prefer a sequential scan when an index
would not pay.

\section{The cost-based optimiser: the pipeline}
\label{sec:optpipeline}

\code{SirixOptimizer} extends brackit's \code{TopDownOptimizer} with a ten-stage pipeline run inline,
each stage rewriting the AST in place, bounded by a \textbf{50\,ms circuit-breaker}
(\code{OPTIMIZATION\_TIMEOUT\_MS}) that returns the best partial plan if a stage overruns:

\dcap{The optimiser pipeline (\texttt{SirixOptimizer})}
\begin{diagramS}
1 JqgmRewriteStage      fixpoint (≤10 iters): join fusion, predicate-into-join, predicate-into-access
2 CostBasedStage        annotate bindings: PATH_CARDINALITY, TOTAL_NODE_COUNT, PREFER_INDEX + costs
3 JoinReorderStage      DPhyp / GOO join reordering over a JoinGraph
4 MeshPopulationStage   build the mesh: seq-vs-index and join-order plan alternatives
5 MeshSelectionStage    pick the cheapest alternative per equivalence class; swap join children
6 IndexDecompositionStage   index-driven join restructure / intersection joins
7 CostDrivenRoutingStage    propagate PREFER_INDEX=false → INDEX_GATE_CLOSED to descendants
8,9 Vectorized{Detection,Routing}   DISABLED — deferred to brackit's vectorised SPI (Ch.8)
10 IndexMatching        rewrite to real index scans (JsonCASStep / JsonPathStep / JsonObjectKeyNameStep)
\end{diagramS}

Stages~1 and 10 bracket the cost-based core: the first normalises the query graph (fusing adjacent
joins, pushing predicates toward the access path) to a fixpoint; the last turns the cost decisions into
actual \code{IndexExpr} scans through the path-rewrite walkers of the previous section.

One query through the pipeline ties the stages together.

\dcap{A query traced through the optimiser stages}
\begin{diagramS}
query:  for $b in jn:doc('db','store').book[] where $b.price gt 50 return $b.title

 1 JqgmRewrite     push the predicate toward the access path:  book[] [price gt 50] → title
 2 CostBased       card(book)=100 000 ; sel(price>50)=0.33 → ~33 000 ; a CAS index on price exists
                   indexCost ≈ 9e3  <  seqCost ≈ 1.1e5   → annotate PREFER_INDEX on the book binding
 3 JoinReorder     no joins in this query → skip
 4 MeshPopulation  alternatives for book[]:  { sequential scan ,  CAS-index scan }
 5 MeshSelection   cheapest per equivalence class → CAS-index scan
10 IndexMatching   rewrite to a JsonCASStep over (price path, value ≥ 50), deref the book, project title

result: a CAS range scan feeds ~33 000 books into the projection, not a 100 000-node document walk
\end{diagramS}

\section{Cardinality, selectivity, and histograms}

Base \textbf{cardinality} is read from the path summary: \code{SirixStatisticsProvider} sums
\code{PathNode.getReferences()} over the path-class records matching a binding's path, and the total
node count is the document root's descendant count. Cardinality \emph{propagates} through the tuple
pipeline: a \code{ForBind} multiplies by the binding cardinality, a \code{Selection} scales by its
predicate selectivity, a group-by divides by the distinct count, and a join multiplies by the join
selectivity.

\textbf{Selectivity} starts from defaults (\textbf{1\,\% equality, 33\,\% range, 25\,\% \code{like}})
and compounds exactly as a modern optimiser does: an \code{and} multiplies its children's
selectivities sorted ascending and \emph{exponentially damped},
$s_{(0)}\cdot s_{(1)}^{1/2}\cdot s_{(2)}^{1/3}\cdots$ (the SQL-Server backoff, which avoids
under-counting correlated conjuncts); an \code{or} uses inclusion--exclusion $P(A){+}P(B){-}P(A)P(B)$; a
\code{not} is $1-s$. When a \textbf{histogram} exists for a field (PostgreSQL-style: 64 buckets plus a
top-10 most-common-value list, equi-width or equi-depth) the estimate sharpens: an equality on a
most-common value uses its exact frequency, a non-MCV value uses
$(1-f_{\text{mcv}})/d_{\text{non-mcv}}$, and a range sums the overlapped buckets by linear interpolation
plus the exact frequencies of any in-range MCVs. Histograms are built by an \code{ANALYZE}-like
stride-sampler (10\,k values), but \textbf{deferred}: a first-seen predicate plans on the default
selectivity and the histogram is collected \emph{after} the writing transaction closes (so the
optimiser never opens a session inside an active write), benefiting later queries.

A range estimate shows the bucket arithmetic the description above compresses. Take a histogram on
\code{/catalog/book/[]/price} --- 64 equi-width buckets over $[0,640]$ (width~10), one million values ---
and the predicate \code{price between 125 and 200}:

\dcap{Histogram range selectivity: sum whole buckets, interpolate the partial one}
\begin{diagramS}
predicate: price between 125 and 200       bucket b covers [10b, 10b+10),  bucket(p) = floor(p/10)

  bucket 12  [120,130)  count 20 000   in range from 125 -> interpolate (130-125)/10 = 0.5  -> 10 000
  buckets 13..19  [130,200)  fully in range,  counts sum to                                    175 000
  bucket 20  [200,210)  excluded (200 is the open upper edge)                                        0
  (no most-common value lies in [125,200], so no exact MCV frequency is added)

  matches ~ 10 000 + 175 000 = 185 000 of 1 000 000   ->   selectivity 0.185
  vs the 0.33 range DEFAULT: the histogram tightens the estimate, and the cost model re-decides on 0.185
\end{diagramS}

\noindent The partial bucket at the low edge is scaled by the fraction of its width that the range covers
(linear interpolation, which assumes values are uniform \emph{within} a bucket), the interior buckets are
summed whole, and any in-range most-common value would be added at its exact stored frequency. The
sharpened $0.185$ then flows into the cost model exactly where the default $0.33$ would have, often
flipping the index-versus-scan decision of \S\ref{sec:joinorder}.

A worked estimate shows the exponential backoff doing its job on a three-way conjunction.

\dcap{Selectivity of an AND, with the SQL-Server exponential damping}
\begin{diagramS}
predicate:  status eq "ACTIVE"  and  price gt 100  and  region eq "EU"      over 1 000 000 rows

per-predicate selectivity:
  status eq "ACTIVE"  → histogram MCV hit, exact frequency      0.40
  price  gt 100       → range default (no histogram yet)        0.33
  region eq "EU"      → equality default                        0.01

combine the AND: sort ascending, damp the k-th by a 1/(k+1) exponent (SQL-Server backoff):
  s = 0.01 · 0.33^(1/2) · 0.40^(1/3)  =  0.01 · 0.574 · 0.737  ≈ 0.0042
  (the naive independent product 0.01·0.33·0.40 = 0.0013 would UNDER-count correlated conjuncts)

estimate ≈ 0.0042 · 1e6 = 4 200 matches → handed to the cost model's index-vs-scan decision
\end{diagramS}

The damping deliberately makes the conjunction \emph{less} selective than the independence assumption
would, because real predicates on one record are usually correlated, so a naive product systematically
under-estimates result sizes and over-eagerly picks an index whose per-row overhead a larger result no
longer repays. The same estimate, sharpened where a histogram exists, is exactly the \code{card} and
\code{selectivity} the next section's cost model multiplies through.

Selectivity prices one predicate; \emph{propagation} carries an estimate along a chain of operators. A
two-binding join shows three of the propagation rules --- bind (multiply), filter (scale), join (multiply
by join selectivity) --- firing in order:

\dcap{Cardinality propagating through ForBind → Selection → Join}
\begin{diagramS}
query:  for $b in jn:doc('db','cat').store.book[]  for $a in jn:doc('db','cat').store.author[]
        where $b.price gt 50 and $b.authorId eq $a.id  return { $b, $a }

  ForBind $b    references(/store/book/[])   = 100 000                       → 100 000
  ForBind $a    references(/store/author/[]) =   5 000   (cartesian: ×5 000) → 5.0e8
  Selection     price gt 50    scale by 0.33 (range default)                 → 1.65e8
  Join          authorId = id  join selectivity 1/distinct(id) = 1/5000      → 33 000

  the optimiser hands the cost model card = 33 000 for the join's output, NOT the 5.0e8
  intermediate — which is why the join is reordered so the 5 000-row author side BUILDS and the
  100 000-row book side PROBES (\S\ref{sec:joinorder}), keeping the hash table small
\end{diagramS}

\noindent Each operator applies exactly one rule to the running estimate, so the number reaching the cost
model is the post-filter, post-join cardinality rather than the raw cross product. The only \emph{exact}
input in the whole chain is the path summary's per-path reference count; everything after it is an
estimate, which is why a wrong reference count or a missing histogram propagates into every downstream
decision.

\section{The cost model: sequential scan versus HOT-index scan}

\code{JsonCostModel} prices the two access paths in abstract I/O{+}CPU units, with a random page
$4\times$ a sequential one, CPU at \code{0.01}/tuple, 1024 records/page, HOT fan-out~16 and 512
entries/leaf:

\dcap{The two costed access paths (\texttt{JsonCostModel})}
\begin{diagramS}
sequential scan:  pages = ⌈N / 1024⌉ ;   cost = pages·1.0  +  N·0.01
index scan:       trieDepth   = ⌈ ln(card) / ln 16 ⌉
                  traversalIO = trieDepth · 0.2 · 4.0            (0.2 = cold-trie fraction)
                  dataIO      = ⌈ (card · selectivity) / 512 ⌉ · 4.0
                  cost = 3.0  +  traversalIO  +  dataIO  +  (card · selectivity)·0.01
\end{diagramS}

\code{CostBasedStage} computes both and sets \code{PREFER\_INDEX} when the index cost is lower, the
crossover a cost model exists to find. For a 10~M-row document with 10 matching rows
($\text{selectivity}=10^{-6}$) the index plan is on the order of $10^4\times$ cheaper; at 90\,\%
selectivity the sequential scan wins, because the index would touch nearly every page at the $4\times$
random-I/O rate. The GUI surfaces the whole pipeline (Figure~\ref{fig:plans}).

\begin{figure}[t]
\centering
\begin{minipage}{0.485\linewidth}\centering
  \includegraphics[width=\linewidth]{plan-parsed.png}\\[2pt]{\small (a) Parsed plan}
\end{minipage}\hfill
\begin{minipage}{0.485\linewidth}\centering
  \includegraphics[width=\linewidth]{plan-optimized.png}\\[2pt]{\small (b) Optimised plan}
\end{minipage}
\caption{The query-plan visualiser in the SirixDB Web GUI: the parsed abstract syntax tree (a) and the
cost-based optimiser's rewritten, index-aware plan (b) for the same JSONiq query.}
\label{fig:plans}
\end{figure}

A concrete predicate shows the cost model deciding.

\dcap{A worked optimisation: sequential scan vs index for one predicate}
\begin{diagramS}
query:  count( jn:doc('db','t').rows[] where .status eq "ACTIVE" )

  N = TOTAL_NODE_COUNT = 10 000 000 ;  a CAS index on /rows/[]/status exists
  selectivity("ACTIVE") from the histogram ≈ 0.05  →  ≈ 500 000 matches

  sequential:  ⌈10M/1024⌉·1 + 10M·0.01        ≈ 9 766 + 100 000     ≈ 1.1e5
  index:       3 + ⌈ln(10M)/ln16⌉·0.2·4 + ⌈500k/512⌉·4 + 500k·0.01
               ≈ 3 + 4.8 + 3 908 + 5 000      ≈ 8.9e3

  index cost << sequential  →  CostBasedStage sets PREFER_INDEX ;  IndexMatching emits a CAS scan
\end{diagramS}

The same model, run at a high selectivity, shows the index \emph{losing} --- the crossover the prose
above asserts, now worked through rather than stated:

\dcap{The other side of the crossover: at 90\% selectivity the sequential scan wins}
\begin{diagramS}
same query and index, but selectivity("common") ≈ 0.90  →  9 000 000 matches

  sequential:  ⌈10M/1024⌉·1 + 10M·0.01            ≈ 9 766 + 100 000      ≈ 1.1e5   (unchanged)
  index:       3 + ⌈ln(10M)/ln16⌉·0.2·4 + ⌈9M/512⌉·4 + 9M·0.01
               ≈ 3 + 4.8 + 70 316 + 90 000        ≈ 1.6e5

  index cost > sequential  →  CostBasedStage keeps the sequential scan
  why: dataIO now charges the 4x random rate on ~every leaf (⌈9M/512⌉ ≈ 17 579 random reads), while the
       per-match CPU (·0.01) is the same on either path — so the index only pays off when few pages move

  solving index(s) = sequential for the selectivity s:
     7.8 + (10M/512)·4·s + 10M·0.01·s = 109 766   →   7.8 + 178 125·s = 109 766   →   s ≈ 0.62
  below ~62% selectivity the CAS scan wins; above it, the full sequential scan does
\end{diagramS}

\noindent The two cases bracket the decision: at $5\%$ the index is $\sim$12$\times$ cheaper, at $90\%$ it
is $\sim$1.5$\times$ \emph{dearer}, and the model places the tipping point near $62\%$. That single number
is the whole reason a cost model, and not a static ``use the index if one exists'' rule, governs the
choice --- the right plan depends on how much of the document the predicate actually admits.

\section{Join ordering and the mesh}
\label{sec:joinorder}

Joins are reordered cost-first. \code{CostBasedJoinReorder} collects a connected group of inner joins
into a \code{JoinGraph} (a dense bitmask adjacency whose edges carry join selectivities) and runs
\code{AdaptiveJoinOrderOptimizer}: \textbf{DPhyp}~\cite{dphyp} (the Moerkotte--Neumann connected-subgraph\,/\,%
complement-pair dynamic program, $O(n\cdot 3^n)$ against the naïve $O(4^n)$) up to 20 relations, with
a greedy \code{GOO} ($O(n^3)$: repeatedly merge the cheapest connected pair) fallback beyond that.
\code{restructureJoinTree} then rewrites the AST so the larger input probes and the smaller builds at
each join. These alternatives feed a \textbf{mesh} of \emph{equivalence classes}: each class holds the
original plan plus seq/index or join-order \code{PlanAlternative}s (an \code{(AST, cost)} pair), and
\code{MeshSelectionStage} walks post-order, copying the cheapest alternative's decisions onto the live
tree.

A three-relation chain shows why the search pays off. The \emph{result} size is the same whatever order
the joins run in; what differs, sometimes by orders of magnitude, is the size of the \emph{intermediate}
the next join must build against.

\dcap{DPhyp join ordering: same answer, very different intermediates}
\begin{diagramS}
relations:   O orders 1 000 000     C customers 10 000     R regions 10
join graph:  O —[sel 1e-5]— C —[sel 1e-3]— R          (no direct O—R predicate)

DPhyp builds {O,C,R} from a connected subgraph and its complement, and costs both splits:

  {C,R} >< {O}:  card(C><R) = 1e4·10·1e-3 = 100        then >< O → 100·1e6·1e-5 = 1 000
                intermediate = 100          ✓ keep this plan
  {O,C} >< {R}:  card(O><C) = 1e6·1e4·1e-5 = 100 000    then >< R → 1e5·10·1e-3 = 1 000
                intermediate = 100 000      ✗ 1000× larger, same final answer

final card = 1 000 either way; DPhyp memoises the cheaper (build C><R first), and
restructureJoinTree puts the smaller side on build, the larger on probe, at each TableJoin.
\end{diagramS}

The dynamic program memoises the best plan for every relation subset, so the $O(n\cdot 3^n)$ cost buys
the provably cheapest ordering up to twenty relations; past that the \code{GOO} fallback gives up
optimality for an $O(n^3)$ greedy merge that is still far better than source order. The chosen ordering
is one \code{PlanAlternative} in the mesh, weighed against the seq/index alternatives for the same
relations before \code{MeshSelectionStage} commits it.

\section{Statistics caching, plan caching, and explainability}

A bitemporal optimisation falls out for free: the \code{StatisticsCatalog} keys statistics on
\code{(database, resource, path, revision)}, and because historical revisions are immutable their
statistics \textbf{never expire} (LRU eviction only); only the latest revision's entries carry a TTL
and are invalidated on write. The whole compile is plan-cached in an LRU (128 entries) keyed on
\code{queryText\,@\,indexSchemaVersion}, where a global index-schema version is bumped on every index
create/drop so stale plans miss; a cache hit returns a \emph{deep copy} of the plan, since ASTs are
mutable. The optimiser's reasoning is inspectable: the \code{sdb:explain} function and the REST
\code{plan} parameter expose the parsed, optimised, and \emph{candidate} (mesh) plans
(Figure~\ref{fig:plans}), the last serialised straight from the mesh's equivalence classes.

Running the same query three times shows the cache key doing its job:

\dcap{Plan cache: miss, hit, then a schema-version invalidation}
\begin{diagramS}
Q = "for $b in jn:doc('db','books').store.book[] where $b.price > 50 return $b"     indexSchemaVersion = 5

run 1   key = (Q @ 5)   MISS -> parse + optimise -> plan P1 (sequential scan) -> execute ; insert into LRU
run 2   key = (Q @ 5)   HIT  -> return a DEEP COPY of P1 (ASTs are mutable) -> execute    (no recompile)
create CAS index on price  ->  global indexSchemaVersion  5 -> 6   (every plan cached @5 is now unreachable)
run 3   key = (Q @ 6)   MISS -> recompile -> plan P2 (CAS index scan, cheaper) -> execute ; insert
\end{diagramS}

\noindent The version-tuple key is what makes index creation correct without walking the cache: bumping a
single global counter strands every stale plan behind a key nobody will form again, so run~3 cannot
accidentally reuse the pre-index P1. And the hit returns a \emph{deep copy} because the optimiser mutates
the plan it is handed --- sharing the cached instance would corrupt it for the next caller.

\begin{honesty}
\small\textbf{Caveat.} The pipeline, cost model, histograms/MCV, DPhyp/GOO reordering, the mesh
search, and the bitemporal statistics catalogue all \emph{ship} as wired, executed code. The genuine
limitations are that histogram collection is \emph{deferred} (a first-seen predicate plans on the
default selectivity until the histogram is built after the write closes), string fields get a
distinct-count but no value histogram, and the mesh's own header still calls it scaffolding for a
fuller search; the broader JQGM/Cascades vision in \pp{docs/COST_BASED_QUERY_OPTIMIZER_PLAN.md} remains
aspirational.
\end{honesty}

\section{The XML face: XQuery over the same store}
\label{app:xml}

The optimiser completes the query story for JSON; but the engine has a second face. The report leads with
JSON, yet \S\ref{sec:xmljson} made the point that the storage, versioning, durability, and indexing
machinery is node-family-agnostic: an XML resource is the same engine with
\code{ELEMENT}\,/\,\code{ATTRIBUTE}\,/\,\code{TEXT} nodes instead of objects, arrays, and fused records. It
is worth showing that XML face directly, so the dual-engine claim is concrete rather than asserted, the
more so because one capability --- FMES structural import --- is wired only on the XML side today
(\S\ref{sec:diff}).

\dcap{Embedded Java: an XML resource is the JSON flow with an XML node family}
\begin{diagramS}
Databases.createXmlDatabase(new DatabaseConfiguration(path));
try (var db = Databases.openXmlDatabase(path)) {
  db.createResource(ResourceConfiguration.newBuilder("lib").useDeweyIDs(true).build());
  try (var session = db.beginResourceSession("lib")) {
    try (var wtx = session.beginNodeTrx()) {                          // same single write permit
      wtx.insertSubtreeAsFirstChild(
            XmlShredder.createStringReader("<library><book><title>A</title></book></library>"));
      wtx.commit("initial");                                         // -> revision 1
    }
    try (var rtx = session.beginNodeReadOnlyTrx(1)) {                 // same lock-free read
      rtx.moveToFirstChild();                                        // <library>; XML adds
      // rtx.moveToAttribute(...) / rtx.moveToNamespace(...) for the markup-only axes
    }
  }
}
\end{diagramS}

\dcap{XQuery: navigation, time travel, diff, and the XML-only FMES import}
\begin{diagramS}
(* XPath navigation over the stored document *)
xml:doc('db','lib')/library/book[1]/title/text()

(* time travel: open the resource as of a wall-clock instant, then query it *)
xml:open('db','lib', xs:dateTime("2026-01-15T00:00:00Z"))/library/book

(* the edit script between two revisions, in the same JSON change format as jn:diff *)
xml:diff('db','lib', 1, 2)

(* fold an external document in as an identity-preserving structural delta (FMES; XML-only) *)
xml:import('db','lib', doc("/path/to/new-library.xml"))

(* store-level sdb: functions work on XML items too *)
sdb:item-history(xml:doc('db','lib')/library/book[1])
\end{diagramS}

The listings differ from the JSON capstone only where the data model genuinely differs: the write feeds
an \code{XMLEventReader} from \code{XmlShredder} instead of a \code{JsonReader} from \code{JsonShredder},
navigation is XPath with the extra attribute and namespace axes, and the query entry points carry the
\code{xml:} prefix. Everything underneath, the page trie, the sliding snapshot, the dual-beacon commit,
the write permit, the change-tracking indexes, and the temporal \code{sdb:} functions, is byte-for-byte
the same code path as for JSON. The one asymmetry is deliberate and noted in the body: \code{xml:import}
runs the FMES fast-match/edit-script algorithm to land an external document as a minimal structural
delta that preserves node identity, and that binding is XML-only for now, with the JSON port
(\code{JsonFMSE}) implemented but not yet wired to a function.

The two node families line up concern for concern; only the markup-specific kinds have no JSON analogue.

\begin{table}[h]
\centering\small\setlength{\tabcolsep}{6pt}\renewcommand{\arraystretch}{1.25}
\begin{tabularx}{\linewidth}{@{}l X X@{}}
\toprule
\textbf{Concern} & \textbf{XML node family} & \textbf{JSON node family}\\
\midrule
container & \code{ELEMENT} & \code{OBJECT}, \code{ARRAY}\\
named field & child \code{ELEMENT} or \code{ATTRIBUTE} & object key, fused into \code{OBJECT\_NAMED\_*}\\
scalar value & \code{TEXT} & \code{STRING}/\code{NUMBER}/\code{BOOLEAN}/\code{NULL} value\\
markup-only & \code{COMMENT}, \code{PROCESSING\_INSTRUCTION}, \code{NAMESPACE} & none\\
extra axes & \code{moveToAttribute}, \code{moveToNamespace} & none\\
ordering & document order & array order; object key order preserved\\
identity & DeweyID {+} stable node key & DeweyID {+} stable node key (identical)\\
query face & \code{xml:}, XPath/XQuery & \code{jn:}, JSONiq\\
\bottomrule
\end{tabularx}
\caption{The two node families compared. Container, field, scalar, ordering, and identity correspond one
to one; XML's markup-only kinds (comments, processing instructions, namespaces) and its attribute and
namespace axes are the only constructs with no JSON counterpart. Everything below the node family, the
page trie, sliding snapshot, commit protocol, indexes, and \code{sdb:} temporal functions, is shared.}
\label{tab:nodefamilies}
\end{table}

\section{A worked bitemporal scenario}
\label{app:bitemporal-scenario}

The minimal example of \S\ref{sec:bitemporal-worked} showed the two axes; this one adds the step a real
application needs, \emph{reconciling overlapping facts}, and shows the \code{sdb:} valid-time accessors
earning their place. Consider an HR resource configured with valid-time paths, recording Alice's salary,
then a backdated correction.

\dcap{Two recordings of a salary, the second a backdated raise}
\begin{diagramS}
resource "hr" (validTimePaths = validFrom/validTo); top level is an array of salary facts

revision r1  recorded 2026-01-05
   [ { emp:"Alice", salary:50000, validFrom:"2026-01-01", validTo:"2026-12-31" } ]

revision r2  recorded 2026-07-10   (a raise, backdated to June, is recorded)
   [ { emp:"Alice", salary:50000, validFrom:"2026-01-01", validTo:"2026-12-31" },
     { emp:"Alice", salary:55000, validFrom:"2026-06-01", validTo:"2026-12-31" } ]
\end{diagramS}

A plain \code{jn:valid-at} for 2026-08-01 now returns \emph{both} facts, because both intervals contain
that instant; the application must decide which is in force. The natural rule, the most recently effective
fact wins, is a one-line FLWOR over the accessors:

\dcap{Reconciling overlapping valid-time facts with the accessors}
\begin{diagramS}
(* the EFFECTIVE salary believed valid on 2026-08-01, as of now: latest validFrom wins *)
let $facts := jn:valid-at('co','hr', xs:dateTime("2026-08-01T00:00:00Z"))
return (
  for $f in $facts
  order by sdb:valid-from($f) descending
  return $f.salary
)[1]
⇒ 55000        (* the backdated raise supersedes the original on the overlap *)

(* but as of a MARCH transaction time, the raise had not been recorded — both axes at once *)
let $then := jn:open-bitemporal('co','hr',
                 xs:dateTime("2026-03-01T00:00:00Z"),     (* transaction time *)
                 xs:dateTime("2026-08-01T00:00:00Z"))     (* valid time *)
return (for $f in $then order by sdb:valid-from($f) descending return $f.salary)[1]
⇒ 50000        (* in March we believed Alice's August salary was 50000 *)
\end{diagramS}

The contrast, 55000 now versus 50000 as-of-March, for the \emph{same} valid instant, is the bitemporal
question stated precisely: not ``what was Alice's salary on 2026-08-01'' but ``what did we \emph{believe},
\emph{and when}, Alice's salary on 2026-08-01 to be''. The valid axis is reconciled in the query language
with \code{sdb:valid-from}; the transaction axis is chosen by which immutable revision
\code{jn:open-bitemporal} opens. Neither rewrites the other, and the original 50000 fact is never
destroyed, the raise is a new fact recorded alongside it, which is what makes the as-of-then answer
reconstructable at all.

\chapter{Vectorised and Columnar Execution}
\label{ch:vector}

\noindent\textit{(\pp{io/sirix/page/pax/*}; \pp{io/sirix/index/projection/*};
\pp{io/sirix/query/scan/SirixVectorizedExecutor.java}.)}

\section{Model and a naming hazard}

Once the optimiser of Chapter~\ref{ch:query} has chosen a plan, an analytical one can run on a
columnar fast path instead of the row-at-a-time iterator. For patterns such as aggregates,
filtered counts, and group-by-count over JSON arrays,
SirixDB can leave the row-at-a-time iterator for a \textbf{batch-oriented columnar path}
processing fixed 1024-row batches with SIMD. A naming hazard to dispel: \code{VectorPage} is
the root of an \emph{embedding / k-nearest-neighbour vector-search} index, a different
subsystem; the columnar machinery lives in \code{pax} and \code{projection}.

\section{Columnar storage}

Two facilities exist. \textbf{PAX regions} (\pp{io/sirix/page/pax/}) compress
numeric/boolean/string field groups \emph{within} a \code{KeyValueLeafPage}, grouping a parent
key's children into contiguous, SIMD-friendly columns. \textbf{Projection-index leaf pages}
(\pp{io/sirix/index/projection/ProjectionIndexLeafPage.java}) are a secondary, declared
projection of up to 1024 rows of a fixed field schema:

\dcap{\code{ProjectionIndexLeafPage} --- per-column encodings}
\begin{diagramS}
ProjectionIndexLeafPage (per column, by kind):
  NUMERIC_LONG : [min][max] zone-map · valueBitWidth · valueBase (frame-of-reference) · packedValues
  BOOLEAN      : packed bitset
  STRING_DICT  : localDict(lengths, concatenated UTF-8) · dictIdBitWidth · packedDictIds
  v2 tail      : per-column flags · presence bitmaps · "PIX2" magic   (v1 readers ignore trailing bytes)
\end{diagramS}

These are orthogonal to the row-oriented primary storage and rebuilt independently.

A worked row group shows the encodings and the byte budget (sizes illustrative). Take 10 rows, three
columns:

\dcap{Encoding a 10-row group into a ProjectionIndexLeafPage}
\begin{diagramS}
dept   [Eng Sales Eng Ops Sales Eng Ops Sales Eng Ops]    salary/k [60 55 58 52 57 61 54 56 59 53]    active [T T F T T T F T T F]

STRING_DICT dept:  localDict {Eng:0, Sales:1, Ops:2}  (3 distinct -> dictIdBitWidth = ceil(log2 3) = 2)
   dictIds [0 1 0 2 1 0 2 1 0 2]  ->  10 x 2 bits = 20 bits = 3 B packed   (+ ~8 B for the dictionary)
NUMERIC_LONG salary:  zone-map [min 52][max 61] ;  valueBase 52  (frame-of-reference)
   deltas [8 3 6 0 5 9 2 4 7 1] ;  max 9 -> valueBitWidth 4 ;  10 x 4 = 40 bits = 5 B packed
BOOLEAN active:  packed bitset 1101110110  ->  10 bits = 2 B

group total ~ 11 + 7 + 2 + footer ~ 25 B for 10x3 cells, vs ~200 B row-oriented; and the zone-map [52,61]
lets a scan SKIP the whole group when a predicate (say salary > 70) cannot possibly match it
\end{diagramS}

\noindent Three ideas do the compression: a per-group \emph{dictionary} turns repeated strings into 2-bit
ids, \emph{frame-of-reference} stores each number as a small delta from the group minimum (so 8 numbers
near 52--61 cost 4 bits each, not 64), and \emph{bit-packing} writes exactly \code{ceil(log2 range)} bits
per cell. The zone-map is the bonus: a one-comparison group skip is what makes the SIMD scan of
\S\ref{sec:vecbench} read only the groups a predicate can touch.

\begin{honesty}
\small\textbf{Sub-leaf chunking (and remaining scale debt).} A projection leaf is split into
fixed-size 4\,KiB chunks, each a \emph{separate} HOT slot, so the sliding-snapshot merge shares
unchanged chunks across revisions: a single-row update rewrites only the changed chunk(s), not
the whole $\sim$20~KB leaf, aligning the projection with the \code{SLIDING\_SNAPSHOT} contract the
framework gives the CAS/PATH/NAME indexes. The open task~\#57 work is a \emph{further} optimisation:
moving chunk bytes into dedicated chunk pages (\code{ChunkDirectory}\,\textrightarrow\,%
\code{PageReference}) so HOT slots carry $\sim$16-byte references rather than 4\,KiB payloads,
pushing the per-leaf scale ceiling well past 100~M records.
\end{honesty}

\section{SIMD execution and the executor}

Comparison kernels (\pp{io/sirix/page/pax/NumberRegionSimd.java}) use the \textbf{Java Vector
API} (\code{jdk.incubator.vector}):

\dcap{Algorithm 8.1 --- Vectorised \texttt{count(values OP threshold)}}
\begin{diagram}
  thr ← broadcast(threshold);  count ← 0
  for i in start … end-LANES step LANES:
      v    ← LongVector.fromMemorySegment(payload, off(i), LITTLE_ENDIAN)
      mask ← v.compare(OP, thr)            // VectorMask, one bit per lane
      count += mask.trueCount()            // popcount
  count += scalar_tail(...)
  return count
\end{diagram}

\dcap{A worked vectorised count: \texttt{count(values > 30)} over ten longs, 4-lane species}
\begin{diagramS}
payload (off-heap, little-endian longs):   12  47  30   5  | 88  31   9  64  | 20  50
threshold thr = broadcast(30);  OP = GT (strictly greater);  count = 0

iter i=0  v ← [12, 47, 30,  5]   v.compare(GT,thr) = [F, T, F, F]   trueCount = 1   count = 1
iter i=4  v ← [88, 31,  9, 64]   v.compare(GT,thr) = [T, T, F, T]   trueCount = 3   count = 4
tail i=8  (2 < LANES left)        20>30=F   50>30=T                  scalar      +1   count = 5

result: 5 lanes-or-rows pass.  Two SIMD steps cover 8 of 10 values; only the n mod 4 = 2 tail is scalar.
\end{diagramS}

A group-by-count adds exactly one step to that skeleton: a SIMD predicate builds a selection mask, then
the surviving lanes scatter-add into a dense per-group counter indexed by the \emph{dictionary-encoded}
grouping key --- never allocating a group object per row, which is what the row-at-a-time interpreter does.

\dcap{A worked vectorised group-by-count: active rows per \texttt{dept}, one 8-row morsel}
\begin{diagramS}
dictionary-encoded columns   (dept: 0=Eng  1=Sales  2=Ops):
  dept    [ 0  1  0  2 | 1  0  2  1 ]
  active  [ 1  1  0  1 | 0  1  1  1 ]

step 1  SIMD filter active == 1  → selection mask [T T F T | F T T T]   (one compare, two 4-lane ops)
step 2  scatter-add the surviving lanes into counts[dept]   (counts starts {0,0,0}):
          rows 0,1,3 pass:  counts[0],counts[1],counts[2] += 1     → {1,1,1}
          rows 5,6,7 pass:  counts[0],counts[2],counts[1] += 1     → {2,2,2}
result  {Eng:2, Sales:2, Ops:2}   — returned as the {groupKey,count} envelope

vs interpreter: it hashes the dept STRING and looks up / allocates a group accumulator PER ROW; the vector
path indexes a 3-slot int[] by the small encoded key, so the group step is branch-light array arithmetic
\end{diagramS}

\noindent That difference --- a masked scatter-add into a dense \code{int[]} versus a per-row string hash
and accumulator lookup --- is the mechanism behind the order-of-magnitude group-by speedup the benchmark
records (\S\ref{sec:vecbench}); dictionary encoding is what turns the grouping key into an array index in
the first place.

A \code{BETWEEN} fuses two thresholds into one pass: \code{lo $\le$ v} AND \code{v $\le$ hi} ANDs two
masks before the single \code{trueCount}, so a range predicate costs the same vector traffic as a
one-sided comparison. Conjunctions across columns AND masks; bit-packed widths up to $\sim$56 are
gathered and unpacked, wider widths fall to scalar. The \code{SirixVectorizedExecutor} fans scans across a worker pool (each
with its own read transaction), takes a \textbf{projection-index fast path} when a covering
index carries the requested column (the status note's ``$\sim$27~s slot-walk\,$\to\sim$1~ms''; measured
at smaller scale in \S\ref{sec:vecbench}, where a million-row interpreter aggregate of 8.5--16.7~s becomes
a sub-millisecond warm repeat), extracts conjunctive per-column predicates (fusing
\code{>}\,{+}\,\code{<} into one \code{BETWEEN}), and re-enters brackit's sequence model through
a thin envelope (\code{\{groupKey,count\}} arrays, scalar counts, or \code{[count,sum,min,max]}).
Results are cached per \code{(session, revision)}.

\section{Status, correctness, and limitations}

The vectorised path is \textbf{functional and differential-tested for its supported query shapes}
(not a work-in-progress) but deliberately \emph{narrow} and \emph{fail-closed}. Running end-to-end
through brackit's \code{VectorizedExecutor} SPI, with genuine \code{jdk.incubator.vector} SIMD
(\code{LongVector}/\allowbreak\code{VectorMask}) on the numeric path: \code{count}/\allowbreak\code{sum}/\allowbreak\code{min}/\allowbreak\code{max}/%
\code{avg}, single- and multi-key group-by-count, filtered counts and aggregates, and count-distinct.
Two rewrites sharpen it: conjunctive \code{>}\,{+}\,\code{<} on one column \textbf{fuse into a single
\code{BETWEEN}} evaluated with a dual-mask SIMD kernel, and a floating-point literal against an integral
column is \textbf{rewritten to compare exactly} on the integral SIMD column. These sit alongside an optional
projection-index byte-scan fast path (when a covering index is installed) and opt-in
\textbf{morsel-driven parallelism}.

Three scoping points. (i) The \emph{optimiser-side} vectorised detection/routing stages are
\textbf{disabled}: they competed with brackit's own dispatch (rewriting the AST before brackit's
executor ran), so routing is left entirely to brackit's SPI and the dead Sirix-side stages remain
in-tree but unused. (ii) \textbf{Compiled-predicate bytecode generation} (a hidden class built per
distinct predicate via the Java~25 \code{java.lang.classfile} API) is a HotSpot-only JIT-warm-up
optimisation: it is gated off under the GraalVM native image (runtime class definition is unavailable
there) and falls back to the interpreter, \emph{producing identical results either way}. (iii)
Unsupported cases \textbf{fail closed rather than risk a wrong answer}: non-\code{double}-exact numerics
(\code{BigDecimal}/\allowbreak\code{BigInteger}/\allowbreak\code{float}), string and boolean aggregation, and order-ambiguous
group-by mixtures throw and defer to the row-at-a-time interpreter. String \emph{predicate} matching is
scalar (\code{String.equals} per row; the dictionary path hashes once per distinct value for
group-by/count-distinct), and \code{StringRegion}'s SIMD docstrings are aspirational; the
implementation is scalar.

Correctness is guarded by \textbf{interpreter-differential testing}:
\code{TypedGroupByDifferentialTest}, \code{PredicateTreeCountIntegrationTest},
\code{RangeFusionPassTest}, \code{FpCmpIntegralRewriteTest}, and \code{MorselFallbackTest}
(morsel-on $=$ morsel-off $=$ ground truth) compare the vectorised result against the row-at-a-time
engine, and \code{NestedSameNamePathScopingTest} regresses a real cardinality bug where local-name-only
indexing conflated a top-level \code{age} with a nested \code{pet.age} (fixed by path-filtered slot
scans). A heavier \code{PredicateTreeDifferentialTest} (shred 200~K records, 20 randomised predicate
shapes) exists but is disabled by default, and a \code{-Dsirix.scan.diag} mode emits per-page
fast-path/slow-path counters to bisect any silent divergence.

\section{Measured: vectorised vs interpreted, and an OLAP comparison}
\label{sec:vecbench}

The speed-up the columnar path buys was previously quoted only from an internal status note; this
section measures it. A JSON array of $N$ records \code{\{dept, city, age, active, amount\}} is stored in
SirixDB and four aggregate shapes, group-by-count, a filtered count, \code{sum}{+}\code{avg}, and a
two-key group-by, run both through brackit's row-at-a-time interpreter and through the
\code{SirixVectorizedExecutor}. The two paths are gated equal by the interpreter-differential tests above
(\code{TypedGroupByDifferentialTest} et al.), so the vectorised numbers are correct, not merely fast.

\begin{table}[h]
\centering\footnotesize\setlength{\tabcolsep}{5pt}\renewcommand{\arraystretch}{1.25}
\begin{tabular}{@{}l r r r r@{}}
\toprule
& \multicolumn{2}{c}{\textbf{$N=100\,000$}} & \multicolumn{2}{c}{\textbf{$N=1\,000\,000$}}\\
\cmidrule(lr){2-3}\cmidrule(lr){4-5}
\textbf{query} & \textbf{interp.} & \textbf{vec.\ (cold/warm)} & \textbf{interp.} & \textbf{vec.\ (cold/warm)}\\
\midrule
group-by-count & 40.8\,ms & 6.8 / \textbf{0.26}\,ms & 8\,705\,ms & 2\,701 / \textbf{0.19}\,ms\\
filtered count & 39.8\,ms & 143 / \textbf{0.23}\,ms & 8\,462\,ms & 191 / \textbf{0.24}\,ms\\
sum {+} avg & 65.5\,ms & 0.59 / \textbf{0.23}\,ms & 16\,713\,ms & 0.53 / \textbf{0.24}\,ms\\
two-key group-by & 53.8\,ms & 4.6 / \textbf{0.40}\,ms & 8\,704\,ms & 9.1 / \textbf{0.21}\,ms\\
\bottomrule
\end{tabular}
\caption{SirixDB interpreter vs vectorised executor (same data, same machine). \emph{Warm} is a
repeated query (the per-(session,revision) result cache hits, sub-millisecond regardless of $N$);
\emph{cold} is a fresh executor (a real SIMD scan, but conflating a one-time JIT, cold buffer pages, and
$-$ for the filtered count $-$ a hidden compiled-predicate class). The honest headline is interpreter vs
warm: a multi-second aggregate over a million rows becomes a \textbf{sub-millisecond} repeat, which is
the measured form of the status note's ``$\sim$27\,s\,$\to\sim$1\,ms''.}
\label{tab:vecbench}
\end{table}

\textbf{Against a real OLAP engine.} The same four aggregates over the same $10^6$ rows, on DuckDB~1.5.2
(its \code{read\_json} into native columns) and PostgreSQL~17 (\code{jsonb} rows, \code{rec->>'field'}):

\dcap{The four aggregates over $10^6$ rows, four engines (ms; lower is better)}
\begin{diagramS}
query             SirixDB interp.   SirixDB vec(warm)   DuckDB 1.5.2   PostgreSQL 17 (jsonb)
group-by-count       8705              0.19                8              180
filtered count       8462              0.24                0.4            68
sum + avg           16713              0.24                1              64
two-key group-by     8704              0.21                8              269
\end{diagramS}

\begin{honesty}
\small\textbf{Caveat, and the honest verdict.} \textbf{DuckDB is the clear winner for ad-hoc analytics}:
a purpose-built vectorised OLAP engine reading into typed columns answers these in 0.4--8\,ms, and that
is the right tool for the job. PostgreSQL's \code{jsonb} path pays a per-row extraction tax ($\sim$10--30$\times$
DuckDB). SirixDB's \emph{interpreter} is the slowest by far, it is a versioned node store, not a column
engine. SirixDB's vectorised \emph{warm} number is sub-millisecond, but it includes a result cache, so it
is the cost of a \emph{repeated} aggregate, not a fresh scan that beats DuckDB; the fair fresh-scan number
is the cold column (0.5--2\,700\,ms, dominated by one-time column build and JIT). The point of the
vectorised path is therefore not to out-run DuckDB but to make analytics over a \emph{versioned,
time-travellable} store, which DuckDB is not, tolerable, turning SirixDB's own 8.5--16.7\,s interpreter
scan into milliseconds for the shapes it supports. Numbers are single-machine, fully cached, one shape
family.
\end{honesty}

%================================================================
\chapter{The Enterprise I/O Stack: io\_uring, S3, Kafka}
\label{ch:enterprise}

\noindent\textit{(\pp{/tmp/sirix-hft/sirix-enterprise-*}.)}

\section{io\_uring via the Java FFM (Panama) API, no JNI}

Each backend in this chapter swaps the storage mechanism under the \code{IOStorage} SPI of
Chapter~\ref{ch:ondisk}, while the engine above stays oblivious. The enterprise io\_uring
backend drives Linux io\_uring \textbf{directly from Java with no
JNI}, via \code{java.lang.foreign} (\pp{FFIIOUring.java}): cache-line-aligned SQE/CQE layouts,
\code{VarHandle}/\allowbreak\code{MethodHandle} accessors, and the
\code{io\_uring\_setup}/\allowbreak\code{io\_uring\_enter}/\allowbreak\code{io\_uring\_register} syscalls, with a
striped ring pool ($\approx$2$\times$cores rings, queue depth 512), best-effort registered
files/buffers, and block-aligned buffers for O\_DIRECT. The storage class implements the same
\code{IOStorage} SPI and, since the V0 port, writes a \textbf{byte-identical V0 on-disk format}
(the same 64-byte superblock, dual uber-page beacons at 4096/8192, and 32-byte little-endian revision
records), so a resource round-trips between the io\_uring and FileChannel backends (a cross-backend
compatibility test asserts exactly this); it differs only in mechanism (MemorySegment I/O, async
submission, O\_DIRECT alignment, striped locking). Its own status document \emph{claims} io\_uring is
\textbf{production-ready and 20--40\,\% faster than FileChannel}; that figure is the backend's own, not
independently benchmarked in this report, and on Linux~6.8 a kernel bug that returns zero bytes for async
ring reads currently forces the path to run \emph{synchronously}, so the async overlap the claim rests on
is not realised on that kernel. The backend stays opt-in (FileChannel is the default storage type). (One documented V0 caveat: gappy compaction
policies that drop \emph{intermediate} revisions can leave checksummed revision records that no longer
open, so non-contiguous retention is not yet safe.)

\begin{property}[Backend-independent on-disk format]\label{prop:backendid}
A resource written by one storage backend opens byte-identically under any other, because every backend
emits the same \code{V0} on-disk format above the \code{IOStorage} SPI. A cross-backend round-trip test
between io\_uring and FileChannel asserts exactly this.
\end{property}

\section{The linked-commit chain and its falsification}

The writer's durability barrier was first \emph{designed} as an io\_uring \textbf{linked-SQE chain}:

\dcap{The linked-SQE commit chain (one submission)}
\begin{diagram}
[ fdatasync(data)  IORING_OP_FSYNC | IORING_FSYNC_DATASYNC | IOSQE_IO_LINK ]
   └▶ [ pwrite(SECONDARY beacon)  RWF_DSYNC | IOSQE_IO_LINK ]
        └▶ [ pwrite(PRIMARY beacon)  RWF_DSYNC ]        # commit acks iff this CQE.res == slot size
\end{diagram}

submitted in one syscall, with a feature-probe and a \code{-Dsirix.iouring.linkedCommitChain}
kill-switch falling back to a serial fdatasync {+} two O\_DSYNC writes of identical durability.

The finding, from the \textbf{W1 commit-throughput benchmark}: the linked chain is
\textbf{0.91--0.97$\times$ the serial path on a single resource}. A commit's barriers are an
inherently serial dependency chain with nothing to overlap, and io\_uring may punt the fsync to
a kernel worker (io-wq), so the chain was first disabled by default and then, in the V0 rewrite,
\textbf{removed outright}: the commit now uses serial write-through O\_DSYNC file descriptors, and the
W1 ordering hazard is the \emph{documented reason} a linked \code{write}\,\textrightarrow\,\code{fsync}\,\textrightarrow\,\code{write}
chain is excluded (its only theoretical advantage was higher \emph{aggregate} throughput across many
concurrent resources). The same work
surfaced and fixed two real bugs (a reader stripe-lock deadlock; an ENOSPC from an unreclaimed
preallocated tail), which informed the (separate, FileChannel-side) preallocated-commit profile
of \S\ref{sec:prealloc}. This is a case where the disciplined benchmark \emph{killed} an
appealing idea, exactly its job.

\begin{figure}[t]
\centering
\includegraphics[width=\linewidth]{flame-iouring.png}
\caption{CPU flame graph of the io\_uring read/scan path under load. The async-submission and
striped-ring machinery are visible; the W1 finding is that for a \emph{single-resource commit}
this concurrency has nothing to overlap, so the V0 commit uses the serial write-through path.}
\label{fig:iouring}
\end{figure}

\section{The io\_uring reader}

The reader (\pp{FFMIOUringReader.java}) does zero-copy, optionally busy-polled, batched page reads
under per-thread stripe locks, with an O\_DIRECT-aware aligned/unaligned split; a batched
descendant-prefetch path reads \emph{N} pages in two syscalls (a header batch then a data batch)
rather than $2N$. It implements the same \code{AbstractReader} contract as the FileChannel reader,
so the rest of the engine is oblivious to the mechanism.

\dcap{Batched descendant prefetch: N pages in two syscalls, not 2N}
\begin{diagramS}
prefetch the N child pages under a node (each on disk as [4B len][compressed body] at a known offset)

  naive FileChannel:  for each page: pread(len) then pread(body)                  → 2N syscalls
  io_uring batch:     submit N "read length" SQEs ─▶ one io_uring_enter ─▶ N CQEs → syscall 1
                      (every body size now known)
                      submit N "read body"   SQEs ─▶ one io_uring_enter ─▶ N CQEs → syscall 2
                      completions land zero-copy in registered buffers; busy-poll can skip the IRQ

⇒ 2 syscalls for N pages instead of 2N; the deep submission queue keeps the NVMe device busy
\end{diagramS}

A concrete \emph{N}=4 prefetch shows the submission-queue/completion-queue cycle the diagram abstracts:

\dcap{One batched prefetch, N=4: SQE submission and CQE completion}
\begin{diagramS}
the node's 4 children sit at data-file offsets  [ 10240, 25600, 40960, 57344 ]

phase 1 — lengths:  fill 4 SQEs, then one io_uring_enter(submit=4, wait=4):
  SQE0 {op:READ, fd:data, off:10240, len:4, buf:reg[0]}     CQE0 {res:4}   reg[0] = 4096
  SQE1 {op:READ, fd:data, off:25600, len:4, buf:reg[1]}     CQE1 {res:4}   reg[1] = 8192
  SQE2 {op:READ, fd:data, off:40960, len:4, buf:reg[2]}     CQE2 {res:4}   reg[2] = 12288
  SQE3 {op:READ, fd:data, off:57344, len:4, buf:reg[3]}     CQE3 {res:4}   reg[3] = 6144     <- syscall 1

phase 2 — bodies:  fill 4 SQEs with the now-known lengths, one io_uring_enter(submit=4, wait=4):
  SQE0 {op:READ, off:10244, len:4096, buf:reg[0]}  ...  SQE3 {off:57348, len:6144, buf:reg[3]}
  CQE0..3 {res = body byte counts}   bodies land zero-copy in the registered buffers          <- syscall 2

the ring head/tail advance by 4 on each enter; in busy-poll mode CQEs are reaped without a completion IRQ
\end{diagramS}

\noindent The two-phase split exists because a page's body length is in its own 4-byte header: the first
batch learns all \emph{N} lengths at once, the second reads all \emph{N} bodies at once, so a subtree
prefetch of any width is two \code{io\_uring\_enter} calls rather than $2N$ \code{pread}s --- and the
registered buffers mean the bytes are never copied out of the kernel. Figure~\ref{fig:iouringring} shows
why a single \code{io\_uring\_enter} can carry a whole batch: the queues are shared memory.

\begin{figure}[t]
\centering
\begin{tikzpicture}[font=\sffamily\small,>=Latex,
  side/.style={draw=slate!60,fill=layerbg,rounded corners=2pt,align=center,inner sep=4pt,minimum width=2.0cm,minimum height=2.4cm},
  cell/.style={minimum width=0.55cm,minimum height=0.5cm,inner sep=0},
  sqe/.style={cell,fill=ember!12,draw=ember!55},
  cqe/.style={cell,fill=nodenew!14,draw=nodenew!60}]
  \node[draw=slate!30,fill=slate!4,rounded corners=3pt,minimum width=6.1cm,minimum height=2.8cm] at (5.25,0) {};
  \node[font=\scriptsize\itshape,slate] at (5.25,1.6) {shared memory (no copy)};
  \node[side] (app) at (0,0) {application\\(SirixDB\\reader)};
  \node[side] (ker) at (10.5,0) {kernel\\(io\_uring)};
  \node[font=\scriptsize,anchor=east] at (3.55,0.95) {SQ:};
  \foreach \i in {0,...,4}{ \node[sqe] at (3.9+0.6*\i,0.95) {}; }
  \node[font=\scriptsize,anchor=east] at (3.55,-0.55) {CQ:};
  \foreach \i in {0,...,4}{ \node[cqe] at (3.9+0.6*\i,-0.55) {}; }
  \draw[->,ember!75] (1.0,0.95) -- (3.05,0.95) node[midway,above,font=\scriptsize]{fill SQEs};
  \draw[->,ember!75] (6.5,0.95) -- (9.5,0.95) node[midway,above,font=\scriptsize]{kernel reads};
  \draw[->,nodenew!70!black] (9.5,-0.55) -- (6.5,-0.55) node[midway,below,font=\scriptsize]{post CQEs};
  \draw[->,nodenew!70!black] (3.05,-0.55) -- (1.0,-0.55) node[midway,below,font=\scriptsize]{app reaps};
  \node[font=\scriptsize,align=center,color=ember!80!black] at (5.25,-1.75) {io\_uring\_enter(): one syscall submits \emph{and} reaps a whole batch};
\end{tikzpicture}
\caption{Why io\_uring batches. The submission queue (SQEs) and completion queue (CQEs) are ring buffers in
memory \emph{shared} between the application and the kernel, so the application enqueues many requests and
dequeues many completions without a syscall per operation; a single \code{io\_uring\_enter} submits and
reaps an entire batch. With the read buffers \emph{registered} in advance, completed data lands in them
zero-copy, and in busy-poll mode the application reaps CQEs without even a completion interrupt --- which is
what turns the $2N$ \code{pread}s of a wide prefetch into two enters.}
\label{fig:iouringring}
\end{figure}

\section{The S3 object-storage backend}

The S3 backend (\pp{sirix-enterprise-s3}) maps a resource onto object storage, but, crucially,
\textbf{not by mirroring the V0 file format}. It defines its own remapped layout of immutable
\emph{segment} objects, one mutable \emph{revisions} object, and a JSON \emph{manifest}, while
keeping the page-record framing (\code{[4B len][compressed page]}) and the alignment constants of
the local format.

\dcap{S3 layout and the segment\,\textrightarrow\,manifest\,\textrightarrow\,revisions commit}
\begin{diagram}
Object layout (per resource, under <prefix>/<database>/<resource>/):
  manifest.json   authoritative segment map: virtualBases[] → s3Keys[], + latestRevision
  revisions       [0,1024) = latest uber page;  rev r → [8B segOffset][8B epochMs] at r*16+1024
  segments/segment-N.dat   immutable ~16 MB batch of many [4B len][compressed page] records

Commit  (atomic acknowledge point = the revisions object, written LAST):
  1. flush + PUT the durable data segment(s)
  2. PUT a candidate manifest.json referencing only now-durable segments
  3. PUT revisions  (carries the new uber page)         ← commit acknowledged here
  crash between steps ⇒ orphaned, unreferenced objects — never a torn or visible revision
\end{diagram}

A page is no longer one object: pages accumulate in an off-heap buffer and flush as
\pp{segments/segment-N.dat} once they cross a threshold (default 16~MB), so a commit may span
several immutable segment objects. The data file's virtual address space is preserved as a
\textbf{virtual offset} that becomes the \code{PageReference} key (offsets begin at 1024 so key~0
stays the uber-page sentinel). The \textbf{manifest} (\code{S3Manifest}: parallel
\code{virtualBases[]}/\allowbreak\code{s3Keys[]}/\allowbreak\code{segmentSizes[]} arrays {+} \code{latestRevision}) is the
authoritative map from virtual offset to \code{(object key, in-segment offset)}, resolved by binary
search; segment discovery is \emph{never} an S3 \code{LIST}, which makes the backend immune to S3's
eventual list-consistency. The commit order (segment, then candidate manifest, then the revisions
object \emph{last}) is the object-storage analogue of the dual-beacon protocol: the uber page in
the revisions object is the commit point exactly as the primary beacon's write-return is on local
V0, so a crash can only orphan unreferenced objects, never expose a torn revision. Manifest-load
failure is deliberately \emph{fatal} (404-not-found means a genuinely new resource; a read error or
truncation throws) so a transient fault can never start the resource on an empty manifest that would
orphan all prior data.

\dcap{Resolving a virtual offset to an S3 object through the manifest}
\begin{diagramS}
read a page whose PageReference key (virtual offset) = 33 554 432

manifest (parallel sorted arrays; the manifest IS the directory, never an S3 LIST):
   virtualBases = [ 1024,  16 778 240,  33 554 432, … ]
   s3Keys       = [ seg-0, seg-1,       seg-2,      … ]
   segmentSizes = [ ~16MB, ~16MB,       ~16MB,      … ]

resolve:  binary-search virtualBases for the largest base ≤ key → index 2 (base 33 554 432)
   object         = s3Keys[2] = segments/segment-2.dat
   in-segment off = key − virtualBases[2] = 0
   ranged GET segment-2.dat bytes [off .. off+len]  (must return HTTP 206) → [4B len][compressed page]

because discovery reads the manifest rather than listing the bucket, the backend is immune to S3's
eventual list-consistency.
\end{diagramS}

The write side, spanning a segment boundary, shows the manifest growing and the commit ordering holding:

\dcap{A multi-segment commit (revision 6 writes ${\sim}$20 MB; segment threshold 16 MB)}
\begin{diagramS}
before:  virtualBases=[1024, 16 778 240]   s3Keys=[seg-0, seg-1]   latestRevision=5

commit revision 6 (buffers ~20 MB of pages off-heap):
  pages buffer; at 16 MB the buffer FLUSHES -> PUT segments/segment-2.dat   (now durable)
            virtualBases += 33 554 432 ;  s3Keys += seg-2
  remaining ~4 MB buffers; at commit it FLUSHES -> PUT segments/segment-3.dat   (now durable)
            virtualBases += 50 331 648 ;  s3Keys += seg-3
  PUT candidate manifest.json -> virtualBases=[1024, 16 778 240, 33 554 432, 50 331 648]
                                 s3Keys=[seg-0 .. seg-3]   latestRevision=6
  PUT revisions  (the rev-6 uber page)                          <- commit acknowledged HERE

a crash before that final PUT leaves seg-2/seg-3 as ORPHANS (no live manifest references them),
never a visible or torn revision 6 — segments are immutable, the revisions object is the single gate
\end{diagramS}

Durability and signing use no AWS SDK: a hand-rolled \code{AwsV4Signer}
(\code{AWS4-HMAC-SHA256}, validated against AWS's published test vector) over the JDK
\code{HttpClient} drives PUT, GET, ranged GET (which \emph{must} return 206), HEAD, and DELETE;
there is no \code{LIST} and no multipart upload. Reads retry on 429/5xx with exponential backoff and
full jitter, self-correct clock-skew 403s from the response \code{Date}, and a follower whose
manifest is stale self-heals by reloading \pp{manifest.json} once and retrying. It implements the
same \code{IOStorage} SPI as the FileChannel/io\_uring backends (\code{ServiceLoader}-discovered,
provider priority~200), differing only in mechanism: a blocking ranged GET replaces a \code{pread},
and durability is S3's 2xx-on-PUT rather than \code{fdatasync}.

\begin{honesty}
\small\textbf{Caveat.} The core
write\,\textrightarrow\,commit\,\textrightarrow\,read\,\textrightarrow\,reopen path is verified, the
SigV4 client matches AWS's published test vector, and an audit took the backend from never-run (49
issues, 16 critical) to a green suite (75 tests) fixing four bug classes: manifest-load data-loss,
page-layout collisions at virtual offset~0, non-atomic commit, and a static-striped-buffer reader race.
Two earlier blockers are now \textbf{closed}: S3 is a first-class \code{StorageType.S3} (no reflective
registration needed), and a real \code{LocalStack}/Testcontainers test exercises the client against the
actual S3 API. Genuinely still open: the \textbf{L2 SSD cache and segment compactor are shipped but
unwired} (reads pay full S3 latency, segment count grows unbounded); uploads are single-part only (no
multipart); there is no conditional-\code{If-Match} PUT for multi-writer safety; metrics are logged at
close rather than exported; and there is no eager revision-index load for cold-JVM time travel. So the
core is sound and integration-tested; what remains is performance and operational hardening.
\end{honesty}

\section{The Kafka replication backend}

The Kafka module (\pp{sirix-enterprise-kafka}) is a \textbf{change-data-capture / replication
sidecar, not a storage backend}; it does not implement \code{IOStorage}. It streams
\emph{committed-revision notifications} from a write \emph{leader} to read \emph{followers} so reads
scale horizontally and survive node loss, while the S3 backend remains the durable source of truth.

\dcap{Leader\,\textrightarrow\,Kafka\,\textrightarrow\,follower replication flow}
\begin{diagram}
Leader (the single writer for a resource)
  commit → S3Writer: segment → manifest → revisions      (durable on S3; revisions = commit point)
         → ReplicationHooks fires a CommitNotification
         → KafkaReplicationListener builds a TransactionEvent  (S3 coordinates only, NOT pages):
              { db, resource, revision, segmentObjectKey, uberPageOffset, uberPageHash, … }
         → publish to topic  sirix.transactions.{db}   (partition key = resource ⇒ per-resource order)
                                   │  [Kafka]
                                   ▼
Follower(s)
  TransactionConsumer  (auto-commit off; dedup: skip any revision ≤ appliedRevision)
         → CacheWarmer: loadManifestFromS3() ; liveManifest.replaceWith(latest)
         → the follower's read address space is now byte-identical to the leader's
         → commit the Kafka offset only AFTER the apply is durable   (at-least-once, idempotent)
\end{diagram}

The streamed unit is a fixed-schema \code{TransactionEvent}, a \emph{pointer} to a committed
revision (database, resource, revision number, the S3 \code{segmentObjectKey}, the uber-page offset,
and a 16-byte uber-page hash), serialised as a hand-rolled binary record directly over a
\code{MemorySegment} (a 64-byte header with a \code{"SIRX"} magic, then length-prefixed UTF-8
fields). \textbf{Page bytes never traverse Kafka}; only the S3 coordinates do. There is one topic
per database, partitioned by resource name, giving per-resource ordering; leadership is just the
Kafka consumer-group assignment for that partition (which coincides with SirixDB's single writer for
the resource), so no external ZooKeeper/etcd coordination is needed. The S3 writer stays decoupled
from Kafka through a process-wide \code{ReplicationHooks} SPI.

The central design point is that \textbf{a follower reloads the leader's authoritative
\pp{manifest.json} from S3 rather than rebuilding it locally} from the event's arithmetic: because the
manifest already encodes the segment map, \code{latestRevision}, and the effect of compaction,
adopting it wholesale (\code{replaceWith}) makes the follower's address space byte-identical to the
leader's and immune to divergence under out-of-order, duplicate, or post-compaction delivery. The
event is thus a \emph{trigger and durability check}, not the payload. Delivery is at-least-once with
idempotent in-order apply: the producer runs with idempotence and \code{acks=all}; the consumer
disables auto-commit and advances the offset only after a batch durably applies, rewinding
\emph{all} partitions in a failed batch; and the dedup key is the revision number (events
$\le$ the applied revision are skipped). Leader failover is ordinary Kafka rebalance, and a dropped
``last commit'' notification self-heals on the next read, since a manifest miss triggers a one-shot
reload from S3.

A delivery trace shows the idempotent apply absorbing the messy cases Kafka permits:

\dcap{A follower under duplicate and out-of-order delivery}
\begin{diagramS}
follower state:  appliedRevision = 40

deliver rev 41   41 > 40  -> apply: reload manifest from S3, replaceWith ; appliedRevision=41 ; commit offset
deliver rev 41   41 <= 41 -> DUPLICATE, skip (no apply) ; commit offset           (at-least-once absorbed)
deliver rev 43   42 unseen: 43 > 41 -> apply: reload manifest -> it authoritatively names the latest
                 revision -> appliedRevision = 43
deliver rev 42   42 <= 43 -> late DUPLICATE, skip                                 (manifest already past it)

invariant: the follower's address space = whatever manifest.json says, NOT an in-order replay of events,
so a gap, a reorder, or a repeat all converge to the leader's manifest — the event is a trigger, not data
\end{diagramS}

\noindent This is why the event carries only S3 coordinates: since the follower adopts the leader's
manifest wholesale rather than reconstructing state from a sequence of deltas, the delivery order of the
\emph{triggers} cannot corrupt it. The strongest property a replication layer can have --- convergence
regardless of duplication or reordering --- falls out of making the durable manifest, not the message
stream, the source of truth.

\begin{honesty}
\small\textbf{Caveat.} Replication is wired end-to-end and tested (39 Kafka tests; an audit took
the original apply-asynchronously, at-most-once follower from 62 issues, 28 critical, to the idempotent
manifest-reload design that collapsed roughly ten of them). It is the \emph{least} mature of the three
backends: verification is still \emph{entirely in-process} (\code{MockProducer}/\allowbreak\code{MockConsumer}),
and a real-broker Testcontainers test, engine-lifecycle auto-bootstrap (the replication manager is still
wired by hand), \emph{enforced} write-leadership fencing (a \code{KafkaLeaderElection} \emph{tracker}
now exists, but producing is unconditional under the single-writer invariant rather than fenced),
follower snapshot bootstrap past topic retention, \code{AdminClient} topic provisioning, and lag/metrics
observability are all still open. The SSD-prefetch arm of the cache warmer is stubbed; today it only
reloads the manifest and verifies durability.
\end{honesty}

%================================================================
\chapter{The APIs: Embedded and REST}
\label{ch:rest}

\noindent\textit{(\pp{io/sirix/api/*}, \pp{io/sirix/access/*}; \pp{bundles/sirix-rest-api/}, Kotlin {+} Vert.x.)}

Everything built so far is reached two ways: \emph{embedded}, through a Java API of resource sessions
and transactions, and over the network, through a Vert.x \textbf{REST} server built on that same API.
This chapter covers both, taking the embedded transaction and concurrency model first, then the REST
surface that wraps it.

\section{The session and transaction model}

A database is opened through the static \code{Databases} factory:
\code{createJsonDatabase(DatabaseConfiguration)} lays out the on-disk folder structure, and
\code{openJsonDatabase(path)} returns a \code{Database<JsonResourceSession>}, on which
\code{createResource(ResourceConfiguration)} and then \code{beginResourceSession(name)} open a
\code{JsonResourceSession}. The session is the unit of concurrency control, and it enforces the defining
contract: \textbf{one read-write transaction per resource, but any number of concurrent read-only
transactions}.

Enforcement is a per-resource \code{Semaphore} with a single permit, shared across all sessions for the
same resource path through a process-wide registry (so the rule holds across threads).
\code{beginNodeTrx()} \code{tryAcquire}s it, waiting up to 5\,s, then throwing
\code{SirixUsageException} (``please close the running read-write transaction first''), and releases
it only when the writer closes. \textbf{Read-only opens take no lock at all}: they are merely registered
in a concurrent map, so $N$ readers run in parallel. This is lock-free snapshot isolation made concrete
(Property~\ref{prop:snapshot}): a read-only transaction is pinned to one immutable committed revision,
and a concurrent writer appends a \emph{new} revision without ever mutating the one a reader sees.

\dcap{One writer, many readers, no read locks}
\begin{diagramS}
per-resource Semaphore(1), shared JVM-wide via WriteLocksRegistry (keyed by resource path)

  writer (at most one)   beginNodeTrx() ─tryAcquire(5s)─▶ build revision N+1 in the TIL
                         commit ─▶ publish N+1, release the permit
  readers (any number)   beginNodeReadOnlyTrx(r) ─▶ pinned to immutable revision r, no lock

   … rev N-2   rev N-1   rev N   ◀── readers pin any committed revision; the writer only appends
\end{diagramS}

\textbf{Time travel is just a choice of read-only transaction}: \code{beginNodeReadOnlyTrx()} reads the
latest revision, \code{beginNodeReadOnlyTrx(int)} a specific revision number, and
\code{beginNodeReadOnlyTrx(Instant)} the revision closest to a wall-clock instant (resolved by binary
search over commit timestamps). A writer ends with \code{commit(message?, timestamp?)} (returning the
new revision) or \code{rollback()}, and can \code{revertTo}/\allowbreak\code{truncateTo} an earlier revision;
\textbf{auto-commit} overloads (\code{beginNodeTrx(maxNodes)} / \code{maxTime}) flush periodically during
a large load, an \code{AfterCommitState} deciding whether the writer survives the flush.

\dcap{A time-travel read: select a revision root, then traverse an immutable snapshot}
\begin{diagramS}
beginNodeReadOnlyTrx(Instant "2026-01-15T00:00Z"):
   binary-search the commit-timestamp index → revision 23 (the closest at or before the instant)
   pin RevisionRootPage 23   (immutable; no lock taken)
   moveTo / axes traverse the revision-23 page trie exactly as for the latest revision
   a concurrent writer appending revision 24 never touches revision 23's pages
\end{diagramS}

\code{revertTo} and \code{truncateTo} look similar but do opposite things to history.
\code{revertTo(r)} is \textbf{non-destructive}: it re-bases the writer on revision~$r$'s content, so the
next commit appends a \emph{new} head whose starting point is $r$ while every revision after $r$ is
kept. This is how editing the past works, since an updating query that targets an older revision reverts
to it first, and it gives a resource's history a branching, append-only shape rather than a single line.
\code{truncateTo(r)} is \textbf{destructive}: it cuts the data and revisions files back to $r$, discards
everything after it, and is also the recovery primitive that rolls a resource back to its last good
revision after an interrupted commit (Ch.~\ref{ch:ondisk}).

\dcap{\texttt{revertTo} appends a branch; \texttt{truncateTo} deletes}
\begin{diagram}
history:    r1 ─ r2 ─ r3 ─ r4        (r4 is the latest)

revertTo(r2); edit; commit  ⇒  a new head whose base is r2's content:
   r1 ─ r2 ─ r3 ─ r4
          └──────── r5     (r3, r4 kept; r5 is the new head, branched from r2)

truncateTo(r2)  ⇒  destructive: r3 and r4 removed, r2 becomes the head again
\end{diagram}

\section{The transaction stack and the commit pipeline}

A transaction is a thin node-layer cursor over a deeper stack. A read-only \code{JsonNodeReadOnlyTrx}
(or its writing subtype) holds a \emph{page-level engine}, a \code{StorageEngineReader} for reads and a
\code{StorageEngineWriter} for writes (\pp{io/sirix/access/trx/page/}), and that engine drives the
copy-on-write page trie (Ch.~\ref{ch:pds}) over the storage SPI (Ch.~\ref{ch:ondisk}). The node layer
never touches bytes. A \code{moveTo(nodeKey)} splits the key into a page key and a slot offset
(\code{nodeKey >> 10} and \code{nodeKey \& 1023}), fetches that \code{KeyValueLeafPage} from the engine
through the buffer cache, and binds a zero-copy flyweight node onto the page's \code{MemorySegment} at
that slot. A write runs the other way: an \code{insert\ldots} call has the node factory allocate a
record in the writer's \emph{modified} page inside the Transaction Intent Log (Ch.~\ref{ch:versioning}),
so all uncommitted state lives in the TIL until commit. Per-resource concurrency is enforced one level
up, in the resource session, which owns the single write \code{Semaphore} and the open-transaction
registries and builds a fresh node factory, path-summary writer, hasher and node-to-revisions index for
each writer. The same stack serves both data models: the \code{Xml-} and \code{Json-} transaction types
are parallel specialisations of one \code{AbstractNodeTrxImpl}, so everything beneath the node-kind
operations is shared.

\code{commit()} turns the accumulated intent log into a durable revision through a fixed pipeline. The
node layer flushes deferred work and runs hooks, then hands off to the page layer, which serialises the
dirty pages, walks the copy-on-write subtree to assign disk offsets, and makes the result durable
through the dual-beacon protocol of Ch.~\ref{ch:ondisk}.

\dcap{The commit pipeline: \texttt{AbstractNodeTrxImpl.commit} \textrightarrow\ the page layer}
\begin{diagram}
node layer (commitInternal):
  1  flushPendingStats()        apply the deferred per-path statistics (one CoW pass per path)
  2  pre-commit hooks
  3  storageEngineWriter.commit(message, timestamp) ──┐
                                                      │  page layer (NodeStorageEngineWriter):
                                                      ├─ serialise dirty KeyValueLeafPages in parallel (> 4 pages)
                                                      ├─ uberPage.commit(): recurse the dirty page subtree,
                                                      │                     write pages, assign disk offsets
                                                      ├─ writeUberPageReference(): fsync data + dual beacons  (Ch. 4)
                                                      ├─ serialise the index catalogue
                                                      └─ clear the intent log, free off-heap pages
  4  publish the new RevisionRootPage as the resource's latest revision
  5  if storeDiffs: serialise the per-commit JSON change file
  6  re-instantiate for the next revision, or mark COMMITTED   (per AfterCommitState)
  7  post-commit hooks
\end{diagram}

Two details are easy to miss. Secondary indexes are not rebuilt at commit. Each node mutation calls
\code{notifyChange(\ldots)} synchronously during the write, and the path, CAS and name change-listeners
write their own records into the same intent log, so the index updates land in the same atomic page set
as the data. The path summary is the one deferred case: its per-path statistics are buffered during the
transaction and applied in a single copy-on-write pass per path at step~1
(Property~\ref{prop:defstats}). \code{rollback()} is the cheap inverse, discarding the intent log and
rewriting the previous beacon with no data writes, then re-instantiating the writer over the prior
revision.

\section{Auto-commit and bulk import}

A bounded transaction is the easy case; the hard cases are a never-ending stream of small edits and a
single enormous load, and both are handled by \emph{auto-commit}. Passing a node-count or a time bound to
\code{beginNodeTrx} makes the writer auto-committing: a count-bounded writer flushes whenever its
modification count crosses the threshold, and a time-bounded one flushes on a background
\code{ScheduledExecutorService}. What happens to the writer at each flush is governed by an
\code{AfterCommitState}. \code{KEEP\_OPEN}, the default, rebuilds the page engine for the next revision so
the same writer keeps going. \code{KEEP\_OPEN\_ASYNC} avoids even that: it rotates the Transaction Intent
Log in $O(1)$ and flushes the old pages on a background thread, producing \emph{no} intermediate
revisions (it is restricted to the \code{FILE\_CHANNEL} backend and count-based bounds). \code{CLOSE}
ends the writer after the flush.

Bulk loading leans on these. \code{insertSubtreeAs\ldots} runs the streaming shredder under a
\textbf{bulk-insert} flag that suppresses the per-node bookkeeping that would otherwise dominate. Hash
and descendant-count maintenance in particular has two mutually exclusive modes: inside one auto-commit
epoch the writer adapts hashes incrementally per insert, and otherwise it performs a single post-order
hash pass over the whole imported subtree at the end, turning an $O(\text{nodes}\times\text{depth})$
ancestor-walk into one traversal. With the asynchronous bulk fast path layering $O(1)$ intent-log
rotation on top, a multi-hundred-million-node import runs in bounded memory without a commit per record.

\dcap{Importing $10^8$ nodes under a 4.5M-node count bound: the three \texttt{AfterCommitState}s}
\begin{diagramS}
AfterCommitState      at each auto-commit flush ...                         intermediate revisions
  KEEP_OPEN (default)  rebuild the page engine; the same writer continues     ~22 (one per flush)
  KEEP_OPEN_ASYNC      rotate the TIL in O(1); flush old pages on a bg thread  NONE (one final revision)
                       (FILE_CHANNEL backend + count bound only)
  CLOSE                end the writer after the flush                          one, then stop

orthogonal: the bulk-insert flag suppresses per-node diff work, and hashes are either adapted
incrementally within an epoch or computed in ONE post-order pass at the end (not the O(nodes·depth) walk).
\end{diagramS}

The decisive option for a pure load is \code{KEEP\_OPEN\_ASYNC}: because it rotates the intent log in
constant time and never publishes an intermediate revision, a hundred-million-node import leaves the
history with a \emph{single} new revision and the writer never blocks on a synchronous flush, the
memory stays bounded by the rotation, not by the document.

\section{Navigation, mutation, and configuration}

Both transaction kinds are \textbf{cursors} (\code{NodeCursor}): \code{moveTo(nodeKey)},
\code{moveToFirstChild}/\allowbreak\code{moveToRightSibling}/\allowbreak\code{moveToParent}/\ldots, with \code{getKind},
\code{getName}, \code{getValue}, \code{getDescendantCount}, \code{getDeweyID}, and JSON type predicates
(\code{isObject}/\allowbreak\code{isArray}/\allowbreak\code{isStringValue}/\ldots). The \code{Axis} family (\code{ChildAxis},
\code{DescendantAxis}, \code{LevelOrderAxis}, and the temporal axes behind Ch.~\ref{ch:query}'s
\code{jn:} functions) are \code{long} iterators layered on a cursor. A \code{JsonNodeTrx} adds the write
surface: \code{insertObjectRecordAsFirstChild(key, value)} and its fused-primitive variant, typed
scalar inserts (\code{insertStringValueAs\ldots}, \code{insertNumberValueAs\ldots}) in all four
positions, bulk \code{insertSubtreeAsFirstChild(\ldots)}, \code{setStringValue}/\allowbreak\code{setObjectKeyName},
\code{remove()}, and genuine \code{moveSubtreeTo\ldots}/\allowbreak\code{copySubtreeAs\ldots}. A resource's
behaviour is fixed at creation by the \code{ResourceConfiguration} builder: \code{hashKind}
(\code{ROLLING}/\allowbreak\code{POSTORDER}/\allowbreak\code{NONE}), \code{useDeweyIDs}, \code{versioningApproach} (default
\code{SLIDING\_SNAPSHOT}), \code{maxNumberOfRevisionsToRestore}, \code{storageType},
\code{byteHandlerPipeline}, \code{buildPathSummary}, \code{storeDiffs}, and \code{storeNodeHistory} (the
change-tracking index of \S\ref{sec:changetracking}).

\dcap{The embedded API end-to-end (open \textrightarrow\ write \textrightarrow\ time-travel read)}
\begin{diagramS}
Databases.createJsonDatabase(new DatabaseConfiguration(path));
try (var db = Databases.openJsonDatabase(path)) {
  db.createResource(ResourceConfiguration.newBuilder("res").useDeweyIDs(true).build());
  try (var session = db.beginResourceSession("res")) {
    try (var wtx = session.beginNodeTrx()) {              // the one writer
      wtx.insertObjectAsFirstChild();
      wtx.insertObjectRecordAsFirstChild("name", new StringValue("sirix"));
      wtx.commit("initial");                              // -> revision 1
    }
    try (var rtx = session.beginNodeReadOnlyTrx(1)) {     // one of N readers, pinned to rev 1
      rtx.moveToFirstChild();
      var axis = new DescendantAxis(rtx, IncludeSelf.YES);
      while (axis.hasNext()) { axis.nextLong(); /* getKind()/getName()/getValue() */ }
    }
  }
}
\end{diagramS}

\section{The REST server}

The network face of the engine is a single Vert.x server exposing the bitemporal store over HTTP,
the substrate the Web GUI (Ch.~\ref{ch:gui}) is built on and the integration point for any non-embedded
client.

The server is one Vert.x \code{CoroutineVerticle} (\code{SirixVerticle}). It serves plaintext HTTP
only when explicitly configured; otherwise \textbf{HTTP/2 over TLS} with ALPN and PEM key/cert,
listening on port 9443 by default. Shutdown is graceful: the HTTP server closes with a timeout and
then \emph{all open databases are closed} to flush writes. A single router carries global handlers
(a configurable-origin CORS handler; a security-headers handler setting
\code{X-Content-Type-Options: nosniff}, \code{X-Frame-Options: DENY}, and a restrictive CSP;
metrics; and unauthenticated \code{/health} and \code{/metrics} routes), then one
authenticate-then-handle pair per business route. A central failure handler logs full detail
server-side but returns only a sanitised \code{\{statusCode, message\}} body; stack traces never
leak.

SirixDB's node layer is synchronous and blocking, so every database operation is pushed off the
event loop with Vert.x's \code{executeBlocking} (Figure~\ref{fig:restflow}). The decisive performance lever is \emph{which}
form is used: the default \code{executeBlocking(callable)} runs \textbf{ordered}, serialising all
blocking tasks queued on a context FIFO, so concurrent requests sharing a context execute
one-at-a-time on the worker pool, the root cause of the $\sim$20$\times$ latency-under-concurrency
regression. The remedy is the two-argument \code{executeBlocking(callable, /*ordered=*/false)} form,
which lets independent requests run in parallel across the pool; \S\ref{sec:histpath} above measures
its effect (the p95 concurrency ratio falling from $\sim$16$\times$ to $\sim$3$\times$). \emph{Caveat:}
a source review finds the lever is not yet applied uniformly: several call sites still take the
ordered default, so passing \code{ordered=false} consistently is a remaining mechanical step, not
a finished one.

\begin{figure}[t]
\centering
\begin{tikzpicture}[font=\sffamily\small,>=Latex,
  proc/.style={draw=slate!55,fill=white,rounded corners=2pt,align=center,inner sep=3pt,
               minimum height=0.95cm,minimum width=2.15cm,font=\footnotesize},
  io/.style={draw=slate!70,fill=slate!12,rounded corners=2pt,align=center,inner sep=3pt,
             minimum height=0.95cm,minimum width=1.8cm,font=\footnotesize},
  worker/.style={draw=ember!70,fill=white,rounded corners=2pt,align=center,inner sep=4pt,
                 minimum height=1.0cm,minimum width=6.5cm,font=\footnotesize},
  e/.style={->,slate},
  be/.style={->,ember,thick}]

  % swim-lane bands (drawn first, behind nodes)
  \path[fill=slate!7,rounded corners=3pt] (-1.4,1.55) rectangle (11.75,3.0);
  \path[fill=ember!7,rounded corners=3pt] (-1.4,-0.55) rectangle (11.75,0.9);
  \node[anchor=north west,font=\scriptsize\itshape,text=slate!72!black] at (-1.33,2.97)
       {EVENT LOOP --- non-blocking};
  \node[anchor=north west,font=\scriptsize\itshape,text=ember!72!black] at (-1.33,0.85)
       {WORKER POOL --- blocking node layer};

  % event-loop pipeline
  \node[io]   (req)   at (0.0,2.2)  {HTTP/2 req\\:9443 TLS};
  \node[proc] (gh)    at (2.5,2.2)  {global handlers\\CORS \textbullet\ headers\\\textbullet\ metrics};
  \node[proc] (auth)  at (5.0,2.2)  {authenticate\\Bearer JWT\\(Keycloak)};
  \node[proc] (route) at (7.6,2.2)  {route by depth\\role \textbullet\ validate};
  \node[io]   (resp)  at (10.45,2.2) {HTTP/2\\response};

  \draw[e] (req)--(gh); \draw[e] (gh)--(auth); \draw[e] (auth)--(route);

  % the boundary: dispatch to the worker pool
  \node[worker] (work) at (5.5,0.2)
       {acquire resource session $\to$ run trx / query\\(synchronous, blocking node layer) $\to$ serialise result};
  \draw[be] (route.south) -- (7.6,0.7);
  \node[fill=white,inner sep=1.5pt,font=\scriptsize] at (7.6,1.2)
       {\code{executeBlocking(\ldots,\,ordered=false)}};
  \draw[e] (work.east) -- node[right,font=\scriptsize,text=slate!75!black,inner sep=2pt]{result} (resp.south);
\end{tikzpicture}
\caption{The REST request lifecycle and the event-loop/worker-pool boundary. Routing, CORS and security
headers, JWT authentication, role and path-segment validation all run \emph{non-blocking} on a Vert.x
event loop; the synchronous, blocking node-layer work is dispatched to the worker pool via
\code{executeBlocking}. The \code{ordered=false} argument is the decisive lever: it lets independent
requests run in parallel across the pool, whereas the \code{ordered=true} default serialises them on a
per-context FIFO --- the $\sim$20$\times$ latency-under-concurrency regression of \S\ref{sec:histpath}.}
\label{fig:restflow}
\end{figure}

\section{Endpoints}

Routing is path-depth based over \code{/:database/:resource}, with content negotiation selecting
the JSON vs XML resource kind (\code{application/json} vs \code{application/xml}). Every path
segment is validated (rejecting blank, over-long, \code{..}, slashes, and NUL) as an
anti-traversal guard.

\begin{table}[h]
\small\setlength{\tabcolsep}{5pt}
\renewcommand{\arraystretch}{1.2}
\begin{tabularx}{\linewidth}{@{}l l X@{}}
\toprule
\textbf{Method \& path} & \textbf{Role} & \textbf{Purpose}\\
\midrule
\code{GET /} & VIEW & list databases, or run a \emph{global} JSONiq/XQuery if \code{query} is given\\
\code{PUT /:db} & CREATE & create a database {+} its first resource (JSON or XML body)\\
\code{POST /:db} & CREATE & create multiple resources from \code{multipart/form-data} uploads\\
\code{GET /:db/:resource} & VIEW & read a resource (or run a resource-scoped query); \code{HEAD} returns the ETag only\\
\code{POST /:db/:resource} & MODIFY & insert/update at a node (\code{nodeId} {+} \code{insert} mode)\\
\code{DELETE /:db/:resource} & DELETE & delete the resource, or a subtree when \code{nodeId} is given\\
\code{GET /:db/:resource/history} & VIEW & the revision list\\
\code{GET /:db/:resource/diff} & VIEW & diff between \code{first-revision} and \code{second-revision}\\
\code{GET /:db/:resource/pathSummary} & VIEW & the path-summary structure\\
\bottomrule
\end{tabularx}
\end{table}

Reads accept a rich query-parameter surface drawn directly from the handlers: temporal selectors
(\code{revision}, \code{revision-timestamp}, and \code{start}/\allowbreak\code{end} variants of each),
sub-document scoping (\code{nodeId}), the pagination cursors below, the serialisation controls
(\code{withMetaData}, \code{prettyPrint}), and, for queries, result-window bounds
(\code{startResultSeqIndex}/\allowbreak\code{endResultSeqIndex}), commit metadata
(\code{commitMessage}/\allowbreak\code{commitTimestamp}), and a \code{plan} flag that returns the
parsed/optimised/candidate query plan (the data behind the GUI's plan visualiser,
Figure~\ref{fig:plans}). Creation accepts \code{hashType} (default \code{NONE}), \code{useDeweyIDs},
and a \code{maxNodes} auto-commit threshold; bitemporal validity paths
(\code{validFromPath}/\allowbreak\code{validToPath}) are accepted too.

\section{Serialisation and pagination}

Output is produced by \code{JsonSerializer} (a node or subtree) or \code{JsonRecordSerializer}
(top-level pagination). The \code{withMetaData} parameter selects three modes
(\code{nodeKey}, \code{nodeKeyAndChildCount}, or full metadata), wrapping each value as
\code{\{key, metadata, value\}}; a \emph{fused} \code{OBJECT\_NAMED\_*} node emits its children as
a bare array directly under \code{value} (no inner record wrapper), the wire-shape subtlety the GUI
normaliser of Ch.~\ref{ch:gui} absorbs. To serve 100~GB+ resources, four controls compose: a depth
cap (\code{maxLevel}), a per-node fan-out cap (\code{maxChildren}), a total node budget
(\code{numberOfNodes}), and cursor-style top-level paging (\code{startNodeKey} {+}
\code{nextTopLevelNodes}) by which the client walks sibling pages from the last-loaded key. Query
results paginate independently through a sequence window. Read/\code{HEAD} responses carry an
\textbf{ETag} equal to the node hash (unless \code{hashType==NONE}).

\dcap{The \texttt{withMetaData} modes, and the fused wire shape}
\begin{diagramS}
source { "price": 12 }  →  one OBJECT_NAMED_NUMBER record (node key 3):
  (none)               → { "price": 12 }
  nodeKey              → { "key":"price", "metadata":{"nodeKey":3}, "value":12 }
  nodeKeyAndChildCount → { "key":"price", "metadata":{"nodeKey":3,"childCount":0}, "value":12 }

a fused container { "tags": ["a","b"] }  →  OBJECT_NAMED_ARRAY (key 4):
  → { "key":"tags", "metadata":{"nodeKey":4,"childCount":2},
      "value": [ {"metadata":{"nodeKey":5},"value":"a"}, {"metadata":{"nodeKey":6},"value":"b"} ] }
  the array sits DIRECTLY under value — there is no inner {metadata,value} wrapper for the fused key
\end{diagramS}

That last line is the whole reason the GUI carries a single wire-shape normaliser (Ch.~\ref{ch:gui}): a
fused key's children are a bare array under its \code{value}, so a parser written for the legacy
key-node-plus-value-node shape would mis-read every fused object until it re-wraps the bare value with the
entry's own metadata.

The four controls compose into cursor pagination over an arbitrarily large array. Take \code{store.books}
with 10\,000 elements, each holding a nested \code{reviews} array:

\dcap{Windowed pagination over a 10 000-element array (maxLevel=2, maxChildren=100)}
\begin{diagramS}
GET /db/store?nodeId=5&maxLevel=2&maxChildren=100&startNodeKey=0
  -> books[0..99] serialised; descent STOPS at level 2 (reviews appear as childCount, not expanded)
     response metadata:  nextTopLevelNodes = 1050      (the node key of books[100])

GET /db/store?nodeId=5&maxLevel=2&maxChildren=100&startNodeKey=1050
  -> books[100..199] ;  nextTopLevelNodes = 2050       (books[200])
  ... repeat, feeding each response's nextTopLevelNodes back in as the next startNodeKey ...

each request reads ONE window: <= 100 top-level siblings, <= 2 levels deep — bounded work and bounded
response regardless of the 10 000-element (or 100 GB) total; the client stitches the pages by the cursor
\end{diagramS}

\noindent The cursor is just the next sibling's node key, so paging never re-reads or re-counts what came
before --- the server seeks straight to \code{startNodeKey} and emits one window. A depth cap plus a
fan-out cap plus a sibling cursor is what lets a single screen of the explorer (Ch.~\ref{ch:gui}) browse a
resource far larger than memory.

\section{Authentication and authorisation}

Access is OAuth2/OIDC via \textbf{Keycloak}: the verticle auto-discovers the realm, and a
\code{/token} route brokers password, refresh-token, and authorisation-code grants. Every business
request must carry a \code{Bearer} JWT or it is rejected with 401. Authorisation uses two role
namespaces: a \emph{global} realm role (\code{view}, \code{create}, \code{modify}, \code{delete})
and a \emph{per-database} role \code{"\$database-\$role"}. Per request, the handler authenticates the
token and then checks the database-scoped role \emph{or} the global role, returning 403 on neither.
The enforcement is defence-in-depth: embedded JSONiq cannot bypass it because the query store
wrapper (\code{JsonSessionDBStore}) re-checks the role on each store operation
(\code{lookup}\,\textrightarrow\,VIEW, \code{create}\,\textrightarrow\,CREATE,
\code{drop}\,\textrightarrow\,DELETE), and an \emph{updating} query is gated at runtime
(\code{PermissionCheckingQuery} requires the MODIFY role before applying updates). The JWT subject
becomes the Sirix \code{User} recorded as the commit author.

\dcap{Two-tier role resolution for one request}
\begin{diagramS}
POST /mydb/employees   (a MODIFY operation)   Bearer JWT for user "alice"

  alice's realm roles: { "view", "mydb-modify" }       (global view + a database-scoped modify on mydb)
  required role for POST /:db/:resource = MODIFY

  check 1 — database-scoped:  has "mydb-modify"?  YES → authorised
  (else check 2 — global:     has "modify"?  ;  neither → 403 Forbidden)

  defence in depth: an embedded JSONiq update is RE-checked at runtime by PermissionCheckingQuery
  (MODIFY required before any update applies), so a query cannot slip past the route-level gate.
\end{diagramS}

\section{Writes and the streaming shredder}

A \code{POST} to a resource carries the JSON/XML fragment; \code{nodeId} selects the target and
\code{insert} selects the mode (\code{ASFIRSTCHILD}, \code{ASRIGHTSIBLING}, \code{ASLEFTSIBLING};
XML adds \code{REPLACE}). Writes are optimistically concurrency-controlled: the handler requires the
client's \code{ETag} to equal the current node hash, else it reports that the resource may have
changed in the meantime. Responses up to a few thousand nodes are serialised inline; larger writes
return 200 {+} ETag and a follow-up paginated \code{GET}. Updating queries reach the engine through
\code{sdb:}/\allowbreak\code{jn:} functions (e.g.\ \code{replace json value of sdb:select-item(...)}).

Large uploads use a \textbf{streaming shredder} (\code{KotlinJsonStreamingShredder}): a Vert.x JSON
parser on the event loop feeds a bounded channel (capacity 100\,000) consumed on a worker thread;
when the channel is full the parser \emph{and the HTTP request are paused} for true back-pressure,
and the producer is resumed every 50\,000 events. The decisive correctness fix (the alpha15
deadlock) is that the consumer resumes the producer \emph{before} suspending on an empty channel
(a low-water fallback), so a tail shorter than the resume threshold cannot wedge. Bulk-insert mode
skips per-node diff work and hashes are adapted incrementally via auto-commit, which is what makes a
300~M-node import feasible without exhausting memory.

\dcap{Streaming-shredder back-pressure (and the lost-wakeup the fix avoids)}
\begin{diagramS}
event loop (producer)                       bounded channel (cap 100 000)        worker thread (consumer)
  JsonParser.handler ─ trySend(event) ─────▶ [#############       ] ───────────▶ wtx.insert{Object,Array,Value}
  channel full → trySend fails:                                                  drains events, blocking inserts
     parser.pause() + httpRequest.pause()  ◀── apply back-pressure
  runOnContext: parser.resume()+request.resume() ◀───────────── resumeProducer() once drained to 50 000 left

  alpha15 fix: the consumer calls resumeProducer() BEFORE it suspends on an empty channel (low-water
  fallback), so a tail shorter than the 50 000 resume threshold can never wedge the producer asleep.

  every ~4.5M nodes: writer auto-commits in bulk-insert mode (no per-node diff) → memory stays bounded
\end{diagramS}

\section{Observability and status}

A Micrometer/Prometheus layer exposes \code{/metrics}: an active-request gauge plus request-duration
and request-count series tagged by method, the matched \emph{route pattern} (low cardinality), and
status; internal cache hit/miss/eviction gauges are bridged into the same registry.

A scrape makes those series concrete, and shows what an operator reads off them:

\dcap{A \texttt{/metrics} snapshot (Prometheus exposition; values illustrative)}
\begin{diagramS}
sirix_requests_active 3
sirix_request_duration_seconds_count{method="POST",route="/{db}/{resource}",status="200"}  42
sirix_request_duration_seconds_sum{method="POST",route="/{db}/{resource}",status="200"}     0.245
sirix_request_duration_seconds_count{method="GET",route="/{db}/{resource}",status="200"}    1204
sirix_request_duration_seconds_sum{method="GET",route="/{db}/{resource}",status="200"}      12.5
sirix_cache_hits_total{cache="record_page"}    98750
sirix_cache_misses_total{cache="record_page"}   1250

read off it:  POST mean = 0.245 / 42  = 5.8 ms      GET mean = 12.5 / 1204 = 10.4 ms
              record-page hit rate    = 98750 / (98750 + 1250) = 98.75%        3 requests in flight now
\end{diagramS}

\noindent The route tag is the \emph{pattern} \code{/\{db\}/\{resource\}}, never the concrete path, so a
busy database cannot explode the metric cardinality; and because the cache gauges share the registry, a
single scrape diagnoses both a slow route (rising \code{duration\_sum/count}) and a thrashing buffer pool
(falling hit rate) without attaching a profiler.

\begin{honesty}
\small\textbf{Caveat.} Besides the \code{ordered=false} lever above, a runtime authorisation
hop inside query evaluation is flagged \code{// FIXME} in the source; the dev entry point hardcodes
an absolute config path (the packaged/native image instead auto-loads its config from the working
directory); a client secret sits in the checked-in demo config (env-overridable); and XML write
support is parallel but thinner than JSON (e.g.\ it does not thread commit metadata into subtree
inserts).
\end{honesty}

\section{A worked trace: one update, end to end}
\label{sec:worked}

To tie the layers together, follow a single sub-document update from the wire to durability and back. A
client issues the JSONiq statement \code{replace json value of jn:doc('db','res').store.book[0].price
with 19.99}.

\textbf{Parse and plan.} brackit parses the statement into an AST, and the cost-based optimiser
(\S\ref{sec:optpipeline}) rewrites the navigation to reach a path index where one exists, leaving an
updating expression whose target resolves to a single fused \code{OBJECT\_NAMED\_NUMBER} node.

\textbf{Bind to the writer.} Evaluating the target produces a \code{JsonDBObject} item. Because the
statement is updating, the binding lazily upgrades the resource's read view to its one write transaction,
acquiring the per-resource write permit (the single-writer rule of this chapter); had the query targeted
an older revision, it would \code{revertTo} it first and branch the history.

\textbf{Apply, in the node layer.} \code{replace json value of} resolves to \code{JsonDBObject.replace},
which, seeing a number-to-number change, calls \code{setNumberValue(19.99)} on the \code{JsonNodeTrx}.
The node factory writes the new value into the writer's \emph{modified} \code{KeyValueLeafPage} inside
the Transaction Intent Log; the rolling hash of the touched node and its ancestors is recomputed in the
same pass, and the path and CAS index listeners record the value change into the same intent log.
Nothing has touched disk yet.

\textbf{Commit.} \code{sdb:commit}, or the implicit commit when the query store closes, runs the commit
pipeline shown earlier in this chapter: the deferred path statistics flush, the dirty pages serialise in
parallel, the copy-on-write subtree is written and assigned disk offsets, and
\code{writeUberPageReference} makes the new revision durable through the dual-beacon protocol of
Chapter~\ref{ch:ondisk}. The result is one new revision in which every page \emph{not} on the change's
root-to-leaf path is still shared with its predecessor.

\textbf{Read it back, then and now.} A later reader opening the new revision sees \code{19.99}; opening
the \emph{previous} revision, with no lock and no reconstruction, still sees the old price, because that
revision's pages were never mutated. And \code{sdb:item-history} on the price node returns exactly the
revisions in which it changed, read straight from the \code{RECORD\_TO\_REVISIONS} index
(\S\ref{sec:changetracking}) without comparing a single pair of trees.

\section{A worked trace: a windowed read of a large resource}
\label{sec:readtrace}

The update trace ended with a read; this one follows a read that must \emph{not} materialise the whole
document. Suppose a \code{res} of tens of gigabytes and a client that wants one screen's worth of the
subtree under node~42, as of revision~17:

\dcap{A windowed GET and the bounded slice it returns}
\begin{diagramS}
GET /db/res?revision=17&nodeId=42&maxLevel=2&maxChildren=50
           &withMetaData=nodeKeyAndChildCount&startNodeKey=0

  route by path depth → (db, res) ; validate Bearer JWT ; check the read role
  beginNodeReadOnlyTrx(17)         no write permit taken — readers never lock (Ch.10)
  moveTo(42)                       O(h) trie navigation to the subtree root
  serialise subtree, but bounded:  stop at depth maxLevel=2, at most maxChildren=50 per node,
                                   each record tagged {nodeKey, childCount} (withMetaData)
  emit a continuation cursor:      nextTopLevelNodes = [last serialised top-level node key]
\end{diagramS}

Three things make this cheap. The open is \textbf{lock-free}: \code{beginNodeReadOnlyTrx(17)} pins the
immutable revision-17 trie and takes no semaphore, so it never contends with a concurrent writer
appending revision~18. The traversal is \textbf{bounded at the source}: \code{maxLevel},
\code{maxChildren}, and \code{numberOfNodes} stop the serialiser inside the engine, so a 30\,GB subtree
costs one screen of work, not 30\,GB of serialisation. And pagination is \textbf{just another node key}:
the response carries \code{nextTopLevelNodes}, and the next request passes it back as \code{startNodeKey},
so the client walks a huge resource window by window with each request resolved by the same $O(h)$
navigation as the first. A follow-up that needs a deeper slice raises \code{maxLevel} or re-roots at a
child's \code{nodeId}; a query rather than a navigation swaps the path parameters for
\code{query=\ldots}, runs through brackit with the same revision-pinned read transaction, and is paged by
\code{startResultSeqIndex}/\allowbreak\code{endResultSeqIndex} instead. That windowed, revision-pinned
contract is the whole external surface of the engine; the rest of this chapter exercises it --- first end
to end through both faces, then pinned down by the guarantees it rests on.

\section{A complete worked example}
\label{app:session}

The chapter has traced each face's mechanics in turn; one short scenario now runs end to end through both,
so the pieces are visible together: create a resource, write a document, correct one value in a second
revision, then read the past, diff, and list a node's history. The two listings do the same thing through
the embedded Java API and over REST/JSONiq.

\dcap{Embedded Java: create, write two revisions, then read the past without a lock}
\begin{diagramS}
Databases.createJsonDatabase(new DatabaseConfiguration(path));         // lay out the on-disk folders
try (var db = Databases.openJsonDatabase(path)) {
  db.createResource(ResourceConfiguration.newBuilder("res")
                       .useDeweyIDs(true).hashKind(HashType.ROLLING).build());
  try (var session = db.beginResourceSession("res")) {

    try (var wtx = session.beginNodeTrx()) {                           // acquires the one write permit
      wtx.insertSubtreeAsFirstChild(
            JsonShredder.createStringReader("{\"price\":12}"));        // stream-shred, no DOM
      wtx.commit("initial");                                          // -> revision 1
    }                                                                  // permit released on close

    try (var wtx = session.beginNodeTrx()) {
      wtx.moveTo(priceNodeKey).trx().setNumberValue(19.99);            // correct the value
      wtx.commit("price correction");                                 // -> revision 2
    }

    try (var r1 = session.beginNodeReadOnlyTrx(1)) { /* sees 12    */ } // no lock taken
    try (var r2 = session.beginNodeReadOnlyTrx(2)) { /* sees 19.99 */ } // concurrent readers OK
  }
}
\end{diagramS}

\dcap{The same over REST and JSONiq}
\begin{diagramS}
# create the resource and write revision 1 (streaming shredder); Bearer JWT on every call
PUT /db/res        Content-Type: application/json        { "price": 12 }

# correct the value as an atomic JSONiq update -> revision 2
POST /              { "query": "replace json value of jn:doc('db','res').price with 19.99" }

# read the past: revision 1 still shows 12, with no reconstruction
GET  /db/res?revision=1            ->   { "price": 12 }
GET  /db/res                       ->   { "price": 19.99 }     # latest

# what changed between the two revisions, and when did this node ever change?
POST /   { "query": "jn:diff('db','res', 1, 2)" }              # an edit script: update price
POST /   { "query": "jn:all-times(jn:doc('db','res').price)" } # 12 then 19.99, oldest to newest
\end{diagramS}

Three things in this small run are the thesis in miniature. The second writer waited for nothing but the
first writer's \code{close} (one permit per resource), and the two readers took no lock at all. Reading
revision~1 after revision~2 committed is an ordinary $O(h)$ traversal of an immutable snapshot, not an
undo or a log replay, because revision~2 only \emph{appended} the copied price node and a new
revision-root, leaving revision~1's pages untouched (Appendix~\ref{app:format}). And ``what changed''
and ``when did it change'' are answered by the stored diff and the \code{RECORD\_TO\_REVISIONS} index
(\S\ref{sec:changetracking}), not by comparing trees. Everything else in the report is detail beneath
these three moves.

\section{Worked REST examples}
\label{app:restsessions}

The embedded walkthrough of \S\ref{app:session} has a network twin. These are the concrete
HTTP exchanges for the common operations against the Vert.x server of this chapter; every
request carries \code{Authorization: Bearer <jwt>} (omitted for brevity), and the server speaks HTTP/2
over TLS on port 9443 by default.

\dcap{Create, write, and read back over REST}
\begin{diagramS}
# create a database + first resource by streaming a JSON body (the shredder ingests it)
PUT /db/store        Content-Type: application/json
     { "book": [ { "title":"A", "price":12 }, { "title":"B", "price":8 } ] }
←  200 OK   ETag: "<rootHash>"

# correct a value with an atomic JSONiq update → a new revision
POST /
     { "query": "replace json value of jn:doc('db','store').book[0].price with 19.99" }
←  200 OK

# read the latest, then read the past — same shape, no reconstruction
GET  /db/store                       ←  { "book":[ {"title":"A","price":19.99}, {"title":"B","price":8} ] }
GET  /db/store?revision=1            ←  { "book":[ {"title":"A","price":12},    {"title":"B","price":8} ] }
\end{diagramS}

\dcap{History, diff, and a windowed slice over REST}
\begin{diagramS}
# the revision list (number + commit timestamp + author)
GET /db/store/history                ←  [ {"revision":2,"timestamp":"…","author":"…"}, {"revision":1, …} ]

# the edit script between two revisions (the stored change format, §7.4)
GET /db/store/diff?first-revision=1&second-revision=2
←  { "database":"db","resource":"store","old-revision":1,"new-revision":2,
     "diffs":[ {"update":{"nodeKey":…,"path":"/book/[0]/price","type":"number","value":19.99}} ] }

# one bounded window of a large resource (depth + fan-out caps + metadata + cursor)
GET /db/store?nodeId=1&maxLevel=2&maxChildren=50&withMetaData=nodeKeyAndChildCount&startNodeKey=0
←  { …bounded subtree…, nextTopLevelNodes:[<lastKey>] }    # pass back as startNodeKey to page on
\end{diagramS}

\dcap{A resource-scoped query and its plan}
\begin{diagramS}
# run a predicate query scoped to one resource; ask for the optimised plan alongside
POST /db/store?plan=optimized
     { "query": "jn:doc('db','store').book[][?$$.price gt 10].title" }
←  200 OK   { "result":[ "A" ],
              "plan": { …the cost-based optimiser's index-aware plan, \S\ref{app:brackitops}… } }
\end{diagramS}

\noindent Every one of these is the same engine the embedded API reaches; the REST layer adds only
routing, authentication (\S\ref{app:rest}), streaming ingestion, and windowed serialisation over the
identical transactions, which is why a result read over HTTP is byte-for-byte what an embedded reader
would see at the same revision.

\section{Transaction isolation: guarantees and anomalies}
\label{app:isolation}

The examples took for granted that the API behaves predictably when a writer and many readers overlap;
this is the place to say exactly what it guarantees. The transaction model above gave the
\emph{mechanism} --- a per-resource \code{Semaphore} with one permit, and lock-free, revision-pinned
readers; the question now is what \emph{isolation level} that mechanism delivers, and why, against the
standard catalogue of concurrency anomalies.

\textbf{The claim.} \emph{Per resource, every execution is serializable}, not merely snapshot-isolated.
The argument is short. All write transactions on a resource run in a total order, because the single
write permit forces them to: revision $r{+}1$ cannot begin its commit until revision $r$'s writer has
released the permit, so the revision sequence \emph{is} a serial order of the writes. Every read
transaction observes one immutable committed revision (Property~\ref{prop:snapshot}) and writes nothing,
so it can be placed in that serial order immediately after the writer that produced the revision it
reads. No execution can therefore diverge from a serial schedule. The one anomaly that ordinarily
separates snapshot isolation from serializability, \emph{write skew} (two concurrent writers each
reading what the other is about to change), is impossible here for the same reason lost updates are:
there is never a second concurrent writer.

\begin{property}[Per-resource serializability]\label{prop:serializable}
For a single resource, any interleaving of one writer and $N$ readers is equivalent to a serial schedule
in which the readers are ordered immediately after the writers whose revisions they observe.
\end{property}

\begin{table}[h]
\centering\small\setlength{\tabcolsep}{5pt}\renewcommand{\arraystretch}{1.25}
\begin{tabularx}{\linewidth}{@{}l c X@{}}
\toprule
\textbf{Anomaly} & \textbf{Possible?} & \textbf{Why}\\
\midrule
Dirty read & no & a writer's uncommitted changes live in its private Transaction Intent Log; readers see only published revisions\\
Dirty write & no & one write permit, so no second writer can overwrite uncommitted state\\
Lost update & no & single writer per resource; there is no concurrent writer to lose an update to\\
Non-repeatable read & no & a read transaction is pinned to one immutable revision; re-reading a node returns the same value\\
Phantom read & no & the whole revision trie is immutable; a re-run query cannot see nodes that did not exist in its snapshot\\
Read skew & no & all reads in a transaction come from one consistent snapshot\\
Write skew & no & requires two concurrent writers; the single permit forbids it\\
\bottomrule
\end{tabularx}
\caption{Standard concurrency anomalies against SirixDB's single-writer, snapshot-reader model. Every
anomaly is precluded by construction, which is what makes per-resource execution serializable
(Property~\ref{prop:serializable}).}
\label{tab:anomalies}
\end{table}

\textbf{Scope and cost.} The guarantee is \emph{per resource}: a transaction targets one resource, and
SirixDB does not offer a multi-resource atomic commit, so a unit of work spanning two resources is two
serializable transactions, not one. The cost of buying serializability this cheaply is the write-throughput
ceiling of Chapter~\ref{ch:eval}: a single resource serializes its writers rather than running them
concurrently and reconciling conflicts afterward. The design's wager, argued in the conclusion, is that
document edits are usually authored serially and that real write concurrency lives \emph{across}
resources, where these per-resource serial orders compose without coordination because no transaction
spans them.

\section{The concurrency and threading model}
\label{app:threads}

Serializability is the guarantee; the threading model is how it is delivered without surrendering
concurrency everywhere else. The report's threads and executors have been introduced where each belonged;
gathering them into one map makes the shape clear. The organising fact is that \emph{writes to one resource
are serial} (one permit) while everything else --- reads, eviction, prefetch, query parallelism, network
handling --- runs concurrently around that serial spine.

{\small
\begin{xltabular}{\linewidth}{@{}>{\RaggedRight\arraybackslash}p{0.27\textwidth} >{\RaggedRight\arraybackslash}p{0.21\textwidth} >{\RaggedRight\arraybackslash}X@{}}
\toprule
\textbf{Component} & \textbf{Thread / executor} & \textbf{Role and bound}\\
\midrule
\endhead
Write serialization & per-resource \code{Semaphore(1)} & the single-writer spine; a second writer waits 5\,s then fails fast (this chapter)\\
Readers & caller threads & any number, lock-free, revision-pinned; register in the epoch tracker (\S\ref{sec:reclaim})\\
Cache eviction & clock-sweeper threads (one per page cache) & wake every 100\,ms, second-chance sweep $\sim$10\% of entries, skip guarded pages (\S\ref{sec:reclaim})\\
Commit serialisation & a parallel pool & dirty \code{KeyValueLeafPage}s serialise in parallel when more than four (\S\ref{sec:commit})\\
Async bulk flush & background thread & \code{KEEP\_OPEN\_ASYNC} flushes the rotated TIL off the writer thread (this chapter)\\
Timed auto-commit & a scheduled executor & flushes a time-bounded writer on its interval (this chapter)\\
Prefetch & virtual threads & fire-and-forget next-sibling-leaf prefetch (HOT range cursor); parallel revision opens for \code{sdb:item-history} (\S\ref{sec:changetracking})\\
Query parallelism & \code{ForkJoinPool} & morsel-driven workers, each a pipeline clone over a 2048-tuple morsel (\S\ref{sec:setoriented})\\
Pipeline parallelism & a producer thread & \code{Parallelizer} fills 2000-tuple buffers for consumer threads (\S\ref{sec:setoriented})\\
REST request handling & Vert.x event loop {+} worker pool & one \code{CoroutineVerticle}; blocking calls dispatched \code{executeBlocking(\ldots, ordered=false)} so requests run in parallel (this chapter)\\
Streaming ingest & a \code{Dispatchers.IO} consumer & drains the shredder channel and does the blocking inserts; back-pressured by the bounded channel (this chapter)\\
Replication (enterprise) & Kafka producer / consumer threads & leader publishes \code{TransactionEvent}s; followers apply idempotently (Ch.~\ref{ch:enterprise})\\
\bottomrule
\end{xltabular}}

\noindent The safety of all this concurrency around a serial writer reduces to three guarantees proved or
described earlier: a reader sees an immutable snapshot (Property~\ref{prop:snapshot}), eviction cannot
free a frame a reader has guarded (\S\ref{sec:reclaim}), and the single permit makes per-resource
execution serializable (Property~\ref{prop:serializable}). Everything in the table is free to run
concurrently precisely because none of it can violate those three.

With both faces of the engine described, exercised end to end, and pinned to their isolation and threading
guarantees, the windowed, revision-pinned contract is complete. The browser client of the next chapter is
built entirely on it.

\chapter{The Web GUI}
\label{ch:gui}

\noindent\textit{(SolidJS {+} Vite {+} TypeScript.)}

The Web GUI is the browser client over the REST API of Chapter~\ref{ch:rest}.
The front-end uses SolidJS fine-grained signals/stores for UI state and \textbf{TanStack Query}
for server state, with structured query keys enabling surgical cache invalidation after a
commit and \code{staleTime: Infinity} on immutable revision diffs. Vite proxies \code{/api} to
the SirixDB REST server of Ch.~\ref{ch:rest}; \textbf{Monaco} is dynamically imported on the query page with a
JSONiq Monarch grammar and a ``Midnight Ember'' theme; Kobalte {+} Tailwind provide accessible,
themed primitives.

The pages map onto the temporal model (Figure~\ref{fig:gui}): an \textbf{Explorer} that lazily
expands the node tree via \code{maxLevel}/\allowbreak\code{maxChildren}/\allowbreak\code{startNodeKey} cursors
(handling 100~GB+ resources, with windowed rendering and a path-summary view), a read-only
historical mode keyed on \code{?revision=N}, and a pending-changes model that batches
inserts/updates/deletes into a \textbf{single atomic JSONiq commit}
(\code{replace json value of sdb:select-item(...)}, \code{insert}, \code{delete},
\code{sdb:commit}); a \textbf{Query} editor with templates, envelope stripping, and
parsed/optimised/candidate \textbf{query-plan} visualisation; a \textbf{History} page with
time-travel (floor-semantics revision-at-timestamp) and revision comparison; a \textbf{Diff}
page that merges consecutive diffs client-side for non-consecutive revision pairs (cancelling
insert{+}delete, folding update{+}delete) with a revision scrubber; and a
\textbf{Sunburst/Treemap} view built from the path summary or paginated data.

\begin{figure}[p]
\centering
\begin{minipage}{0.48\linewidth}\centering\includegraphics[width=\linewidth]{gui-explorer.png}\\[1pt]{\footnotesize (a) Explorer --- lazy node tree}\end{minipage}\hfill
\begin{minipage}{0.48\linewidth}\centering\includegraphics[width=\linewidth]{gui-query.png}\\[1pt]{\footnotesize (b) Query --- Monaco JSONiq editor}\end{minipage}\\[7pt]
\begin{minipage}{0.48\linewidth}\centering\includegraphics[width=\linewidth]{gui-history.png}\\[1pt]{\footnotesize (c) History --- time travel}\end{minipage}\hfill
\begin{minipage}{0.48\linewidth}\centering\includegraphics[width=\linewidth]{gui-diff.png}\\[1pt]{\footnotesize (d) Diff --- revision comparison}\end{minipage}\\[7pt]
\begin{minipage}{0.48\linewidth}\centering\includegraphics[width=\linewidth]{gui-sunburst.png}\\[1pt]{\footnotesize (e) Sunburst --- path-summary view}\end{minipage}\hfill
\begin{minipage}{0.48\linewidth}\centering\includegraphics[width=\linewidth]{gui-treemap.png}\\[1pt]{\footnotesize (f) Treemap --- structural weight}\end{minipage}
\caption{The SirixDB Web GUI. Each view is a projection of the temporal model: lazy windowed
exploration (a) over 100\,GB+ resources, a JSONiq editor with plan visualisation (b), time
travel by revision or timestamp (c), client-merged semantic diffs (d), and path-summary-driven
sunburst/treemap visualisations (e,~f).}
\label{fig:gui}
\end{figure}

The \code{SirixApiClient} (\pp{src/lib/api.ts}) wraps the REST endpoints and tolerantly parses
responses (multi-object, raw-text, and single-named-scalar-unwrap fallbacks). A
well-documented subtlety is the \textbf{wire format}: SirixDB serialises object entries either
as legacy nested \code{\{metadata,value\}} wrappers or as \emph{fused} \code{OBJECT\_NAMED\_*}
records whose children sit as a bare array; a single normaliser (\pp{src/lib/sirixTree.ts})
re-wraps the fused shape so one parser handles both, avoiding the class of bug where duplicated
parsers drift. Auth is Keycloak JWT with epoch-guarded refresh (a refresh resolving after
logout cannot resurrect the session) and a \code{waitForAuth} gate that blocks data requests
until a token is ready; \code{VITE\_DEMO\_MODE} yields a read-only cloud-sandbox build.

\section{Lazy, windowed exploration of very large resources}

The Explorer must render a tree that may not fit in memory, let alone the DOM, so it never loads a
resource whole. Four REST controls compose into a windowing protocol (Ch.~\ref{ch:rest}): a depth cap
(\code{maxLevel}), a per-node fan-out cap (\code{maxChildren}), a total node budget
(\code{numberOfNodes}), and a top-level cursor (\code{startNodeKey} {+} \code{nextTopLevelNodes}). The
client expands a subtree on demand by re-fetching the clicked node with a deeper \code{maxLevel}, and it
pages long sibling lists by walking the cursor from the last loaded key, so memory and render cost track
the \emph{visible} window rather than the document. Each fetch is keyed in TanStack Query by
\code{(resource, revision, nodeKey, window)}; because a committed revision is immutable, those entries
are cached with \code{staleTime: Infinity} and never re-validated.

\dcap{A windowed cursor walk: paging a 100\,000-child root, 50 at a time}
\begin{diagramS}
request 1   GET /db/res?nodeId=1&maxLevel=1&maxChildren=50&startNodeKey=0
            ◀ 50 children (keys 2..51)  + nextTopLevelNodes=[51]      cache key (res,rev,node 1,win 0)
request 2   GET /db/res?nodeId=1&maxLevel=1&maxChildren=50&startNodeKey=51
            ◀ next 50 (keys 52..101)    + nextTopLevelNodes=[101]     cache key (res,rev,node 1,win 1)
expand one  GET /db/res?nodeId=37&maxLevel=2&maxChildren=50          deepen ONLY the clicked subtree

memory and DOM ∝ the visible window, never the 100 000-child document;
immutable committed revisions ⇒ every window is cached staleTime:Infinity and never re-fetched
\end{diagramS}

\section{Client-side diff merging}

The REST \code{/diff} endpoint compares two revisions, but the revision scrubber lets a user compare any
pair, including non-consecutive ones. Rather than ask the server for every pairwise diff, the client
fetches the consecutive per-commit diffs (the stored change format of \S\ref{sec:diff}) and merges them.
The merge is a small algebra over operations on the same node: an \code{insert} followed later by a
\code{delete} of the same subtree cancels; an \code{update} followed by a \code{delete} folds to a
\code{delete}; successive \code{update}s collapse to the last. The result is the minimal edit script
between the two chosen revisions, computed without a server round-trip per pair, and rendered with
insert/delete/update colouring against the node tree.

\dcap{Client-side diff merge: composing consecutive edit scripts}
\begin{diagram}
diff(r5 → r6):  insert X ;  update Y = 3
diff(r6 → r7):  delete X ;  update Y = 7

merge (compose on the same node, then cancel / fold):
   X:  insert then delete    → cancels (X is absent at r5 and gone by r7)
   Y:  update 3 then update 7 → keep the last → update Y = 7

diff(r5 → r7) = [ update Y = 7 ]      (computed in the client, no extra server round-trip)
\end{diagram}

A longer, non-consecutive span exercises the rest of the algebra, including \code{replace} (a keyed
delete-plus-insert that the merge must follow across the key change):

\dcap{Merging three consecutive diffs into one r5 → r8 script}
\begin{diagram}
diff(r5 → r6):  insert E (nodeKey 10) ;  update price = 12
diff(r6 → r7):  replace E (nodeKey 10 -> 15) ;  update price = 19
diff(r7 → r8):  delete E (nodeKey 15)

merge, following each node through its key changes:
   E:  insert(10) -> replace(10->15) -> delete(15)   the inserted element is ultimately deleted -> CANCELS
   price:  update 12 -> update 19 -> (unchanged)     -> keep the last -> update price = 19

diff(r5 -> r8) = [ update price = 19 ]    — E never appears: it was born and died inside the span
\end{diagram}

\noindent The \code{replace} is the subtle case: because it renumbers the node (10\,$\to$\,15), the merge
threads the chain by identity across the key change before deciding it cancels. Folding the stored
per-commit diffs this way means comparing two revisions $k$ apart costs the client $k$ small local
compositions and \emph{zero} extra server round-trips, which is what makes scrubbing the revision slider
feel instant.

\section{Reactive state and surgical invalidation}

Two kinds of state meet in the client. \emph{UI} state (selection, expansion, the pending-changes set)
lives in SolidJS signals and stores, whose fine-grained reactivity updates only the affected DOM nodes.
\emph{Server} state lives in TanStack Query under structured keys. The pending-changes model collects a
user's inserts, updates, and deletes locally and flushes them as one atomic JSONiq commit
(\S\ref{sec:worked}); on success the client invalidates exactly the query keys under the affected
resource and revision, so the views that changed refetch while the rest stay put. The effect is an
editing experience that feels local while every committed change is a real, atomic, server-side
revision.

%================================================================
\chapter{Evaluation}
\label{ch:eval}

Having described the whole system, this chapter measures it against PostgreSQL, reproducing the
project's own measurements (\pp{docs/COMPARISON_POSTGRES.md},
\pp{docs/BENCHMARKS.md}) with their caveats~\cite{sirixdocs}. The summary is stated first: \textbf{in the
small-document, single-writer regime, PostgreSQL wins most raw numbers.} The value of the
evaluation is the \emph{analysis} of where each system's advantage actually lies.

\section{Methodology}

The evaluation is built around one principle: a versioned store should be judged on the \emph{shape} of
its cost and the capabilities it enables, not on raw throughput against an engine optimised for the
opposite workload. Three commitments follow. First, every comparison is \textbf{same-machine}: SirixDB
embedded and PostgreSQL server-side loops run on identical hardware, and PostgreSQL is given its most
favourable configuration (no client round-trips). Second, durability is \textbf{verified, not assumed}:
each system's commit path is measured against the disk's own \code{fsync} floor (\code{pg\_test\_fsync},
\code{fio}), so a throughput figure is reported only alongside evidence that the writes were genuinely
durable, since a faster commit rate that quietly weakened durability would make the comparison
meaningless. Third, results are \textbf{reproduced} from the project's own measurement scripts rather
than re-run ad hoc, and both systems are cross-checked to have produced identical histories (the
counter-sum invariant below). Where a number comes from a different machine, it is reported as a
\emph{shape}, a ratio or a trend, never as an absolute. The threats-to-validity section states plainly
what this methodology does \emph{not} establish.

\section{SirixDB vs PostgreSQL (same machine, durability verified)}
\label{sec:w1}

Setup: one $\sim$2.4~KB JSON document, 5\,000 single-field durable updates
$\rightarrow$ 5\,001 retained versions, same i7-12700H / NVMe / ext4. SirixDB embedded
(GraalVM JDK 25, \code{FILE\_CHANNEL}, \code{SLIDING\_SNAPSHOT}); PostgreSQL 17
(\code{synchronous\_commit=on}, \code{fsync=on}), server-side plpgsql loops (no client
round-trips, PostgreSQL's most favourable setup). Durability is \emph{verified}, not
assumed: \code{pg\_test\_fsync} measured \textbf{4\,778 fdatasync/s (209\,$\mu$s)} on the same
volume, and PostgreSQL's 4\,015 commits/s is 84\,\% of that floor, i.e.\ it is fsync-bound
and well tuned. SirixDB's per-commit protocol (write-ahead fsync {+} three FUA
write-throughs) is, if anything, \emph{stronger} than a single fdatasync. Both runs reproduce
the identical 5\,001-version history (cross-checked: counter sum $=$ 12\,502\,500 on both).

\begin{table}[h]
\centering\footnotesize\setlength{\tabcolsep}{4pt}\renewcommand{\arraystretch}{1.3}
\begin{tabular}{@{}>{\RaggedRight}p{2.9cm} >{\RaggedRight}p{2.7cm} >{\RaggedRight}p{2.5cm} >{\RaggedRight}p{2.6cm} >{\RaggedRight}p{3.0cm}@{}}
\toprule
\textbf{Workload} & \textbf{SirixDB (full)} & \textbf{SirixDB (lean)} & \textbf{PostgreSQL 17} & \textbf{Winner}\\
\midrule
\textbf{W1} 5\,000 durable single-field commits & 375 commits/s (2.66\,ms) & 429 commits/s (peak 555) & \textbf{4\,015} commits/s (0.25\,ms) & PostgreSQL \textbf{5.5--10.7$\times$}\\
\textbf{W2} 1\,000 random point-in-time full reads & 75.7\,$\mu$s/read & 74.5\,$\mu$s & 17.5\,$\mu$s batched \textperiodcentered\ $\sim$104\,$\mu$s per-stmt & PG batched 4.3$\times$ --- \textbf{SirixDB wins per-statement}\\
\textbf{W3} list 5\,001 timestamps & 4.57\,ms & 3.36\,ms & 1.99\,ms & PG $\sim$2$\times$ (but see \S\ref{sec:histpath})\\
\textbf{W4} one field across all 5\,001 versions & 55/49\,ms & 49/48\,ms & 6.9\,ms & PG $\sim$7$\times$ (dense field)\\
\textbf{W5} storage for full history & 16.4\,MiB (full) & 11.8\,MiB (lean) & \textbf{4.66\,MiB} & PG 2.5--3.5$\times$\\
\textbf{W6} diff of adjacent versions & \textbf{0.30\,ms} node-level patch & 0.40\,ms & 0.15\,ms top-level only & sub-ms tie; \textbf{SirixDB on capability}\\
\bottomrule
\end{tabular}
\end{table}

\textbf{Why PostgreSQL wins this regime, precisely.} Its per-commit work is one compact WAL
record {+} one fdatasync; SirixDB writes a CoW page-tree (several pages) plus an fsync barrier
and (in the protocol measured) three FUA writes, hashes, and, in the full config, a
per-commit diff file. For history \emph{scans} (W3/W4), a heap scan over 5\,001
pglz-compressed rows beats opening 5\,001 revision contexts ($\sim$10--18\,$\mu$s each). For
\emph{storage} at 2.4~KB, a compressed full copy (836\,B/row) beats SirixDB's per-revision
metadata floor ($\sim$2.4--3.4~KB/version: revision root {+} indirect pages {+} record
fragment), since the floor exceeds a compressed full copy when the document is barely larger than
the floor.

\textbf{Where SirixDB wins, and the cost-shape argument.} Per-statement read latency from an
application is \textbf{75.7\,$\mu$s embedded vs $\sim$104\,$\mu$s} for client-server PostgreSQL
even over a local unix socket; there is no batching trick for an app that needs one document
\emph{now}. Versioning is a first-class capability, not a pattern to assemble: numbered
\emph{and} timestamped revisions, stable node identity across versions, time-travel axes,
per-node history indexes, audit-grade rolling hashes, and \textbf{node-level semantic diffs in
0.3\,ms}, for which the PostgreSQL side must ship both full documents to the application and
diff there. Most importantly, the \textbf{cost shape}: SirixDB's per-version cost is
$O(\text{changed nodes}) + \text{fixed metadata}$, \emph{independent of document size};
PostgreSQL's is $O(\text{document})$ per version (a full jsonb copy in heap \emph{and} WAL). At
2.4~KB this is PostgreSQL's win; the storage/write-amplification crossover sits in
the tens-of-KB document range, \emph{now measured} on 100~KB\,/\,1~MB\,/\,10~MB documents
(\S\ref{sec:crossover}): the storage crossover is already past 100~KB and the commit-speed crossover
is near 1~MB.

The project's own positioning conclusion is reproduced verbatim in spirit: \textbf{do not pitch
SirixDB as ``faster than PostgreSQL for keeping history of small documents'', it is not.}
The pitch is the cost shape, the capability set, and embedding.

\section{History-path improvements (different machine; shape, not absolutes)}
\label{sec:histpath}

A reworked history read path (\code{getHistoryTimestamps()} from the resident in-memory
\code{RevisionIndex}, a \code{long[]} {+} one \code{arraycopy}) serves all 5\,001 timestamps in
\textbf{$\sim$0.05\,ms, $\sim$40$\times$ faster} than the PostgreSQL heap scan and the prior
SirixDB path. For a field that changes \emph{rarely} (the common real shape),
\code{scanValueRuns} reads \textbf{one} record and reconstructs the value across all 5\,001
versions in \textbf{$\sim$0.05\,ms vs PostgreSQL's $\sim$9.8\,ms ($\sim$200$\times$)}, the
$O(\text{changes})$ cost shape made visible. For a field that changes on \emph{every} commit,
the change set is everything, so all paths still read 5\,001 records and PostgreSQL's
compressed heap scan still leads ($\sim$8$\times$); closing that dense case needs a deferred
fragment-chain single-pass scan, left as future work.

\section{Concurrency and the revision-depth pathology (and its fix)}
\label{sec:concurrency}

A REST concurrency benchmark validated that moving blocking work off Vert.x's \emph{ordered}
worker queue (\code{ordered = false}) cut the p95(c{=}16)/p95(c{=}1) ratio from a
queue-serialised $\sim$16$\times$ to \textbf{$\sim$3$\times$}, with read throughput scaling
1$\rightarrow$16 by 4.85$\times$ before CPU saturation. It also surfaced a pathology: reading
the \emph{latest} revision of a same-sized document became $\sim$4$\times$ slower (c{=}1) and
$\sim$10--12$\times$ lower throughput (c{=}16) once the resource carried $\sim$1\,400 revisions,
and \emph{fully recovered} after re-seeding to 6 revisions on the same JVM, ruling out
aging/GC. Root cause (\S\ref{sec:evolution}): every storage open eagerly reloaded the
per-revision index, $O(\text{revisions})$ per request, and a hidden $O(R^2)$ \code{access(2)}
walk. After the fix, read performance at 1\,900 revisions \textbf{exceeds} the fresh-resource
baseline (c{=}16: 1\,042 $\rightarrow$ \textbf{18\,361 req/s}, p99 245\,ms $\rightarrow$
\textbf{1.84\,ms}), and the large-history build is flat at $\sim$570 commits/s through 10\,000
commits. Random-revision access is \textbf{position-independent}: trx open{+}read is
$\sim$18\,$\mu$s warm whether the target is revision 1, 5\,000, or 10\,000, a direct
dividend of structural sharing.

\textbf{Lock-free reads, measured under JMH.}\label{par:lockfreeread} The report repeatedly claims reads
take no lock and so scale with cores (Property~\ref{prop:snapshot}); this measures it with the harness of
\S\ref{sec:commitphases}. On the 20-core host, $T$ benchmark threads each hold their own read-only
transaction over one shared, fully-cached 1\,MB resource and do random point reads; aggregate throughput,
with 99.9\,\% confidence intervals:

\dcap{Aggregate point-read throughput vs.\ reader threads (JMH, lock-free; mean $\pm$ CI)}
\begin{diagramS}
threads   reads/s            speedup
  1        7.8 M  +/- 0.2     1.00
  2       11.3 M  +/- 0.2     1.44     rising, not flat: there is no read lock to serialise on
  4       16.2 M  +/- 2.1     2.08
  8       24.5 M  +/- 0.2     3.15
  16      27.4 M  +/- 0.3     3.52     flattening: shared buffer pool + memory bandwidth, not a lock
\end{diagramS}

The shape is the point. A read \emph{lock} would pin throughput near the single-thread 7.8\,M/s however
many threads piled on; instead it climbs monotonically to $3.5\times$ at sixteen, the signature of
genuinely parallel, lock-free reads (a reader takes no write permit and registers only an epoch slot,
\S\ref{sec:reclaim}). The scaling is \emph{sub}-linear and flattens past eight threads, but the ceiling is
the shared resources every reader touches --- the JVM-wide buffer pool, the off-heap frames, and memory
bandwidth on a fully-cached document --- not a serialisation point. (A hand-rolled tight-loop measurement
reported a higher $4.75\times$ at sixteen; as with the write scaling, JMH's warmup and forks give the
trustworthy, more conservative figure, and the report quotes the JMH one.) This measures read concurrency
only; \emph{write} concurrency within a resource is, by design, one writer
(Property~\ref{prop:serializable}), so the scaling axis for writes is across resources, not threads.

That across-resource axis is itself measurable, and the claim deserves a number measured \emph{properly}.
Under JMH (the harness of \S\ref{sec:commitphases}, which forks fresh JVMs and warms each before timing),
each thread owns a \emph{distinct} resource and commits a stream of small durable updates (8 changed slots,
1\,MB resource, \code{FILE\_CHANNEL}, fsync per commit); aggregate throughput at increasing thread counts,
with 99.9\,\% confidence intervals over 10 samples:

\dcap{Write scaling across resources (JMH, 8-slot durable commits; mean $\pm$ CI)}
\begin{diagramS}
threads (= resources)   commits/s        speedup
  1                      114  +/- 7        1.00      single-resource rate ~ 8.7 ms/commit
  2                      214  +/- 10       1.87
  4                      362  +/- 33       3.17
  8                      553  +/- 73       4.83      SUB-linear: per-thread latency rises 8.7 -> 14.5 ms
\end{diagramS}

The scaling is real but \emph{sub}-linear: eight independent writers deliver $4.8\times$ the throughput of
one, not eight, because concurrent durable commits contend on what they share --- the one disk their fsyncs
hit, the single JVM-wide buffer pool, and the off-heap allocator --- so per-thread commit latency climbs
from 8.7\,ms at one writer to 14.5\,ms at eight. The single-resource rate (8.7\,ms/commit, lean
configuration) agrees with the commit-latency benchmark of \S\ref{sec:commitphases} rather than with the
bulk-ingest one: a small commit is light serialisation plus the durability barrier, not the $\sim$83\%
serialisation of a full rewrite.

This number carries a methodology lesson, because getting it wrong was instructive. Two earlier hand-rolled
attempts reported \emph{super}-linear scaling ($14.6\times$, then $8.9\times$ after a warmup sweep was
added); both were artifacts --- a cold single-thread baseline inflated by JIT, and a smaller working set
--- that a hand-rolled timing loop cannot reliably exclude. JMH's forks and warmup iterations remove
exactly that class of error, and its verdict, sub-linear, is the one to trust; the convergence of three
successively more careful measurements ($14.6\to8.9\to4.8\times$) is itself the argument for using the
harness. The qualitative takeaway survives the correction: the single-writer-per-resource rule gives up
little, because writing many resources at once recovers \emph{most} of the parallelism a multi-writer
design would chase within one resource --- just not, as the rigorous measurement shows, more than all of
it.

\begin{figure}[p]
\centering
\begin{minipage}{\linewidth}\centering\includegraphics[width=\linewidth]{flame-cpu.png}\\[1pt]{\footnotesize (a) CPU flame graph}\end{minipage}\\[8pt]
\begin{minipage}{\linewidth}\centering\includegraphics[width=\linewidth]{flame-wall.png}\\[1pt]{\footnotesize (b) Wall-clock flame graph}\end{minipage}\\[8pt]
\begin{minipage}{\linewidth}\centering\includegraphics[width=\linewidth]{flame-alloc.png}\\[1pt]{\footnotesize (c) Allocation flame graph}\end{minipage}
\caption{Representative flame graphs from the profiling that drove the optimisations of
Chapters~\ref{ch:ondisk} and~\ref{ch:vector}: CPU time (a), wall-clock (b, which exposed the
quadratic \code{access(2)} pathology of \S\ref{sec:evolution}), and allocation (c). The
methodology (profile, attribute cost, fix the dominant term, re-measure) is the same one
that turned the large-history build from 48.4\,s to 20.5\,s.}
\label{fig:flame}
\end{figure}

\section{Cost-shape analysis}

The numbers above are specific to one machine and one workload; the claim they support is general and
worth stating as a model. Let a resource hold $D$ nodes and let a transaction change $c$ of them.

\textbf{Write.} An update-in-place engine pays roughly $O(c)$ to mutate the touched records plus its
write-ahead-log overhead. SirixDB pays $O(c + h)$, where $h = \lceil\log_{1024} D\rceil$ is the trie
height, because the change also copies the indirect pages along each touched leaf's root-to-leaf path,
and it pays a fixed per-revision metadata floor $m$ (a new revision-root page, the dual beacons, the
index deltas). For a small document $h$ and $m$ dominate $c$, and the in-place engine wins; as $c$ grows
the copied path amortises over many changed records and the floor becomes negligible, so the two
converge. That crossover is now \emph{located} (\S\ref{sec:crossover}): near 1\,MB for commit speed and
already past 100\,KB for storage, with SirixDB $\sim$24$\times$ faster and $\sim$19$\times$ smaller at
10\,MB.

\textbf{Storage.} An in-place engine that keeps no history stores $O(D)$. SirixDB after $k$ commits
stores $O(D + \sum_i c_i)$: the initial document plus the sum of all changes, never $O(kD)$. The price
of retaining all history is therefore proportional to the \emph{total change}, not to the document size
times the number of revisions, which is what makes deep histories affordable.

\textbf{Read.} A current-version point read is $O(h)$ for both. A \emph{historical} read is where the
shapes diverge: reconstructing an old version in a log-based system is $O(\text{log replayed})$ or a
validity-interval join, whereas in SirixDB it is a direct $O(h)$ traversal of a snapshot already on
disk, touching at most $R$ fragments per page. The temporal questions, ``what did this node look like at
revision $r$'' and ``in which revisions did it change'', are $O(h)$ and $O(\text{changes})$
respectively, with no scan of unrelated history.

A single set of numbers makes the storage shape concrete. Take a resource of $D=10^{6}$ nodes, and
$k=100$ commits that each change $c=100$ nodes:

\dcap{Storage and cost for $10^6$ nodes, 100 commits of 100 changes each}
\begin{diagramS}
trie height        h = ceil(log_1024 10^6) = 2          (at most 2 indirect levels + leaf)

snapshot-per-commit (naive)   k · D            = 100 · 10^6 = 100 000 000 node-slots   (100× the data)
SirixDB (sliding snapshot)    D + Σ c_i        = 10^6 + 100·100 = 1 010 000 node-slots (≈ +1% for full history)
in-place, no history          D                = 1 000 000 node-slots                  (keeps nothing)

write per commit              O(c + h) + floor = 100 + 2 copied pages + (1 revroot, 2 beacons, index deltas)
read current node             O(h)             = 2 indirections + ≤ R fragment reads
read node @ revision r        O(h)             = same shape; no log replay, no validity-interval join
"when did it change?"         O(changes)       = a RECORD_TO_REVISIONS lookup, not a tree-pair scan
\end{diagramS}

\noindent Retaining a hundred revisions of complete history costs about one percent over the live data,
because the bill is the \emph{total change} ($\sum c_i$), not the document size times the revision count
($kD$) that a snapshot-per-commit store would pay. That single contrast, $+1\%$ versus $100\times$ for
the same retained history, is the cost shape the whole design exists to produce.

The honest conclusion is the one the measurements show: for a single small document updated in place,
pay-as-you-go beats pay-the-floor; for data whose history must be kept and queried, paying $O(c)$ per
revision and reading the past in $O(h)$ is the better cost shape. The evaluation establishes the first
empirically and the second both empirically, through the history-read results, and by the argument here.

\section{Storage growth, measured: the change is the bill}
\label{sec:storagegrowth}

The cost-shape claim most worth seeing as a curve is that retained-history storage grows with the
\emph{total change} $\sum c_i$, not the document size times the revision count $kD$. A 1\,MB document is
committed 100 times, each commit changing 20 scattered values, and the total on-disk size is sampled along
the way; the ``snapshot-per-commit'' column is the $kD$ a copy-everything store would pay for the same
history.

\begin{table}[h]
\centering\footnotesize\setlength{\tabcolsep}{8pt}\renewcommand{\arraystretch}{1.25}
\begin{tabular}{@{}r r r@{}}
\toprule
\textbf{after commit} & \textbf{SirixDB (MB)} & \textbf{snapshot-per-commit (MB)}\\
\midrule
0 (base) & 0.73 & 0.73\\
1        & 0.75 & 1.47\\
20       & 1.92 & 15.4\\
40       & 3.42 & 30.1\\
60       & 4.90 & 44.8\\
80       & 6.40 & 59.5\\
100      & \textbf{7.97} & 74.1\\
\bottomrule
\end{tabular}
\caption{On-disk size as 100 commits of 20 scattered changes accumulate on a 1\,MB document. SirixDB grows
linearly in the commit count at a \emph{shallow} slope ($\sim$0.07\,MB/commit, the changed pages plus fixed
metadata); a snapshot-per-commit store grows ten times steeper to $\sim$9$\times$ the size for identical
history.}
\label{tab:storagegrowth}
\end{table}

The shape is the result. SirixDB's footprint rises \emph{linearly} in the number of commits but at a
shallow slope --- about 0.07\,MB per commit, the cost of that commit's changed pages plus its fixed
metadata --- reaching 8.0\,MB of full history over a 0.73\,MB base. The snapshot-per-commit line rises ten
times steeper, a full 0.73\,MB copy each commit, to 74\,MB: $\sim$9$\times$ larger for the \emph{same}
retained revisions. Those two slopes \emph{are} the two cost shapes, $O(\sum c_i)$ against $O(kD)$. The
absolute ratio depends on how much each commit touches (here 20 scattered slots, each dragging its leaf
fragment and root path), but the slope being change-proportional and revision-count-linear rather than
document-proportional is exactly what makes deep histories affordable, and is the measured form of the
contrast the cost-shape diagram states in the limit.

\section{The versioning strategies, measured (and the window knob)}
\label{sec:versioningbench}

The cost-shape argument is generic, but SirixDB also exposes the trade-off \emph{internally}: the same
logical change can be persisted under four selectable versioning strategies
(\S\ref{sec:strategyspace}) governed by one tuning knob, the reconstruction window
$R=\textit{revisionsToRestore}$. The choice of \code{SLIDING\_SNAPSHOT} with $R=3$ as the default has so
far rested on two \emph{invariants} of \S\ref{sec:slidealg} --- bounded reconstruction
(Property~\ref{prop:bounded}, a read combines at most $R$ fragments) and even amortisation
(Property~\ref{prop:amort}, preservation costs $\approx 1/R$ extra writes per revision). Invariants bound
the shape but not the constant; this section measures it. A single 2\,MB JSON document is ingested under
each strategy and then updated by 100 durable commits, each changing 32 scattered numeric slots. Three
quantities are recorded: write latency per commit, total on-disk storage for the 101 retained revisions,
and a \emph{cold full scan} --- a fresh session open (SirixDB's in-process page caches empty, the OS file
cache warm) that visits every node once, so each leaf page is reconstructed from its fragments exactly
once. The cold scan thus isolates the CPU cost of fragment combination, which is the price a delta
strategy pays back on every read.

\begin{table}[h]
\centering\footnotesize\setlength{\tabcolsep}{7pt}\renewcommand{\arraystretch}{1.25}
\begin{tabular}{@{}l r r r@{}}
\toprule
\textbf{Strategy} & \textbf{write ms/commit} & \textbf{storage (MB)} & \textbf{cold scan (ms)}\\
\midrule
\code{FULL}             & 122.7 & 31.5 & \textbf{207}\\
\code{DIFFERENTIAL}     & 70.8  & 31.9 & 212\\
\code{INCREMENTAL}      & 55.4  & 17.6 & 225\\
\code{SLIDING\_SNAPSHOT} & \textbf{41.4} & \textbf{13.9} & 252\\
\bottomrule
\end{tabular}
\caption{The four versioning strategies on one 2\,MB document, 100 commits of 32 scattered slot
changes, $R=3$, \code{FILE\_CHANNEL}, hashing off, durable commits. Lower is better throughout;
\textbf{bold} marks the best in each column. The strategies trace the read/write-amplification curve:
\code{FULL} reads fastest but writes and stores the most; \code{SLIDING\_SNAPSHOT} writes and stores the
least at a bounded read tax.}
\label{tab:versioning-strategies}
\end{table}

Table~\ref{tab:versioning-strategies} traces the full amplification curve. \code{FULL} rewrites every
touched leaf page in its entirety, so it pays the most to write ($\sim$123\,ms/commit) and stores the most
(31.5\,MB) yet reads fastest (207\,ms), because every page is self-contained --- one load, no merge.
\code{DIFFERENTIAL} writes the cumulative delta against the last full page; with scattered changes that
delta grows until it rivals re-storing the page, so its storage (31.9\,MB) matches \code{FULL} while its
two-fragment reads stay cheap (212\,ms). \code{INCREMENTAL} writes only the change since the previous
revision, cheaper again (55\,ms, 17.6\,MB), but its reconstruction chain is unbounded between the periodic
full dumps. \code{SLIDING\_SNAPSHOT} writes the least (41\,ms/commit) and stores the least (13.9\,MB) by
deferring preservation across the window, at the cost of merging up to $R$ fragments per page on read
(252\,ms). That is a $\sim$21\% read tax over \code{FULL} in exchange for $3.0\times$ cheaper commits and
$2.3\times$ smaller storage \emph{for identical retained history} --- the trade the default is built to
make, since reads hit the combined-page cache (\S\ref{sec:caching}) on the second touch whereas the write
and storage costs are paid on every commit, forever.

The window $R$ is the dial that slides a \emph{single} strategy along that same curve, and sweeping it
turns the two invariants into a measured law.

\begin{table}[h]
\centering\footnotesize\setlength{\tabcolsep}{7pt}\renewcommand{\arraystretch}{1.25}
\begin{tabular}{@{}r r r r@{}}
\toprule
\textbf{$R$} & \textbf{write ms/commit} & \textbf{storage (MB)} & \textbf{cold scan (ms)}\\
\midrule
1  & 66.2 & 31.9 & 209\\
2  & 46.3 & 18.2 & 231\\
3  & 39.9 & 13.9 & 227\\
5  & 35.2 & 10.7 & 255\\
10 & 31.4 & 8.9  & 271\\
\bottomrule
\end{tabular}
\caption{Sweeping the sliding-snapshot window $R$ on the same workload. As $R$ grows the write cost and
storage fall monotonically (a wider window defers more preservation and re-stores fewer redundant slots)
while the cold-scan reconstruction cost rises (more fragments to merge per page) --- the read/write
amplification trade-off of Properties~\ref{prop:bounded} and~\ref{prop:amort}, now in numbers.}
\label{tab:versioning-window}
\end{table}

Three things in Table~\ref{tab:versioning-window} are worth drawing out. First, the endpoints behave
exactly as the algorithm dictates: at $R=1$ the window holds a single fragment, so every fragment must be
self-contained and the strategy \emph{collapses to} \code{FULL} --- its storage (31.9\,MB) and read cost
(209\,ms) match the \code{FULL} row of Table~\ref{tab:versioning-strategies} to within noise. Second, the
write and storage curves are monotone in $R$: widening the window from 1 to 10 makes commits
$2.1\times$ cheaper (66\,$\to$\,31\,ms) and the store $3.6\times$ smaller (31.9\,$\to$\,8.9\,MB), the
direct payoff of even amortisation, while the cold scan climbs $\sim$30\% (209\,$\to$\,271\,ms) as bounded
reconstruction touches more fragments. Third, the gains are sharply concave: most of the write and storage
benefit is already captured by $R=3$ (commits down to 40\,ms, storage to 13.9\,MB), and pushing on to
$R=10$ buys only a further $\sim$5\,MB and $\sim$8\,ms while adding $\sim$60\,ms of read latency. The
default sits at that knee, not at an extreme. (The cold scan carries a run-to-run variance of order 10\%
--- the $R=3$ scan reads 227\,ms here and 252\,ms in Table~\ref{tab:versioning-strategies}, the same
configuration measured twice --- so it is the monotone \emph{trend}, not any single adjacent pair, that
the sweep establishes.)

\dcap{The knob in one line: $R$ trades read amplification for write/storage amplification}
\begin{diagramS}
  R = 1  ── self-contained pages ── writes/stores like FULL,    reads in 1 fragment   (max write amp.)
  R = 3  ── default knee        ── ~3x cheaper writes, ~2.3x smaller, ~+10-20% read   (balanced)
  R = 10 ── deep deferral       ── cheapest writes & storage,   reads merge up to 10  (max read amp.)

  bounded reconstruction (Prop. 3.1): read cost grows with R     (more fragments per page)
  even amortisation      (Prop. 3.2): write/storage shrink ~1/R  (preservation deferred across window)
\end{diagramS}

\noindent The measurement closes the loop opened in \S\ref{sec:slidealg}: the two properties are not only
provable invariants but the visible axes of a tunable trade-off, and the shipped default ($R=3$,
\code{SLIDING\_SNAPSHOT}) is the empirically located balance point between them rather than a guess.

\section{What the structural features cost at ingest}
\label{sec:ingestcost}

The versioning strategy is one write-time knob; the \emph{structural} features are the others. Node
hashing, the path summary, and DeweyID labelling are each optional
(\code{ResourceConfiguration.Builder}), and the report has so far described what they \emph{buy} without
pricing what they \emph{cost}. This section prices them under JMH. A fixed 4\,MB JSON document
($\approx$308\,000 nodes) is shredded into a fresh resource (created and removed per measured invocation,
so the database never accumulates) under a lean baseline (hashing off, no path summary, no DeweyIDs, diffs
off), then with each feature switched on individually, then with all three together. JMH forks a fresh JVM
per configuration, so the cross-configuration JIT-ordering bias is removed by construction; each row is the
mean of 12 samples with a 99.9\,\% confidence interval, the commit durable (\code{FILE\_CHANNEL}, fsync).

\begin{table}[h]
\centering\footnotesize\setlength{\tabcolsep}{8pt}\renewcommand{\arraystretch}{1.25}
\begin{tabular}{@{}l r r r r@{}}
\toprule
\textbf{Configuration} & \textbf{ingest ms} & \textbf{$\pm$CI} & \textbf{MB/s} & \textbf{overhead}\\
\midrule
baseline (lean)        & \textbf{181}  & 5.7  & \textbf{22.1} & ---\\
\quad+ path summary    & 217  & 7.6  & 18.4 & $+20\%$\\
\quad+ rolling hashes  & 283  & 7.8  & 14.1 & $+57\%$\\
\quad+ DeweyIDs        & 285  & 7.3  & 14.0 & $+58\%$\\
all three (full)       & 500  & 12.9 & 8.0  & $+176\%$\\
\bottomrule
\end{tabular}
\caption{Per-feature write-time cost of shredding one 4\,MB document ($\approx$308\,000 nodes), durably
committed, under JMH (\code{Mode.AverageTime}, mean of 12 samples). Each optional structural feature trades
ingest throughput for a query-time capability; the path summary is cheapest because it is maintained per
\emph{distinct path}, the others per \emph{node}. (A hand-rolled interleaved run gave $+23/63/68/194\%$ ---
the same ranking and magnitudes, so this is the one benchmark JMH \emph{confirmed} rather than corrected.)}
\label{tab:ingestcost}
\end{table}

The lean baseline shreds at $\sim$22\,MB/s ($\sim$1.7\,M nodes/s), and the features price out in a
sensible order. The \textbf{path summary} is the cheapest at $+20\%$ because it is maintained per
\emph{distinct path}, not per node: this document has only a handful of distinct paths
(\code{/rows/[]/id}, \code{.../name}, \ldots), so the summary trie is tiny and seldom grows --- a deeply
heterogeneous document with thousands of distinct paths would pay more, which is precisely why the path
summary is a configuration choice rather than a default. \textbf{Rolling hashes} ($+57\%$) and
\textbf{DeweyIDs} ($+58\%$) each add a genuine per-node computation: a 64-bit hash folded up the ancestor
chain for every node, and an \code{int[]} label plus its lazily-encoded byte form for every node.
Enabling \textbf{all three} nearly triples ingest time ($+176\%$, throughput $22\to8$\,MB/s), and the
stack is mildly super-additive --- three per-node bookkeeping passes competing for the same memory
bandwidth cost a little more together than apart. The lesson is the one the design intends: the lean
configuration used by the demo and by the large-document benchmarks (\S\ref{sec:crossover}) is the fast
path, and each structural feature is a deliberate, separately-priced trade of ingest throughput for a
later capability --- hashing for $O(1)$ subtree equality and hash-guided diffs (\S\ref{sec:diff}), the
path summary for index-aware planning (\S\ref{ch:pathsummary}), DeweyIDs for arithmetic document order
and constant-time ancestor tests (\S\ref{sec:deweyids}). None is free, and none is mandatory.

\section{The hash-guided diff, measured (and an honest refinement)}
\label{sec:diffbench}

\S\ref{sec:diff} claims the structural diff costs ``proportional to the changed region, not the document''
through hash pruning. That is half true, and the measurement says which half. A document is ingested under
\code{ROLLING} hashing with \code{storeDiffs} \emph{off} (so the diff is recomputed, not read from the
stored per-commit delta), a revision changes $K$ scattered values, and \code{DiffFactory} (mode
\code{HASHED}) recomputes the edit script. Two counters are recorded: wall time, and the number of
\code{diffListener} callbacks --- one per node the walk \emph{visits}, where a \code{SAMEHASH} callback is
a pruned, never-descended subtree.

\begin{table}[h]
\centering\footnotesize\setlength{\tabcolsep}{7pt}\renewcommand{\arraystretch}{1.25}
\begin{tabular}{@{}l r r r r@{}}
\toprule
\textbf{Document (nodes)} & \textbf{diff ms (JMH)} & \textbf{callbacks} & \textbf{of which SAMEHASH} & \textbf{changed}\\
\midrule
100\,KB ($\sim$8.3\,K)  & 0.20 & 1\,429   & 1\,407   & 10\\
1\,MB ($\sim$79\,K)     & 1.61 & 13\,290  & 13\,268  & 10\\
10\,MB ($\sim$765\,K)   & 25.1 & 127\,496 & 127\,474 & 10\\
\bottomrule
\end{tabular}
\caption{Recomputed hash-guided diff of a 10-value change, varying document size (JMH, mean of 12 samples;
$\pm$CI $\le$0.43\,ms). The callback count (the nodes the walk visits) grows with the document but stays
far below the total node count --- the prune skips descent into each unchanged subtree, and
\code{SAMEHASH}$\approx$\code{callbacks} shows nearly every visited node was dismissed by one comparison.
Time tracks callbacks at $\sim$120--200\,ns each.}
\label{tab:diffbench-size}
\end{table}

\begin{table}[h]
\centering\footnotesize\setlength{\tabcolsep}{7pt}\renewcommand{\arraystretch}{1.25}
\begin{tabular}{@{}r r r r@{}}
\toprule
\textbf{$K$ changed} & \textbf{diff ms (JMH)} & \textbf{callbacks} & \textbf{changed emitted}\\
\midrule
1    & 7.32  & 51\,717 & 1\\
10   & 7.75  & 51\,762 & 10\\
100  & 8.38  & 52\,212 & 100\\
1000 & 11.08 & 56\,657 & 995\\
\bottomrule
\end{tabular}
\caption{Same 4\,MB document ($\sim$310\,K nodes), varying the number of changed values $K$ (JMH, mean of
12 samples, $\pm$CI $\le$0.22\,ms). The callback count is essentially fixed (the array's breadth); only the
time rises, the linear cost of descending into and emitting the $K$ changed subtrees.}
\label{tab:diffbench-k}
\end{table}

The hash prune is real and large: a 10\,MB document of $\sim$765\,000 nodes diffs in $\sim$127\,000
callbacks, not 765\,000 --- the walk pruned the descent into every unchanged row's subtree (an unhashed
\code{DiffOptimized.NO} walk visits each node), and \code{sameHash}$\approx$\code{callbacks} confirms nearly
every visited node was dismissed by a single comparison. But the callback count tracks the \emph{breadth}
the walk enumerates --- here the one wide array's child list --- which grows with the document. So the
prune is \textbf{vertical} (skip unchanged subtrees) and not \textbf{horizontal} (siblings are still walked
child by child): for a wide, flat fan-out the recompute is $O(\text{breadth})$, not
$O(\text{changed region})$. Table~\ref{tab:diffbench-k} isolates the other term: at a fixed document the
callback count barely moves while time rises only modestly ($7.3\to11.1$\,ms as $K$ goes $1\to1000$), the
linear cost of the changed subtrees layered on the fixed breadth scan.

The honest conclusion refines the claim rather than simply confirming it: a hash-guided diff is
$O(\text{siblings enumerated along the changed levels} + \text{changed region})$. It eliminates the
document-sized term \emph{below} unchanged nodes but not the fan-out \emph{beside} them --- which is
exactly why the default keeps a precomputed per-commit delta (\code{storeDiffs}) for the common
adjacent-revision case and reserves this recompute path for arbitrary revision pairs.

\section{Large documents: the measured crossover (vs PostgreSQL and DuckDB)}
\label{sec:crossover}

The small-document benchmark (\S\ref{sec:w1}) is PostgreSQL's home turf, and the cost-shape argument
above predicts the picture inverts as documents grow. That prediction was previously unmeasured; this
section measures it. The same harness is run on documents of \textbf{100\,KB, 1\,MB, and 10\,MB} (scaled
by a large \code{items} array), with the \emph{same tiny update}, the first field, \code{counter}, set
to a new value, committed durably 100 times, so 101 full versions are retained. SirixDB is embedded
(\code{SLIDING\_SNAPSHOT}, \code{FILE\_CHANNEL}), in both its full-feature and lean configurations;
PostgreSQL~17 keeps history with the same \code{jsonb}{+}trigger pattern as~\S\ref{sec:w1}
(\code{synchronous\_commit=on}, server-side loop); DuckDB~1.5.2 retains the 101 versions in a
\code{JSON} column. The update is $O(1)$ in all systems' \emph{intent}; what differs is what each one
\emph{does} with it.

\begin{table}[h]
\centering\footnotesize\setlength{\tabcolsep}{6pt}\renewcommand{\arraystretch}{1.25}
\begin{tabular}{@{}l r r r@{}}
\toprule
\textbf{Document} & \textbf{SirixDB (full)} & \textbf{SirixDB (lean)} & \textbf{PostgreSQL 17}\\
\midrule
\multicolumn{4}{@{}l}{\emph{ms per single-field durable commit (lower is better)}}\\
100\,KB & 9.97 & 7.41 & \textbf{2.09}\\
1\,MB   & 10.66 & \textbf{7.39} & 19.25\\
10\,MB  & 10.13 & \textbf{8.18} & 193.9\\
\bottomrule
\end{tabular}
\caption{Per-commit latency for a tiny update to a growing document. SirixDB's cost is \emph{flat}
($\sim$8--10\,ms) because it writes the changed node plus fixed metadata; PostgreSQL's grows
$\sim$linearly ($2\to19\to194$\,ms) because each \code{jsonb\_set} rewrites the whole document and the
trigger stores a full copy. The commit-speed crossover sits near 1\,MB; at 10\,MB lean SirixDB is
\textbf{$\sim$24$\times$} faster.}
\label{tab:crossover-commit}
\end{table}

\begin{table}[h]
\centering\footnotesize\setlength{\tabcolsep}{6pt}\renewcommand{\arraystretch}{1.25}
\begin{tabular}{@{}l r r r r@{}}
\toprule
\textbf{Document} & \textbf{SirixDB (full)} & \textbf{SirixDB (lean)} & \textbf{PostgreSQL 17} & \textbf{DuckDB 1.5.2}\\
\midrule
\multicolumn{5}{@{}l}{\emph{total on-disk storage for 101 retained versions (lower is better)}}\\
100\,KB & 1.07\,MB & \textbf{0.63\,MB} & 2.26\,MB & 11\,MB\\
1\,MB   & 2.62\,MB & \textbf{1.49\,MB} & 20\,MB & 98\,MB\\
10\,MB  & 17.6\,MB & \textbf{9.99\,MB} & 193\,MB & 965\,MB\\
\bottomrule
\end{tabular}
\caption{History storage for 101 versions. SirixDB stores the document \emph{once} plus per-version
deltas; PostgreSQL and DuckDB store (compressed) full copies per version. At 10\,MB lean SirixDB is
\textbf{$\sim$19$\times$} smaller than PostgreSQL and \textbf{$\sim$97$\times$} smaller than DuckDB.}
\label{tab:crossover-storage}
\end{table}

\textbf{What the numbers say.} The defining row is SirixDB's: its per-commit time barely moves from
100\,KB to 10\,MB ($\sim$8--10\,ms), because a single-field update copies one leaf plus the root-to-leaf
path and a fixed metadata floor, \emph{independent of document size} (Ch.~\ref{ch:pds}); the per-commit
diff is hash-pruned, so it too does not scan the document. PostgreSQL pays the opposite shape: every
\code{jsonb\_set} produces a fresh full document and the trigger appends a full copy, so its per-commit
cost climbs from 2\,ms to 194\,ms across the three sizes, and its retained history from 2\,MB to 193\,MB.
The \textbf{commit-speed crossover is near 1\,MB} (below it PostgreSQL's WAL{+}fsync is unbeatable; above
it SirixDB's delta-write wins, by $\sim$24$\times$ at 10\,MB), and the \textbf{storage crossover is
earlier still}, already past 100\,KB. DuckDB, an analytical engine pressed into a versioned-document
role it is not built for, keeps the largest history of all, it has no structural sharing and its JSON
column barely compresses near-identical large values; the comparison shows only that it offers no
versioning advantage, not that it is a poor analytical engine (raw OLAP scans are its strength and a
different axis from anything measured here).

\begin{honesty}
\small\textbf{Caveat.} This regime is deliberately the mirror image of \S\ref{sec:w1}, and its limits are
the same in reverse: one document, one writer, \emph{101} versions (not thousands), and a single tiny
update shape. \textbf{Full-document reads do not cross over}, serialising a 10\,MB document is
$\sim$210\,ms whatever the store, because the bytes have to be produced; SirixDB's win is on
\emph{writes}, \emph{storage}, and \emph{sub-document} reads, not on regenerating a whole large document.
PostgreSQL's storage is after \code{CHECKPOINT} and excludes WAL; DuckDB's is a single appended table.
SirixDB is the same dev build as \S\ref{sec:w1}; PostgreSQL~17 and DuckDB~1.5.2 are GA releases. The
result is not ``SirixDB is faster than PostgreSQL'', it is that the \emph{cost shape} predicted in
\S\ref{sec:w1}'s analysis is real and measurable: flat-in-document-size for SirixDB, linear for the
full-copy stores, with both crossovers now located rather than extrapolated.
\end{honesty}

\section{Against the other history-retainers: Datomic and XTDB}
\label{sec:retainers}

PostgreSQL and DuckDB do not retain queryable history natively, they were full-copy baselines.
\textbf{Datomic} (accumulate-only datoms) and \textbf{XTDB} (bitemporal) do, which makes them SirixDB's
real peers. To compare the \emph{cost shape} of retaining history, a flat record with $K$ fields has one
field updated 1\,000 times, retaining the full version log, and the per-update latency is measured as $K$
(the record size) grows. The question is whether a single-field update costs $O(1)$ or $O(\text{record})$.

\begin{table}[h]
\centering\footnotesize\setlength{\tabcolsep}{5pt}\renewcommand{\arraystretch}{1.25}
\begin{tabular}{@{}l r r r r r@{}}
\toprule
\textbf{record} & \textbf{SirixDB} & \textbf{Datomic} & \textbf{XTDB v1} & \textbf{XTDB v2} & \textbf{PostgreSQL}\\
 & \emph{(durable)} & \emph{(mem)} & \emph{(mem)} & \emph{(PG-wire)} & \emph{(durable)}\\
\midrule
$\sim$1\,KB (30 f) & 7.4 & \textbf{0.22} & 0.37 & 8.6 & 0.4\\
$\sim$10\,KB (300 f) & 11.1 & \textbf{0.20} & 2.55 & 13.1 & ---\\
$\sim$100\,KB (3000 f) & 10.0 & \textbf{0.20} & 24.5 & --- & 2.1\\
\bottomrule
\end{tabular}
\caption{ms per single-field update, full history retained, as the record grows. \textbf{SirixDB} (delta
page) and \textbf{Datomic} (one datom) are \emph{flat}, $O(1)$ in record size; \textbf{XTDB v1} (a full
document version per update) and \textbf{PostgreSQL} (a full \code{jsonb} copy) \emph{grow linearly},
$O(\text{record})$; \textbf{XTDB v2}'s columnar storage is mostly flat with a high fixed per-transaction
cost. PostgreSQL's 1\,MB/10\,MB points (19/194\,ms) are in \S\ref{sec:crossover}.}
\label{tab:retainers}
\end{table}

The history-retainers split into two camps, and the split is the structural point. \textbf{SirixDB and
Datomic update in $O(1)$}: SirixDB copies one leaf plus its root-to-leaf path, Datomic asserts one datom,
and neither cost depends on how big the surrounding record is, so both stay flat from 1\,KB to 100\,KB.
\textbf{XTDB v1 and PostgreSQL update in $O(\text{record})$}: each stores a fresh full version of the
whole document (XTDB v1's content-addressed document, PostgreSQL's \code{jsonb} copy), so XTDB v1 climbs
$0.37\to24.5$\,ms and PostgreSQL $0.4\to2.1$\,ms (and on to 194\,ms at 10\,MB) as the record grows. XTDB
v2's columnar engine is the interesting middle: its per-update barely moves with size ($8.6\to13$\,ms),
but a high fixed transaction cost dominates. SirixDB earns its place in the first camp.

\begin{honesty}
\small\textbf{Caveat.} This is a \emph{cost-shape} comparison, not an absolute-latency leaderboard, because
the systems' durability and drive models differ and are \emph{not} normalised: SirixDB commits durably
(fsync); Datomic ran on \code{datomic:mem} (in-memory, synchronous, \emph{not} durable); XTDB v1 ran as an
in-memory node with asynchronous submission; XTDB v2 ran as the latest stable Docker node driven over its
PostgreSQL wire protocol; PostgreSQL ran \code{synchronous\_commit=on}. So Datomic's 0.2\,ms and SirixDB's
7--11\,ms are not comparable as latencies, SirixDB is paying for a durable fsync per commit that the
in-memory engines skip. What \emph{is} comparable, and durability-independent, is the \textbf{slope}: flat
for SirixDB and Datomic, rising for XTDB v1 and PostgreSQL. All five retain full, queryable history; what
this measures is the price each pays to do so per update as records grow.
\end{honesty}

\section{Reading the past is not slower: depth-independence, measured}
\label{sec:histreadparity}

The retainer comparison measured the \emph{write} side of history; this section measures the \emph{read}
side of the same claim. A central promise of the design (\S\ref{sec:slidealg}, and the cost-shape
``Read'' argument above) is that a read at revision $r$ is a direct $O(h)$ snapshot traversal touching at
most $R$ fragments per page (Property~\ref{prop:bounded}), with \emph{no log replay and no
validity-interval join} --- so read latency must not grow with how far into the past one reads. A
log-structured or event-sourcing store cannot make that promise: reconstructing an old state there means
replaying forward from the nearest checkpoint, so cost rises with distance into the past. The claim is
directly testable. A 1\,MB document is committed 100 times under \code{SLIDING\_SNAPSHOT} ($R=3$),
yielding a 102-revision history, and then each of revisions $\{1,25,50,75,100\}$ is read in a
\textbf{fresh JVM} (SirixDB's in-process page caches empty, the OS file cache warm) in two ways: a
\emph{cold} full traversal that reconstructs every leaf page of that revision from its fragments exactly
once, and a \emph{warm} stream of random point reads once those pages are resident.

\begin{table}[h]
\centering\footnotesize\setlength{\tabcolsep}{8pt}\renewcommand{\arraystretch}{1.25}
\begin{tabular}{@{}r r r@{}}
\toprule
\textbf{revision read (of 102)} & \textbf{cold reconstruct (ms)} & \textbf{warm reads/s}\\
\midrule
1 (oldest)   & \textbf{291} & 7.83\,M\\
25           & 416 & 7.77\,M\\
50           & 421 & 7.76\,M\\
75           & 413 & \textbf{8.09\,M}\\
100 (newest) & 402 & 7.56\,M\\
\bottomrule
\end{tabular}
\caption{Reading any revision of a 102-revision history. \emph{Warm} point-read throughput is flat at
$\sim$7.8\,M reads/s regardless of which revision is open; \emph{cold} reconstruction is bounded and does
not grow with depth ($\sim$410\,ms from revision~25 on), with the oldest revision the cheapest of all
because its pages are still pristine single fragments. Reading the deep past costs no more than reading
the present.}
\label{tab:histreadparity}
\end{table}

The numbers show depth-independence on both axes. Warm point-read throughput is \textbf{flat at
$\sim$7.8\,M reads/s} across the whole history --- reading the oldest revision is no slower than reading
the newest, because once a page is resident every read is the same $O(h)$ trie walk irrespective of
revision (and the single-thread figure matches the per-thread rate of the concurrency scan,
\S\ref{sec:concurrency}, a reassuring cross-check). The cold reconstruction cost is the more revealing
measurement: it is \textbf{bounded and does not grow with depth} --- $\sim$410\,ms whether the target is
revision 25, 50, 75, or 100 --- exactly the behaviour Property~\ref{prop:bounded} predicts, since
reconstruction touches at most $R$ fragments per page no matter how deep the history. The one outlier is
instructive rather than anomalous: revision~1 reconstructs \emph{fastest} of all (291\,ms), because the
oldest snapshot's pages are pristine single fragments, with no sliding-snapshot deltas yet accumulated
against them. The deepest past is therefore not merely un-penalised but the \emph{cheapest} to
materialise, and from the revision at which the $R$-window saturates onward the cost is constant. That
flat line is the read-side counterpart to the flat write-side slope of \S\ref{sec:retainers}: the entire
purpose of materialising every revision as a navigable snapshot, rather than as a delta to be replayed
forward, is that time travel costs the same wherever it lands.

\section{Asymptotic cost summary}
\label{app:costs}

This section gathers the cost of the engine's core operations in one place; each row is established in
the chapter cited and is repeated here only for reference. Throughout, $N$ is the number of nodes in a
resource, $f=1024$ the page-trie fan-out, $h=\lceil\log_f N\rceil$ the trie height, $R$ the
sliding-snapshot window (\code{revisionsToRestore}), $D$ the number of dirty pages in a commit's intent log,
$d$ a node's tree depth, $\delta$ the size of the changed region between two revisions, $p$ a path's
length in segments, $k$ a result's cardinality, $n$ the number of rows a scan touches, and $W$ the SIMD
lane count of the host vector species.

{\small
\begin{xltabular}{\linewidth}{@{}>{\RaggedRight\arraybackslash}p{0.235\textwidth} >{\RaggedRight\arraybackslash}p{0.18\textwidth} >{\RaggedRight\arraybackslash}X@{}}
\toprule
\textbf{Operation} & \textbf{Cost} & \textbf{Basis}\\
\midrule
\endhead
Point lookup by node key & $O(h)=O(\log_f N)$ & descend Uber\,$\to$\,RevisionRoot\,$\to$\,Indirect$^\ast$\,$\to$\,leaf; $h\le 5$ for any $N<10^{15}$ at $f{=}1024$ (Ch.~\ref{ch:pds})\\
Record access within a leaf & $O(1)$ / $O(L)$ & direct slot offset on the slotted page; $O(L)$ in key-byte length $L$ on the HOT walk (\S\ref{sec:slotted}, Ch.~\ref{ch:hot})\\
Insert or update one node & $O(h+d)$ & copy-on-write of the $h$ pages on the root path, plus an $O(d)$ hash re-spine up the ancestors; DeweyID careting is $O(1)$ amortised (Ch.~\ref{ch:pds})\\
Commit & $O(D)$ & write the $D$ dirty intent-log pages and maintain affected indexes; $D\propto\delta$, not $N$ (Ch.~\ref{ch:ondisk})\\
Time-travel open by revision \# & $O(1)$ & index the revision array and read its RevisionRootPage (Ch.~\ref{ch:rest})\\
Time-travel open by instant & $O(\log R_{\text{tot}})$ & binary search over the $R_{\text{tot}}$ commit timestamps (Property~\ref{prop:histcost})\\
Reconstruct a leaf at revision $r$ & $O(R)$ & merge at most $R$ page fragments back to the last full dump (sliding snapshot, Ch.~\ref{ch:versioning})\\
Diff, adjacent revisions & $O(\delta)$ & read the stored per-commit delta (default-on \code{storeDiffs}) (\S\ref{sec:diff})\\
Diff, arbitrary revisions & $O(\delta)$ / $O(N)$ & hash-pruned lock-step descent; degrades to $O(N)$ only under \code{HashType.NONE} (\S\ref{sec:diff})\\
Secondary index point/range scan & $O(L+k)$ / $O(\log m+k)$ & HOT backend in key-byte length $L$; RBTREE backend in tree size $m$; $k$ matches returned (Ch.~\ref{ch:hot}, Ch.~\ref{ch:pathsummary})\\
Path-summary lookup & $O(p)$ & descend the summary trie one node per path segment (\S\ref{sec:pathsummary})\\
DeweyID order comparison & $O(\min\text{ divisions})$ & lexicographic over tier-encoded divisions; siblings differ in the last (\S\ref{sec:deweyids})\\
Vectorised aggregate / count & $O(n/W + n\bmod W)$ & $\lfloor n/W\rfloor$ SIMD steps plus a scalar tail; the projection-index fast path is an $O(n)$ byte-scan ($\sim$1\,ms for $10^8$ rows) (Ch.~\ref{ch:vector})\\
FMES import match & $O(ND)$ \dots\ $O(n^2)$ & Myers $O(ND)$ LCS per same-kind bucket then a greedy sweep; near-linear when most nodes match, quadratic in the unmatched leaves (\S\ref{sec:diff})\\
Buffer-pool hit / eviction & $O(1)$ & concurrent-map probe; MVCC-aware clock-sweep eviction is amortised $O(1)$ (\S\ref{sec:buffering})\\
\bottomrule
\end{xltabular}}

\noindent One theme runs down the cost column: \textbf{work is proportional to the region touched, not
to the document}. A point lookup, an insert, a commit, and a diff are all sublinear or
delta-proportional; the only genuinely $O(N)$ entries are a full scan with no applicable index and a
diff run with hashing switched off, both of which the system avoids by default. The $f{=}1024$ fan-out
is what makes the logarithmic terms behave like constants in practice: $\log_{1024}$ of a trillion
nodes is four, so ``$O(\log_f N)$'' is, concretely, ``at most five page touches'' for any database that
will ever be built. The reconstruction bound $O(R)$ is the price the sliding snapshot pays for its
write savings, and it is exactly why $R$ is a per-resource knob (Ch.~\ref{ch:versioning}): a
read-mostly resource lowers $R$ to make every historical read cheaper, a write-mostly one raises it to
write less per commit.

\section{Threats to validity}

Single machine, co-located client{+}server, fully cached, one document shape, single writer;
``cold'' drops only in-process caches (OS page cache stays warm); SirixDB is a dev build with
JIT warm-up tax, PostgreSQL a GA release. The small-document regime (\S\ref{sec:w1}) deliberately
favours PostgreSQL. The \textbf{large-document crossover is now measured} (\S\ref{sec:crossover}) rather
than extrapolated, but it too is a single document, one writer, and 101 versions; the
\emph{many-document}, \emph{many-version}, and \emph{concurrent-writer} regimes remain unmeasured. The
discipline the project held to, \emph{measure the crossover, do not extrapolate it in public}, is what
this revision discharges.

\chapter{Related Work}
\label{ch:related}

With the system described and measured, this chapter situates its design in the literature it draws on.

\textbf{Persistent data structures.} SirixDB's structural sharing is the database-scale
analogue of purely functional persistent structures (Okasaki~\cite{okasaki}, \emph{Purely Functional
Data Structures}): a ``modification'' produces a new version sharing the unchanged majority. The
page trie is a wide-fan-out persistent trie; the sliding snapshot is a page-granular
versioning policy layered on it.

\textbf{Copy-on-write filesystems.} WAFL~\cite{wafl}, ZFS, and Btrfs popularised CoW B-trees with atomic
root-pointer swaps and snapshotting; SirixDB applies the same root-swap-on-commit discipline at
the document/page level, with dual beacons playing the role ZFS's überblock array plays:
redundant, checksummed roots tolerant of torn writes.

\textbf{Versioned and temporal stores.} Git and Fossil version content with content-addressed
structural sharing but are not queryable databases; bitemporal relational features (SQL:2011
system-versioned tables~\cite{sql2011}, Datomic's accumulate-only model) offer time travel at the row/datom
level. SirixDB differs in offering \emph{sub-document}, node-level history and diffs over
hierarchical JSON/XML, with stable node identity across revisions, capabilities the W6/W4
results show have no native relational equivalent. More recent versioned-data systems (Dolt's
git-style versioned tables, TerminusDB's immutable layered graph, XTDB's bitemporal store) share the
accumulate-only philosophy, but SirixDB is distinguished by node-granular history and semantic diffing
over hierarchical documents.

\textbf{History-retention architectures, read off the benchmark.} The measured cost shapes of
\S\ref{sec:retainers} are a fingerprint of \emph{where each retainer writes the change}, and reading
them back illuminates the design space. \textbf{Datomic} stores accumulate-only \emph{datoms}
(entity-attribute-value-transaction); a single-field update is literally one datom, so its cost is
$O(1)$ in the surrounding record, the flat $\sim$0.2\,ms measured, and its history is the immutable log
itself. The price is the model: a flat entity-attribute-value store does not natively carry nested
documents or offer a sub-document semantic diff. \textbf{XTDB v1} is a content-addressed \emph{document}
store, bitemporal in valid and transaction time; a \code{put} writes a fresh \emph{whole-document}
version, so its per-update cost grows with the document, the $0.37\to24.5$\,ms slope measured, the same
$O(\text{record})$ shape as PostgreSQL's full-\code{jsonb}-copy pattern, traded for first-class bitemporal
queries. \textbf{XTDB v2} rebuilds that on a columnar Apache Arrow engine with SQL; its per-update barely
moves with record size (the columnar write is small) but carries a high fixed transaction cost, the
flat-but-elevated $8.6\to13$\,ms measured. \textbf{DuckDB}, by contrast, is not a history store at all,
it is the OLAP engine that won the aggregate benchmark (\S\ref{sec:vecbench}) and keeps no queryable past.
\textbf{SirixDB} sits in Datomic's $O(1)$-update camp by a different route, a copy-on-write page trie
whose single-field update copies one leaf and its root path, but unlike Datomic it retains the
\emph{hierarchical document} with stable node identity, and unlike XTDB it writes a \emph{delta} rather
than a whole version. The distinguishing combination, node-granular $O(1)$ delta updates, bitemporal
revisions \emph{and} valid time, stable node keys across history, and a sub-document semantic diff, is
what none of the others offers in one system; the benchmark's two camps (delta/datom $O(1)$ versus
full-version $O(\text{record})$) are the structural reason why.

\begin{table}[h]
\centering\footnotesize\setlength{\tabcolsep}{4pt}\renewcommand{\arraystretch}{1.25}
\begin{tabularx}{\linewidth}{@{}l l l l l X@{}}
\toprule
\textbf{System} & \textbf{History} & \textbf{Node id} & \textbf{Time travel} & \textbf{Diff} & \textbf{Storage model}\\
\midrule
SirixDB & node-level & stable & revision / instant & sub-document & versioned CoW page trie\\
PostgreSQL (SQL:2011) & row-level & n/a & system / valid time & none & heap {+} WAL\\
Datomic & datom-level & entity id & as-of & none & accumulate-only log\\
git / Fossil & file-level & n/a & commit & textual & content-addressed CoW\\
Dolt & row / table & primary key & as-of & row-level & Merkle-DAG CoW\\
XTC (Kaiserslautern) & none & SPLID & none & sub-document & SPLID-labelled node store\\
BaseX (native XML) & none & pre/dist & none & external & pre/dist preorder table\\
\bottomrule
\end{tabularx}
\caption{SirixDB against related versioned and temporal systems. Its distinguishing combination is
node-level history with stable identity, time travel by revision or wall-clock instant, and
sub-document semantic diffs over hierarchical JSON and XML.}
\label{tab:related}
\end{table}

\textbf{Storage engines and durability.} Against the log-structured-merge mainstream
(LevelDB/RocksDB) and the WAL-and-heap mainstream (PostgreSQL), SirixDB has \emph{no separate
WAL}; its append-only data files are the entire on-disk story. The commit-latency analysis
(Ch.~\ref{ch:ondisk},~\ref{ch:eval}) is in dialogue with the TU-Munich/LeanStore/Umbra line
(Haubenschild et al.~\cite{haubenschild} on rethinking logging/recovery; Haas \& Leis~\cite{haas} on what modern NVMe can do):
SirixDB adopts their core NVMe insight: group commit is a spinning-disk tax, and the lever on
flash is removing per-commit barrier overhead, not amortising fsync. That is precisely what
the preallocated profile does. It does \emph{not} adopt their parallelism mechanism
(distributed per-worker logging with remote-flush-avoidance), because a SirixDB resource has a
single writer; its analogue of their parallelism is \emph{across} resources.

\textbf{Log-structured merge trees.} The LSM mainstream (LevelDB, RocksDB, Cassandra) shares SirixDB's
refusal to overwrite in place: writes land in an in-memory memtable, flush to immutable on-disk runs, and
background \emph{compaction} merges runs to reclaim space and bound read fan-out. The cost shapes diverge,
though. An LSM trades \emph{read} and \emph{write amplification} (a key may sit in several levels, so a
point read probes bloom filters across them, and compaction rewrites the same data several times) for
sequential write throughput, and it \emph{discards} superseded versions when levels merge. SirixDB's
copy-on-write page trie is also append-only, but its compaction analogue is the sliding snapshot's bounded
carry-forward (\S\ref{sec:slidealg}), a point read is a single $O(h)$ trie walk with no per-level probe,
and superseded versions are \emph{retained} as the whole point of the structure. The slogan contrast is
that an LSM appends in order to \emph{reclaim}, whereas SirixDB appends in order to \emph{remember}.

\textbf{Event sourcing and CQRS.} A SirixDB revision history resembles an event log, each commit an
immutable, ordered fact. The architectural difference is where current and past state live. Event-sourcing
systems reconstruct state by \emph{replaying} events forward from the nearest snapshot, so reading an
arbitrary past state costs work proportional to the events since that checkpoint, and the
command-query-responsibility-segregation (CQRS) pattern exists precisely to maintain a separate
materialised read model beside the log. SirixDB instead materialises \emph{every} revision as a navigable
snapshot in the shared page DAG, so time travel is an $O(h)$ traversal rather than an $O(\text{replay})$
fold, the depth-independence the historical-read benchmark measures directly (\S\ref{sec:histreadparity}).
In those terms SirixDB collapses the CQRS split: the log and the read model are the same versioned store,
and ``replay'' is just opening a different root.

\textbf{Lakehouse table formats.} Apache Iceberg, Delta Lake, and Apache Hudi add snapshot isolation and
time travel to columnar data lakes through a metadata log of \emph{snapshots} over immutable
Parquet/ORC files. This is structurally the same idea as SirixDB's S3 backend, immutable data objects
plus an authoritative manifest naming the live set (\S\ref{sec:storagespi}), but pitched at \emph{file
and table} granularity: a row update rewrites a whole data file (copy-on-write tables) or appends a
delete-vector and a new file (merge-on-read), and the time-travel unit is a table snapshot. SirixDB
pushes the identical immutable-objects-plus-manifest discipline down to \emph{page} granularity over a
tree, so its delta is one leaf and its time-travel unit is a single node at a revision, finer by orders of
magnitude. The trade is symmetric: the lake formats win bulk analytic scans over flat columnar data, the
axis on which DuckDB beat SirixDB (\S\ref{sec:vecbench}), while SirixDB wins sub-document update cost,
node identity, and fine-grained history.

\textbf{Indexing.} HOT (Binna et al.~\cite{binna}) and ART (Leis et al.~\cite{leisart}) are the lineage of the in-memory
trie; SirixDB's HOT adapts the single-key-leaf reference design to multi-key leaves, which is
the source of the invariant subtleties of Ch.~\ref{ch:hot}. The path summary is in the
tradition of XML path indexes (DataGuides, A(k)-indexes) repurposed for JSON and extended with
per-path statistics for cost-based optimisation.

\textbf{Columnar execution.} The PAX layout (Ailamaki et al.~\cite{ailamaki}) and the C-Store~\cite{cstore}/MonetDB columnar
tradition inform the projection-index pages and SIMD scan; frame-of-reference {+} bit-packing
and dictionary encoding are standard lightweight column codecs. SirixDB's twist is layering an
\emph{optional, declared} columnar projection over a row-oriented, versioned primary store.

\textbf{Query processing.} The embedded brackit engine is set-oriented in the sense of B\"achle's
thesis~\cite{baechle}, and its parallel and vectorised operators follow the now-standard lines:
\emph{morsel-driven parallelism} (Leis et al.~\cite{morsel}) for NUMA-aware, load-balanced fan-out, and a
columnar group-by kernel in the vein of the column-at-a-time tradition above. The cost-based join
ordering is the connected-subgraph/complement-pair dynamic program \emph{DPhyp} (Moerkotte \&
Neumann~\cite{dphyp}), with a greedy fallback past twenty relations. The structural diff and the
identity-preserving import rest on classic tree-differencing: the Fast-Match/Edit-Script algorithm of
Chawathe et al.~\cite{chawathe} over a longest-common-subsequence core (Myers~\cite{myers}). SirixDB's
contribution is not these algorithms but their composition over a \emph{versioned, immutable} store, where
a diff can short-circuit to a hash comparison and an import lands as one new revision.

\textbf{Approximate nearest-neighbour search.} The vector index is a faithful Hierarchical Navigable
Small World graph (Malkov \& Yashunin~\cite{hnsw}), the dominant in-memory ANN structure, novel here only
in being persisted through the copy-on-write page trie so that the proximity graph is itself versioned
and time-travellable, a capability a standalone ANN library does not offer. Its honest status (no
query-language binding yet, explicit population, float32, the most experimental member of the index
family) is recorded in Chapter~\ref{ch:pathsummary}.

\textbf{Buffer management and off-heap memory.} The page guards, frame-version validation, and
size-classed off-heap slot allocator follow the LeanStore line (Leis et al.~\cite{leanstore}) of
pointer-swizzling, optimistic-latch buffer managers, adapted to a single-writer, MVCC-snapshot store
where the guard's role is to keep a frame alive under a lock-free reader rather than to coordinate
concurrent writers.

\textbf{Node labeling schemes.} Assigning every node a stable, order-encoding identifier is a research
area in its own right. Range-based labels (start/end/level intervals) answer ancestor queries in
constant time but must relabel on insertion; prefix-based path labels avoid relabelling at the cost of
longer identifiers. ORDPATH (O'Neil et al.~\cite{ordpath}), shipped in Microsoft SQL Server, is the
prefix scheme SirixDB follows: careted-in odd/even components give insert-without-relabel, and a
compressed binary form whose byte order \emph{is} document order. H\"arder, Haustein, Mathis and
Wagner~\cite{haerderlabels} reconsidered the family for dynamic XML and refined it into stable path
labeling identifiers (SPLIDs), the immediate ancestor of SirixDB's \code{SirixDeweyID}
(\S\ref{sec:deweyids}).

The design space splits cleanly on one axis (Table~\ref{tab:labeling}). \emph{Range-based} labels
encode each node as a \code{(start, end, level)} interval, so an ancestor test is a constant-time
containment check, but an insertion that runs out of integer space between two siblings forces a
renumbering of part of the document. \emph{Prefix-based} (Dewey) labels make each child extend its
parent's label, so the ancestor test becomes a prefix check and document order falls out of a
lexicographic comparison, at the cost of longer, variable-length identifiers. The naive Dewey scheme
still relabels siblings on an insert; the ORDPATH contribution is the \emph{caret}, an even component
that adds depth without adding a level, which removes relabelling entirely; H\"arder's SPLID work
generalised the spacing and the binary, order-preserving encoding for dynamic XML in XTC. SirixDB sits
at the end of this line, and the reason it chose a prefix scheme over the constant-size interval scheme
is exactly its headline property: \textbf{stable node identity across revisions}. A versioned store
inserts on nearly every commit, and a label scheme that renumbered existing nodes on insert would
destroy the cross-revision identity that history, diffs, and temporal navigation all depend on.

\begin{table}[h]
\centering\footnotesize\setlength{\tabcolsep}{4pt}\renewcommand{\arraystretch}{1.25}
\begin{tabularx}{\linewidth}{@{}l l c l l X@{}}
\toprule
\textbf{Scheme} & \textbf{Label form} & \textbf{Insert w/o relabel} & \textbf{Ancestor} & \textbf{Order} & \textbf{Exemplar}\\
\midrule
Range / interval & (start, end, level) & no & containment $O(1)$ & compare start & many early XML DBs\\
pre/dist & (pre, dist, size) & no & containment $O(1)$ & pre value & BaseX\\
Dewey / path & \code{1.2.3} & partial & prefix $O(d)$ & lexicographic & XPath Accelerator\\
ORDPATH & careted prefix & yes & prefix $O(d)$ & byte order & MS SQL Server\\
SPLID & spaced prefix & yes & prefix $O(d)$ & byte order & XTC (Kaiserslautern)\\
\code{SirixDeweyID} & tier-coded prefix & yes & prefix $O(d)$ & unsigned bytes & SirixDB\\
\bottomrule
\end{tabularx}
\caption{Node-labeling schemes. The prefix family (Dewey, ORDPATH, SPLID, \code{SirixDeweyID}) trades a
constant-size interval label for prefix-encoded structure and, from ORDPATH onward, insertion without
relabelling. SirixDB's \code{SirixDeweyID} is a tier-coded, prefix-free binary refinement of SPLID whose
unsigned byte order is document order (Theorem~\ref{thm:dewey}).}
\label{tab:labeling}
\end{table}

\textbf{Native XML databases and the query heritage.} SirixDB descends, by way of Treetank at the
University of Konstanz (Kramis~\cite{kramis}; Graf et al.~\cite{treetank,graf}), from the
native-XML-DBMS generation: systems such as Natix, BaseX, eXist, and Theo H\"arder's XTC at TU
Kaiserslautern, which worked out node-level storage, node labeling, and fine-grained locking for
tree-structured data. Their labeling approaches diverged, and it is the divergence SirixDB inherits a
side of: XTC used SPLIDs, the prefix/DeweyID family SirixDB follows (\S\ref{sec:deweyids}), whereas
\textbf{BaseX} stores nodes in a \emph{preorder table} addressed by a \code{pre} value, with each record
holding a \code{dist} (the distance to its parent) and a subtree \code{size}, a range/structural
encoding rather than a prefix label, compact and fast to navigate but order-sensitive on insertion. The
query engine SirixDB embeds, \emph{brackit}, is itself from that line: a
set-oriented XQuery/JSONiq processor designed for physical data independence and parallelism in
B\"achle's doctoral work~\cite{baechle}. SirixDB's contribution is to carry these XML-era results, namely
versioned node storage, prefix labeling, and a compositional XQuery engine, into a JSON/JSONiq,
off-heap, multi-version setting.

%================================================================
\chapter{Conclusion and Future Work}
\label{ch:conclusion}

SirixDB derives a broad surface of useful properties (time travel, structural sharing,
lock-free snapshot reads, sub-document history, hash-guided diffs, append-only durability with
no WAL) from one disciplined commitment to immutability and copy-on-write, executed through a
fan-out-1024 page trie, a sliding-snapshot page-versioning algorithm with bounded
reconstruction and even amortisation, and a self-describing, dual-beacon, checksum-protected
on-disk format whose recovery is validated under SIGKILL and power-loss fault injection. Above
the store sit a HOT-based navigation/index layer, a path summary with incremental statistics, a
brackit JSONiq/XQuery engine with an index-aware cost-based optimiser, and an emerging
columnar/SIMD execution path; alongside it, pluggable FileChannel/mmap/io\_uring/S3/Kafka
back-ends and a SolidJS GUI built explicitly for the temporal model. The evaluation is honest
about the regime where a relational engine wins, and precise about the cost-shape and
capability axes where SirixDB does. In the terms of \S\ref{sec:principles}, the preceding chapters have
worked the five design principles out in full; this one revisits the choices behind them and the price
each pays.

\section{Design rationale: the choices and their costs}

Each of SirixDB's defining decisions buys a capability at a stated price, and it is worth setting them
out together.

\textbf{Immutability and copy-on-write, not update-in-place.} Never overwriting a page is what makes
every past revision a first-class, directly readable object and makes snapshot reads lock-free. The
price is a per-revision metadata floor and write amplification along the root-to-leaf path of each
change; the evaluation shows this losing to an update-in-place engine in the small-document regime and
winning on the temporal-capability and history-read axes. The bet is that for data whose \emph{history}
matters, paying $O(\text{change})$ per revision is the right cost shape.

\textbf{The sliding snapshot, not full or plain incremental versioning.} Of the four page-versioning
policies, only the sliding snapshot bounds both the read cost (at most $R$ fragments) and the write
spikes (no periodic full rewrite), at the cost of a small, evenly amortised carry-forward. It is the one
strategy that keeps both tails flat, which is why it is the default.

\textbf{Off-heap, flyweight nodes, not a heap object graph.} Holding pages as native
\code{MemorySegment}s keeps the engine's largest memory consumer off the garbage collector, gives
zero-copy reads, and makes the in-memory and on-disk images identical, so that a commit is mostly a byte
copy. The price is manual frame lifecycle and a dependence on the Foreign Function and Memory API.

\textbf{One writer per resource, many readers.} A single-writer rule makes the concurrency model simple
to reason about and the read path lock-free, and it sidesteps write--write conflict resolution entirely.
The price is that a single resource does not scale writes across cores; SirixDB's answer is parallelism
\emph{across} resources, not within one.

\textbf{Prefix labels (DeweyIDs), not range labels.} Encoding a node's position as an ORDPATH-style path
gives insert-without-relabel and document order by byte comparison, at the cost of longer identifiers
than a compact integer would need. For a store whose whole point is stable identity across edits, never
having to renumber is the property that matters.

\textbf{Embedding brackit, not building a query engine.} Reusing a proven, set-oriented XQuery/JSONiq
engine with a clean store SPI let SirixDB add only what is genuinely its own, the temporal functions and
an index-aware cost-based optimiser, instead of re-implementing a language runtime. The price is living
within another system's extension points, which the optimiser stages and the store binding are shaped to
fit.

The common thread is that none of these choices is free; the report has tried to name the cost of each
rather than present them as unalloyed wins.

\section{Discussion: questions the design invites}

A few questions come up often enough to answer directly.

\textbf{Why the JVM, and not C or Rust?} The performance-sensitive parts that would normally argue for a
native language, off-heap memory and SIMD, are reachable from modern Java without JNI: the Foreign
Function and Memory API gives direct \code{MemorySegment} access (and the io\_uring backend calls the
kernel through it), and the Vector API gives real SIMD. In return SirixDB keeps the JVM's tooling, the
GraalVM native-image option, and, decisively, brackit, a mature JVM query engine it would otherwise have
to replace.

\textbf{Why copy-on-write, and not an LSM tree?} Log-structured-merge storage is built to forget: its
compaction merges and discards superseded versions, which is the opposite of what a history store wants.
Copy-on-write retains every version as a by-product of never overwriting, so the discipline that looks
like overhead for a key-value store is exactly the right primitive here.

\textbf{Why a single writer per resource?} Document edits are usually authored serially, and a
single-writer rule removes write--write conflict resolution, MVCC garbage, and reader locks in one
stroke, while reads scale without limit because they see immutable snapshots. The scalability a single
resource gives up is recovered across resources, which is where real concurrency in a document store
tends to live.

\textbf{Why not just version JSON in git?} Git versions content with structural sharing, but it is a
file store, not a database: it has no node identity, no secondary indexes, no query language, and no way
to ask ``what did this field look like at time $t$'' without checking out and re-parsing. SirixDB keeps
git's sharing discipline while adding the machinery that makes history \emph{queryable}.

\textbf{Is retaining all history not ruinously expensive?} Storage grows with the total volume of
\emph{change}, not with the document size multiplied by the number of revisions (as the cost-shape
analysis of Chapter~\ref{ch:eval} shows); a field touched once in a thousand revisions costs one extra
record, not a thousand copies. Where even that is too much, the versioning window and compaction
policies bound it, at the cost of the oldest history.

\textbf{Why store everything off-heap instead of on the Java heap?} Two reasons compound. A page's
in-memory form is byte-identical to its on-disk image, so reading is a slice and writing is a copy with
no serialise/deserialise step, which is only possible off-heap; and a multi-gigabyte page cache on the
Java heap would dominate garbage-collection pause times. The FFM \code{MemorySegment} gives the raw
access without JNI, and the custom frame allocator (\S\ref{sec:allocator}) adds stable addresses and
syscall-free recycling on top. The cost is manual lifetime management, which is exactly what the page
guards exist to make safe.

\textbf{Why fuse object keys into their value node?} The overwhelmingly common JSON shape is
``\code{key: scalar}'', and storing it as two records (a key node and a value node) doubles the node
count, the index entries, and the hash work for the most frequent case. Fusing them into one
\code{OBJECT\_NAMED\_*} record halves all three. The price is a wire-shape subtlety, a fused container's
children hang directly under its value, which one normaliser in the serialiser and one in the GUI absorb;
the report judged that a worthwhile trade and the measurements bear it out.

\textbf{Why a set-oriented engine under a pull interface, rather than one or the other?} A pure
tuple-at-a-time iterator is simple but caps throughput; a pure vectorised engine is fast but rigid. By
keeping brackit's logical algebra a pull-based \code{Cursor} while letting its operators consume tuple
\emph{batches}, the same query plan can run serially, under morsel-driven parallelism, or through a
vectorised columnar kernel without being recompiled, B\"achle's physical data independence. The engine
pays for that flexibility with a more intricate operator contract, but it is what lets the SIMD path of
Chapter~\ref{ch:vector} slot in as an SPI rather than a fork.

\textbf{Is the valid-time axis indexed, like the transaction axis?} Yes, though it costs more to
maintain. The transaction axis falls out of the immutable revision structure for free; the valid axis is
accelerated by a persistent \code{VALIDTIME} Relational-Interval-Tree that stabs a query instant in
$O(h)$, with a CAS-index narrowing path and the $O(n)$ field filter as the fallback when neither index is
present (\S\ref{sec:bitemporal-worked}). Because valid bounds are ordinary JSON fields, the index is an
explicit per-resource opt-in rather than a free by-product of the structure --- the model stays simple,
and the acceleration is there once a workload declares it needs it.

\textbf{Could the diff prune siblings, not just subtrees?} The measured diff (\S\ref{sec:diffbench}) prunes
\emph{vertically} --- an unchanged subtree costs one hash comparison --- but still walks a node's children
one by one, so diffing a revision of a million-element array enumerates the whole array even if one element
changed. A \emph{horizontal} prune is conceivable: a Merkle structure over runs of child hashes (or the
existing change-tracking index, which already records the changed subtree-roots per revision) could skip
unchanged sibling ranges. SirixDB does not do this in the general walk; it instead sidesteps the cost for
the overwhelmingly common adjacent-revision case by reading the precomputed per-commit delta
(\code{storeDiffs}), and reserves the breadth-scanning recompute for arbitrary revision pairs. The honest
answer is ``yes in principle, and partly available through the change index, but not in the lock-step walk
today,'' and the breadth measurement is precisely what would justify building it.

\textbf{Why three storage backends behind one SPI, rather than one?} The page layer reads and writes only
through \code{IOStorage}, so the engine is oblivious to which backend sits underneath, and the three that
ship occupy different points of a portability/latency/scale triangle. \code{FILE\_CHANNEL} is the portable,
durable default that runs anywhere a JVM does; \code{IO\_URING} trades portability for NVMe throughput,
reaching the kernel's asynchronous interface through the Foreign-Function API with no JNI
(\S\ref{sec:storagespi}); and \code{S3} trades local-disk latency for elastic, shared object storage that
doubles as the replication source of truth. One SPI means a resource can move between them without touching
a line of the page, versioning, or query code --- the same property that let the io\_uring and S3 backends
be added as \code{ServiceLoader} providers rather than forks.

\textbf{So is the approach sound --- and when is it the wrong tool?} Sound, on the evidence this report
gathers, but deliberately specialised, and the two are best stated together. The foundation is a
persistent copy-on-write trie with structural sharing (Chapter~\ref{ch:pds}), a lineage it shares with
purely functional data structures, copy-on-write filesystems, and Datomic (Chapter~\ref{ch:related}); on
that base the sliding snapshot bounds reconstruction to at most $R$ fragments
(Property~\ref{prop:bounded}), time travel is an $O(1)$ index into the revisions file rather than a scan or
a log replay, and reads are lock-free immutable snapshots. Those are the central claims, and they survive
measurement: the cost \emph{shape} is flat in document size where an update-in-place store is linear, with
the crossover now measured rather than argued (\S\ref{sec:crossover}), and reads scale without contention
(\S\ref{par:lockfreeread}). That is the honest case --- not that the engine is fast everywhere, but that
its structure delivers the one capability it is built for, bitemporal node-granular history at bounded read
and write cost, which neither a document store nor SQL system-versioned tables provide.

The caveats are about \emph{fit} and \emph{maturity}, not about the design being wrong, and they mark where
it should not be reached for. First, it is a \emph{specialist}: the append-only ``remember everything''
premise is a feature for audit, temporal, and versioned-document workloads and a tax for those that never
look back --- the same crossover study that shows SirixDB ahead on versioned documents shows update-in-place
PostgreSQL and columnar DuckDB ahead on their own ground of small-document churn and ad-hoc analytics
(\S\ref{sec:w1}, Chapter~\ref{ch:eval}). Second, write throughput on a single resource is \emph{capped} by
the one-writer rule (Property~\ref{prop:serializable}): concurrency is recovered \emph{across} resources,
not within one, and the measured write-across-resources scaling is sub-linear, not free
(\S\ref{sec:concurrency}), so a single write-hot document is a ceiling. Third, the implementation is
\emph{younger than the design} --- correctness rests on the test suites named throughout rather than on
years of production hardening, and the report draws its still-open boundaries (the surviving HOT routing
invariant of Chapter~\ref{ch:hot}, the deliberately narrow vectoriser of Chapter~\ref{ch:vector}, the
staged enterprise pieces gathered just below) in the open rather than blurring them. The guidance that
falls out is one line: reach for SirixDB when history, time travel, audit, or node-level diffing are
first-class requirements; reach for something else when the need is a general-purpose document store, heavy
single-dataset write throughput, or a mature ad-hoc analytics engine.

\section{Limitations}
\label{sec:limitations}

The caveats throughout this report flag local boundaries; gathered together, the substantive limitations
are these. \textbf{The small-document write regime} is a genuine loss to an update-in-place engine, as
Chapter~\ref{ch:eval} shows; the cost-shape crossover for larger documents is now \emph{measured}
(\S\ref{sec:crossover}: SirixDB overtakes PostgreSQL on commit speed near 1\,MB and on storage past
100\,KB), and concurrent \emph{read} scaling is now measured (\S\ref{par:lockfreeread}: lock-free,
$\sim$2$\times$/$\sim$3$\times$ at 2/4 threads), though the many-document and
concurrent-\emph{writer} regimes remain unmeasured. \textbf{Within a resource there is one writer}, so a single hot
document does not scale writes across cores. \textbf{The vectorised executor is deliberately narrow}:
it is fail-closed to a small set of numeric, columnar query shapes, with string-predicate SIMD and
non-\code{double} numeric columns still open. \textbf{The HOT index backend} carries one surviving,
reads-correct routing-invariant deviation (Chapter~\ref{ch:hot}) and remains opt-in behind the
red-black-tree default. \textbf{The JSON FMES import} is library-only, not wired to a function, and its
delete phase does not yet cover the fused structural kinds. \textbf{Histogram collection is deferred}, so
a first-seen predicate plans on a default selectivity until statistics are built. And several
\textbf{enterprise pieces}, the public S3 selector wiring, multi-writer fencing, and a broker-backed
Kafka replication test, are staged rather than finished. None of these undermines the core claims; each
is a boundary the report has tried to draw plainly rather than blur.

The open frontier, stated plainly:

\begin{itemize}[leftmargin=1.4em]
\item \textbf{HOT}: the single remaining marginal routing invariant and its encoding rework
  (AND-encoding / proactive disc-bit extension / recompute-on-CoW).
\item \textbf{Vectorisation}: SIMD string-predicate matching (currently scalar), broadening
  numeric support beyond \code{double}-exact columns, and the projection-index chunk-page refactor
  (task~\#57) to push the per-leaf scale ceiling.
\item \textbf{Query}: the broader JQGM/Cascades optimiser vision
  (\pp{docs/COST_BASED_QUERY_OPTIMIZER_PLAN.md}), and eager (non-deferred) histogram collection so a
  first-seen predicate is not planned on the default selectivity.
\item \textbf{History scans}: the deferred fragment-chain single-pass scan that would close
  the dense-field W4 case.
\item \textbf{Storage}: a big-endian
  (e.g.\ \code{qemu-s390x}) test run to certify the now-completed little-endian field pinning before a
  format freeze --- on little-endian hardware the pinning is a no-op, so only a big-endian host
  actually exercises it.
\item \textbf{Durability}: for a release candidate of the new default commit profile,
  validation on power-loss-protected enterprise NVMe and a torn-block (not merely
  process-death) fault injector.
\item \textbf{Evaluation}: the large-document crossover (100~KB / 1~MB / 10~MB) is now measured
  (\S\ref{sec:crossover}); still open are the concurrency, many-document, and many-version regimes.
\item \textbf{Enterprise back-ends}: wiring the S3 L2 SSD cache and segment compactor,
  a public \code{StorageType.S3} selector and multi-writer fencing; a real-broker
  (Testcontainers) Kafka replication test and engine-lifecycle auto-bootstrap of the replicator.
\item \textbf{Serving}: applying the unordered \code{executeBlocking} lever uniformly across
  the REST handlers, and broadening XML write parity with JSON.
\end{itemize}

None of these undercut the central claim. They are the ordinary frontier of a system that has
chosen a harder, more principled foundation than most, and has, by its own measurements,
paid for that choice honestly.

\section{Future directions, worked through}
\label{sec:futuredirections}

Several adjacent designs are worth working through, because each is a \emph{design} that the rest of the
system already constrains. The first --- valid-time acceleration --- has since been \textbf{built}
(\code{1.0.0-beta3}) and is described below as shipped; the rest are \textbf{future work}, presented as a
design space rather than built behaviour --- but the existing architecture narrows that space enough that
the shape of each solution, and its cost, can be reasoned about now, and each connects to a measured
result earlier in the report.

\subsection*{Accelerating the valid-time axis (shipped)}

The transaction-time axis is indexed for free: every revision is a materialised snapshot, so ``as of
revision~$r$'' is an $O(h)$ open (\S\ref{sec:histreadparity}). The valid-time axis is not free --- valid
bounds are ordinary JSON fields --- so ``which facts hold at instant~$t$'' was for a long time an $O(n)$
filter over the candidate records (\S\ref{sec:bitemporal-worked}). It is now accelerated by a persistent
index, and the engine's grain chose its shape from three candidates. A \textbf{CAS index on
\code{valid-from} alone} is the cheapest to build but answers the wrong question: ``valid at~$t$'' is the
\emph{stabbing} query $\textit{from}\le t<\textit{to}$, a two-sided interval-containment test a
single-attribute ordered index cannot satisfy without scanning every record with \code{valid-from}$\le t$
and checking \code{valid-to} per row --- so it survives only as a \emph{narrowing} fast-path. A
\textbf{segment tree} over the endpoint-sorted timeline is read-optimal but rebuild-heavy, poorly suited to
a store that inserts on nearly every commit. The shape that fits SirixDB's copy-on-write grain, and the one
that shipped (\code{1.0.0-beta3}), is a \textbf{Relational-Interval-Tree}: a persistent \code{VALIDTIME}
index that maps each record's $[\textit{from},\textit{to})$ through a monotone interval domain and stabs a
query instant in $O(h)$, returning candidate record-object keys directly. \code{jn:valid-at} and
\code{jn:open-bitemporal} consult it first, re-verify each candidate against its exact bounds (so the
result is provably the linear scan's), and fall back to the CAS-narrowing path or the field filter when no
interval index exists. Its price is maintenance --- every commit that changes a valid bound updates the
index, a second $O(h)$ atop the one the document already pays --- incurred only when bounds change. The
\emph{model} was complete before; the \emph{acceleration} is now built, gated behind an explicit
per-resource index so workloads that never stab pay nothing.

\subsection*{Horizontal diff pruning}

The measured diff (\S\ref{sec:diffbench}) prunes \emph{vertically} --- an unchanged subtree is one hash
comparison --- but walks siblings one by one, so a one-element change in a million-element array still
enumerates the array, the $O(\text{breadth})$ cost the benchmark made precise. A \textbf{horizontal} prune
would skip unchanged sibling \emph{runs}, and two ingredients the engine already has make it approachable.
First, the \code{CHANGED\_NODES} index records, per revision, the changed subtree-roots
(\S\ref{sec:changetracking}); intersecting two revisions' changed-sets yields the sibling positions that
\emph{can} differ before the walk starts, turning enumeration into a lookup. Second, a per-node
child-hash-run summary (a small Merkle structure over contiguous child hashes) would let the lock-step walk
binary-search to the first differing child rather than scan to it. The cost is that summary's maintenance on
every child insert or delete --- write-time bookkeeping traded for read-time diff speed, the same trade the
path summary and the change indexes already make. For the common \emph{adjacent}-revision diff the question
is moot, since the stored per-commit delta (\code{storeDiffs}) already sidesteps the walk; the horizontal
prune earns its keep only for frequent \emph{arbitrary}-pair diffs over wide arrays, a narrower need that
keeps it future work.

\subsection*{The dense-field history scan}

The history-read results (\S\ref{sec:histpath}) showed the $O(\text{changes})$ shape win decisively for a
field that changes \emph{rarely} --- one record reconstructs the whole history --- but lose to
PostgreSQL's compressed heap scan for a field that changes on \emph{every} commit, because the current
path reconstructs each of the $R$ revisions' pages independently. The deferred optimisation is a
\textbf{single-pass fragment-chain scan}: a field that changes every revision has, on its leaf page, a
fragment chain whose successive links \emph{are} its successive values, so one ordered walk of that chain
(the \code{PageFragmentKey} chain of \S\ref{sec:caching}, already on disk) yields all $R$ values without
reconstructing $R$ separate pages. The expected shape is still $O(R)$ \emph{reads}, but with one decode
pass and sequential I/O in place of $R$ random reconstructions --- the constant-factor and cache-locality
win the dense case needs. It is future work because it optimises an already-correct path, not a missing
one: the dense history is \emph{readable} today, just not at the throughput a bulk temporal analytic would
want.

\subsection*{Bounding the S3 segment count}

The S3 backend (\S\ref{sec:storagespi}) writes immutable $\sim$16\,MB segment objects under an
authoritative manifest; over a long-lived resource the segment count and the manifest's parallel arrays
grow without bound, so a cold read's manifest load and binary search creep upward and the object count
itself becomes an operational cost. A \textbf{compactor} --- shipped but unwired --- would merge runs of
small or sparsely-referenced segments into fewer larger ones and rewrite the manifest, the live-set
discipline making it safe: write the merged segment, publish a manifest referencing it, and only then mark
the superseded segments collectable, so a crash mid-compaction orphans objects but never breaks a readable
manifest (the ordering the commit protocol already uses). The open questions are policy, not mechanism ---
\emph{when} to compact (a segment-count threshold, an age window, a read-amplification estimate) --- and a
replication follower observes the result without a torn view for free, since it adopts whatever manifest it
loads (\S\ref{sec:storagespi}). The mechanism fits the architecture; only the trigger policy is the
deferred design decision.

\subsection*{Scan-over-compressed pages}

The compression measurement (\S\ref{sec:names}) drew a sharp line: the \emph{inner} bake-off
(frame-of-reference, dictionary, bit-packing) is lightweight and byte-addressable, so it preserves the
engine's \mbox{in-memory~$=$~on-disk} property; the \emph{outer} LZ4 layer buys roughly $2\times$ more ratio
but forfeits it, since an LZ4 page must be decompressed before any node is read. The principled direction
is therefore \emph{not} to lean harder on block compression but to push the inner codecs toward
\textbf{scan-over-compressed}, in the manner of the TU~Munich Data~Blocks design~\cite{datablocks}: a
predicate (\code{price > 100}) and even an aggregate can run directly over frame-of-reference and
bit-packed columns without materialising the decompressed page, using the zone-maps the PAX format already
stores (Chapter~\ref{ch:vector}) to skip whole blocks and SIMD to evaluate the survivors in their packed
form. The projection-index columns are already most of the way there --- they are bit-packed and
zone-mapped --- and the vectorised executor already consumes the packed representation in place rather
than a widened copy: \code{NumberRegionSimd} evaluates predicates and aggregates directly over the
bit-packed bytes, and the projection-index scan is zero-copy. The
payoff is the one the Data~Blocks paper claims: the storage stays small \emph{and} directly scannable, so
the engine never has to choose between compression ratio and its defining zero-copy access --- and never
\emph{needs} a general block compressor like LZ4 at all. The remaining work --- extending the same
treatment to the row-oriented leaf's inner codecs, and to string-predicate SIMD and non-\code{double}
columns --- is a performance capability being completed on a representation that is already correct and
already compact, not a missing one.

\subsection*{Native compilation for the cold analytical query}

An analytical query --- a scan, an aggregate, an ad-hoc temporal question --- is often run \emph{once}, so
it pays the JVM's JIT warm-up in full and never amortises it; the vectorised benchmark
(\S\ref{sec:vecbench}) measures exactly that, a \emph{cold} column dominated by one-time JIT and column
build before the warm path falls to sub-millisecond. The textbook answer is ahead-of-time compilation
(a GraalVM native image, removing warm-up) with profile-guided optimisation (recovering the speculation
the JIT would otherwise perform). A probe on the shipping toolchain (Oracle GraalVM for JDK~25) makes the
trade-off concrete --- and it inverts the obvious worry. The Java Vector API \emph{survives} native
compilation: built with \code{-H:+VectorAPISupport}, the columnar SIMD kernel
(\code{LongVector.fromMemorySegment}, Chapter~\ref{ch:vector}) is intrinsified to machine SIMD at
$\sim$0.24\,ns per element, and the native image even auto-vectorises a plain \code{long[]} reduction
($\sim$0.22\,ns) --- both within $\sim$3--4$\times$ of the JIT. The blocker is narrow and specific: the
\emph{scalar} \code{MemorySegment} accessor --- used for every field read and write by the flyweight-node
model (Chapter~\ref{ch:pds}), and by the column \emph{hydration} that feeds the SIMD --- is \emph{not}
lowered to an inlined load under AOT. Under HotSpot it is fully optimised (\textbf{JMH}: $\sim$0.09\,ns per
element, indistinguishable from the array and SIMD paths), yet under the native image it runs at
$\sim$32\,ns (heap- and arena-backed segments alike, and \code{-O3} does not close it): \emph{two orders of
magnitude} slower, on the one access pattern the whole engine is built on, even as arrays and the Vector
API \emph{in the same binary} stay fast. A native binary would thus speed the SIMD \emph{tip} of an
analytical query while crippling the off-heap body that feeds it, and every other path in the engine
besides. The lever is therefore real but gated --- not on SirixDB but on the ahead-of-time toolchain
closing the FFM scalar-access gap (accessor inlining and the auto-vectorisation HotSpot already does); PGO
recovers branch and call speculation, not this codegen gap. It is recorded with its measurement so the
condition for revisiting it is explicit: viable the day a native image reads an off-heap \code{long} as
fast as HotSpot does.

Five of these six share a pattern worth naming: in each, the \emph{capability} exists and is correct, and what is
deferred is an \emph{acceleration} whose cost is write-time bookkeeping (or executor work) justified only by
a query pattern frequent enough to amortise it. That is the same cost-shape discipline the rest of the
report applies to features that \emph{are} built --- which is why these are presented as principled
deferrals, not gaps. The sixth, native compilation, is deferred for a different reason --- a gap in the
ahead-of-time toolchain rather than a cost SirixDB controls --- and is recorded with its measurement so
the condition for revisiting it stays explicit.

%================================================================
\printbibliography[heading=bibintoc,title={References},prenote=refsnote]

%================================================================\appendix

\chapter{Reference: Notation, Glossary, and Tables}
\label{app:reference}

\section{Notation and parameters}
\label{app:notation}


The quantitative parameters and named constants used throughout the report, with the value the engine
ships and where it is set.

{\small
\begin{xltabular}{\linewidth}{@{}>{\RaggedRight\arraybackslash}p{0.22\textwidth} >{\RaggedRight\arraybackslash}p{0.16\textwidth} >{\RaggedRight\arraybackslash}X@{}}
\toprule
\textbf{Symbol / constant} & \textbf{Value} & \textbf{Meaning and origin}\\
\midrule
\endhead
$f$ (page fan-out) & $1024$ & references per \code{IndirectPage}; fixes trie height $h=\lceil\log_f N\rceil$ (Ch.~\ref{ch:pds})\\
$h$ (trie height) & $\le 5$ & levels Uber\,$\to$\,leaf; $\le 5$ for $N<f^{5}\approx1.1\times10^{15}$\\
$R$ (\code{revisionsToRestore}) & default $3$ & sliding-snapshot window; also the period of full dumps in the differential strategy (\code{revision \% R == 0}) (Ch.~\ref{ch:versioning})\\
slot page capacity & varies & records per \code{KeyValueLeafPage}, bounded by the page byte budget, not a fixed count (\S\ref{sec:slotted})\\
\code{FMESF} & $0.5$ & leaf-match threshold: a Levenshtein ratio above it pairs two leaves (\S\ref{sec:diff})\\
\code{FMESTHRESHOLD} & $0.5$ & inner-node match threshold: fraction of shared matched descendants (\S\ref{sec:diff})\\
DeweyID tier prefixes & 0, 10, 110, 1110, 1111 & unary-style selectors for tiers $0$–$4$ of the division encoding (\S\ref{sec:deweyids})\\
\code{maxDivisionValue[3]} & $270{,}549{,}119$ & upper bound of tier-3 division values before tier-4 is used (\S\ref{sec:deweyids})\\
hash word & $64$ bits & rolling node hash (\code{HashType.ROLLING} default) propagated up the ancestor spine (Ch.~\ref{ch:pds})\\
checksum & XXH3-64 & per-page and per-beacon integrity tag in the durable-commit protocol (Ch.~\ref{ch:ondisk})\\
write permit & $1$ & per-resource \code{Semaphore} permit: one writer, $N$ lock-free readers (Ch.~\ref{ch:rest})\\
$W$ (SIMD lanes) & $4$ or $8$ & \code{LongVector} lanes of the preferred species (AVX2 / AVX-512) (Ch.~\ref{ch:vector})\\
projection chunk & $4$\,KiB & sub-leaf chunk size for the columnar projection index (Ch.~\ref{ch:vector})\\
\bottomrule
\end{xltabular}}

\section{Glossary}
\label{app:glossary}


\begin{description}[leftmargin=0pt,style=nextline,font=\bfseries\sffamily\color{slate}]
\item[AfterCommitState] the policy for a writer at an auto-commit flush: \code{KEEP\_OPEN},
  \code{KEEP\_OPEN\_ASYNC} (rotate the TIL, no intermediate revision), or \code{CLOSE}.
\item[Beacon] one of two redundant on-disk copies of the uber page (offsets 4096/8192), each
  \code{[len][payload][XXH3]}, in separate filesystem blocks.
\item[BiNode] an ephemeral two-child HOT node produced by a split, integrated as a single-mask
  compound node rather than a distinct class.
\item[Bitemporal] tracking two independent time axes: transaction (system) time and user-defined
  valid time.
\item[Brackit] the embedded, set-oriented XQuery/JSONiq engine SirixDB binds its store to.
\item[Cardinality] the optimiser's estimated number of nodes in a path class.
\item[Caret] an even DeweyID division that orders a node without adding a tree level, used to wedge a
  new node between two siblings.
\item[CAS index] content-and-structure index, keyed on a (value, type, path-class) triple.
\item[Clock sweeper] the background second-chance cache-eviction thread.
\item[CoW] copy-on-write.
\item[DeweyID] order-preserving hierarchical node identifier.
\item[Discriminative bit] a key bit on which a HOT node routes between its children.
\item[Division] one component of a DeweyID, one per level from the root.
\item[DPhyp] the Moerkotte--Neumann connected-subgraph\,/\,complement-pair join-ordering dynamic
  program ($O(n\cdot 3^n)$), used up to 20 relations before the greedy GOO fallback.
\item[Epoch tracker] the global record of the oldest revision any live reader is pinned to, consulted
  by MVCC-aware cache eviction.
\item[Eytzinger layout] a breadth-first array order that makes binary search cache-friendly.
\item[FFM] the Foreign Function and Memory API (Project Panama), used for off-heap memory.
\item[FMES] Fast-Match\,/\,Edit-Script, the Chawathe tree-correction algorithm used to import a
  document as an identity-preserving structural delta.
\item[Flyweight node] a thin cursor that reads and writes its fields directly against a page's
  \code{MemorySegment}.
\item[Fragment] a partial on-disk version of a logical page, combined at read time.
\item[Frame of reference] a column codec storing values as small offsets from a per-block base.
\item[Frame slot] a fixed-address off-heap \code{MemorySegment} from the size-classed
  \code{FrameSlotAllocator} pool, recycled with a seqlock version counter.
\item[FSST] a fast static symbol table, a per-page string-compression scheme.
\item[FUA] Force Unit Access; a write flag forcing data through to stable media.
\item[HNSW] Hierarchical Navigable Small World; the layered proximity graph backing the vector $k$-NN
  index.
\item[HOT] Height-Optimised Trie.
\item[Histogram] a per-field value distribution (64 buckets {+} a top-10 most-common-value list) that
  sharpens the optimiser's selectivity estimates.
\item[HyperLogLog] a probabilistic distinct-count sketch kept per path class in PathStats.
\item[io\_uring] the Linux asynchronous I/O interface a storage backend drives through the FFM API
  without JNI.
\item[JSONiq] the JSON query language, a JSON-native dialect of XQuery, that brackit implements.
\item[MemorySegment] an FFM handle to a contiguous off-heap region; a page is one segment.
\item[Mesh] the optimiser's set of equivalence classes, each holding plan alternatives with costs;
  \code{MeshSelectionStage} copies the cheapest onto the live tree.
\item[Morsel-driven parallelism] splitting a scan into fixed-size morsels processed by a worker pool;
  opt-in on the vectorised path.
\item[MVCC] multi-version concurrency control; here, immutable per-revision snapshots.
\item[NodeReferences] an index posting list, backed by a Roaring bitmap of node keys.
\item[Overflow record] a value too large for an inline slot, stored on its own page and reached by a
  \code{PageReference}.
\item[Page guard] a reader's pin on a cached page (a refcount plus a frame-version check) that bars
  eviction while held and detects frame reuse via \code{FrameReusedException}.
\item[Parallelizer] the brackit operator that parallelises a pipeline, running the upstream in a
  producer thread that fills tuple buffers for consumer threads.
\item[Pipeline breaker] an operator (\code{GroupBy}, \code{OrderBy}, join) that must see all input
  before emitting; kept outside a morsel boundary.
\item[ORDPATH] the insert-friendly prefix node-labeling scheme that DeweyIDs follow.
\item[PAX] Partition Attributes Across; an in-page columnar layout.
\item[Path class] an equivalence class of document nodes sharing a path; a node in the path summary.
\item[Path summary] the compact tree of every distinct path class in a resource.
\item[PathStats] per-path-class statistics: count, sum, min/max, a HyperLogLog sketch, and a page
  bitmap.
\item[PCR] path-class reference; a node's path identity in the path summary.
\item[PEXT/PDEP] parallel bit-extract/deposit instructions used by HOT to compress and expand
  partial keys.
\item[Posting list] the set of node keys an index entry points to.
\item[Preallocated commit] the journal-free, low-latency commit profile that writes into preallocated
  file space behind two data barriers.
\item[Projection index] a declared columnar secondary index over a fixed field schema.
\item[Revision] an immutable committed version of a resource.
\item[Roaring bitmap] a compressed integer-set representation used for posting lists and page maps.
\item[Selectivity] the fraction of a path class a predicate is estimated to keep.
\item[Set orientation] brackit's execution model in which operators process tuple \emph{batches}
  (\code{Tuple[]}) under a pull interface, not one tuple at a time (B\"achle).
\item[Shredder] the streaming parser that turns a JSON or XML document into stored nodes.
\item[Sliding snapshot] the default page-versioning strategy: bounded-window reconstruction
  with out-of-window preservation.
\item[Snapshot isolation] the read guarantee that a transaction sees one immutable committed revision
  throughout its lifetime.
\item[Sparse path] the on-path key bits a HOT node keeps for a child via \code{PEXT}.
\item[Storage SPI] the \code{IOStorage}\,/\,\code{Reader}\,/\,\code{Writer} service interface over which
  all backends plug in.
\item[SPLID] stable path labeling identifier; the TU Kaiserslautern refinement of ORDPATH that
  DeweyIDs descend from.
\item[Superblock] the 64-byte validated header at the start of each on-disk file.
\item[TIL] Transaction Intent Log, the in-memory uncommitted page set.
\item[Tombstone] a marker that a node is deleted in a revision, surfaced by \code{sdb:is-deleted}.
\item[Transaction time] when the system recorded a fact; in SirixDB, a revision's commit time.
\item[Uber page] the per-resource root descriptor pointing at the current revision root.
\item[Valid time] when a fact is true in the modeled world; set via valid-from/valid-to.
\item[ValidTimeConfig] the per-resource declaration of the \code{validFrom}/\code{validTo} JSON paths
  the bitemporal functions read for the valid axis.
\item[Vectorised executor] the SIMD, columnar scan path for the supported aggregate and group-by query
  shapes.
\item[Volcano model] the pull-based, one-tuple-at-a-time iterator execution model.
\item[XXH3] the 64-bit hash used for page checksums and node hashing.
\end{description}

\section{Index-type reference}
\label{app:indextypes}

Every index in the system, internal or user-declared, is an \code{IndexType} (a one-byte tag on its
backing page tree). They divide into three groups: \emph{structural} indexes that the engine always
maintains, \emph{change-tracking} indexes that are on by default and make history a lookup
(\S\ref{sec:changetracking}), and \emph{secondary} indexes a user declares for acceleration
(Chapter~\ref{ch:pathsummary}).

{\small
\begin{xltabular}{\linewidth}{@{}>{\ttfamily\small}p{0.21\textwidth} c >{\RaggedRight\arraybackslash}p{0.135\textwidth} >{\RaggedRight\arraybackslash}X@{}}
\toprule
\textnormal{\textbf{IndexType}} & \textbf{tag} & \textbf{Group} & \textbf{What it maps / accelerates}\\
\midrule
\endhead
REVISIONS & 0 & structural & the revision-root directory (revision \textrightarrow\ root page)\\
DOCUMENT & 1 & structural & the primary node index (node key \textrightarrow\ record), the document itself\\
CHANGED\_NODES & 2 & change-track & per revision, the set of changed subtree-roots (a CoW trie off the revision root)\\
RECORD\_TO\_REVISIONS & 3 & change-track & node \textrightarrow\ ascending list of revisions that touched it ($O(k)$ history)\\
PATH\_SUMMARY & 4 & structural & the tree of distinct path classes {+} per-path \code{PathStats}\\
PATH & 5 & secondary & nodes by their path class (declared path index)\\
CAS & 6 & secondary & content-and-structure: (value, type, path) \textrightarrow\ nodes\\
NAME & 7 & secondary & object keys / element names \textrightarrow\ nodes\\
DEWEYID\_TO\_RECORDID & 8 & identity & DeweyID \textrightarrow\ node key, present when \code{useDeweyIDs}; order-preserving range scans\\
VECTOR & 9 & secondary & HNSW $k$-NN over float embeddings (\S\ref{sec:setoriented} naming hazard; experimental)\\
PROJECTION & 10 & secondary & declared columnar covering index feeding the SIMD scan (Ch.~\ref{ch:vector})\\
\bottomrule
\end{xltabular}}

\noindent The structural and change-tracking groups (tags 0--4, plus 8 when DeweyIDs are on) are
maintained unconditionally by the write transaction; the secondary group (tags 5--7, 9, 10) is created
on demand and kept current by per-type listeners (\S\ref{sec:catalogue}). All of them are ordinary
copy-on-write page trees, so each is versioned, time-travellable, and recovered by exactly the same
machinery as the document, an index as of revision~17 is opened the same way the document is.

\section{Configuration reference}
\label{app:config}

\subsection{Per-resource configuration}

A resource's behaviour is fixed at \code{createResource} time through the
\code{ResourceConfiguration} builder and persisted into the resource directory as JSON. The options
that matter for this report:

\dcap{\code{ResourceConfiguration} options and their defaults}
\begin{diagramS}
storageType                   FILE_CHANNEL (default) · MEMORY_MAPPED · FILE · IN_MEMORY · IO_URING · S3
versioningApproach            SLIDING_SNAPSHOT (default) · INCREMENTAL · DIFFERENTIAL · FULL
maxNumberOfRevisionsToRestore the reconstruction window R for the snapshot strategies
hashKind                      ROLLING (default) · POSTORDER · NONE
useDeweyIDs                   store DeweyID labels                            (default off)
buildPathSummary              maintain the path summary                       (default on)
storeDiffs                    write the per-commit JSON change file           (default on)
storeNodeHistory              maintain the RECORD_TO_REVISIONS change index   (default on)
indexBackendType              RBTREE (default) · HOT
byteHandlerPipeline           a compression / encryption handler chain        (default identity)
deweyIdSiblingDistance        inter-sibling gap for new DeweyIDs              (default 16)
\end{diagramS}

\subsection{JVM system properties}

A handful of \code{-D} properties tune behaviour that is not per-resource:

\dcap{Selected \texttt{-D} system properties}
\begin{diagramS}
compression / caching
  -Dsirix.compression               none (default) | lz4    the outer ByteHandler pipeline default
  -Dsirix.cache.recordPage          override the record-page cache share of the off-heap budget
  -Dsirix.cache.recordPageFragment  override the fragment-cache share
  -Dsirix.cache.page                override the metadata page-cache share

FileChannel backend / commit profile
  -Dsirix.filechannel.bufferBytes   pooled direct read buffer size (default 128 KiB)
  -Dsirix.fadvise                   sequential | random | none (default)  per-channel kernel readahead hint
  -Dsirix.commit.preallocated       pre-grow the data file so a commit writes in place (force(false))
  -Dsirix.commit.preallocChunkBytes prealloc growth chunk (default 64 MiB)
  -Dsirix.commit.bufferedBeacons    buffer both beacons behind a single fsync (needs preallocated)
  -Dsirix.revision.file.data.cache.size   RevisionFileData cache entries (default ~1M)

concurrency
  -Ddisable.single.threaded.check   relax the single-writer thread assertion (embedded multi-thread use)
\end{diagramS}

\noindent The per-page codec bake-off (\S\ref{sec:storagespi} and Chapter~\ref{ch:ondisk}) runs
regardless of the \code{-Dsirix.compression} default, choosing the smallest of the built-in codecs per
leaf; the property only changes the outer pipeline applied on top. The
\code{commit.preallocated}/\allowbreak\code{bufferedBeacons} pair is the low-latency commit profile measured in
Chapter~\ref{ch:eval}: it trades a pre-grown file for the removal of the per-commit metadata-journal
fsync, and is off by default because it changes the file-growth pattern an operator may be monitoring.

\section{Temporal-function reference}
\label{app:funcs}

The \code{jn:}, \code{xml:}, and \code{sdb:} function families --- all \emph{added by SirixDB} (registered
by the \code{sirix-query} binding, not part of brackit) --- layer time travel, history, and metadata access
over brackit's XQuery/JSONiq engine (Chapter~\ref{ch:query}). The principal functions:

\renewcommand{\arraystretch}{1.3}\footnotesize\setlength{\tabcolsep}{6pt}
\begin{longtable}{@{}>{\ttfamily}p{5.4cm} >{\RaggedRight}p{10.3cm}@{}}
\toprule
\textnormal{\textbf{Function}} & \textbf{Purpose}\\
\midrule\endhead
jn:doc(coll, res[, rev]) & open a resource, optionally at a given revision number\\
jn:open(coll, res, instant) & open a resource as of a wall-clock instant\\
jn:open-revisions(coll, res, from, to) & open the resource across a revision range\\
jn:valid-at(coll, res, validTime) & facts valid at an instant, in the latest revision\\
jn:open-bitemporal(coll, res, txTime, validTime) & facts valid at an instant, as of a transaction time (both axes)\\
jn:load / jn:store & load or store a document\\
jn:past(item) / jn:future(item) & the item in earlier / later revisions\\
jn:previous(item) / jn:next(item) & the item in the adjacent revision\\
jn:first(item) / jn:last(item) & the item in its first / last revision\\
jn:first-existing / jn:last-existing & first / last revision in which the item exists\\
jn:all-times(item) & the item in every revision it appears in\\
jn:diff(coll, res, r1, r2) & the edit script between two revisions\\
sdb:item-history(item) & the item in each revision that touched it\\
sdb:nodekey(item) & the stable node key\\
sdb:hash(item) & the node's rolling hash\\
sdb:path(item) & the node's path-summary path\\
sdb:revision(item) & the revision number the item was read at\\
sdb:timestamp(item) & the revision's commit timestamp\\
sdb:author-id / sdb:author-name & the commit author\\
sdb:child-count / sdb:descendant-count & structural counts\\
sdb:valid-from / sdb:valid-to & the item's valid-time bounds\\
sdb:is-deleted(item) & whether the item is tombstoned in this revision\\
sdb:select-item(item, nodeKey) & move to a node by key\\
sdb:select-parent(item) & move to the parent node\\
sdb:commit(item[, msg[, dateTime]]) & commit the writer, returning the new revision number\\
sdb:rollback(item) & discard uncommitted changes\\
\bottomrule
\end{longtable}

\section{Brackit physical-operator reference}
\label{app:brackitops}

The set-oriented pipeline of \S\ref{sec:physops} is built from the operators below
(\pp{io/brackit/query/operator/}). Each is created over a single tuple or a \code{Tuple[]} batch
(\code{Operator.create}); the table notes which operators are \emph{pipeline breakers} (they must see the
whole input before emitting), and which spill to disk under memory pressure.

{\small
\begin{xltabular}{\linewidth}{@{}>{\ttfamily\small}p{0.255\textwidth} >{\RaggedRight\arraybackslash}p{0.135\textwidth} >{\RaggedRight\arraybackslash}X@{}}
\toprule
\textnormal{\textbf{Operator}} & \textbf{Role} & \textbf{Behaviour}\\
\midrule
\endhead
Start & pipeline source & seeds the tuple stream from a single binding\\
ForBind & bind (iterate) & evaluates a sequence per input tuple and emits one output tuple per item; optional position binding\\
LetBind & bind (eager) & evaluates a sequence once per input tuple and appends it whole as one binding\\
Select & filter & passes tuples whose predicate is true; emits a dead marker under iteration-nesting\\
Count & enumerate & appends a 1-based position counter, reset at group boundaries\\
NLJoin & join & nested loop, $O(\lvert L\rvert\lvert R\rvert)$; \code{leftJoin} adds left-outer padding; best for tiny or non-equality right input\\
TableJoin & join (breaker) & hash join; builds the table on the smaller input, probes with the larger; \code{HashJoinTable} for equality, \code{SortedJoinTable} for order; mixed key types numerically promoted; left-outer supported\\
GroupBy & group (breaker) & three cursors: \code{AllGroupBy} (no key), \code{SequentialGroupBy} (pre-clustered input), \code{HashGroupBy} (\code{LinkedHashMap}, SIMD key compare); pluggable count/sum/min/max/avg\\
SpillableGroupBy & group (breaker) & radix hash-aggregation spill (DuckDB-style): quarter-heap budget, overflow groups written raw to 64 partitions and re-aggregated, recursively repartitioned under a depth-rotated hash; reused probe keeps the resident-hit path key-allocation-free\\
OrderBy & sort (breaker) & sorts via \code{TupleSort} with a quarter-heap budget; \textbf{spills} sorted runs and external-merge-sorts beyond it\\
Parallelizer & parallelism & runs the upstream in a producer thread filling 2000-tuple \code{Tuple[]} buffers; consumers dequeue whole buffers (\code{BlockingParallelizer} = synchronising form)\\
MorselPipeline & parallelism & morsel-driven: a serial \code{MorselSource} hands 2048-tuple \code{Morsel}s to ForkJoinPool workers, each a pipeline clone; breakers kept outside the morsel boundary\\
VectorizedGroupByExec & vectorised & streaming group-by over reusable 2048-element columnar arrays; in-place aggregation\\
ParallelGroupByExec & vectorised & parallel byte-level group-by/filter/aggregate (count, sum, min, max, avg) over 8\,MiB mmap windows with SIMD field-search; engages only above a $\sim$100\,MB input\\
\bottomrule
\end{xltabular}}

\noindent The translator's \code{SequentialPipelineStrategy} decides the execution mode. It first tries
the optimiser's vectorised-expression annotations (count-distinct, group-by-count, filtered count,
sorted scan, pure aggregate); failing that, if morsel parallelism is enabled and the pipeline is
\emph{morsel-safe} (it begins with a data-driven \code{ForBind} and contains no breaker), it wraps the
root in a \code{MorselPipeline}; otherwise it runs the plain serial \code{Cursor} chain. A process-wide
and a thread-local hook let a host such as SirixDB install its own \code{VectorizedExecutor} (the SIMD
numeric kernel of Chapter~\ref{ch:vector}) for the shapes it can accelerate.

\section{REST API reference}
\label{app:rest}

The Vert.x server (Chapter~\ref{ch:rest}) exposes the store over HTTP/2, routing by path depth over
\code{/:database/:resource}; every business request carries a Bearer JWT. Beyond the endpoint table of
that chapter, reads and writes accept a rich parameter surface.

\subsection{Read and query parameters}

\dcap{Read / query parameters}
\begin{diagramS}
revision, revision-timestamp           select a revision by number or by wall-clock instant
start-revision, end-revision           a revision range (each with a -timestamp variant)
nodeId                                 scope to a sub-document (a node key)
maxLevel, maxChildren, numberOfNodes   depth / fan-out / total-node caps for large resources
startNodeKey, nextTopLevelNodes        the top-level pagination cursor
withMetaData                           nodeKey | nodeKeyAndChildCount | full metadata
prettyPrint                            pretty or compact serialisation
startResultSeqIndex, endResultSeqIndex result-window bounds for a query
query                                  a JSONiq/XQuery to run (global or resource-scoped)
plan                                   return the parsed / optimised / candidate query plan
commitMessage, commitTimestamp         metadata attached to a write
\end{diagramS}

\subsection{Write parameters and modes}

\dcap{Write parameters}
\begin{diagramS}
nodeId                       the target node key
insert                       ASFIRSTCHILD | ASRIGHTSIBLING | ASLEFTSIBLING   (XML adds REPLACE)
hashType                     NONE | ROLLING | POSTORDER   (at creation; REST default NONE)
useDeweyIDs                  store DeweyID labels                            (at creation)
maxNodes                     auto-commit threshold for a large load
validFromPath, validToPath   the document paths carrying valid-time bounds
\end{diagramS}

\noindent Note the layering: the REST \code{PUT}/create endpoint defaults \code{hashType} to
\code{NONE} when the parameter is omitted, whereas the \emph{embedded} \code{ResourceConfiguration}
defaults to \code{ROLLING} (\S\ref{app:config}); a deployment that wants audit-grade hashing over
REST must pass \code{hashType=ROLLING} explicitly at creation.

An \code{ETag} equal to the node hash provides optimistic concurrency: a write must present the current
ETag, or it is told the resource may have changed in the meantime.

\section{Node-kind reference}
\label{app:nodekinds}

Every record carries a one-byte \code{NodeKind} tag. The \emph{data} kinds, the ones a document is made
of, are below; the engine additionally defines internal kinds for index nodes (\code{PATH},
\code{CASRB}, \code{PATHRB}, \code{NAMERB}), change tracking (\code{REVISION\_REFERENCES\_NODE}), the
name dictionary (\code{HASH\_ENTRY}), DeweyID mapping, projection/vector indexes, and a
\code{DELETE} tombstone, none of which a user authors directly.

{\small
\begin{xltabular}{\linewidth}{@{}>{\ttfamily\small}p{0.265\textwidth} c >{\RaggedRight\arraybackslash}X@{}}
\toprule
\textnormal{\textbf{NodeKind}} & \textbf{tag} & \textbf{Role}\\
\midrule
\endhead
\multicolumn{3}{@{}l}{\emph{JSON}}\\
JSON\_DOCUMENT & 31 & the document root\\
OBJECT & 24 & an object container (anonymous; no path node)\\
ARRAY & 25 & an array container\\
STRING\_VALUE & 30 & a bare string scalar\\
NUMBER\_VALUE & 28 & a bare numeric scalar\\
BOOLEAN\_VALUE & 27 & a bare boolean scalar\\
NULL\_VALUE & 29 & a bare null scalar\\
OBJECT\_NAMED\_STRING & 50 & \emph{fused} key{+}string value in one record\\
OBJECT\_NAMED\_NUMBER & 49 & \emph{fused} key{+}number value\\
OBJECT\_NAMED\_BOOLEAN & 48 & \emph{fused} key{+}boolean value\\
OBJECT\_NAMED\_NULL & 51 & \emph{fused} key{+}null value\\
OBJECT\_NAMED\_OBJECT & 52 & \emph{fused} key whose value is an object (children sit under it directly)\\
OBJECT\_NAMED\_ARRAY & 53 & \emph{fused} key whose value is an array\\
\midrule
\multicolumn{3}{@{}l}{\emph{XML}}\\
XML\_DOCUMENT & 9 & the document root\\
ELEMENT & 1 & an element node\\
ATTRIBUTE & 2 & an attribute (reached by \code{moveToAttribute})\\
NAMESPACE & 13 & a namespace binding (\code{moveToNamespace})\\
TEXT & 3 & character data\\
COMMENT & 8 & a comment node\\
PROCESSING\_INSTRUCTION & 7 & a processing instruction\\
\bottomrule
\end{xltabular}}

\noindent The six \code{OBJECT\_NAMED\_*} kinds are the fusion optimisation
(Ch.~\ref{ch:pds}): the common ``\code{key: value}'' shape becomes \emph{one} record instead of an
object-key node plus a separate value node, halving the node count for typical JSON. The trade is the
wire-shape subtlety the GUI's parsers handle (Ch.~\ref{ch:gui}): a fused container's children sit as a
bare array directly under its value, with no intermediate key-node wrapper. The scalar value kinds
(\code{STRING\_VALUE}\,\ldots) appear bare only as direct array elements or document roots; inside an
object they are fused into the \code{OBJECT\_NAMED\_*} parent.

\section{Index-management function reference}
\label{app:indexfuncs}

Secondary indexes (\S\ref{sec:catalogue}, \S\ref{app:indextypes}) are created and inspected from the
query language through the \code{jn:} family below. Creating an index returns the document node;
\code{find-*} returns the index number; once an index exists the optimiser rewrites matching predicate
and path queries onto it automatically (Ch.~\ref{ch:query}), so explicit \code{scan-*} calls are rarely
needed in application code.

{\small
\begin{xltabular}{\linewidth}{@{}>{\ttfamily\footnotesize}p{0.46\textwidth} >{\RaggedRight\arraybackslash}X@{}}
\toprule
\textnormal{\textbf{Function}} & \textbf{Effect}\\
\midrule
\endhead
jn:create-cas-index(\$doc, \$type?, \$paths?) & build a content-and-structure index (optional value type, optional path set)\\
jn:create-path-index(\$doc, \$paths?) & build a path index over the given paths (or index-all)\\
jn:create-name-index(\$doc, \$names?) & build a name index over object keys / element names\\
jn:find-cas-index(\$doc, \$type, \$path) & the index number of a matching CAS index, or $-1$\\
jn:find-path-index(\$doc, \$path) & the index number of a matching path index\\
jn:find-name-index(\$doc, \$name) & the index number of a matching name index\\
jn:scan-cas-index(\$doc, \$idx, \$key, \ldots) & low-level point scan of a CAS index by number\\
jn:scan-cas-index-range(\$doc, \$idx, \$lo, \$hi, \ldots) & low-level inclusive range scan of a CAS index\\
jn:scan-path-index(\$doc, \$idx, \$pcr) & low-level scan of a path index by path-class reference\\
jn:scan-name-index(\$doc, \$idx, \$name) & low-level scan of a name index by name\\
\bottomrule
\end{xltabular}}

\noindent The same functions exist under the \code{xml:} prefix for XML resources. A representative
session creates a typed CAS index and then lets an ordinary predicate use it:

\dcap{Create an index, then let the optimiser reach for it}
\begin{diagramS}
jn:create-cas-index(jn:doc('db','res'), 'xs:decimal', '/store/book/[]/price')   (* returns the doc node *)

(* thereafter this is rewritten onto the CAS index automatically (\S\ref{app:indextypes}): *)
jn:doc('db','res').store.book[][?$$.price ge 10 and $$.price le 50]
\end{diagramS}

\section{Storage-backend reference}
\label{app:backends}

The pluggable backends behind the \code{IOStorage} SPI (\S\ref{sec:storagespi}) are summarised here on
the axes that matter for choosing one: how reads and writes work, what durability they give, whether
reads are zero-copy, and what they are for. The three built-ins ship in the core; \code{IO\_URING} and
\code{S3} arrive through a \code{ServiceLoader} from the enterprise modules.

{\small
\begin{xltabular}{\linewidth}{@{}>{\ttfamily\footnotesize}p{0.135\textwidth} >{\RaggedRight\arraybackslash}X >{\RaggedRight\arraybackslash}p{0.31\textwidth}@{}}
\toprule
\textnormal{\textbf{Backend}} & \textbf{Read / durability} & \textbf{Use \& notes}\\
\midrule
\endhead
FILE\_CHANNEL & pread into a pooled 128\,KiB direct buffer; fsync {+} dual beacon over 3 channels (data buffered, revisions O\_SYNC, beacon O\_DSYNC) & the default; the two-file V0 layout the report describes\\
MEMORY\_MAPPED & slice a mapped \code{MemorySegment} (zero-copy), \code{madvise} SEQUENTIAL; writes reuse \code{FileChannelWriter}, so durability is identical to FILE\_CHANNEL & read-heavy, warm resources; generational arena remap on growth\\
IN\_MEMORY & \code{ConcurrentHashMap} get/put; no durability & tests / ephemeral resources only; persists nothing\\
IO\_URING & zero-copy, optionally busy-polled, batched (N pages in 2 syscalls); NVMe FUA, same dual-beacon contract & enterprise; high-concurrency NVMe via the FFM io\_uring SPI (Ch.~\ref{ch:enterprise})\\
S3 & ranged GET (206) resolved through the manifest, never a LIST; commit point is the revisions object (segment\,\textrightarrow\,manifest\,\textrightarrow\,revisions) & enterprise; immutable segment objects, own remapped layout, \code{AwsV4Signer}, no AWS SDK\\
\bottomrule
\end{xltabular}}

\noindent Orthogonal to the backend is the \textbf{byte-handler pipeline} that wraps each page payload:
\code{DeflateCompressor} (stream zlib), \code{FFILz4Compressor} (FFM-native LZ4, zero-copy over
\code{MemorySegment}s, \code{FAST}/\code{HIGH\_COMPRESSION}), and \code{Encryptor} (Tink streaming AEAD),
composed compress-then-encrypt; the pipeline is zero-copy only if every handler in it is
(\S\ref{sec:storagespi}). And \textbf{Kafka} is not in this table because it is not an \code{IOStorage}
at all: it is a replication sidecar (Ch.~\ref{ch:enterprise}) that streams committed-revision pointers
from a leader to followers, with S3 remaining the durable source of truth. Whichever backend serves the
bytes, the page layer above sees only \code{IOStorage}, which is why a resource can be moved between
backends without the engine noticing (Property~\ref{prop:backendid}).

\section{JSONiq cookbook}
\label{app:cookbook}


A few representative queries make the temporal surface concrete.

\dcap{Time travel and history}
\begin{diagramS}
(* the document as of revision 17 *)
jn:doc('db','res', 17)

(* the document as of a wall-clock instant *)
jn:open('db','res', xs:dateTime("2026-01-15T00:00:00Z"))

(* every value a field ever held, oldest to newest *)
jn:all-times(jn:doc('db','res').store.book[0].price)

(* the revisions in which one node changed *)
sdb:item-history(jn:doc('db','res').store.book[0])
\end{diagramS}

\dcap{Diff, update, and bitemporal}
\begin{diagramS}
(* the edit script between revisions 5 and 9 *)
jn:diff('db','res', 5, 9)

(* an atomic sub-document update (one new revision) *)
replace json value of jn:doc('db','res').store.book[0].price with 19.99

(* insert a field, then commit explicitly *)
insert json {"discount": 0.1} into jn:doc('db','res').store.book[0]
sdb:commit(jn:doc('db','res'))

(* facts currently believed valid at an instant (latest revision) *)
jn:valid-at('db','res', xs:dateTime("2026-03-01T00:00:00Z"))

(* facts believed valid at an instant AS OF a past transaction time (both axes) *)
jn:open-bitemporal('db','res', xs:dateTime("2026-01-20T00:00:00Z"),
                               xs:dateTime("2026-03-01T00:00:00Z"))
\end{diagramS}

\dcap{Aggregation over a covered column}
\begin{diagramS}
(* count and average over a projection-indexed field; runs on the SIMD path *)
let $rows := jn:doc('db','t').rows[]
return { "n":   count($rows[.status eq "ACTIVE"]),
         "avg": avg($rows[.status eq "ACTIVE"] ! .amount) }
\end{diagramS}

\dcap{Secondary indexes: create, then let the optimiser use them}
\begin{diagramS}
(* a path index over two paths; returns the document node *)
jn:create-path-index(jn:doc('db','res'), ('/store/book', '/store/book/[]/price'))

(* a content-and-structure (CAS) index, typed, over one path *)
jn:create-cas-index(jn:doc('db','res'), 'xs:decimal', '/store/book/[]/price')

(* a name index over object keys, then look an index up by number *)
jn:create-name-index(jn:doc('db','res'))
jn:find-cas-index(jn:doc('db','res'), 'xs:decimal', '/store/book/[]/price')
(* once created, an ordinary predicate query is rewritten onto the index automatically:
   jn:doc('db','res').store.book[][?$$.price gt 10]   -> CAS-index range scan *)
\end{diagramS}

\dcap{Branching and rewriting history}
\begin{diagramS}
(* edit the past: an updating query targeting an older revision revertTo()s it first,
   appending a NEW head branched from that revision (revisions after it are kept) *)
replace json value of jn:doc('db','res', 3).store.book[0].price with 9.99

(* read back the two heads: the original latest and the branched one *)
jn:doc('db','res')          (* the most recent committed revision *)
jn:doc('db','res', 3)       (* the point the branch was based on *)
\end{diagramS}

\dcap{Structural and navigational metadata}
\begin{diagramS}
(* a node's stable key, path-summary path, and structural counts *)
let $b := jn:doc('db','res').store.book[0]
return { "key":         sdb:nodekey($b),
         "path":        sdb:path($b),
         "children":    sdb:child-count($b),
         "descendants": sdb:descendant-count($b),
         "hash":        sdb:hash($b) }

(* jump to a node by its stable key, then to its parent *)
sdb:select-parent(sdb:select-item(jn:doc('db','res'), 42))
\end{diagramS}

\dcap{Revision-spanning analytics: a field's whole history as a time series}
\begin{diagramS}
(* (revision, commit-time, value) for one field across every revision it appears in *)
for $v in jn:all-times(jn:doc('db','res').store.book[0].price)
return { "rev":   sdb:revision($v),
         "at":    sdb:timestamp($v),
         "price": $v }

(* aggregate over history without materialising a single tree: the peak price ever recorded *)
max( jn:all-times(jn:doc('db','res').store.book[0].price) )
\end{diagramS}

\dcap{Comparing two points in time: as-of-then vs.\ as-of-now}
\begin{diagramS}
(* what changed for one node between revision 5 and the latest, structurally *)
let $then := jn:doc('db','res', 5).store.book[0]
let $now  := jn:doc('db','res').store.book[0]
return { "wasPrice": $then.price, "isPrice": $now.price,
         "delta": $now.price - $then.price }

(* the node's value at the revision closest to a wall-clock instant *)
jn:open('db','res', xs:dateTime("2026-03-01T00:00:00Z")).store.book[0].price
\end{diagramS}

\dcap{The XQuery / XML face of the same patterns}
\begin{diagramS}
(* XPath navigation + time travel over an XML resource (App.~O) *)
xml:open('db','lib', xs:dateTime("2026-01-15T00:00:00Z"))/library/book[1]/title/text()

(* fold an external XML document in as an identity-preserving structural delta (FMES; XML-only) *)
xml:import('db','lib', doc("/path/to/new-library.xml"))

(* store-level temporal accessors work on XML items too *)
sdb:item-history(xml:doc('db','lib')/library/book[1])
\end{diagramS}

\dcap{Aggregations and grouping (set-oriented / vectorised where covered)}
\begin{diagramS}
(* group-by-count over a field; runs on the columnar SIMD path when a projection index covers it *)
let $r := jn:doc('db','t').rows[]
for $row in $r
group by $s := $row.status
return { "status": $s, "n": count($row) }

(* a filtered multi-aggregate in one pass *)
let $a := jn:doc('db','t').rows[][?$$.amount gt 0]
return { "n": count($a), "sum": sum($a ! .amount), "max": max($a ! .amount) }
\end{diagramS}

\chapter{On-Disk Binary Format}
\label{app:format}


This appendix consolidates the wire format described in prose across Ch.~\ref{ch:pds} (the slotted leaf
page) and Ch.~\ref{ch:ondisk} (the superblock, beacons, and commit) into one byte-level reference. A
resource is two append-only files, \pp{sirix.data} and \pp{sirix.revisions}. Integers are
\textbf{little-endian} throughout --- the superblock, revision records, and the page-record length
prefixes and primitives alike --- so the bytes are host-independent; the superblock's endianness word is
kept as a defensive gate that rejects a foreign-endian host at open. The only \textbf{big-endian} fields
are the canonical page-hash word and the HOT trie keys, where byte order is deliberately sort order.

\dcap{Page-kind tags (the one-byte \texttt{PageKind} prefix)}
\begin{diagramS}
tag   page kind           role
  1   KeyValueLeafPage    a slotted leaf of records (primary data and index leaves)
  3   UberPage            the per-resource root descriptor
  4   IndirectPage        one trie level (an adaptive reference container)
  5   RevisionRootPage    the root of one revision, with refs to the document index and auxiliaries
 12   HOTLeafPage         a HOT index leaf (when the HOT backend is selected)
the metadata pages (NamePage, PathSummaryPage, CASPage, PathPage, DeweyIDPage, OverflowPage) carry
their own tags and follow the same tag-then-body framing.
\end{diagramS}

\section{File geometry and the superblock}

\begin{diagram}
sirix.data                                sirix.revisions
  0       superblock (role=DATA, 64 B)      0      superblock (role=REVISIONS, 64 B)
  4096    PRIMARY beacon slot   (4 KiB)     4096   REVISIONS_RECORDS_START;
  8192    SECONDARY beacon slot (4 KiB)            slot r at 4096 + r*32, 32 bytes:
  12288   DATA_REGION_START, page records          [u64 offset][u64 epochMs]
          [u32 len][payload], 8-aligned ▶          [u64 checksum][u64 pageHash]
\end{diagram}

\begin{diagram}
Superblock (64 bytes, little-endian, at offset 0 of both files)
  off  size  field
  0    8     magic  "SIRIXDB!"  (8 ASCII bytes)
  8    4     layout version (= 0)
  12   1     role  (0 = DATA, 1 = REVISIONS)            13  3   reserved / flags
  16   4     endianness check word (= 0x01020304)
  20   4     beacon-slot size  (DATA 4096; REVISIONS 0)
  24   8     primary-beacon offset (DATA 4096; REVISIONS 0)
  32   8     content-start  (DATA 12288; REVISIONS 4096)
  40   16    reserved (future resource UUID)
  56   8     XXH3-64 over bytes [0, 56)
\end{diagram}

\section{Beacons and revision records}

\begin{diagram}
Beacon slot — one 4096-byte block; PRIMARY @4096, SECONDARY @8192 (separate blocks):
  [u32 len][ UberPage payload … ][u64 XXH3 over payload][zero padding to 4096]

Revision record — 32 bytes, slot r at REVISIONS_RECORDS_START + r*32:
  [u64 rootPageOffset][u64 epochMillis][u64 checksum][u64 rootPageHash]
  checksum = XXH3 over the first 24 bytes (hash present) or 16 (legacy; rootPageHash == 0)
\end{diagram}

A worked trace shows how the two files grow, and what is appended versus overwritten, across three
commits. Only the two beacon slots are ever rewritten in place; the data region and the revision-record
slots are strictly append-only.

\dcap{File growth across three commits (offsets illustrative, geometry exact)}
\begin{diagramS}
                              sirix.data                              sirix.revisions
bootstrap (rev 0, empty)      [..12288] superblock+beacons            [0] superblock
                              12288: initial pages, RevRootPage@R0    slot 0 @4096: {offset R0, ts, xxh3, hash}
                              beacons@4096/8192 → point at R0

commit rev 1 (write 5 nodes)  APPEND leaf + indirect path + RevRoot   APPEND slot 1 @4128:
                              at the current end-of-data → R1         {offset R1, ts, xxh3, hash}
                              OVERWRITE both 4KiB beacons → R1        (revisions file grew by 32 bytes)

commit rev 2 (change 1 node)  APPEND only the 1 changed leaf, its     APPEND slot 2 @4160:
                              2 root-path indirect pages, RevRoot     {offset R2, ts, xxh3, hash}
                              → R2 ; OVERWRITE both beacons → R2
                              (pages NOT on that path stay shared with rev 1; the file only grew)

read rev 1 after rev 2        beacons name R2, but slot 1 still names R1 → open RevisionRootPage at R1,
                              traverse rev-1's immutable pages; rev 2's appended pages are simply never visited
\end{diagramS}

\noindent The invariant the trace exposes is that a commit never moves or rewrites a byte the previous
revision depended on: it appends the copied root-to-leaf path and a fresh revision-root, appends one
32-byte revision record, and swaps the two beacons. That is why an interrupted commit can damage at most
the tail and the two beacon slots, never an acknowledged revision (Property~\ref{prop:torn}), and why
reading any past revision is a direct traversal rather than an undo.

\section{Page framing and page kinds}

A page record is \code{[u32 len][payload]}, 8-byte aligned (the UberPage fills its whole 4\,KiB beacon
slot; the RevisionRootPage is 256-aligned). The payload, after the byte-handler pipeline, begins with a
one-byte page-kind tag and a one-byte \code{BinaryEncodingVersion} (\code{V0}):

\begin{diagramS}
[u8 pageKind][u8 V0][ body ]      Page-kind tags (1 byte; tag 7 is unused):
  1 KeyValueLeafPage   4 IndirectPage      8 CASPage       11 DeweyIDPage      14 BitmapChunkPage
  2 NamePage           5 RevisionRootPage  9 OverflowPage  12 HOTLeafPage      15 VectorPage
  3 UberPage           6 PathSummaryPage  10 PathPage      13 HOTIndirectPage  16 ProjectionPage
\end{diagramS}

\section{Reference containers}

Every reference-bearing page embeds a child-reference container, written as a \textbf{1-byte sub-tag}
then the body (\code{RevisionRootPage} is the lone exception: it omits the sub-tag and always uses a
bitmap container). A single \code{PageReference}, and the three container layouts (sub-tags 0/1/2,
promoting $4 \to \text{bitmap} \to \text{full}$ as a level fills):

\begin{diagramS}
PageReference   [u8 fragCount][ fragment × (i32 rev, i64 key) ][i64 key][u8 hashPresent][hash:8?]
   (hash = XXH3-64; databaseId/resourceId are NOT stored — re-injected on open)
0 ReferencesPage4       [u8 size][ size × PageReference ][ size × i16 offset ]
1 BitmapReferencesPage  [i16 bitmapLen][bitmap][ cardinality × PageReference ]   (position from set bits)
2 FullReferencesPage    [i16 bitmapLen][bitmap]; per set bit ascending: [i64 key][PageReference body]
\end{diagramS}

\section{The structural pages}

Each follows the \code{[u8 tag][u8 V0]} prefix; ``delegate'' is a reference container. Widths:
\code{i32}/\code{i64} are fixed little-endian, \code{varint}/\code{varlong} stop-bit, \code{u8} one byte,
and an \emph{array} is an \code{i32} count then the elements.

\begin{diagramS}
UberPage          [i32 revisionCount]
IndirectPage      [delegate]
RevisionRootPage  [delegate — NO sub-tag][i32 revision][i64 maxNodeKeyDoc][i64 maxNodeKeyChanged]
                  [i64 maxNodeKeyRecordToRev][i64 revisionTimestamp][bool msg? → i32 len + bytes]
                  [u8 levelDoc][u8 levelChanged][u8 levelRecToRev][bool user? → utf8 name, utf8 uuid]
NamePage          [delegate][i32+i64 maxNodeKeys][i32+i64 maxHotPageKeys][i32 numberOfArrays]
                  [i32+u8 currentMaxLevels][i32 liveEntries; per dict: i32 roaringLen + bytes]
PathSummaryPage   [delegate][i32+i64 maxNodeKeys][i32+u8 currentMaxLevels]
CAS/Path/Projection [delegate][i32+i64 maxNodeKeys][i32+i64 maxHotPageKeys][i32+u8 currentMaxLevels]
VectorPage        [delegate][i32+i64 maxNodeKeys][i32+u8 currentMaxLevels]
DeweyIDPage       [delegate][i64 maxNodeKey][u8 currentMaxLevel]            (scalars, not arrays)
OverflowPage      [i32 dataLen][dataLen bytes]                             (one opaque overflow value)
\end{diagramS}

The \code{maxNodeKey}/\code{maxHotPageKey} arrays size each index tree's trie, \code{currentMaxLevels}
its current height. \code{NamePage} additionally persists the five interned-name dictionaries
(attributes, elements, namespaces, processing-instructions, JSON object keys) through the delegate's
referenced pages, as Roaring64 bitmaps of live entry keys.

\section{The HOT pages}

\begin{diagramS}
HOT_LEAF_PAGE     [varint pageKey][i32 revision][u8 indexType][i16 commonPrefixLen][commonPrefix]
  full dump       [i32 entryCount (bit31 = completeDump)][i32 usedSlotMem][ count × i32 slotOffset ][slotMem]
  sparse dump     [i32 dirtyCount][i32 dirtyUsed][ dirtyCount × i32 packedOffset ][packed]
  one slot        [u16 suffixLen][suffix][u16 valueLen][value]
HOT_INDIRECT_PAGE [varint pageKey][i32 revision][u8 height][u8 nodeType][u8 layoutType][i32 numChildren]
  single-mask     [i16 initialBytePos][i64 bitMask][i16 mostSigBit]
  multi-mask      [i16 mostSigBit][i16 nExtractBytes][positions][ ceil(n/8) × i64 extractionMask ]
  partialKeys     numChildren × (u8 | i16 | i32   — width by disc-bit count)
  children        numChildren × ( i64 key, u8 fragCount, fragCount × (i32 rev, i64 key) )
  trailer         [256-byte childIndex]   (single-mask + multi-node only)
BITMAP_CHUNK_PAGE [varint pageKey][ opaque BitmapChunkPage blob ]
\end{diagramS}

\section{The slotted leaf page (\texttt{KeyValueLeafPage})}

The leaf page is the format's centrepiece, a slotted page that is \emph{in-memory == on-disk},
column-compressed then byte-compressed. Its frame and the 160-byte header it copies raw:

\begin{diagramS}
[u8 0x01][u8 V0][varlong recordPageKey][i32 revision][u8 indexType]
[160-byte header+bitmap region][encoded body][i32 regionTableLen|0 + PAX columns][overlong entries]

header+bitmap (160 B):  0 8 recordPageKey · 8 4 revision · 12 2 populatedCount · 14 4 heapEnd
  18 4 heapUsed · 22 1 indexType · 23 1 flags(bit0 deweyIDs, bit1 FSST) · 24 8 reserved
  32 128 slotBitmap (16 × u64 = 1024 bits; bit set = slot populated)
\end{diagramS}

The \textbf{encoded body} carries the records as a column store, then one compression codec:

\begin{diagramS}
[i32 populatedCount][i32 onDiskHeapSize][u8 templateCount]
 templateCount>0 :  [u8 structFlags][i32 templatePoolBytes][i32 compLen][u8 codec][ blob ]
 templateCount=0 :  [i32 0][i32 compLen][u8 codec][ blob ]
structFlags: 0x01 hashElision · 0x02 parentKeyCol · 0x04 pathNodeKeyCol · 0x08 valueElision · 0x10 nameKeyElision
\end{diagramS}

The decompressed \code{blob} is, in order: (1) a \textbf{compact directory}, \code{populatedCount}
$\times$ 4 bytes \code{[dataLength:3][kindId:1]}; (2) the \textbf{offset-table template pool} (the
deduplicated per-record field-offset tables); (3) per-slot \textbf{template ids}; then, gated by
\code{structFlags}, optional columns: (4) a \textbf{zero-hash bitmap} (slots whose all-zero hash was
stripped), (5) a \textbf{parent-key column}, (6) a dictionary-encoded \textbf{path-node-key column},
(7) a \textbf{value-elision} and (8) a \textbf{name-key-elision} table (fused \code{OBJECT\_NAMED\_*}
values/names lifted into PAX regions); and finally (9) the \textbf{record heap}. A heap record (the
\code{dataLength} bytes a slot points at) is:

\begin{diagramS}
[u8 nodeKind][ offset table: fieldCount × u8 ][ data: delta-varint / varint / fixed ][ payload ][ deweyId ][u16 len]
\end{diagramS}

\noindent The one-byte offset table gives $O(1)$ random field access, and the trailing \code{u16} length
lets a reader find the record's extent back-to-front. The stored DeweyID is delta-encoded against the
previous record's label (a one-byte shared-prefix length plus the divergent suffix), so a run of
siblings costs only a byte or two each. The 8\,KB slot \emph{directory}
(\code{[heapOffset:4][dataLength:3][kindId:1]} $\times$ 1024) and the preservation bitmap are
reconstructed on read, never stored.

\section{Node records, by kind}

A node's fields and encodings (delta-varint encodes a neighbour as \code{zigzag(field − ownNodeKey)+1},
so 0 marks a null neighbour and adjacent siblings cost one byte; \code{hash}/\code{descendantCount} are
omitted when \code{HashType==NONE}):

\begin{diagramS}
OBJECT (24)          parentKey, rightSib, leftSib, firstChild, lastChild  (delta-varint);
                     prevRev, lastModRev (signed varint); childCount? (signed varint);
                     hash:8 + descendantCount (signed varint)
ARRAY (25)           = OBJECT, plus a pathNodeKey (delta-varint) after lastChild
OBJECT_NAMED_NUMBER  parentKey, rightSib, leftSib (delta-varint); nameKey (signed varint);
  (fused, 49)        pathNodeKey (delta-varint); prevRev, lastModRev (signed varint); hash:8;
                     number payload [u8 type][data]: 0 Double8 · 1 Float4 · 2 int · 3 long · 4 BigInt · 5 BigDec
STRING_VALUE (30)    parentKey, rightSib, leftSib (delta-varint); prevRev, lastModRev (signed varint);
                     payload [u8 fsstCompressed][signed-varint len][bytes]    (no hash field)
ELEMENT (XML, 1)     parentKey, rightSib, leftSib, firstChild, lastChild, pathNodeKey (delta-varint);
                     prefixKey, localNameKey, uriKey, prevRev, lastModRev (signed varint); childCount?;
                     hash:8 + descendantCount; payload [varint attrCount][attrKeys…][varint nsCount][nsKeys…]
\end{diagramS}

\section{Varints and compression codecs}

\textbf{Varint} is LEB128 (7 data bits/byte, high bit continues, little-endian); \textbf{signed varint}
zigzags first ($(v\!\ll\!1)\oplus(v\!\gg\!63)$); \textbf{delta-varint} is \code{zigzag(key − base) + 1}
(0 = null). \code{recordPageKey} and \code{BigInteger} lengths instead use Chronicle stop-bit varlongs.
The body is then run through one of four codecs (\code{codec} id selects the decoder); the smallest
encoding wins:

\begin{diagramS}
0 ZeroRunByteCodec   [0xFF][varint size][ literal runs (0x00–0x7F) · zero runs (0x80–0xFE, 0xFF+varint) ]
2 ByteRunCodec       [0xFE][varint size][ literal · zero-run · const-run (+1 value byte) · long runs ]
3 SirixLZ77Codec     [0xFD][varint size][ LZ4-block tokens: (litNib<<4 | matchNib), literals, u16 dist ]
1 FFILz4Compressor   [i32 size][payload]   — size < 0 ⇒ (−size) literal bytes STORED uncompressed
\end{diagramS}

\noindent Page hashing is XXH3-64 over the compressed bytes, stored big-endian in the parent
\code{PageReference}; the per-page \code{V0} byte is the only format-version hook. The page-record fields
are pinned little-endian (this hash word and the HOT keys are the deliberate big-endian exceptions), so
the format is host-independent; the pinning also corrected a latent big-endian fault in the
word-at-a-time varint decode, which had read the leading byte in host order and would have decoded
backwards on a big-endian machine.

\chapter{Selected Algorithms and Formal Properties}
\label{app:algorithms}


This appendix collects, in pseudocode, several algorithms described narratively in the body.

\dcap{Hash-guided diff (\texttt{DiffFactory})}
\begin{diagramS}
diff(old, new):
  if hash(old) == hash(new):  emit SAMEHASH ; return          // prune the whole subtree
  if kind or value differs:   emit UPDATE / REPLACE
  align children by node key:
     child only in new   → emit INSERT
     child only in old   → emit DELETE
     matched child pair  → diff(child_old, child_new)         // recurse only where hashes differ
\end{diagramS}

\dcap{FMES fast match (\texttt{FMSEAlgorithm.match}), per NodeKind bucket}
\begin{diagramS}
for each NodeKind present in both trees:
   pairs ← longestCommonSubsequence(old_nodes, new_nodes, cmp)   // Myers O(ND) diff
   add each LCS pair to the matching ; remove them
   greedily pair any remaining equal nodes
// leaf match: similarity > FMESF = 0.5    inner match: |common| / max(|x|,|y|) >= 0.5
\end{diagramS}

\dcap{Access-path cost (\texttt{JsonCostModel})}
\begin{diagramS}
seqCost(N)           = ceil(N / 1024) * 1.0  +  N * 0.01
indexCost(card, sel) = 3.0
                     + ceil(ln(card)/ln 16) * 0.2 * 4.0   // trie-traversal I/O
                     + ceil(card*sel / 512) * 4.0         // matched-leaf data I/O
                     + (card*sel) * 0.01                  // per-result CPU
prefer index  when  indexCost < seqCost
\end{diagramS}

\dcap{Sliding-snapshot read combine (\texttt{combineRecordPages}); cf.\ the write path, Alg.~3.1}
\begin{diagramS}
combineRecordPages(chain F = [f0 (newest) … f_m], window R):   // the READ path
  complete ← empty page
  for f in f0 … f_{min(R-1, m)}:           // the newest R fragments only
      for slot s populated in f:
          if s not yet in complete:  complete[s] ← f[s]    // newest writer wins; never overwritten
  return complete                          // fully materialised; Property 3.1 guarantees ≤ R reads
\end{diagramS}

\dcap{MVCC-aware clock-sweep eviction (\texttt{ClockSweeper}); cf.\ \S\ref{sec:buffering}}
\begin{diagramS}
evict(targetBytes):
  watermark ← epochTracker.oldestLiveReaderRevision()      // global MVCC epoch
  freed ← 0
  while freed < targetBytes:
      e ← hand.advance()                   // circular scan over resident fragments
      if e.referenced:                     // second-chance bit set by the last hit
          e.referenced ← false ;  continue // spare it for one sweep
      if liveReaderCouldNeed(e, watermark): // a snapshot ≤ watermark may still read e
          continue                         // MVCC-safe: never reclaim below the watermark
      reclaim(e) ;  freed += e.size         // free the off-heap MemorySegment
\end{diagramS}

\dcap{DeweyID careting insert (\texttt{SirixDeweyID}); cf.\ \S\ref{sec:deweyids}, default distance 16}
\begin{diagramS}
newSibling(prev, next):                    // insert a node between siblings prev < next
  if next is none:        return prev with last division += 16        // append on the right
  mid ← (prev.last + next.last) / 2
  if mid is odd and prev.last < mid < next.last:
        return prev[:−1] ++ [mid]          // an odd value fits between them
  else: return prev ++ [even caret] ++ [16]  // no gap: descend one level via an even caret
  // odd = real node, even = caret (adds depth, not a level); no existing label is ever rewritten
\end{diagramS}

\section{Algorithms: concurrency and search}
\label{app:algorithms2}

The first part of this appendix collects the storage and query algorithms; this part collects the
concurrency and search algorithms developed in Chapters~\ref{ch:pds}, \ref{ch:ondisk},
\ref{ch:hot}, and~\ref{ch:pathsummary}, in pseudocode.

\dcap{Page-guard acquire / release (\texttt{KeyValueLeafPage}, \texttt{PageGuard}); cf.\ \S\ref{sec:reclaim}}
\begin{diagramS}
acquireGuard():
  flags ← STATE_FLAGS.volatileGet()
  if flags has CLOSED:               return false
  if flags has ORPHANED and guardCount ≤ 0:  return false      // no resurrection at zero guards
  guardCount.incrementAndGet() ;     return true

read with guard:
  g.versionAtFix ← page.version                                // recorded AFTER the guard is held
  … traverse page.MemorySegment …
release (g.close):
  v ← page.version ;  page.releaseGuard()
  if v ≠ g.versionAtFix:  throw FrameReusedException           // frame recycled under us → reload
\end{diagramS}

\dcap{Optimistic frame-slot read (\texttt{FrameSlotAllocator}); even/odd seqlock; cf.\ \S\ref{sec:allocator}}
\begin{diagramS}
read(classIdx, slotIdx):
  v0 ← slotVersion.acquireGet(classIdx, slotIdx)
  if (v0 & 1) ≠ 0:  retry                       // odd = a writer (alloc/release) is in progress
  … copy the slot bytes …
  v1 ← slotVersion.acquireGet(classIdx, slotIdx)
  if v1 ≠ v0:  retry                            // version moved → the slot was recycled mid-read
  // else the bytes are a consistent snapshot of this slot generation
\end{diagramS}

\dcap{HNSW $k$-NN search (\texttt{HnswGraph.searchKnn}); cf.\ \S\ref{sec:setoriented} naming hazard}
\begin{diagramS}
searchKnn(query, k, efSearch):
  if graph empty: return []
  ep ← entryPoint
  for layer in maxLevel … 1:                    // greedy descent through the upper layers
      ep ← greedilyHopToClosest(ep, query, layer)
  R  ← searchLayer(ep, query, ef = efSearch, layer = 0)   // beam search at the base
  return the k nearest of R (skip tombstoned nodes)

searchLayer(ep, query, ef, layer):              // candidate min-heap C, result max-heap W, visited set
  push ep into C and W
  while C not empty:
      c ← C.popClosest()
      if dist(c) > W.farthest():  break         // early termination
      for n in neighbours(c, layer) not visited:
          if dist(n,query) < W.farthest() or |W| < ef:
              push n into C and W ;  trim W to ef
  return W
\end{diagramS}

\dcap{Morsel-driven dispatch (\texttt{MorselPipeline}); cf.\ \S\ref{sec:setoriented}}
\begin{diagramS}
run(upstream):
  source ← MorselSource(upstream)               // serial producer; nextMorsel() is synchronized
  spawn N = poolSize workers, each:
      loop:
          m ← source.nextMorsel()               // grab a 2048-tuple batch (null ⇒ done)
          if m is null: break
          process m through a PER-WORKER clone of the non-breaker pipeline
          deposit results into the shared output queue
  consumer drains the output queue              // breakers (GroupBy/OrderBy) sit OUTSIDE this boundary
\end{diagramS}

\section{List of properties and theorems}
\label{app:properties}


The formal statements the report relies on, gathered for reference. Each is stated, and where noted
proved or test-gated, at the cross-reference given.

{\small
\begin{xltabular}{\linewidth}{@{}l >{\RaggedRight\arraybackslash}p{0.26\textwidth} >{\RaggedRight\arraybackslash}X@{}}
\toprule
\textbf{\#} & \textbf{Name} & \textbf{Statement}\\
\midrule
\endhead
Prop.~\ref{prop:snapshot} & Snapshot immutability & the pages reachable from a committed root are fixed for all time, so a revision-$N$ read is a consistent, lock-free snapshot\\
Thm.~\ref{thm:dewey} & Order preservation & \code{compareTo} on two DeweyIDs agrees with \code{compareUnsigned} on their encoded bytes; document order \emph{is} byte order\\
Prop.~\ref{prop:norelabel} & Insertion without relabelling & inserting a node assigns a fresh DeweyID and rewrites no existing label, preserving every order/ancestor/sibling relation\\
Prop.~\ref{prop:bounded} & Bounded reconstruction & a logical page is reconstructable from at most $R$ fragments at any revision (sliding snapshot)\\
Prop.~\ref{prop:amort} & Even amortisation & each record is preserved at most once per window pass, spreading carry-forward cost ($\approx 1/R$ extra writes/rev) instead of bunching it\\
Prop.~\ref{prop:torn} & Torn-block safety & whichever beacon survives a torn-block crash names a durable tail, because the tail is fsynced before any beacon and the copies occupy distinct blocks\\
Prop.~\ref{prop:defstats} & Deferred-stats cost & maintaining per-path statistics over a transaction touching $V$ values across $P$ paths costs $O(P)$ CoW operations, not $O(V)$\\
Prop.~\ref{prop:histcost} & History-lookup cost & listing the $k$ revisions in which a node changed costs $O(k)$ from \code{RECORD\_TO\_REVISIONS}, independent of total revisions or document size\\
Prop.~\ref{prop:backendid} & Backend-independent format & a resource round-trips byte-identically across the io\_uring and FileChannel backends (cross-backend test)\\
Prop.~\ref{prop:pdi} & Physical data independence & one logical \code{Operator} plan runs serial, parallel, morsel, or vectorised --- and over any storage backend --- chosen at translation through SPIs, with no plan rewrite (B\"achle; Ch.~\ref{ch:query})\\
Prop.~\ref{prop:serializable} & Per-resource serializability & any interleaving of one writer and $N$ readers on a resource is equivalent to a serial schedule (the single write permit excludes write skew)\\
\bottomrule
\end{xltabular}}

\noindent The through-line is that each property is local and checkable: immutability and torn-block
safety are structural consequences of copy-on-write and the commit ordering; the DeweyID theorem and the
cost properties are gated by the tests named in their home sections; and per-resource serializability
follows from the single-writer permit. None is asserted without either a proof sketch in place or a test
that exercises it.

\chapter{Reproduction, Benchmarking, and Operations}
\label{app:benchmethod}


So that every measured number in Chapters~\ref{ch:vector} and~\ref{ch:eval} is traceable, this appendix
records the harnesses, versions, and conditions. All runs were single-machine, the client and server
co-located, warm caches, on the host of \S\ref{sec:w1} (Linux~6.8). SirixDB ran embedded on OpenJDK~25
with the Vector API and FFM enabled (\code{-Xmx} 4--8\,GB, \code{FILE\_CHANNEL}, \code{SLIDING\_SNAPSHOT});
the comparison engines ran as noted below.

\textbf{The three benchmarks.}
\begin{itemize}[leftmargin=1.4em]
\item \textbf{Small-document W1--W6} (\S\ref{sec:w1}): a Java driver against an embedded SirixDB and a
  PostgreSQL~17 \code{jsonb}{+}trigger history table; a single $\sim$2.4\,KB document, 5\,001 versions.
\item \textbf{Large-document crossover} (\S\ref{sec:crossover}): a single-field update to a
  100\,KB\,/\,1\,MB\,/\,10\,MB document, 100 durable commits, 101 versions. SirixDB embedded (full and lean
  configs); PostgreSQL~17 (Docker, \code{jsonb}{+}trigger, \code{synchronous\_commit=on}, server-side
  loop); DuckDB~1.5.2 (\code{JSON} column, 101 appended versions). Storage measured after
  \code{CHECKPOINT}; per-commit latency verified with \code{clock\_timestamp()}.
\item \textbf{Aggregate / vectorisation} (\S\ref{sec:vecbench}): a JSON array of $N$ records; four
  aggregate shapes run through brackit's interpreter and the \code{SirixVectorizedExecutor} (correctness
  gated by \code{TypedGroupByDifferentialTest} et al.); the same aggregates on DuckDB~1.5.2
  (\code{read\_json} into native columns) and PostgreSQL~17 (\code{jsonb} rows). Vectorised \emph{cold}
  uses a fresh executor (empty per-(session,revision) cache, a real scan), \emph{warm} a kept executor
  (cache hit).
\item \textbf{History-retainer cost shape} (\S\ref{sec:retainers}): a flat $K$-field record, one field
  updated 1\,000 times, full history retained, across SirixDB (durable embedded), \textbf{Datomic} (peer
  1.0.7622, \code{datomic:mem}, in-memory), \textbf{XTDB v1} (1.24.5, in-memory node, async submit),
  \textbf{XTDB v2} (latest stable Docker image, driven over its PostgreSQL wire protocol), and
  PostgreSQL~17. Per-update latency measured as $K$ grows.
\item \textbf{Lock-free read scaling} (\S\ref{par:lockfreeread}): $T$ threads, each holding one read-only
  transaction, doing random point reads over a 1\,MB document for a fixed window; aggregate reads/s at
  $T\in\{1,2,4,8,16\}$ on the 20-core host.
\end{itemize}

\textbf{JMH harnesses.} The per-operation measurements --- commit latency (\S\ref{sec:commitphases}), write
and read scaling (\S\ref{sec:concurrency}), the recomputed diff (\S\ref{sec:diffbench}), and per-feature
ingest cost (\S\ref{sec:ingestcost}) --- are run through the \textbf{Java Microbenchmark Harness}
(JMH~1.37) in the project's \pp{sirix-benchmarks} module rather than a hand-rolled timing loop. JMH forks a
fresh JVM per measurement, runs warmup iterations before timing, and reports a mean with a 99.9\,\%
confidence interval over many samples; the throughput benchmarks take their thread count from a
\code{-Pjmh.threads} parameter so one benchmark sweeps $1\ldots16$ readers or writers. This removes a class
of error a hand-rolled loop cannot reliably exclude --- the cold-JIT \emph{ordering bias} that inflates
whichever configuration runs first. Migrating the scaling benchmarks to JMH revised three figures
\emph{downward} toward the truth (write scaling from a spurious $8.9\times$ to $4.8\times$ at eight threads,
read scaling from $4.75\times$ to $3.5\times$ at sixteen), and the convergence of successively more careful
measurements is itself the argument for using the harness. The hand-rolled drivers are retained for the
multi-process and cross-engine benchmarks (PostgreSQL, DuckDB, Datomic, XTDB) that JMH's single-JVM model
does not fit.

\begin{honesty}
\small\textbf{What is and isn't comparable.} The benchmarks are deliberately narrow (one machine, one
writer, one document/array shape, warm) and each favours a different side: the small-document regime
favours PostgreSQL, the large-document regime SirixDB, the aggregate regime DuckDB. Across systems the
\emph{drive and durability models are not normalised}, an embedded engine pays no IPC, an in-memory store
pays no fsync, a client-server store pays a round trip, so absolute latencies across systems are
indicative, not a leaderboard; the \emph{cost shapes} (flat vs linear in size, interpreter vs vectorised)
are the durability-independent findings the report leans on. SirixDB is a dev build; the comparison
engines are GA/stable releases. Raw per-query logs were captured for each run.
\end{honesty}

\section{Measured results in full}
\label{sec:rawdata}

The complete numbers behind the body's tables, gathered for reference (ms unless noted).

\dcap{Large-document crossover (\S\ref{sec:crossover}): per-commit and storage, 101 versions}
\begin{diagramS}
                     SirixDB(full)  SirixDB(lean)  PostgreSQL 17   DuckDB 1.5.2
ms/commit  100 KB       9.97           7.41           2.09            (n/a)
           1 MB        10.66           7.39          19.25            (n/a)
           10 MB       10.13           8.18         193.9             (n/a)
storage    100 KB     1.07 MB         0.63 MB        2.26 MB          11 MB
           1 MB       2.62 MB         1.49 MB         20 MB           98 MB
           10 MB     17.6 MB          9.99 MB        193 MB          965 MB
\end{diagramS}

\dcap{Aggregate / vectorisation (\S\ref{sec:vecbench}): ms per query}
\begin{diagramS}
N = 1,000,000        interp.   vec.cold   vec.warm   DuckDB   PG(jsonb)
group-by-count        8705      2701.0      0.19       8        180
filtered count        8462       191.2      0.24       0.4       68
sum + avg            16713         0.53      0.24       1         64
two-key group-by      8704         9.06      0.21       8        269
N = 100,000          interp.   vec.cold   vec.warm   DuckDB   PG(jsonb)
group-by-count          40.8      6.79      0.26       5         17.5
filtered count          39.8    143.0       0.23       1         13.2
sum + avg               65.5      0.59      0.23       1         22.7
two-key group-by        53.8      4.64      0.40       7         25.7
\end{diagramS}

\dcap{History-retainer cost shape (\S\ref{sec:retainers}): ms per single-field update}
\begin{diagramS}
record (fields)   SirixDB   Datomic   XTDB v1   XTDB v2   PostgreSQL
  ~1 KB   (30)      7.4       0.22      0.37      8.6        0.4
  ~10 KB  (300)    11.1       0.20      2.55     13.1        ---
  ~100 KB (3000)   10.0       0.20     24.50      ---        2.1
SirixDB/PostgreSQL 1 MB / 10 MB points: 10.7 / 10.1  and  19.3 / 193.9 (from the crossover above).
\end{diagramS}

\dcap{Lock-free read scaling (\S\ref{par:lockfreeread}): aggregate point reads/s, JMH, 20-core host}
\begin{diagramS}
threads     1          2           4           8           16
reads/s   7.8 M      11.3 M      16.2 M      24.5 M      27.4 M
speedup   1.00       1.44        2.08        3.15        3.52
\end{diagramS}

These are single-machine, warm-cache, single-writer measurements; the honesty boxes
beside each body table state what is and is not comparable. The cost \emph{shapes}, flat versus linear in
document size, interpreter versus vectorised, are the robust findings; the absolute cross-system latencies
are indicative given the unnormalised durability and drive models recorded above.



\section{Building, testing, and benchmarking}
\label{app:build}

\subsection{Build}

SirixDB is a Gradle multi-module build (JDK~25 with preview features). The deployable REST jar is the
shadow (``fat'') jar of \code{sirix-rest-api}.

\dcap{Common Gradle targets}
\begin{diagramS}
./gradlew :sirix-core:test            the storage-engine test suite
./gradlew :sirix-rest-api:test        the REST-API tests (run these before any REST change)
./gradlew :sirix-rest-api:shadowJar   the deployable -fat.jar
./gradlew :sirix-benchmarks:jmh       the JMH micro-benchmarks
\end{diagramS}

The SolidJS front-end builds with \code{npm run build} (\code{dev}, \code{typecheck}, \code{lint}, and
\code{format} for development).

\subsection{Fault-injection and differential tests}

Two harnesses guard the hardest properties.

\begin{itemize}[leftmargin=1.4em]
\item \textbf{Crash recovery.} \code{CrashRecoveryInjectionTest} (\code{-Dsirix.crash.run=true}) forks a
  writer, \code{SIGKILL}s it at random instants in its commit loop, and verifies on reopen that the
  database opens, every acknowledged commit is intact, all pages deserialise, and the resource accepts a
  new writer.
\item \textbf{Vectorised correctness.} The columnar/SIMD path is checked by
  \emph{interpreter-differential} testing: each supported query shape runs both through the vectorised
  executor and through the scalar interpreter, and the results must match exactly, which is what lets the
  path stay fail-closed to only the shapes it has proven.
\end{itemize}

A round-trip fidelity sweep over many JSON shapes, and a cross-backend test that round-trips a resource
between the io\_uring and FileChannel backends (Property~\ref{prop:backendid}), round out the regression
gates.


The architecture this report describes runs unchanged whether embedded in a Java process or served over
the network. The following sections record how the public demo is deployed, as a concrete operational reference.

\section{Deployment: the single-container demo}
\label{app:deploy}

The demo runs one SirixDB JVM container behind \textbf{nginx} (TLS termination and static-asset serving
for the SolidJS front-end), a small \textbf{Node.js sandbox-proxy} (request routing, demo-data seeding,
orphan cleanup), and \textbf{Keycloak} (OAuth2/OIDC). A single host suffices because one SirixDB
container, sized for write throughput, already needs most of the machine's memory.

\dcap{The demo deployment topology}
\begin{diagramS}
browser ──HTTPS──▶ nginx ──▶ sandbox-proxy (Node.js) ──▶ SirixDB REST (JVM container, HTTP/2)
                     │                  │                         │
                static SolidJS    seeds demo data           one writer / N readers per resource
                     assets       via an admin token
                                       └──▶ Keycloak (OAuth2 / OIDC; demo user has only the view role)
\end{diagramS}

\section{Operational constraints}

A few constraints are worth recording because they are easy to get wrong:

\begin{itemize}[leftmargin=1.4em]
\item \textbf{Heap.} The JVM container needs an \code{-Xmx} of at least 8\,GB, or large \code{PUT}
  (database-creation) requests stall; \code{-Ddisable.single.threaded.check} is required, and
  \code{-XX:+AlwaysPreTouch} must \emph{not} be set, since it forces the full resident set upfront. One
  such container fits comfortably only on a 16\,GB host.
\item \textbf{Read-only enforcement, in depth.} The Keycloak demo user carries only the \code{view}
  role, the REST layer rejects writes for a role-less user server-side, and the front-end's
  \code{isCloudDemo} flag hides every write control. Demo data is seeded server-side by the proxy with
  admin credentials, never by the demo user.
\item \textbf{Start-up ordering.} The container manager waits for Keycloak to be listening before
  launching SirixDB, so a cold host reboot does not crash-loop the backend on a refused auth connection.
\item \textbf{Availability.} The proxy binds its port before seeding and runs the seed in the
  background, so a slow or failed seed degrades gracefully instead of returning 502 for the whole site.
\end{itemize}

\section{Performance-tuning guide}
\label{app:tuning}

The defaults are chosen for correctness and general-purpose behaviour; this section collects, as
actionable guidance, the levers a deployment can pull and the trade each carries. Every knob is described
in its home chapter; here they are organised by the goal they serve.

\textbf{Sizing memory.} Give the JVM \code{-Xmx} of at least 8\,GB for write workloads (a smaller heap
stalls large \code{PUT}s, \S\ref{app:deploy}), set \code{-Ddisable.single.threaded.check} for embedded
multi-thread callers, and do \emph{not} set \code{-XX:+AlwaysPreTouch} (it forces the full resident set
up front). The page cache is off-heap and budgeted separately; \code{-Dsirix.cache.recordPage} and its
siblings shift the off-heap budget between the record-page, fragment, and metadata caches when a
workload is skewed toward one (\S\ref{sec:buffering}).

\textbf{Lowering commit latency.} The default commit pays a metadata-journal fsync on every file growth.
For a write-heavy, latency-sensitive deployment on a journalling filesystem, enable the preallocated
profile (\code{-Dsirix.commit.preallocated}, optionally with \code{-Dsirix.commit.bufferedBeacons}): it
pre-grows the file so a commit writes in place and skips the journal fsync, the 1.29--1.51$\times$
commit-rate gain measured in Chapter~\ref{ch:eval}. The trade is a changed file-growth pattern an
operator may be monitoring, which is why it is off by default.

\textbf{Tuning reads.} \code{-Dsirix.fadvise=sequential} biases the kernel toward read-ahead (good for
cold analytical scans and projection hydration); \code{random} suppresses it (good for point lookups).
The \code{MEMORY\_MAPPED} backend serves reads zero-copy from a mapped segment and suits read-heavy,
already-warm resources; \code{FILE\_CHANNEL} (default) is the safe general choice; \code{IO\_URING}
(enterprise) is for high-concurrency NVMe (Ch.~\ref{ch:enterprise}).

\textbf{Trading history for throughput.} The per-resource sliding-snapshot window
(\code{revisionsToRestore}) is the central read/write dial: lower it on a read-mostly resource to make
every historical read cheaper, raise it on a write-mostly one to write less per commit
(Ch.~\ref{ch:versioning}). \code{HashType.NONE} buys about 13\,\% more commit throughput by dropping the
rolling hash, at the cost of hash-guided diff and per-node integrity, choose it only when neither is
needed (\S\ref{sec:w1}).

\textbf{Bulk loading.} For a one-shot import, drive \code{insertSubtreeAs\ldots} so the bulk-insert flag
suppresses per-node diff work, and use a count-bounded \code{KEEP\_OPEN\_ASYNC} writer so the intent log
rotates in $O(1)$ and the load leaves a single revision instead of dozens (\S\ref{sec:worked} and the
auto-commit section). Memory then tracks the rotation, not the document.

\textbf{Accelerating analytics.} Declare a \code{PROJECTION} index over the hot field list to put
aggregate and group-by-count queries on the SIMD columnar path (Chapter~\ref{ch:vector}); a covering
projection turns a multi-second interpreter aggregate into a sub-millisecond warm repeat (measured,
\S\ref{sec:vecbench}). Morsel-driven
parallelism is opt-in for the brackit pipeline (\S\ref{sec:setoriented}); the vectorised columnar path is
fail-closed to its supported shapes and otherwise falls back transparently.

\textbf{Scaling concurrency.} Write concurrency lives \emph{across} resources, not within one: a single
resource serialises its writers (Property~\ref{prop:serializable}), so partition write-hot data into
separate resources and let the JVM-wide buffer pool and per-resource write permits do the rest. Reads
never contend and scale freely.

\begin{honesty}
\small\textbf{Caveat.} These are levers, not a tuned configuration. Several interact (the preallocated
profile needs a count-based, not timed, auto-commit; \code{bufferedBeacons} needs \code{preallocated};
\code{KEEP\_OPEN\_ASYNC} is \code{FILE\_CHANNEL}-only), and the headline gains quoted here are from the
specific machines and workloads of Chapter~\ref{ch:eval}, not guarantees. The honest default advice is to
measure on the target workload before turning anything off that protects correctness (hashing, checksum
verification) for throughput.
\end{honesty}

\end{document}
